\chapter{Caaguazú}\label{dept:Caaguazú}
\mapaparaguai{5}{Caaguazú}
\vfill\clearpage
\section{Caaguazu /\allowbreak  pacote 07}

O pacote 07 fecha a leitura documental de , com 22 localidades e o
departamento mais logisticamente central do interior paraguaio.

\begin{itemize}
\tightlist
\item
  Coronel Oviedo é o nó rodoviário nacional (cruzamento PY02/\allowbreak PY08) e
  polo de serviços do interior.
\item
  A base agropecuária e agroindustrial (soja, milho, cana) sustenta alta
  resiliência produtiva regional.
\item
  A leitura local está ancorada nos  de cada localidade
  e no capitulo .
\item
  Lacuna de solar/\allowbreak pluviometria da capital fechada com dados NASA POWER
  (2001-2020).
\end{itemize}

\subsection{Ficha climática de Coronel Oviedo (capital
departamental)}

Fonte: NASA POWER Climatology API, período 2001-2020, coordenadas
-25.44°S /\allowbreak  -56.44°W.

\textbf{Irradiância solar global --- ALLSKY\_SFC\_SW\_DWN (kWh/\allowbreak m²/\allowbreak dia):}

\begin{longtable}[]{@{}
  >{\raggedright\arraybackslash}p{(\linewidth - 22\tabcolsep) * \real{0.0833}}
  >{\raggedright\arraybackslash}p{(\linewidth - 22\tabcolsep) * \real{0.0833}}
  >{\raggedright\arraybackslash}p{(\linewidth - 22\tabcolsep) * \real{0.0833}}
  >{\raggedright\arraybackslash}p{(\linewidth - 22\tabcolsep) * \real{0.0833}}
  >{\raggedright\arraybackslash}p{(\linewidth - 22\tabcolsep) * \real{0.0833}}
  >{\raggedright\arraybackslash}p{(\linewidth - 22\tabcolsep) * \real{0.0833}}
  >{\raggedright\arraybackslash}p{(\linewidth - 22\tabcolsep) * \real{0.0833}}
  >{\raggedright\arraybackslash}p{(\linewidth - 22\tabcolsep) * \real{0.0833}}
  >{\raggedright\arraybackslash}p{(\linewidth - 22\tabcolsep) * \real{0.0833}}
  >{\raggedright\arraybackslash}p{(\linewidth - 22\tabcolsep) * \real{0.0833}}
  >{\raggedright\arraybackslash}p{(\linewidth - 22\tabcolsep) * \real{0.0833}}
  >{\raggedright\arraybackslash}p{(\linewidth - 22\tabcolsep) * \real{0.0833}}@{}}
\toprule\noalign{}
\begin{minipage}[b]{\linewidth}\raggedright
Jan
\end{minipage} & \begin{minipage}[b]{\linewidth}\raggedright
Fev
\end{minipage} & \begin{minipage}[b]{\linewidth}\raggedright
Mar
\end{minipage} & \begin{minipage}[b]{\linewidth}\raggedright
Abr
\end{minipage} & \begin{minipage}[b]{\linewidth}\raggedright
Mai
\end{minipage} & \begin{minipage}[b]{\linewidth}\raggedright
Jun
\end{minipage} & \begin{minipage}[b]{\linewidth}\raggedright
Jul
\end{minipage} & \begin{minipage}[b]{\linewidth}\raggedright
Ago
\end{minipage} & \begin{minipage}[b]{\linewidth}\raggedright
Set
\end{minipage} & \begin{minipage}[b]{\linewidth}\raggedright
Out
\end{minipage} & \begin{minipage}[b]{\linewidth}\raggedright
Nov
\end{minipage} & \begin{minipage}[b]{\linewidth}\raggedright
Dez
\end{minipage} \\
\midrule\noalign{}
\endhead
\bottomrule\noalign{}
\endlastfoot
6.74 & 6.20 & 5.50 & 4.48 & 3.35 & 2.88 & 3.21 & 3.96 & 4.59 & 5.43 &
6.43 & 6.81 \\
\end{longtable}

Média anual: 4.96 kWh/\allowbreak m²/\allowbreak dia. Inclinação recomendada: 25° Norte (anual);
35° inverno; 15° verão.

\textbf{Precipitação --- PRECTOTCORR (mm/\allowbreak mês):}

\begin{longtable}[]{@{}
  >{\raggedright\arraybackslash}p{(\linewidth - 22\tabcolsep) * \real{0.0833}}
  >{\raggedright\arraybackslash}p{(\linewidth - 22\tabcolsep) * \real{0.0833}}
  >{\raggedright\arraybackslash}p{(\linewidth - 22\tabcolsep) * \real{0.0833}}
  >{\raggedright\arraybackslash}p{(\linewidth - 22\tabcolsep) * \real{0.0833}}
  >{\raggedright\arraybackslash}p{(\linewidth - 22\tabcolsep) * \real{0.0833}}
  >{\raggedright\arraybackslash}p{(\linewidth - 22\tabcolsep) * \real{0.0833}}
  >{\raggedright\arraybackslash}p{(\linewidth - 22\tabcolsep) * \real{0.0833}}
  >{\raggedright\arraybackslash}p{(\linewidth - 22\tabcolsep) * \real{0.0833}}
  >{\raggedright\arraybackslash}p{(\linewidth - 22\tabcolsep) * \real{0.0833}}
  >{\raggedright\arraybackslash}p{(\linewidth - 22\tabcolsep) * \real{0.0833}}
  >{\raggedright\arraybackslash}p{(\linewidth - 22\tabcolsep) * \real{0.0833}}
  >{\raggedright\arraybackslash}p{(\linewidth - 22\tabcolsep) * \real{0.0833}}@{}}
\toprule\noalign{}
\begin{minipage}[b]{\linewidth}\raggedright
Jan
\end{minipage} & \begin{minipage}[b]{\linewidth}\raggedright
Fev
\end{minipage} & \begin{minipage}[b]{\linewidth}\raggedright
Mar
\end{minipage} & \begin{minipage}[b]{\linewidth}\raggedright
Abr
\end{minipage} & \begin{minipage}[b]{\linewidth}\raggedright
Mai
\end{minipage} & \begin{minipage}[b]{\linewidth}\raggedright
Jun
\end{minipage} & \begin{minipage}[b]{\linewidth}\raggedright
Jul
\end{minipage} & \begin{minipage}[b]{\linewidth}\raggedright
Ago
\end{minipage} & \begin{minipage}[b]{\linewidth}\raggedright
Set
\end{minipage} & \begin{minipage}[b]{\linewidth}\raggedright
Out
\end{minipage} & \begin{minipage}[b]{\linewidth}\raggedright
Nov
\end{minipage} & \begin{minipage}[b]{\linewidth}\raggedright
Dez
\end{minipage} \\
\midrule\noalign{}
\endhead
\bottomrule\noalign{}
\endlastfoot
131 & 136 & 131 & 152 & 172 & 90 & 75 & 52 & 92 & 187 & 190 & 166 \\
\end{longtable}

Média anual: \textasciitilde1.573 mm/\allowbreak ano. Estação chuvosa Out-Mai; seca
Jun-Set.
\clearpage
\subsection*{Dados Departamentais}
\subsubsection{Desenvolvimento Humano e
Social}

\subsection{IDH Departamental}

\begin{longtable}[]{@{}ll@{}}
\toprule\noalign{}
Indicador & Valor \\
\midrule\noalign{}
\endhead
\bottomrule\noalign{}
\endlastfoot
IDH & 0,715 \\
Classificação IDH & alto \\
Ranking nacional (1=melhor) & 7/\allowbreak 18 \\
Esperança de vida & 73,1 anos \\
Anos de escolaridade média & 8,4 anos \\
RNB per capita (USD PPA) & 9.000 \\
\end{longtable}

\subsection{Indicadores de Pobreza}

\begin{longtable}[]{@{}ll@{}}
\toprule\noalign{}
Indicador & Valor \\
\midrule\noalign{}
\endhead
\bottomrule\noalign{}
\endlastfoot
Taxa de pobreza total (\%) & 28,5\% \\
Taxa de pobreza extrema (\%) & 6,8\% \\
Índice de Gini & 0,447 \\
\end{longtable}

\subsection{Infraestrutura Básica}

\begin{longtable}[]{@{}ll@{}}
\toprule\noalign{}
Indicador & Valor \\
\midrule\noalign{}
\endhead
\bottomrule\noalign{}
\endlastfoot
Acesso a água potável (\%) & 79,8\% \\
Acesso a saneamento (\%) & 78,3\% \\
\end{longtable}

\subsubsection{Segurança Pública}

\subsection{Indicadores}

\begin{longtable}[]{@{}
  >{\raggedright\arraybackslash}p{(\linewidth - 2\tabcolsep) * \real{0.6111}}
  >{\raggedright\arraybackslash}p{(\linewidth - 2\tabcolsep) * \real{0.3889}}@{}}
\toprule\noalign{}
\begin{minipage}[b]{\linewidth}\raggedright
Indicador
\end{minipage} & \begin{minipage}[b]{\linewidth}\raggedright
Valor
\end{minipage} \\
\midrule\noalign{}
\endhead
\bottomrule\noalign{}
\endlastfoot
Taxa de homicídios (por 100k hab) & \textasciitilde6--8 (estimativa,
próxima ou acima da média nacional) \\
Crimes registrados (total/\allowbreak ano) & --- dados não desagregados
publicamente \\
Número de comisarías & \textasciitilde20 (estimativa --- departamento
extenso e populoso) \\
Índice segurança relativo & médio \\
\end{longtable}

\subsubsection{Mercado de Terra Rural}

\subsection{Preços por Tipo de Terra}

\begin{longtable}[]{@{}llll@{}}
\toprule\noalign{}
Tipo & Preço (USD/\allowbreak ha) & Faixa & Tendência \\
\midrule\noalign{}
\endhead
\bottomrule\noalign{}
\endlastfoot
Agrícola alta produtividade & 8.800 & 6.000 -- 12.000 & ↑ \\
Pastagem & 4.000 & 2.800 -- 5.500 & ↑ \\
Mata /\allowbreak  reserva & 1.800 & 1.000 -- 3.000 & → \\
\end{longtable}

\subsection{Características do
Mercado}

\begin{longtable}[]{@{}ll@{}}
\toprule\noalign{}
Parâmetro & Valor \\
\midrule\noalign{}
\endhead
\bottomrule\noalign{}
\endlastfoot
Valorização anual média (\%) & 8--10\% \\
Lote mínimo rural (ha) & 10 \\
Imposto territorial rural (\%) & 1\% sobre valor fiscal \\
Liquidez do mercado & alta \\
\end{longtable}
\clearpage
\newpage
\secaoDiagnostico{Caaguazu}{\mapaConteudo{-25.4500}{-56.0166}}{Caaguazú é o coração logístico e industrial do departamento. Oferece uma
infraestrutura de serviços e uma capacidade de produção de recursos
(madeira, grãos e proteína) que a tornam um ativo de suprimento
inigualável. Contudo, sua força logística é também sua fraqueza
estratégica: a localização na principal rodovia do país a expõe a fluxos
populacionais intensos em cenários de crise urbana na capital ou Ciudad
del Este.

\subsubsection{Por que sim}

\begin{itemize}
\tightlist
\item
  Solo de altíssima produtividade e abundância de recursos industriais
  básicos.
\item
  Infraestrutura de saúde e comércio superior à média do interior.
\item
  Conectividade rápida com as principais cidades do país via asfalto
  duplicado.
\item
  Sem restrições da Lei de Fronteira.
\end{itemize}

\subsubsection{Por que não}

\begin{itemize}
\tightlist
\item
  Alvo estratégico logístico de alta visibilidade nacional.
\item
  Concentração populacional urbana elevada para um refúgio isolado.
\item
  Vulnerabilidade social em cenários de interrupção de suprimentos
  rodoviários.
\end{itemize}

\subsubsection{Recomendação
Técnica}

Altamente indicado como base logística, industrial ou para residências
rurais periféricas (10-20 km de distância do centro). Requer um
planejamento de segurança focado na proteção contra fluxos migratórios
massivos. O GSS de 6.8 reflete a solidez dos recursos e da segurança
institucional da região. Referencia atual de score: GSS 6.8 em
2026-03-05.

\subsubsection{}

\begin{itemize}
\tightlist
\item
  \begin{enumerate}
  \def\labelenumi{\arabic{enumi})}
  \tightlist
  \item
    Dados censitários completos (INE\footnote{Instituto Nacional de Estadística (Instituto Nacional de Estatística), órgão oficial de dados demográficos e censitários.} 2022).
  \end{enumerate}
\item
  \begin{enumerate}
  \def\labelenumi{\arabic{enumi})}
  \setcounter{enumi}{1}
  \tightlist
  \item
    Relatórios de indústria madeireira e agrícola (MIC/\allowbreak MAG).
  \end{enumerate}
\item
  \begin{enumerate}
  \def\labelenumi{\arabic{enumi})}
  \setcounter{enumi}{2}
  \tightlist
  \item
    Mapa de conectividade e duplicação da Ruta PY02 (MOPC).
  \end{enumerate}
\item
  \begin{enumerate}
  \def\labelenumi{\arabic{enumi})}
  \setcounter{enumi}{3}
  \tightlist
  \item
    Registro de segurança pública e governança municipal.
  \end{enumerate}
\end{itemize}}
\begin{table}[ht!]
\small \begin{tabular}{p{3.0cm}p{0.8cm}p{6.0cm}} \toprule
\textbf{Critério} & \textbf{Nota} & \textbf{Justificativa} \\ \midrule
A - Ameaças Estratégicas & 5.5 & Hub logístico tático vital na PY02; visibilidade estratégica elevada \\
B - Risco Social & 6.0 & Densidade metropolitana significativa; risco social em crises de suprimento urbana \\
C - Riscos Naturais & 7.0 & Geologia muito estável e baixo risco de grandes desastres naturais \\
D - Autossuficiência & 8.5 & Recursos agroindustriais de elite: solos férteis, madeira, proteína animal e água \\
E - Institucional & 8.0 & Forte segurança institucional e serviços de saúde consolidados na região \\
\midrule \textbf{GSS FINAL} & \textbf{6.8} & \textbf{Moderadamente Seguro (Ideal como base administrativa e de suprimentos para o projeto).} \\ \bottomrule \end{tabular}
\end{table}
\subsubsection{Vulnerabilidades e
Mitigação}\label{vuln-Caaguazu-vulnerabilidades-e-mitigauxe7uxe3o}

\begin{itemize}
\tightlist
\item
  \textbf{Pressão de Massa:} Por estar no eixo PY02, é o destino
  imediato de refugiados urbanos (Mitigação: residência rural afastada
  10-15 km do eixo asfáltico, com autonomia total).
\item
  \textbf{Dependência de Redes:} Alta dependência de sistemas
  centralizados de energia e logística (Mitigação: implementação de
  backup solar e captação de água de chuva).
\item
  \textbf{Visibilidade Logística:} Alvo de bloqueios rodoviários em
  crises nacionais (Mitigação: mapeamento de rotas rurais vicinais em
  direção a San Pedro ou Guairá).
\end{itemize}
\subsubsection{Dossiê de Campo}\label{dossie-Caaguazu-dossiuxea-de-campo}
\begin{description}[style=nextline,leftmargin=0.5cm,font=\bfseries]
\item[Coordenadas]  25°27'S 56°01'W.
\item[Topografia]  Relevo predominantemente plano com colinas suaves. Altitude \textasciitilde 200m. Área de \textasciitilde 977 km². Geologicamente estável, com histórico de sismos leves intraplaca sem danos estruturais.
\item[Alvos Estratégicos]  Nó Logístico Vital da Ruta PY02 Duplicada (Principal eixo comercial Assunção-CDE); Grandes indústrias de beneficiamento de madeira e moinhos; Subestação regional da ANDE\footnote{Administración Nacional de Electricidade (Administração Nacional de Eletricidade), companhia estatal de energia.}.
\item[Fallout]  Ventos predominantes do Norte (quentes e úmidos) e Leste. Baixo risco de fallout direto.
\item[População]  \textasciitilde 98.200 habitantes (Censo 2022 Final). Segunda cidade mais importante do departamento.
\item[Densidade]  \textasciitilde 100,5 hab/\allowbreak km² (Densidade urbana elevada, indicando vulnerabilidade social e risco de desabastecimento rápido em crises urbanas agudas).
\item[Vias de Acesso]  Localização crítica sobre a Ruta PY02 duplicada. Ponto de intersecção com a Ruta PY13 (Conexão para Canindeyú). Conectividade asfáltica de elite, mas com altíssimo risco de saturação por fluxos migratórios em cenários de ruptura metropolitana.
\item[Serviços]  Hospital Distrital de Caaguazú e forte rede de clínicas privadas. Centro financeiro e comercial robusto com silos e terminais de carga.
\item[Custo de Vida]  Médio-baixo; economia impulsionada pela indústria e comércio transfronteiriço indireto.
\item[Preço da Terra]  US\$ 4.000-7.000/\allowbreak ha (áreas rurais agrícolas); lotes industriais e residenciais urbanos com alta valorização (US\$ 25.000-50.000/\allowbreak lote).
\item[Hidrologia]  Diversos arroios cruzam o distrito (ex: Tapiracuái). Baixo risco de inundações catastróficas sistêmicas, mas vulnerabilidade localizada em bairros baixos durante tempestades severas (média pluvial 1.600 mm/\allowbreak ano).
\item[Clima]  Subtropical Úmido. Verões intensos (>38°C) e invernos com geadas ocasionais.
\item[Energia]  Rede nacional (Itaipu). Ponto estratégico na rede de transmissão nacional leste-oeste.
\item[Água]  Captação via poços artesianos profundos e sistemas de distribuição urbana (ESSAP\footnote{Empresa de Servicios Sanitarios del Paraguay (Empresa de Serviços Sanitários do Paraguai), responsável pelo saneamento e água urbana.}/\allowbreak Juntas).
\item[Qualidade do Solo]  Latossolo Vermelho (Terra Vermelha); muito fértil, profundo e mecanizável. Excelente aptidão para soja, milho e reflorestamento industrial.
\item[Resources Locais]  Grande polo madeireiro, produção de gado de corte e soja. Elevado potencial para autossuficiência alimentar regional, embora a cidade dependa de logística interna para distribuição.
\item[Segurança]  Cidade com dinamismo comercial intenso. Criminalidade urbana comum presente (furtos/\allowbreak roubos). Estabilidade social sustentada pelo setor industrial e agropecuário. Patrulhamento policial reforçado por ser um corredor logístico internacional.
\item[Leis Local: Município consolidado e polo regional ("Capital da Madeira"). Livre de Restrição de Fronteira]  Localizado no coração do país, permitindo plena titularidade para investidores brasileiros.
\end{description}
\textbf{Fonte climática:} NASA POWER Climatology API (período 2001-2020)
\textbf{Fonte luminosa:} estimativa world atlas

\paragraph{Irradiação Solar
(kWh/\allowbreak m²/\allowbreak dia)}\label{irradiauxe7uxe3o-solar-kwhmuxb2dia}

\begin{longtable}[]{@{}
  >{\raggedright\arraybackslash}p{(\linewidth - 24\tabcolsep) * \real{0.0746}}
  >{\raggedright\arraybackslash}p{(\linewidth - 24\tabcolsep) * \real{0.0746}}
  >{\raggedright\arraybackslash}p{(\linewidth - 24\tabcolsep) * \real{0.0746}}
  >{\raggedright\arraybackslash}p{(\linewidth - 24\tabcolsep) * \real{0.0746}}
  >{\raggedright\arraybackslash}p{(\linewidth - 24\tabcolsep) * \real{0.0746}}
  >{\raggedright\arraybackslash}p{(\linewidth - 24\tabcolsep) * \real{0.0746}}
  >{\raggedright\arraybackslash}p{(\linewidth - 24\tabcolsep) * \real{0.0746}}
  >{\raggedright\arraybackslash}p{(\linewidth - 24\tabcolsep) * \real{0.0746}}
  >{\raggedright\arraybackslash}p{(\linewidth - 24\tabcolsep) * \real{0.0746}}
  >{\raggedright\arraybackslash}p{(\linewidth - 24\tabcolsep) * \real{0.0746}}
  >{\raggedright\arraybackslash}p{(\linewidth - 24\tabcolsep) * \real{0.0746}}
  >{\raggedright\arraybackslash}p{(\linewidth - 24\tabcolsep) * \real{0.0746}}
  >{\raggedright\arraybackslash}p{(\linewidth - 24\tabcolsep) * \real{0.1045}}@{}}
\toprule\noalign{}
\begin{minipage}[b]{\linewidth}\raggedright
Jan
\end{minipage} & \begin{minipage}[b]{\linewidth}\raggedright
Fev
\end{minipage} & \begin{minipage}[b]{\linewidth}\raggedright
Mar
\end{minipage} & \begin{minipage}[b]{\linewidth}\raggedright
Abr
\end{minipage} & \begin{minipage}[b]{\linewidth}\raggedright
Mai
\end{minipage} & \begin{minipage}[b]{\linewidth}\raggedright
Jun
\end{minipage} & \begin{minipage}[b]{\linewidth}\raggedright
Jul
\end{minipage} & \begin{minipage}[b]{\linewidth}\raggedright
Ago
\end{minipage} & \begin{minipage}[b]{\linewidth}\raggedright
Set
\end{minipage} & \begin{minipage}[b]{\linewidth}\raggedright
Out
\end{minipage} & \begin{minipage}[b]{\linewidth}\raggedright
Nov
\end{minipage} & \begin{minipage}[b]{\linewidth}\raggedright
Dez
\end{minipage} & \begin{minipage}[b]{\linewidth}\raggedright
Média
\end{minipage} \\
\midrule\noalign{}
\endhead
\bottomrule\noalign{}
\endlastfoot
6.74 & 6.20 & 5.50 & 4.48 & 3.35 & 2.88 & 3.21 & 3.96 & 4.59 & 5.43 &
6.43 & 6.81 & \textbf{4.96} \\
\end{longtable}

\textbf{Inclinação solar recomendada:} 25° N (anual) · 35° N (inverno
jun-ago) · 15° N (verão nov-jan)

\paragraph{Precipitação (mm/\allowbreak mês)}\label{precipitauxe7uxe3o-mmmuxeas}

\begin{longtable}[]{@{}
  >{\raggedright\arraybackslash}p{(\linewidth - 24\tabcolsep) * \real{0.0704}}
  >{\raggedright\arraybackslash}p{(\linewidth - 24\tabcolsep) * \real{0.0704}}
  >{\raggedright\arraybackslash}p{(\linewidth - 24\tabcolsep) * \real{0.0704}}
  >{\raggedright\arraybackslash}p{(\linewidth - 24\tabcolsep) * \real{0.0704}}
  >{\raggedright\arraybackslash}p{(\linewidth - 24\tabcolsep) * \real{0.0704}}
  >{\raggedright\arraybackslash}p{(\linewidth - 24\tabcolsep) * \real{0.0704}}
  >{\raggedright\arraybackslash}p{(\linewidth - 24\tabcolsep) * \real{0.0704}}
  >{\raggedright\arraybackslash}p{(\linewidth - 24\tabcolsep) * \real{0.0704}}
  >{\raggedright\arraybackslash}p{(\linewidth - 24\tabcolsep) * \real{0.0704}}
  >{\raggedright\arraybackslash}p{(\linewidth - 24\tabcolsep) * \real{0.0704}}
  >{\raggedright\arraybackslash}p{(\linewidth - 24\tabcolsep) * \real{0.0704}}
  >{\raggedright\arraybackslash}p{(\linewidth - 24\tabcolsep) * \real{0.0704}}
  >{\raggedright\arraybackslash}p{(\linewidth - 24\tabcolsep) * \real{0.1549}}@{}}
\toprule\noalign{}
\begin{minipage}[b]{\linewidth}\raggedright
Jan
\end{minipage} & \begin{minipage}[b]{\linewidth}\raggedright
Fev
\end{minipage} & \begin{minipage}[b]{\linewidth}\raggedright
Mar
\end{minipage} & \begin{minipage}[b]{\linewidth}\raggedright
Abr
\end{minipage} & \begin{minipage}[b]{\linewidth}\raggedright
Mai
\end{minipage} & \begin{minipage}[b]{\linewidth}\raggedright
Jun
\end{minipage} & \begin{minipage}[b]{\linewidth}\raggedright
Jul
\end{minipage} & \begin{minipage}[b]{\linewidth}\raggedright
Ago
\end{minipage} & \begin{minipage}[b]{\linewidth}\raggedright
Set
\end{minipage} & \begin{minipage}[b]{\linewidth}\raggedright
Out
\end{minipage} & \begin{minipage}[b]{\linewidth}\raggedright
Nov
\end{minipage} & \begin{minipage}[b]{\linewidth}\raggedright
Dez
\end{minipage} & \begin{minipage}[b]{\linewidth}\raggedright
Total/\allowbreak ano
\end{minipage} \\
\midrule\noalign{}
\endhead
\bottomrule\noalign{}
\endlastfoot
131 & 136 & 131 & 152 & 172 & 90 & 75 & 52 & 92 & 187 & 190 & 166 &
\textbf{1577 mm} \\
\end{longtable}

\paragraph{Poluição Luminosa}\label{poluiuxe7uxe3o-luminosa}

\begin{longtable}[]{@{}ll@{}}
\toprule\noalign{}
Parâmetro & Valor \\
\midrule\noalign{}
\endhead
\bottomrule\noalign{}
\endlastfoot
Escala Bortle & 4 --- Céu rural-suburbano \\
Radiância artificial & 3.0 nW/\allowbreak cm²/\allowbreak sr \\
\end{longtable}
\paragraph{Solo (SoilGrids 2.0, média ponderada 0--30
cm)}

\begin{longtable}[]{@{}ll@{}}
\toprule\noalign{}
Parâmetro & Valor \\
\midrule\noalign{}
\endhead
\bottomrule\noalign{}
\endlastfoot
pH (H$_2$O) & 5.22 \\
Carbono orgânico (SOC) & 18.82 g/\allowbreak kg \\
Argila & 27.33 \% \\
Areia & 52.23 \% \\
Silte (calc.) & 20.4 \% \\
Densidade aparente & 1.27 g/\allowbreak cm³ \\
\textbf{Aptidão agrícola} & \textbf{Média} \\
\end{longtable}

Fonte: ISRIC SoilGrids 2.0 via WCS. Coords: -25.45°, -56.0167°. Média
ponderada camadas 0-5, 5-15, 15-30 cm.
\subsubsection{Indicadores Sociais}

\textbf{Fontes:} PNUD 2020, INE\footnote{Instituto Nacional de Estadística (Instituto Nacional de Estatística), órgão oficial de dados demográficos e censitários.} Censo 2022, INE\footnote{Instituto Nacional de Estadística (Instituto Nacional de Estatística), órgão oficial de dados demográficos e censitários.} EPHC 2023, Ministerio
Público 2024\\
\textbf{Nota:} Valores marcados como \emph{(dept.)} referem-se ao
departamento de Caaguazú; valores \emph{(dist.)} são específicos deste
distrito. Dados marcados \emph{(est.)} são estimativas.

\begin{longtable}[]{@{}
  >{\raggedright\arraybackslash}p{(\linewidth - 4\tabcolsep) * \real{0.4231}}
  >{\raggedright\arraybackslash}p{(\linewidth - 4\tabcolsep) * \real{0.2692}}
  >{\raggedright\arraybackslash}p{(\linewidth - 4\tabcolsep) * \real{0.3077}}@{}}
\toprule\noalign{}
\begin{minipage}[b]{\linewidth}\raggedright
Indicador
\end{minipage} & \begin{minipage}[b]{\linewidth}\raggedright
Valor
\end{minipage} & \begin{minipage}[b]{\linewidth}\raggedright
Âmbito
\end{minipage} \\
\midrule\noalign{}
\endhead
\bottomrule\noalign{}
\endlastfoot
IDH (2020) & 0,715 & dept. \\
Ranking IDH nacional & 7/\allowbreak 18 & dept. \\
Esperança de vida & 73,1 anos & dept. \\
Escolaridade média & 8.9 anos & dist. \\
RNB per capita (USD PPA) & 9.000 & dept. \\
Pobreza monetária (\%) & 28,5\% & dept. \\
Pobreza extrema (\%) & 6,8\% & dept. \\
Índice de Gini & 0,447 & dept. \\
Acesso a água potável (\%) & 79,8\% & dept. \\
Acesso a saneamento (\%) & 78,3\% & dept. \\
Taxa de homicídios (est., /\allowbreak 100k hab) & \textasciitilde6--8 (estimativa,
próxima ou acima da média nacional) & dept. \\
Índice de segurança & médio & dept. \\
População (Censo 2022) & 982 hab. & dist. \\
Idade mediana & 29.0 anos & dist. \\
Taxa de fecundidade & N/\allowbreak D & dist. \\
\end{longtable}
\subsubsection{Combustível}

\textbf{Referência:} PETROPAR /\allowbreak  postos locais (2024) \textbf{Tipo de
localidade:} Interior

\begin{longtable}[]{@{}lll@{}}
\toprule\noalign{}
Combustível & USD/\allowbreak litro & Gs/\allowbreak litro (aprox.) \\
\midrule\noalign{}
\endhead
\bottomrule\noalign{}
\endlastfoot
Gasolina 93 oct & 0.98 & 7,252 \\
Gasolina 97 oct (premium) & 1.08 & 7,992 \\
Diesel & 0.91 & 6,734 \\
\end{longtable}

\begin{quote}
Preços podem variar ±5\% conforme posto e sazonalidade. Chaco e interior
remoto apresentam maior variação.
\end{quote}

\subsubsection{Cobertura Celular}

\textbf{Fonte:} CONATEL PY /\allowbreak  operadoras (2024)

\begin{longtable}[]{@{}ll@{}}
\toprule\noalign{}
Parâmetro & Valor \\
\midrule\noalign{}
\endhead
\bottomrule\noalign{}
\endlastfoot
Cobertura 4G população (dept.) & 88\% \\
Cobertura 4G área rural & 72\% \\
Melhor operadora & Tigo \\
Qualidade rural & boa \\
\end{longtable}

\begin{quote}
Para áreas rurais fora do núcleo urbano, recomenda-se chip Tigo como
principal e Personal como backup.
\end{quote}

\subsubsection{Internet}

\textbf{Fonte:} CONATEL /\allowbreak  Speedtest Ookla (2024)

\begin{longtable}[]{@{}ll@{}}
\toprule\noalign{}
Parâmetro & Valor \\
\midrule\noalign{}
\endhead
\bottomrule\noalign{}
\endlastfoot
Velocidade média download & 55 Mbps \\
Domicílios com internet (dept.) & 62\% \\
Tecnologia predominante & rádio/\allowbreak fibra \\
Opção rural & Starlink disponível (\textasciitilde USD 44/\allowbreak mês) \\
\end{longtable}

\subsubsection{Mercado Imobiliário e Terra
Rural}

\textbf{Fonte:} INDERT /\allowbreak  Clasificados.com.py (2024)

\begin{longtable}[]{@{}ll@{}}
\toprule\noalign{}
Tipo & Referência \\
\midrule\noalign{}
\endhead
\bottomrule\noalign{}
\endlastfoot
Terra agrícola alta prod. (USD/\allowbreak ha) & 8,800 \\
Imóvel urbano (USD/\allowbreak m²) & 850 \\
Aluguel 2 quartos (USD/\allowbreak mês) & 350 \\
\end{longtable}

\begin{quote}
Valores de referência departamental. Localidades menores podem ter
preços 20--40\% abaixo da capital departamental.
\end{quote}

\subsubsection{Saúde}

\textbf{Fonte:} MSPBS /\allowbreak  IPS Paraguay (2024)

\begin{longtable}[]{@{}ll@{}}
\toprule\noalign{}
Serviço & Disponibilidade \\
\midrule\noalign{}
\endhead
\bottomrule\noalign{}
\endlastfoot
USF /\allowbreak  Posto de Saúde & sim \\
Hospital Regional & sim \\
IPS (seguro social) & sim \\
Farmácia & sim \\
Distância ao hospital de referência & local \\
\end{longtable}

\textbf{Principais estabelecimentos:} Hospital Distrital de Caaguazú;
IPS Caaguazú; rede de clínicas privadas e farmácias.

\textbf{Observação para imigrantes:} Boa cobertura de saúde para capital
departamental. IPS acessível mediante vínculo formal. Referência
regional para casos complexos: Coronel Oviedo (\textasciitilde40 km).
Cobertura privada recomendada para especialidades.
\newpage
\secaoDiagnostico{Carayao}{\mapaConteudo{-25.1666}{-56.4000}}{Carayaó é um dos maiores refúgios demográficos do departamento de
Caaguazú. Sua maior força reside na vasta extensão territorial com
baixíssima densidade populacional, combinada com uma segurança social
superior e solo apto para policultivos e pecuária. É o local ideal para
quem busca o isolamento tático massivo sem estar geograficamente
desconectado de eixos rodoviários importantes (PY08).

\subsubsection{Por que sim}

\begin{itemize}
\tightlist
\item
  Densidade populacional entre as menores da Região Oriental
  (privacidade total).
\item
  Ambiente social extremamente seguro e coeso.
\item
  Solo de boa qualidade e abundância de águas superficiais para gado.
\item
  Sem restrições da Lei de Fronteira.
\end{itemize}

\subsubsection{Por que não}

\begin{itemize}
\tightlist
\item
  Infraestrutura de saúde local limitada a atendimentos primários.
\item
  Dependência de estradas vicinais que podem sofrer em chuvas severas.
\item
  Pequena escala comercial exige deslocamento para Villarrica ou Coronel
  Oviedo para suprimentos diversos.
\end{itemize}

\subsubsection{Recomendação
Técnica}

Altamente indicado para estabelecimentos de autossuficiência de médio
porte ou estâncias pecuárias que busquem isolamento demográfico. Requer
autonomia em termos de logística básica e saúde. O GSS de 6.7 reflete a
segurança do isolamento em contraste com a limitação institucional
local. Referencia atual de score: GSS 6.7 em 2026-03-05.

\subsubsection{}

\begin{itemize}
\tightlist
\item
  \begin{enumerate}
  \def\labelenumi{\arabic{enumi})}
  \tightlist
  \item
    Dados censitários validados (INE\footnote{Instituto Nacional de Estadística (Instituto Nacional de Estatística), órgão oficial de dados demográficos e censitários.} 2022).
  \end{enumerate}
\item
  \begin{enumerate}
  \def\labelenumi{\arabic{enumi})}
  \setcounter{enumi}{1}
  \tightlist
  \item
    Levantamento territorial oficial do distrito (Portal
    Geoestadístico).
  \end{enumerate}
\item
  \begin{enumerate}
  \def\labelenumi{\arabic{enumi})}
  \setcounter{enumi}{2}
  \tightlist
  \item
    Mapa de solos e hidrografia de Caaguazú (MAG).
  \end{enumerate}
\item
  \begin{enumerate}
  \def\labelenumi{\arabic{enumi})}
  \setcounter{enumi}{3}
  \tightlist
  \item
    Registro de segurança pública departamental.
  \end{enumerate}
\end{itemize}}
\begin{table}[ht!]
\small \begin{tabular}{p{3.0cm}p{0.8cm}p{6.0cm}} \toprule
\textbf{Critério} & \textbf{Nota} & \textbf{Justificativa} \\ \midrule
A - Ameaças Estratégicas & 6.5 & Localização discreta e protegida de rotas principais; alta invisibilidade estratégica \\
B - Risco Social & 7.5 & Baixíssima densidade demográfica; risco social urbano desprezível \\
C - Riscos Naturais & 7.0 & Geologia muito estável e riscos climáticos típicos manejáveis \\
D - Autossuficiência & 7.0 & Bons recursos naturais - gado e água - mas limitado pela logística de insumos técnicos \\
E - Institucional & 5.5 & Ambiente institucional seguro, mas com infraestrutura pública local limitada \\
\midrule \textbf{GSS FINAL} & \textbf{6.7} & \textbf{Moderadamente Seguro (Ideal para quem busca o máximo de isolamento demográfico em uma região segura e produtiva).} \\ \bottomrule \end{tabular}
\end{table}
\subsubsection{Vulnerabilidades e
Mitigação}\label{vuln-Carayao-vulnerabilidades-e-mitigauxe7uxe3o}

\begin{itemize}
\tightlist
\item
  \textbf{Isolamento de Saúde:} Distância de centros cirúrgicos
  (Mitigação: manutenção de estoque médico básico e kit trauma local).
\item
  \textbf{Dependência Logística:} Necessidade de deslocamento para
  Coronel Oviedo para suprimentos técnicos (Mitigação: estoque de
  suprimentos para 60 dias).
\item
  \textbf{Instabilidade Elétrica Rural:} Quedas ocasionais em temporais
  (Mitigação: instalação de sistemas solares fotovoltaicos redundantes).
\end{itemize}
\subsubsection{Dossiê de Campo}\label{dossie-Carayao-dossiuxea-de-campo}
\begin{description}[style=nextline,leftmargin=0.5cm,font=\bfseries]
\item[Coordenadas]  25°10'S 56°24'W.
\item[Topografia]  Relevo predominantemente plano a suavemente ondulado. Localizado no centro-oeste do departamento. Altitude \textasciitilde 160m. Área vasta de \textasciitilde 920 km². Geologicamente muito estável, risco sísmico desprezível.
\item[Alvos Estratégicos]  Áreas de produção pecuária extensiva; Ponto de passagem secundário para o norte de Caaguazú. Ausência de alvos militares ou industriais primários, oferecendo excelente invisibilidade tática e discrição estratégica.
\item[Fallout]  Ventos predominantes do Norte (quentes e úmidos) e Leste. Baixo risco de fallout direto de alvos continentais primários.
\item[População]  10.832 habitantes (Censo 2022 Final). Perfil essencialmente rural disperso em colônias históricas.
\item[Densidade]  \textasciitilde 11,8 hab/\allowbreak km² (Densidade baixíssima, oferecendo isolamento social superior e privacidade estratégica).
\item[Vias de Acesso]  Acesso facilitado pela Ruta PY08 (Dr. Blas Garay). Dependência de estradas vicinais de terra para o interior do distrito, que podem apresentar dificuldades de trafegabilidade em períodos de chuva extrema (média 1.600 mm/\allowbreak ano).
\item[Serviços]  Unidade de Saúde da Família (USF) local. Dependência crítica de Coronel Oviedo (35 km) para serviços hospitalares de alta complexidade e logística avançada.
\item[Custo de Vida]  Muito baixo; economia baseada na pecuária de corte e agricultura de subsistência.
\item[Preço da Terra]  US\$ 2.000-4.000/\allowbreak ha (áreas de pecuária ou mata); US\$ 5.500/\allowbreak ha (áreas com aptidão agrícola mecanizada).
\item[Hidrologia]  Baixo risco de inundações sistêmicas catastróficas. Presença de diversos arroios locais (ex: Arroyo Tobatiry) que favorecem a pecuária. Danos pontuais em acessos rurais durante tempestades de verão.
\item[Clima]  Subtropical Úmido. Verões quentes; invernos com entradas bruscas de ar frio do Sul.
\item[Energia]  Rede nacional (Itaipu); instável em áreas rurais remotas.
\item[Água]  Abundante via poços artesianos e lençol freático superficial produtivo.
\item[Qualidade do Solo]  Solo franco-arenoso a argiloso; fértil e profundo, ideal para pastagens, mandioca, milho e pequena produção de tabaco.
\item[Resources Locais]  Grande produção de gado de corte, leite e grãos para subsistência. Elevado potencial para autossuficiência alimentar básica em nível de propriedade.
\item[Segurança]  Zona rural extremamente pacífica e estável. Baixíssima criminalidade urbana. Coesão social baseada em comunidades rurais tradicionais e laboriosas. Baixa visibilidade para conflitos de massa.
\item[Leis Local: Município histórico (fundado em 1770). Livre de Restrição de Fronteira]  Localizado no interior do país, permitindo plena titularidade para investidores estrangeiros.
\end{description}
\textbf{Fonte climática:} NASA POWER Climatology API (período 2001-2020)
\textbf{Fonte luminosa:} estimativa world atlas

\paragraph{Irradiação Solar
(kWh/\allowbreak m²/\allowbreak dia)}\label{irradiauxe7uxe3o-solar-kwhmuxb2dia}

\begin{longtable}[]{@{}
  >{\raggedright\arraybackslash}p{(\linewidth - 24\tabcolsep) * \real{0.0746}}
  >{\raggedright\arraybackslash}p{(\linewidth - 24\tabcolsep) * \real{0.0746}}
  >{\raggedright\arraybackslash}p{(\linewidth - 24\tabcolsep) * \real{0.0746}}
  >{\raggedright\arraybackslash}p{(\linewidth - 24\tabcolsep) * \real{0.0746}}
  >{\raggedright\arraybackslash}p{(\linewidth - 24\tabcolsep) * \real{0.0746}}
  >{\raggedright\arraybackslash}p{(\linewidth - 24\tabcolsep) * \real{0.0746}}
  >{\raggedright\arraybackslash}p{(\linewidth - 24\tabcolsep) * \real{0.0746}}
  >{\raggedright\arraybackslash}p{(\linewidth - 24\tabcolsep) * \real{0.0746}}
  >{\raggedright\arraybackslash}p{(\linewidth - 24\tabcolsep) * \real{0.0746}}
  >{\raggedright\arraybackslash}p{(\linewidth - 24\tabcolsep) * \real{0.0746}}
  >{\raggedright\arraybackslash}p{(\linewidth - 24\tabcolsep) * \real{0.0746}}
  >{\raggedright\arraybackslash}p{(\linewidth - 24\tabcolsep) * \real{0.0746}}
  >{\raggedright\arraybackslash}p{(\linewidth - 24\tabcolsep) * \real{0.1045}}@{}}
\toprule\noalign{}
\begin{minipage}[b]{\linewidth}\raggedright
Jan
\end{minipage} & \begin{minipage}[b]{\linewidth}\raggedright
Fev
\end{minipage} & \begin{minipage}[b]{\linewidth}\raggedright
Mar
\end{minipage} & \begin{minipage}[b]{\linewidth}\raggedright
Abr
\end{minipage} & \begin{minipage}[b]{\linewidth}\raggedright
Mai
\end{minipage} & \begin{minipage}[b]{\linewidth}\raggedright
Jun
\end{minipage} & \begin{minipage}[b]{\linewidth}\raggedright
Jul
\end{minipage} & \begin{minipage}[b]{\linewidth}\raggedright
Ago
\end{minipage} & \begin{minipage}[b]{\linewidth}\raggedright
Set
\end{minipage} & \begin{minipage}[b]{\linewidth}\raggedright
Out
\end{minipage} & \begin{minipage}[b]{\linewidth}\raggedright
Nov
\end{minipage} & \begin{minipage}[b]{\linewidth}\raggedright
Dez
\end{minipage} & \begin{minipage}[b]{\linewidth}\raggedright
Média
\end{minipage} \\
\midrule\noalign{}
\endhead
\bottomrule\noalign{}
\endlastfoot
6.74 & 6.20 & 5.50 & 4.48 & 3.35 & 2.88 & 3.21 & 3.96 & 4.59 & 5.43 &
6.43 & 6.81 & \textbf{4.96} \\
\end{longtable}

\textbf{Inclinação solar recomendada:} 25° N (anual) · 35° N (inverno
jun-ago) · 15° N (verão nov-jan)

\paragraph{Precipitação (mm/\allowbreak mês)}\label{precipitauxe7uxe3o-mmmuxeas}

\begin{longtable}[]{@{}
  >{\raggedright\arraybackslash}p{(\linewidth - 24\tabcolsep) * \real{0.0704}}
  >{\raggedright\arraybackslash}p{(\linewidth - 24\tabcolsep) * \real{0.0704}}
  >{\raggedright\arraybackslash}p{(\linewidth - 24\tabcolsep) * \real{0.0704}}
  >{\raggedright\arraybackslash}p{(\linewidth - 24\tabcolsep) * \real{0.0704}}
  >{\raggedright\arraybackslash}p{(\linewidth - 24\tabcolsep) * \real{0.0704}}
  >{\raggedright\arraybackslash}p{(\linewidth - 24\tabcolsep) * \real{0.0704}}
  >{\raggedright\arraybackslash}p{(\linewidth - 24\tabcolsep) * \real{0.0704}}
  >{\raggedright\arraybackslash}p{(\linewidth - 24\tabcolsep) * \real{0.0704}}
  >{\raggedright\arraybackslash}p{(\linewidth - 24\tabcolsep) * \real{0.0704}}
  >{\raggedright\arraybackslash}p{(\linewidth - 24\tabcolsep) * \real{0.0704}}
  >{\raggedright\arraybackslash}p{(\linewidth - 24\tabcolsep) * \real{0.0704}}
  >{\raggedright\arraybackslash}p{(\linewidth - 24\tabcolsep) * \real{0.0704}}
  >{\raggedright\arraybackslash}p{(\linewidth - 24\tabcolsep) * \real{0.1549}}@{}}
\toprule\noalign{}
\begin{minipage}[b]{\linewidth}\raggedright
Jan
\end{minipage} & \begin{minipage}[b]{\linewidth}\raggedright
Fev
\end{minipage} & \begin{minipage}[b]{\linewidth}\raggedright
Mar
\end{minipage} & \begin{minipage}[b]{\linewidth}\raggedright
Abr
\end{minipage} & \begin{minipage}[b]{\linewidth}\raggedright
Mai
\end{minipage} & \begin{minipage}[b]{\linewidth}\raggedright
Jun
\end{minipage} & \begin{minipage}[b]{\linewidth}\raggedright
Jul
\end{minipage} & \begin{minipage}[b]{\linewidth}\raggedright
Ago
\end{minipage} & \begin{minipage}[b]{\linewidth}\raggedright
Set
\end{minipage} & \begin{minipage}[b]{\linewidth}\raggedright
Out
\end{minipage} & \begin{minipage}[b]{\linewidth}\raggedright
Nov
\end{minipage} & \begin{minipage}[b]{\linewidth}\raggedright
Dez
\end{minipage} & \begin{minipage}[b]{\linewidth}\raggedright
Total/\allowbreak ano
\end{minipage} \\
\midrule\noalign{}
\endhead
\bottomrule\noalign{}
\endlastfoot
124 & 143 & 125 & 147 & 173 & 90 & 70 & 46 & 94 & 184 & 190 & 168 &
\textbf{1557 mm} \\
\end{longtable}

\paragraph{Poluição Luminosa}\label{poluiuxe7uxe3o-luminosa}

\begin{longtable}[]{@{}ll@{}}
\toprule\noalign{}
Parâmetro & Valor \\
\midrule\noalign{}
\endhead
\bottomrule\noalign{}
\endlastfoot
Escala Bortle & 4 --- Céu rural-suburbano \\
Radiância artificial & 3.0 nW/\allowbreak cm²/\allowbreak sr \\
\end{longtable}
\paragraph{Solo (SoilGrids 2.0, média ponderada 0--30
cm)}

\begin{longtable}[]{@{}ll@{}}
\toprule\noalign{}
Parâmetro & Valor \\
\midrule\noalign{}
\endhead
\bottomrule\noalign{}
\endlastfoot
pH (H$_2$O) & 5.3 \\
Carbono orgânico (SOC) & 20.0 g/\allowbreak kg \\
Argila & 29.8 \% \\
Areia & 49.66 \% \\
Silte (calc.) & 20.5 \% \\
Densidade aparente & 1.27 g/\allowbreak cm³ \\
\textbf{Aptidão agrícola} & \textbf{Média} \\
\end{longtable}

Fonte: ISRIC SoilGrids 2.0 via WCS. Coords: -25.1667°, -56.4°. Média
ponderada camadas 0-5, 5-15, 15-30 cm.
\subsubsection{Indicadores Sociais}

\textbf{Fontes:} PNUD 2020, INE\footnote{Instituto Nacional de Estadística (Instituto Nacional de Estatística), órgão oficial de dados demográficos e censitários.} Censo 2022, INE\footnote{Instituto Nacional de Estadística (Instituto Nacional de Estatística), órgão oficial de dados demográficos e censitários.} EPHC 2023, Ministerio
Público 2024\\
\textbf{Nota:} Valores marcados como \emph{(dept.)} referem-se ao
departamento de Caaguazú; valores \emph{(dist.)} são específicos deste
distrito. Dados marcados \emph{(est.)} são estimativas.

\begin{longtable}[]{@{}
  >{\raggedright\arraybackslash}p{(\linewidth - 4\tabcolsep) * \real{0.4231}}
  >{\raggedright\arraybackslash}p{(\linewidth - 4\tabcolsep) * \real{0.2692}}
  >{\raggedright\arraybackslash}p{(\linewidth - 4\tabcolsep) * \real{0.3077}}@{}}
\toprule\noalign{}
\begin{minipage}[b]{\linewidth}\raggedright
Indicador
\end{minipage} & \begin{minipage}[b]{\linewidth}\raggedright
Valor
\end{minipage} & \begin{minipage}[b]{\linewidth}\raggedright
Âmbito
\end{minipage} \\
\midrule\noalign{}
\endhead
\bottomrule\noalign{}
\endlastfoot
IDH (2020) & 0,715 & dept. \\
Ranking IDH nacional & 7/\allowbreak 18 & dept. \\
Esperança de vida & 73,1 anos & dept. \\
Escolaridade média & 7.9 anos & dist. \\
RNB per capita (USD PPA) & 9.000 & dept. \\
Pobreza monetária (\%) & 28,5\% & dept. \\
Pobreza extrema (\%) & 6,8\% & dept. \\
Índice de Gini & 0,447 & dept. \\
Acesso a água potável (\%) & 79,8\% & dept. \\
Acesso a saneamento (\%) & 78,3\% & dept. \\
Taxa de homicídios (est., /\allowbreak 100k hab) & \textasciitilde6--8 (estimativa,
próxima ou acima da média nacional) & dept. \\
Índice de segurança & médio & dept. \\
População (Censo 2022) & 10.832 hab. & dist. \\
Idade mediana & 31.0 anos & dist. \\
Taxa de fecundidade & N/\allowbreak D & dist. \\
\end{longtable}
\subsubsection{Combustível}

\textbf{Referência:} PETROPAR /\allowbreak  postos locais (2024) \textbf{Tipo de
localidade:} Interior

\begin{longtable}[]{@{}lll@{}}
\toprule\noalign{}
Combustível & USD/\allowbreak litro & Gs/\allowbreak litro (aprox.) \\
\midrule\noalign{}
\endhead
\bottomrule\noalign{}
\endlastfoot
Gasolina 93 oct & 0.98 & 7,252 \\
Gasolina 97 oct (premium) & 1.08 & 7,992 \\
Diesel & 0.91 & 6,734 \\
\end{longtable}

\begin{quote}
Preços podem variar ±5\% conforme posto e sazonalidade. Chaco e interior
remoto apresentam maior variação.
\end{quote}

\subsubsection{Cobertura Celular}

\textbf{Fonte:} CONATEL PY /\allowbreak  operadoras (2024)

\begin{longtable}[]{@{}ll@{}}
\toprule\noalign{}
Parâmetro & Valor \\
\midrule\noalign{}
\endhead
\bottomrule\noalign{}
\endlastfoot
Cobertura 4G população (dept.) & 88\% \\
Cobertura 4G área rural & 72\% \\
Melhor operadora & Tigo \\
Qualidade rural & boa \\
\end{longtable}

\begin{quote}
Para áreas rurais fora do núcleo urbano, recomenda-se chip Tigo como
principal e Personal como backup.
\end{quote}

\subsubsection{Internet}

\textbf{Fonte:} CONATEL /\allowbreak  Speedtest Ookla (2024)

\begin{longtable}[]{@{}ll@{}}
\toprule\noalign{}
Parâmetro & Valor \\
\midrule\noalign{}
\endhead
\bottomrule\noalign{}
\endlastfoot
Velocidade média download & 55 Mbps \\
Domicílios com internet (dept.) & 62\% \\
Tecnologia predominante & rádio/\allowbreak fibra \\
Opção rural & Starlink disponível (\textasciitilde USD 44/\allowbreak mês) \\
\end{longtable}

\subsubsection{Mercado Imobiliário e Terra
Rural}

\textbf{Fonte:} INDERT /\allowbreak  Clasificados.com.py (2024)

\begin{longtable}[]{@{}ll@{}}
\toprule\noalign{}
Tipo & Referência \\
\midrule\noalign{}
\endhead
\bottomrule\noalign{}
\endlastfoot
Terra agrícola alta prod. (USD/\allowbreak ha) & 8,800 \\
Imóvel urbano (USD/\allowbreak m²) & 850 \\
Aluguel 2 quartos (USD/\allowbreak mês) & 350 \\
\end{longtable}

\begin{quote}
Valores de referência departamental. Localidades menores podem ter
preços 20--40\% abaixo da capital departamental.
\end{quote}

\subsubsection{Saúde}

\textbf{Fonte:} MSPBS /\allowbreak  IPS Paraguay (2024)

\begin{longtable}[]{@{}ll@{}}
\toprule\noalign{}
Serviço & Disponibilidade \\
\midrule\noalign{}
\endhead
\bottomrule\noalign{}
\endlastfoot
USF /\allowbreak  Posto de Saúde & sim \\
Hospital Regional & não \\
IPS (seguro social) & não \\
Farmácia & sim \\
Distância ao hospital de referência & \textasciitilde50 km (Coronel
Oviedo) \\
\end{longtable}

\textbf{Principais estabelecimentos:} USF Carayaó; farmácias locais.

\textbf{Observação para imigrantes:} Atendimento primário disponível via
USF. Para internações ou emergências, referência em Coronel Oviedo.
Cobertura privada recomendada; IPS não disponível localmente.
\newpage
\secaoDiagnostico{Coronel Oviedo}{\mapaConteudo{-25.4333}{-56.4333}}{Coronel Oviedo é o centro de gravidade logística do Paraguai. Oferece a
infraestrutura de saúde e serviços mais robusta de todo o interior do
país, garantindo um suporte institucional de elite. Contudo, seu
posicionamento geográfico como o principal cruzamento rodoviário
nacional a torna um alvo tático evidente em qualquer cenário de
instabilidade. Sua viabilidade como refúgio depende de um posicionamento
rural periférico que aproveite os recursos da cidade sem sofrer os
riscos da sua densidade urbana.

\subsubsection{Por que sim}

\begin{itemize}
\tightlist
\item
  Melhor infraestrutura de saúde regional (Novo Hospital Geral).
\item
  Centro logístico absoluto para aquisição de suprimentos e
  equipamentos.
\item
  Geologia estável e solos de altíssima produtividade.
\item
  Sem restrições da Lei de Fronteira.
\end{itemize}

\subsubsection{Por que não}

\begin{itemize}
\tightlist
\item
  Alvo estratégico logístico de altíssima visibilidade nacional.
\item
  Maior risco de agitação social e bloqueios rodoviários em crises.
\item
  Concentração populacional urbana elevada.
\end{itemize}

\subsubsection{Recomendação
Técnica}

Altamente indicado como base logística e de serviços de suporte para o
projeto. Para moradia permanente, recomenda-se a zona rural do distrito
ou municípios vizinhos menos densos (como Nueva Londres ou Carayaó),
mantendo Oviedo como centro de suprimentos. O GSS de 6.3 reflete a força
institucional equilibrada pela alta visibilidade estratégica. Referencia
atual de score: GSS 6.3 em 2026-03-05.

\subsubsection{}

\begin{itemize}
\tightlist
\item
  \begin{enumerate}
  \def\labelenumi{\arabic{enumi})}
  \tightlist
  \item
    Dados censitários validados (INE\footnote{Instituto Nacional de Estadística (Instituto Nacional de Estatística), órgão oficial de dados demográficos e censitários.} 2022).
  \end{enumerate}
\item
  \begin{enumerate}
  \def\labelenumi{\arabic{enumi})}
  \setcounter{enumi}{1}
  \tightlist
  \item
    Relatórios de tráfego nacional e logística (MOPC 2026).
  \end{enumerate}
\item
  \begin{enumerate}
  \def\labelenumi{\arabic{enumi})}
  \setcounter{enumi}{2}
  \tightlist
  \item
    Mapa geológico e de riscos hídricos (SEN\footnote{Secretaría de Emergencia Nacional (Secretaria de Emergência Nacional), órgão governamental de resposta a desastres.}/\allowbreak DMH).
  \end{enumerate}
\item
  \begin{enumerate}
  \def\labelenumi{\arabic{enumi})}
  \setcounter{enumi}{3}
  \tightlist
  \item
    Registro de investimentos em saúde pública (BID/\allowbreak MOPC).
  \end{enumerate}
\end{itemize}}
\begin{table}[ht!]
\small \begin{tabular}{p{3.0cm}p{0.8cm}p{6.0cm}} \toprule
\textbf{Critério} & \textbf{Nota} & \textbf{Justificativa} \\ \midrule
A - Ameaças Estratégicas & 4.5 & Entroncamento logístico nevrálgico nacional; alvo estratégico prioritário de alta visibilidade \\
B - Risco Social & 5.5 & População regional elevada e densidade urbana significativa; risco social em crises \\
C - Riscos Naturais & 7.0 & Geologia muito estável e baixo risco de grandes desastres naturais \\
D - Autossuficiência & 7.5 & Riqueza em recursos: hub de suprimentos, água abundante e solos de elite \\
E - Institucional & 8.0 & Forte segurança institucional e melhor infraestrutura hospitalar do interior \\
\midrule \textbf{GSS FINAL} & \textbf{6.3} & \textbf{Moderadamente Seguro (Ideal como centro de operações administrativas e de saúde para o projeto).} \\ \bottomrule \end{tabular}
\end{table}
\subsubsection{Vulnerabilidades e
Mitigação}\label{vuln-Coronel_Oviedo-vulnerabilidades-e-mitigauxe7uxe3o}

\begin{itemize}
\tightlist
\item
  \textbf{Gargalo Logístico Nacional:} Risco de isolamento por bloqueios
  táticos das Rutas PY02 e PY08 (Mitigação: residência rural afastada
  15-20 km do centro, com rotas de terra alternativas mapeadas).
\item
  \textbf{Pressão Populacional Urbana:} Risco de instabilidade social em
  crises de suprimento (Mitigação: estabelecimento de base rural com
  autonomia alimentar e energética total).
\item
  \textbf{Dependência de Redes:} Alta dependência de sistemas urbanos
  (Mitigação: poço artesiano próprio e sistema solar off-grid).
\end{itemize}
\subsubsection{Dossiê de Campo}\label{dossie-Coronel_Oviedo-dossiuxea-de-campo}
\begin{description}[style=nextline,leftmargin=0.5cm,font=\bfseries]
\item[Coordenadas]  25°26'S 56°26'W.
\item[Topografia]  Relevo predominantemente plano com colinas suaves e vales abertos. Altitude \textasciitilde 170m. Área de \textasciitilde 878 km². Geologicamente muito estável, risco sísmico desprezível (zona intraplaca).
\item[Alvos Estratégicos]  Coração Logístico do Paraguai (Cruzamento vital das Rutas PY02 e PY08); Sede do Governo Departamental; Grande Subestação de Transmissão da ANDE\footnote{Administración Nacional de Electricidade (Administração Nacional de Eletricidade), companhia estatal de energia.}; Novo Hospital Geral de Coronel Oviedo (referência regional).
\item[Fallout]  Ventos predominantes do Norte (quentes e úmidos) e Leste. Risco de fallout indireto em caso de ataques severos a Assunção (\textasciitilde 130 km).
\item[População]  98.323 habitantes (Censo 2022 Final). Cidade mais populosa e centro administrativo de Caaguazú.
\item[Densidade]  \textasciitilde 111,9 hab/\allowbreak km² (Densidade urbana elevada, indicando alta vulnerabilidade social e pressão sobre serviços em crises de massa).
\item[Vias de Acesso]  Ponto nevrálgico absoluto. Conecta Assunção, Ciudad del Este, Concepción e o Sul do país. Infraestrutura asfáltica duplicada na PY02. Extremo risco de paralisia e controle militar em cenários de instabilidade nacional por ser o "funil" do tráfego paraguaio.
\item[Serviços]  Melhor infraestrutura hospitalar do interior do país (Novo Hospital Geral). Polo universitário, comercial e financeiro consolidado.
\item[Custo de Vida]  Médio; economia pujante baseada na prestação de serviços, transporte e agroindústria.
\item[Preço da Terra]  US\$ 5.000-10.000/\allowbreak ha (áreas rurais valorizadas); lotes urbanos e industriais com preços de mercado metropolitano.
\item[Hidrologia]  Diversos arroios na zona urbana (ex: Arroyo Pasito). Baixo risco de inundações catastróficas sistêmicas. Vulnerabilidade localizada em drenagem urbana pluvial durante tempestades de verão.
\item[Clima]  Subtropical Úmido. Verões intensos (>40°C frequentes).
\item[Irradiância solar global — ALLSKY\_SFC\_SW\_DWN (kWh/\allowbreak m²/\allowbreak dia)]  Jan 6.74 | Fev 6.20 | Mar 5.50 | Abr 4.48 | Mai 3.35 | Jun 2.88 | Jul 3.21 | Ago 3.96 | Set 4.59 | Out 5.43 | Nov 6.43 | Dez 6.81. Média anual: 4.96 kWh/\allowbreak m²/\allowbreak dia. Mínimo: Jun (2.88) | Máximo: Dez (6.81). Fonte: NASA POWER Climatology API, período 2001-2020.
\item[Precipitação — PRECTOTCORR (mm/\allowbreak mês)]  Jan 131 | Fev 136 | Mar 131 | Abr 152 | Mai 172 | Jun 90 | Jul 75 | Ago 52 | Set 92 | Out 187 | Nov 190 | Dez 166. Média anual: \textasciitilde 1.573 mm/\allowbreak ano. Estação chuvosa Out-Mai; seca Jun-Set (mínimo em Ago). Fonte: NASA POWER.
\item[Inclinação recomendada para placas solares]  25° voltado para o Norte (anual); 35° no inverno (Jun-Ago); 15° no verão (Nov-Jan). Base: latitude local -25.44°S.
\item[Energia]  Rede nacional (Itaipu). Nodo vital de distribuição elétrica nacional; alta prioridade de manutenção institucional.
\item[Água]  Sistema ESSAP\footnote{Empresa de Servicios Sanitarios del Paraguay (Empresa de Serviços Sanitários do Paraguai), responsável pelo saneamento e água urbana.} e abundância de poços artesianos profundos (Aquífero Guarani).
\item[Qualidade do Solo]  Latossolo Vermelho de alta fertilidade; excelente para grãos, cana-de-açúcar e silvicultura nas áreas periféricas.
\item[Resources Locais]  Grande hub de suprimentos regionais, silos de grãos e usinas de etanol. Elevada capacidade de autossuficiência regional, mas dependente de ordem pública para distribuição.
\item[Segurança]  Centro urbano dinâmico com criminalidade urbana típica (furtos/\allowbreak roubos). Segurança institucional forte através da Direção de Polícia Departamental. Sensível a bloqueios rodoviários e manifestações sociais devido ao papel estratégico do cruzamento rodoviário.
\item[Leis Local]  Sede de todos os poderes estatais regionais. \\ \textbf{Livre de Restrição} de Fronteira. \\ \textit{Aquisição de terra rural proibida para estrangeiros limítrofes.}
\end{description}
\textbf{Fonte climática:} NASA POWER Climatology API (período 2001-2020)
\textbf{Fonte luminosa:} estimativa world atlas

\paragraph{Irradiação Solar
(kWh/\allowbreak m²/\allowbreak dia)}\label{irradiauxe7uxe3o-solar-kwhmuxb2dia}

\begin{longtable}[]{@{}
  >{\raggedright\arraybackslash}p{(\linewidth - 24\tabcolsep) * \real{0.0746}}
  >{\raggedright\arraybackslash}p{(\linewidth - 24\tabcolsep) * \real{0.0746}}
  >{\raggedright\arraybackslash}p{(\linewidth - 24\tabcolsep) * \real{0.0746}}
  >{\raggedright\arraybackslash}p{(\linewidth - 24\tabcolsep) * \real{0.0746}}
  >{\raggedright\arraybackslash}p{(\linewidth - 24\tabcolsep) * \real{0.0746}}
  >{\raggedright\arraybackslash}p{(\linewidth - 24\tabcolsep) * \real{0.0746}}
  >{\raggedright\arraybackslash}p{(\linewidth - 24\tabcolsep) * \real{0.0746}}
  >{\raggedright\arraybackslash}p{(\linewidth - 24\tabcolsep) * \real{0.0746}}
  >{\raggedright\arraybackslash}p{(\linewidth - 24\tabcolsep) * \real{0.0746}}
  >{\raggedright\arraybackslash}p{(\linewidth - 24\tabcolsep) * \real{0.0746}}
  >{\raggedright\arraybackslash}p{(\linewidth - 24\tabcolsep) * \real{0.0746}}
  >{\raggedright\arraybackslash}p{(\linewidth - 24\tabcolsep) * \real{0.0746}}
  >{\raggedright\arraybackslash}p{(\linewidth - 24\tabcolsep) * \real{0.1045}}@{}}
\toprule\noalign{}
\begin{minipage}[b]{\linewidth}\raggedright
Jan
\end{minipage} & \begin{minipage}[b]{\linewidth}\raggedright
Fev
\end{minipage} & \begin{minipage}[b]{\linewidth}\raggedright
Mar
\end{minipage} & \begin{minipage}[b]{\linewidth}\raggedright
Abr
\end{minipage} & \begin{minipage}[b]{\linewidth}\raggedright
Mai
\end{minipage} & \begin{minipage}[b]{\linewidth}\raggedright
Jun
\end{minipage} & \begin{minipage}[b]{\linewidth}\raggedright
Jul
\end{minipage} & \begin{minipage}[b]{\linewidth}\raggedright
Ago
\end{minipage} & \begin{minipage}[b]{\linewidth}\raggedright
Set
\end{minipage} & \begin{minipage}[b]{\linewidth}\raggedright
Out
\end{minipage} & \begin{minipage}[b]{\linewidth}\raggedright
Nov
\end{minipage} & \begin{minipage}[b]{\linewidth}\raggedright
Dez
\end{minipage} & \begin{minipage}[b]{\linewidth}\raggedright
Média
\end{minipage} \\
\midrule\noalign{}
\endhead
\bottomrule\noalign{}
\endlastfoot
6.74 & 6.20 & 5.50 & 4.48 & 3.35 & 2.88 & 3.21 & 3.96 & 4.59 & 5.43 &
6.43 & 6.81 & \textbf{4.96} \\
\end{longtable}

\textbf{Inclinação solar recomendada:} 25° N (anual) · 35° N (inverno
jun-ago) · 15° N (verão nov-jan)

\paragraph{Precipitação (mm/\allowbreak mês)}\label{precipitauxe7uxe3o-mmmuxeas}

\begin{longtable}[]{@{}
  >{\raggedright\arraybackslash}p{(\linewidth - 24\tabcolsep) * \real{0.0704}}
  >{\raggedright\arraybackslash}p{(\linewidth - 24\tabcolsep) * \real{0.0704}}
  >{\raggedright\arraybackslash}p{(\linewidth - 24\tabcolsep) * \real{0.0704}}
  >{\raggedright\arraybackslash}p{(\linewidth - 24\tabcolsep) * \real{0.0704}}
  >{\raggedright\arraybackslash}p{(\linewidth - 24\tabcolsep) * \real{0.0704}}
  >{\raggedright\arraybackslash}p{(\linewidth - 24\tabcolsep) * \real{0.0704}}
  >{\raggedright\arraybackslash}p{(\linewidth - 24\tabcolsep) * \real{0.0704}}
  >{\raggedright\arraybackslash}p{(\linewidth - 24\tabcolsep) * \real{0.0704}}
  >{\raggedright\arraybackslash}p{(\linewidth - 24\tabcolsep) * \real{0.0704}}
  >{\raggedright\arraybackslash}p{(\linewidth - 24\tabcolsep) * \real{0.0704}}
  >{\raggedright\arraybackslash}p{(\linewidth - 24\tabcolsep) * \real{0.0704}}
  >{\raggedright\arraybackslash}p{(\linewidth - 24\tabcolsep) * \real{0.0704}}
  >{\raggedright\arraybackslash}p{(\linewidth - 24\tabcolsep) * \real{0.1549}}@{}}
\toprule\noalign{}
\begin{minipage}[b]{\linewidth}\raggedright
Jan
\end{minipage} & \begin{minipage}[b]{\linewidth}\raggedright
Fev
\end{minipage} & \begin{minipage}[b]{\linewidth}\raggedright
Mar
\end{minipage} & \begin{minipage}[b]{\linewidth}\raggedright
Abr
\end{minipage} & \begin{minipage}[b]{\linewidth}\raggedright
Mai
\end{minipage} & \begin{minipage}[b]{\linewidth}\raggedright
Jun
\end{minipage} & \begin{minipage}[b]{\linewidth}\raggedright
Jul
\end{minipage} & \begin{minipage}[b]{\linewidth}\raggedright
Ago
\end{minipage} & \begin{minipage}[b]{\linewidth}\raggedright
Set
\end{minipage} & \begin{minipage}[b]{\linewidth}\raggedright
Out
\end{minipage} & \begin{minipage}[b]{\linewidth}\raggedright
Nov
\end{minipage} & \begin{minipage}[b]{\linewidth}\raggedright
Dez
\end{minipage} & \begin{minipage}[b]{\linewidth}\raggedright
Total/\allowbreak ano
\end{minipage} \\
\midrule\noalign{}
\endhead
\bottomrule\noalign{}
\endlastfoot
131 & 136 & 131 & 152 & 172 & 90 & 75 & 52 & 92 & 187 & 190 & 166 &
\textbf{1577 mm} \\
\end{longtable}

\paragraph{Poluição Luminosa}\label{poluiuxe7uxe3o-luminosa}

\begin{longtable}[]{@{}ll@{}}
\toprule\noalign{}
Parâmetro & Valor \\
\midrule\noalign{}
\endhead
\bottomrule\noalign{}
\endlastfoot
Escala Bortle & 5 --- Céu suburbano \\
Radiância artificial & 6.0 nW/\allowbreak cm²/\allowbreak sr \\
\end{longtable}
\paragraph{Solo (SoilGrids 2.0, média ponderada 0--30
cm)}

\begin{longtable}[]{@{}ll@{}}
\toprule\noalign{}
Parâmetro & Valor \\
\midrule\noalign{}
\endhead
\bottomrule\noalign{}
\endlastfoot
pH (H$_2$O) & 5.2 \\
Carbono orgânico (SOC) & 18.77 g/\allowbreak kg \\
Argila & 24.38 \% \\
Areia & 55.24 \% \\
Silte (calc.) & 20.4 \% \\
Densidade aparente & 1.3 g/\allowbreak cm³ \\
\textbf{Aptidão agrícola} & \textbf{Média} \\
\end{longtable}

Fonte: ISRIC SoilGrids 2.0 via WCS. Coords: -25.4333°, -56.4333°. Média
ponderada camadas 0-5, 5-15, 15-30 cm.
\subsubsection{Indicadores Sociais}

\textbf{Fontes:} PNUD 2020, INE\footnote{Instituto Nacional de Estadística (Instituto Nacional de Estatística), órgão oficial de dados demográficos e censitários.} Censo 2022, INE\footnote{Instituto Nacional de Estadística (Instituto Nacional de Estatística), órgão oficial de dados demográficos e censitários.} EPHC 2023, Ministerio
Público 2024\\
\textbf{Nota:} Valores marcados como \emph{(dept.)} referem-se ao
departamento de Caaguazú; valores \emph{(dist.)} são específicos deste
distrito. Dados marcados \emph{(est.)} são estimativas.

\begin{longtable}[]{@{}
  >{\raggedright\arraybackslash}p{(\linewidth - 4\tabcolsep) * \real{0.4231}}
  >{\raggedright\arraybackslash}p{(\linewidth - 4\tabcolsep) * \real{0.2692}}
  >{\raggedright\arraybackslash}p{(\linewidth - 4\tabcolsep) * \real{0.3077}}@{}}
\toprule\noalign{}
\begin{minipage}[b]{\linewidth}\raggedright
Indicador
\end{minipage} & \begin{minipage}[b]{\linewidth}\raggedright
Valor
\end{minipage} & \begin{minipage}[b]{\linewidth}\raggedright
Âmbito
\end{minipage} \\
\midrule\noalign{}
\endhead
\bottomrule\noalign{}
\endlastfoot
IDH (2020) & 0,715 & dept. \\
Ranking IDH nacional & 7/\allowbreak 18 & dept. \\
Esperança de vida & 73,1 anos & dept. \\
Escolaridade média & 10.4 anos & dist. \\
RNB per capita (USD PPA) & 9.000 & dept. \\
Pobreza monetária (\%) & 28,5\% & dept. \\
Pobreza extrema (\%) & 6,8\% & dept. \\
Índice de Gini & 0,447 & dept. \\
Acesso a água potável (\%) & 79,8\% & dept. \\
Acesso a saneamento (\%) & 78,3\% & dept. \\
Taxa de homicídios (est., /\allowbreak 100k hab) & \textasciitilde6--8 (estimativa,
próxima ou acima da média nacional) & dept. \\
Índice de segurança & médio & dept. \\
População (Censo 2022) & 98.323 hab. & dist. \\
Idade mediana & 30.0 anos & dist. \\
Taxa de fecundidade & N/\allowbreak D & dist. \\
\end{longtable}
\subsubsection{Combustível}

\textbf{Referência:} PETROPAR /\allowbreak  postos locais (2024) \textbf{Tipo de
localidade:} Capital departamental

\begin{longtable}[]{@{}lll@{}}
\toprule\noalign{}
Combustível & USD/\allowbreak litro & Gs/\allowbreak litro (aprox.) \\
\midrule\noalign{}
\endhead
\bottomrule\noalign{}
\endlastfoot
Gasolina 93 oct & 0.98 & 7,252 \\
Gasolina 97 oct (premium) & 1.08 & 7,992 \\
Diesel & 0.91 & 6,734 \\
\end{longtable}

\begin{quote}
Preços podem variar ±5\% conforme posto e sazonalidade. Chaco e interior
remoto apresentam maior variação.
\end{quote}

\subsubsection{Cobertura Celular}

\textbf{Fonte:} CONATEL PY /\allowbreak  operadoras (2024)

\begin{longtable}[]{@{}ll@{}}
\toprule\noalign{}
Parâmetro & Valor \\
\midrule\noalign{}
\endhead
\bottomrule\noalign{}
\endlastfoot
Cobertura 4G população (dept.) & 88\% \\
Cobertura 4G área rural & 72\% \\
Melhor operadora & Tigo \\
Qualidade rural & boa \\
\end{longtable}

\begin{quote}
Para áreas rurais fora do núcleo urbano, recomenda-se chip Tigo como
principal e Personal como backup.
\end{quote}

\subsubsection{Internet}

\textbf{Fonte:} CONATEL /\allowbreak  Speedtest Ookla (2024)

\begin{longtable}[]{@{}ll@{}}
\toprule\noalign{}
Parâmetro & Valor \\
\midrule\noalign{}
\endhead
\bottomrule\noalign{}
\endlastfoot
Velocidade média download & 55 Mbps \\
Domicílios com internet (dept.) & 62\% \\
Tecnologia predominante & rádio/\allowbreak fibra \\
Opção rural & Starlink disponível (\textasciitilde USD 44/\allowbreak mês) \\
\end{longtable}

\subsubsection{Mercado Imobiliário e Terra
Rural}

\textbf{Fonte:} INDERT /\allowbreak  Clasificados.com.py (2024)

\begin{longtable}[]{@{}ll@{}}
\toprule\noalign{}
Tipo & Referência \\
\midrule\noalign{}
\endhead
\bottomrule\noalign{}
\endlastfoot
Terra agrícola alta prod. (USD/\allowbreak ha) & 8,800 \\
Imóvel urbano (USD/\allowbreak m²) & 977 \\
Aluguel 2 quartos (USD/\allowbreak mês) & 420 \\
\end{longtable}

\begin{quote}
Valores de referência departamental. Localidades menores podem ter
preços 20--40\% abaixo da capital departamental.
\end{quote}

\subsubsection{Saúde}

\textbf{Fonte:} MSPBS /\allowbreak  IPS Paraguay (2024)

\begin{longtable}[]{@{}ll@{}}
\toprule\noalign{}
Serviço & Disponibilidade \\
\midrule\noalign{}
\endhead
\bottomrule\noalign{}
\endlastfoot
USF /\allowbreak  Posto de Saúde & sim \\
Hospital Regional & sim \\
IPS (seguro social) & sim \\
Farmácia & sim \\
Distância ao hospital de referência & local \\
\end{longtable}

\textbf{Principais estabelecimentos:} Hospital Geral de Coronel Oviedo
(referência regional); IPS Coronel Oviedo; diversas clínicas e farmácias
privadas.

\textbf{Observação para imigrantes:} Melhor infraestrutura de saúde do
interior do país. IPS acessível mediante contrato de trabalho formal.
Cobertura privada complementar recomendada para especialidades de alta
complexidade (encaminhamento a Assunção).
\newpage
\secaoDiagnostico{Doctor Cecilio Baez}{\mapaConteudo{-25.0666}{-56.2500}}{Doctor Cecilio Baez (Caaguazu) apresenta perfil operacional consolidado
para comparacao territorial, com leitura conjunta de infraestrutura,
risco hidroclimatico e capacidade de continuidade de servicos
essenciais.

\subsubsection{Por que sim}

\begin{itemize}
\tightlist
\item
  Base institucional de referencia disponivel e rastreavel.
\item
  Estrutura comparativa (A-E/\allowbreak GSS) consistente para decisao relativa.
\item
  Plano de mitigacao acionavel ja definido no dossie.
\end{itemize}

\subsubsection{Por que não}

\begin{itemize}
\tightlist
\item
  Serie oficial distrital granular ainda pode apresentar lacunas
  temporais.
\item
  Sensibilidade do score exige monitoramento continuo de risco e
  conectividade.
\item
  Decisao final depende de validacao de campo e janela temporal recente.
\end{itemize}

\subsubsection{Pre-condicoes Minimas de
Decisao}

\begin{itemize}
\tightlist
\item
  Atualizar indicadores criticos em ciclo trimestral.
\item
  Confirmar trafegabilidade e contingencia hidrica/\allowbreak energetica local.
\item
  Validar checklist de resiliencia domiciliar/\allowbreak logistica antes de decisao
  final.
\end{itemize}

\subsubsection{Recomendação
Técnica}

Uso recomendado como candidato em analise multicriterio, condicionado ao
cumprimento das pre-condicoes acima. Referencia atual de score: GSS 7.3
em 2026-03-05; risco predominante: baixo.

\subsubsection{}

\begin{itemize}
\tightlist
\item
  \begin{enumerate}
  \def\labelenumi{\arabic{enumi})}
  \tightlist
  \item
    Delimitacao territorial validada por georreferencia institucional
    (Fonte: acesso: 2026-03-05).
  \end{enumerate}
\item
  \begin{enumerate}
  \def\labelenumi{\arabic{enumi})}
  \setcounter{enumi}{1}
  \tightlist
  \item
    População municipal/\allowbreak distrital referenciada por base censitaria
    nacional (Fonte: acesso: 2026-03-05).
  \end{enumerate}
\item
  \begin{enumerate}
  \def\labelenumi{\arabic{enumi})}
  \setcounter{enumi}{2}
  \tightlist
  \item
    Comparabilidade intra-departamental apoiada em indicadores
    distritais oficiais (Fonte: acesso: 2026-03-05).
  \end{enumerate}
\item
  \begin{enumerate}
  \def\labelenumi{\arabic{enumi})}
  \setcounter{enumi}{3}
  \tightlist
  \item
    Estado macro de conectividade viaria ancorado em informacoes de
    rodovias (Fonte: acesso: 2026-03-05).
  \end{enumerate}
\item
  \begin{enumerate}
  \def\labelenumi{\arabic{enumi})}
  \setcounter{enumi}{4}
  \tightlist
  \item
    Exposicao hidrometeorologica monitorada por avisos oficiais (Fonte:
    acesso: 2026-03-05).
  \end{enumerate}
\item
  \begin{enumerate}
  \def\labelenumi{\arabic{enumi})}
  \setcounter{enumi}{5}
  \tightlist
  \item
    Historico de contexto meteorologico apoiado em publicacoes tecnicas
    (Fonte: acesso: 2026-03-05).
  \end{enumerate}
\item
  \begin{enumerate}
  \def\labelenumi{\arabic{enumi})}
  \setcounter{enumi}{6}
  \tightlist
  \item
    Capacidade institucional de resposta a eventos apoiada em acoes da
    SEN\footnote{Secretaría de Emergencia Nacional (Secretaria de Emergência Nacional), órgão governamental de resposta a desastres.} (Fonte: acesso: 2026-03-05).
  \end{enumerate}
\item
  \begin{enumerate}
  \def\labelenumi{\arabic{enumi})}
  \setcounter{enumi}{7}
  \tightlist
  \item
    Infraestrutura energetica analisada via operador nacional (Fonte:
    acesso: 2026-03-05).
  \end{enumerate}
\item
  \begin{enumerate}
  \def\labelenumi{\arabic{enumi})}
  \setcounter{enumi}{8}
  \tightlist
  \item
    Referencias de obras e logistica nacional verificadas no MOPC
    (Fonte: acesso: 2026-03-05).
  \end{enumerate}
\item
  \begin{enumerate}
  \def\labelenumi{\arabic{enumi})}
  \setcounter{enumi}{9}
  \tightlist
  \item
    Contexto sociopolitico institucional referenciado por autoridade
    eleitoral (Fonte: acesso: 2026-03-05).
  \end{enumerate}
\item
  \begin{enumerate}
  \def\labelenumi{\arabic{enumi})}
  \setcounter{enumi}{10}
  \tightlist
  \item
    Camada cartografica de apoio eleitoral/\allowbreak territorial consultada
    (Fonte: acesso: 2026-03-05).
  \end{enumerate}
\item
  \begin{enumerate}
  \def\labelenumi{\arabic{enumi})}
  \setcounter{enumi}{11}
  \tightlist
  \item
    Dados publicos complementares para verificacao cruzada (Fonte:
    acesso: 2026-03-05).
  \end{enumerate}
\end{itemize}}
\begin{table}[ht!]
\small \begin{tabular}{p{3.0cm}p{0.8cm}p{6.0cm}} \toprule
\textbf{Critério} & \textbf{Nota} & \textbf{Justificativa} \\ \midrule
A - Ameaças Estratégicas & 8.2 & - \\
B - Risco Social & 6.6 & - \\
C - Riscos Naturais & 7.3 & - \\
D - Autossuficiência & 6.8 & - \\
E - Institucional & 7.2 & - \\
\midrule \textbf{GSS FINAL} & \textbf{7.3} & \textbf{Seguro (bom potencial com preparacao).} \\ \bottomrule \end{tabular}
\end{table}
\subsubsection{Vulnerabilidades e
Mitigação}\label{vuln-Doctor_Cecilio_Baez-vulnerabilidades-e-mitigauxe7uxe3o}

\begin{itemize}
\tightlist
\item
  Exposicao hidrometeorologica controlada (\textless45/\allowbreak 100).
\item
  Conectividade viaria favoravel (\textgreater=70/\allowbreak 100).
\item
  Estoque de contingencia para 14 dias (agua/\allowbreak alimentos/\allowbreak combustivel),
  priorizado quando risco hidro=34/\allowbreak 100.
\item
  Duplicar rotas/\allowbreak logistica quando conectividade viaria=74/\allowbreak 100 for
  \textless70; manter rota secundaria mapeada e testada trimestralmente.
\item
  Plano de energia resiliente com autonomia minima de 72h
  (geracao/\allowbreak backup), revisao semestral.
\end{itemize}
\subsubsection{Dossiê de Campo}\label{dossie-Doctor_Cecilio_Baez-dossiuxea-de-campo}
\begin{description}[style=nextline,leftmargin=0.5cm,font=\bfseries]
\item[Coordenadas]  25°04'S 56°15'W (geocodificado via Nominatim).
\end{description}
\textbf{Fonte climática:} NASA POWER Climatology API (período 2001-2020)
\textbf{Fonte luminosa:} estimativa world atlas

\paragraph{Irradiação Solar
(kWh/\allowbreak m²/\allowbreak dia)}\label{irradiauxe7uxe3o-solar-kwhmuxb2dia}

\begin{longtable}[]{@{}
  >{\raggedright\arraybackslash}p{(\linewidth - 24\tabcolsep) * \real{0.0746}}
  >{\raggedright\arraybackslash}p{(\linewidth - 24\tabcolsep) * \real{0.0746}}
  >{\raggedright\arraybackslash}p{(\linewidth - 24\tabcolsep) * \real{0.0746}}
  >{\raggedright\arraybackslash}p{(\linewidth - 24\tabcolsep) * \real{0.0746}}
  >{\raggedright\arraybackslash}p{(\linewidth - 24\tabcolsep) * \real{0.0746}}
  >{\raggedright\arraybackslash}p{(\linewidth - 24\tabcolsep) * \real{0.0746}}
  >{\raggedright\arraybackslash}p{(\linewidth - 24\tabcolsep) * \real{0.0746}}
  >{\raggedright\arraybackslash}p{(\linewidth - 24\tabcolsep) * \real{0.0746}}
  >{\raggedright\arraybackslash}p{(\linewidth - 24\tabcolsep) * \real{0.0746}}
  >{\raggedright\arraybackslash}p{(\linewidth - 24\tabcolsep) * \real{0.0746}}
  >{\raggedright\arraybackslash}p{(\linewidth - 24\tabcolsep) * \real{0.0746}}
  >{\raggedright\arraybackslash}p{(\linewidth - 24\tabcolsep) * \real{0.0746}}
  >{\raggedright\arraybackslash}p{(\linewidth - 24\tabcolsep) * \real{0.1045}}@{}}
\toprule\noalign{}
\begin{minipage}[b]{\linewidth}\raggedright
Jan
\end{minipage} & \begin{minipage}[b]{\linewidth}\raggedright
Fev
\end{minipage} & \begin{minipage}[b]{\linewidth}\raggedright
Mar
\end{minipage} & \begin{minipage}[b]{\linewidth}\raggedright
Abr
\end{minipage} & \begin{minipage}[b]{\linewidth}\raggedright
Mai
\end{minipage} & \begin{minipage}[b]{\linewidth}\raggedright
Jun
\end{minipage} & \begin{minipage}[b]{\linewidth}\raggedright
Jul
\end{minipage} & \begin{minipage}[b]{\linewidth}\raggedright
Ago
\end{minipage} & \begin{minipage}[b]{\linewidth}\raggedright
Set
\end{minipage} & \begin{minipage}[b]{\linewidth}\raggedright
Out
\end{minipage} & \begin{minipage}[b]{\linewidth}\raggedright
Nov
\end{minipage} & \begin{minipage}[b]{\linewidth}\raggedright
Dez
\end{minipage} & \begin{minipage}[b]{\linewidth}\raggedright
Média
\end{minipage} \\
\midrule\noalign{}
\endhead
\bottomrule\noalign{}
\endlastfoot
6.74 & 6.20 & 5.50 & 4.48 & 3.35 & 2.88 & 3.21 & 3.96 & 4.59 & 5.43 &
6.43 & 6.81 & \textbf{4.96} \\
\end{longtable}

\textbf{Inclinação solar recomendada:} 25° N (anual) · 35° N (inverno
jun-ago) · 15° N (verão nov-jan)

\paragraph{Precipitação (mm/\allowbreak mês)}\label{precipitauxe7uxe3o-mmmuxeas}

\begin{longtable}[]{@{}
  >{\raggedright\arraybackslash}p{(\linewidth - 24\tabcolsep) * \real{0.0704}}
  >{\raggedright\arraybackslash}p{(\linewidth - 24\tabcolsep) * \real{0.0704}}
  >{\raggedright\arraybackslash}p{(\linewidth - 24\tabcolsep) * \real{0.0704}}
  >{\raggedright\arraybackslash}p{(\linewidth - 24\tabcolsep) * \real{0.0704}}
  >{\raggedright\arraybackslash}p{(\linewidth - 24\tabcolsep) * \real{0.0704}}
  >{\raggedright\arraybackslash}p{(\linewidth - 24\tabcolsep) * \real{0.0704}}
  >{\raggedright\arraybackslash}p{(\linewidth - 24\tabcolsep) * \real{0.0704}}
  >{\raggedright\arraybackslash}p{(\linewidth - 24\tabcolsep) * \real{0.0704}}
  >{\raggedright\arraybackslash}p{(\linewidth - 24\tabcolsep) * \real{0.0704}}
  >{\raggedright\arraybackslash}p{(\linewidth - 24\tabcolsep) * \real{0.0704}}
  >{\raggedright\arraybackslash}p{(\linewidth - 24\tabcolsep) * \real{0.0704}}
  >{\raggedright\arraybackslash}p{(\linewidth - 24\tabcolsep) * \real{0.0704}}
  >{\raggedright\arraybackslash}p{(\linewidth - 24\tabcolsep) * \real{0.1549}}@{}}
\toprule\noalign{}
\begin{minipage}[b]{\linewidth}\raggedright
Jan
\end{minipage} & \begin{minipage}[b]{\linewidth}\raggedright
Fev
\end{minipage} & \begin{minipage}[b]{\linewidth}\raggedright
Mar
\end{minipage} & \begin{minipage}[b]{\linewidth}\raggedright
Abr
\end{minipage} & \begin{minipage}[b]{\linewidth}\raggedright
Mai
\end{minipage} & \begin{minipage}[b]{\linewidth}\raggedright
Jun
\end{minipage} & \begin{minipage}[b]{\linewidth}\raggedright
Jul
\end{minipage} & \begin{minipage}[b]{\linewidth}\raggedright
Ago
\end{minipage} & \begin{minipage}[b]{\linewidth}\raggedright
Set
\end{minipage} & \begin{minipage}[b]{\linewidth}\raggedright
Out
\end{minipage} & \begin{minipage}[b]{\linewidth}\raggedright
Nov
\end{minipage} & \begin{minipage}[b]{\linewidth}\raggedright
Dez
\end{minipage} & \begin{minipage}[b]{\linewidth}\raggedright
Total/\allowbreak ano
\end{minipage} \\
\midrule\noalign{}
\endhead
\bottomrule\noalign{}
\endlastfoot
124 & 143 & 125 & 147 & 173 & 90 & 70 & 46 & 94 & 184 & 190 & 168 &
\textbf{1557 mm} \\
\end{longtable}

\paragraph{Poluição Luminosa}\label{poluiuxe7uxe3o-luminosa}

\begin{longtable}[]{@{}ll@{}}
\toprule\noalign{}
Parâmetro & Valor \\
\midrule\noalign{}
\endhead
\bottomrule\noalign{}
\endlastfoot
Escala Bortle & 4 --- Céu rural-suburbano \\
Radiância artificial & 3.0 nW/\allowbreak cm²/\allowbreak sr \\
\end{longtable}
\paragraph{Solo (SoilGrids 2.0, média ponderada 0--30
cm)}

\begin{longtable}[]{@{}ll@{}}
\toprule\noalign{}
Parâmetro & Valor \\
\midrule\noalign{}
\endhead
\bottomrule\noalign{}
\endlastfoot
pH (H$_2$O) & 5.3 \\
Carbono orgânico (SOC) & 17.76 g/\allowbreak kg \\
Argila & 23.77 \% \\
Areia & 51.71 \% \\
Silte (calc.) & 24.5 \% \\
Densidade aparente & 1.32 g/\allowbreak cm³ \\
\textbf{Aptidão agrícola} & \textbf{Média} \\
\end{longtable}

Fonte: ISRIC SoilGrids 2.0 via WCS. Coords: -25.0667°, -56.25°. Média
ponderada camadas 0-5, 5-15, 15-30 cm.
\subsubsection{Indicadores Sociais}

\textbf{Fontes:} PNUD 2020, INE\footnote{Instituto Nacional de Estadística (Instituto Nacional de Estatística), órgão oficial de dados demográficos e censitários.} Censo 2022, INE\footnote{Instituto Nacional de Estadística (Instituto Nacional de Estatística), órgão oficial de dados demográficos e censitários.} EPHC 2023, Ministerio
Público 2024\\
\textbf{Nota:} Valores marcados como \emph{(dept.)} referem-se ao
departamento de Caaguazú; valores \emph{(dist.)} são específicos deste
distrito. Dados marcados \emph{(est.)} são estimativas.

\begin{longtable}[]{@{}
  >{\raggedright\arraybackslash}p{(\linewidth - 4\tabcolsep) * \real{0.4231}}
  >{\raggedright\arraybackslash}p{(\linewidth - 4\tabcolsep) * \real{0.2692}}
  >{\raggedright\arraybackslash}p{(\linewidth - 4\tabcolsep) * \real{0.3077}}@{}}
\toprule\noalign{}
\begin{minipage}[b]{\linewidth}\raggedright
Indicador
\end{minipage} & \begin{minipage}[b]{\linewidth}\raggedright
Valor
\end{minipage} & \begin{minipage}[b]{\linewidth}\raggedright
Âmbito
\end{minipage} \\
\midrule\noalign{}
\endhead
\bottomrule\noalign{}
\endlastfoot
IDH (2020) & 0,715 & dept. \\
Ranking IDH nacional & 7/\allowbreak 18 & dept. \\
Esperança de vida & 73,1 anos & dept. \\
Escolaridade média & 8,4 anos & dept. \\
RNB per capita (USD PPA) & 9.000 & dept. \\
Pobreza monetária (\%) & 28,5\% & dept. \\
Pobreza extrema (\%) & 6,8\% & dept. \\
Índice de Gini & 0,447 & dept. \\
Acesso a água potável (\%) & 79,8\% & dept. \\
Acesso a saneamento (\%) & 78,3\% & dept. \\
Taxa de homicídios (est., /\allowbreak 100k hab) & \textasciitilde6--8 (estimativa,
próxima ou acima da média nacional) & dept. \\
Índice de segurança & médio & dept. \\
População (Censo 2022) & 4.812 hab. & dist. \\
Idade mediana & 30.0 anos & dist. \\
Taxa de fecundidade & N/\allowbreak D & dist. \\
\end{longtable}
\subsubsection{Combustível}

\textbf{Referência:} PETROPAR /\allowbreak  postos locais (2024) \textbf{Tipo de
localidade:} Interior

\begin{longtable}[]{@{}lll@{}}
\toprule\noalign{}
Combustível & USD/\allowbreak litro & Gs/\allowbreak litro (aprox.) \\
\midrule\noalign{}
\endhead
\bottomrule\noalign{}
\endlastfoot
Gasolina 93 oct & 0.98 & 7,252 \\
Gasolina 97 oct (premium) & 1.08 & 7,992 \\
Diesel & 0.91 & 6,734 \\
\end{longtable}

\begin{quote}
Preços podem variar ±5\% conforme posto e sazonalidade. Chaco e interior
remoto apresentam maior variação.
\end{quote}

\subsubsection{Cobertura Celular}

\textbf{Fonte:} CONATEL PY /\allowbreak  operadoras (2024)

\begin{longtable}[]{@{}ll@{}}
\toprule\noalign{}
Parâmetro & Valor \\
\midrule\noalign{}
\endhead
\bottomrule\noalign{}
\endlastfoot
Cobertura 4G população (dept.) & 88\% \\
Cobertura 4G área rural & 72\% \\
Melhor operadora & Tigo \\
Qualidade rural & boa \\
\end{longtable}

\begin{quote}
Para áreas rurais fora do núcleo urbano, recomenda-se chip Tigo como
principal e Personal como backup.
\end{quote}

\subsubsection{Internet}

\textbf{Fonte:} CONATEL /\allowbreak  Speedtest Ookla (2024)

\begin{longtable}[]{@{}ll@{}}
\toprule\noalign{}
Parâmetro & Valor \\
\midrule\noalign{}
\endhead
\bottomrule\noalign{}
\endlastfoot
Velocidade média download & 55 Mbps \\
Domicílios com internet (dept.) & 62\% \\
Tecnologia predominante & rádio/\allowbreak fibra \\
Opção rural & Starlink disponível (\textasciitilde USD 44/\allowbreak mês) \\
\end{longtable}

\subsubsection{Mercado Imobiliário e Terra
Rural}

\textbf{Fonte:} INDERT /\allowbreak  Clasificados.com.py (2024)

\begin{longtable}[]{@{}ll@{}}
\toprule\noalign{}
Tipo & Referência \\
\midrule\noalign{}
\endhead
\bottomrule\noalign{}
\endlastfoot
Terra agrícola alta prod. (USD/\allowbreak ha) & 8,800 \\
Imóvel urbano (USD/\allowbreak m²) & 850 \\
Aluguel 2 quartos (USD/\allowbreak mês) & 350 \\
\end{longtable}

\begin{quote}
Valores de referência departamental. Localidades menores podem ter
preços 20--40\% abaixo da capital departamental.
\end{quote}

\subsubsection{Saúde}

\textbf{Fonte:} MSPBS /\allowbreak  IPS Paraguay (2024)

\begin{longtable}[]{@{}ll@{}}
\toprule\noalign{}
Serviço & Disponibilidade \\
\midrule\noalign{}
\endhead
\bottomrule\noalign{}
\endlastfoot
USF /\allowbreak  Posto de Saúde & sim \\
Hospital Regional & não \\
IPS (seguro social) & não \\
Farmácia & sim \\
Distância ao hospital de referência & \textasciitilde60 km (Coronel
Oviedo) \\
\end{longtable}

\textbf{Principais estabelecimentos:} USF local; farmácias distritais.

\textbf{Observação para imigrantes:} Atendimento primário disponível.
Casos de maior complexidade encaminhados a Coronel Oviedo. Cobertura
privada recomendada; IPS não disponível localmente.
\newpage
\secaoDiagnostico{Doctor Eulogio Estigarribia}{\mapaConteudo{-25.3500}{-55.7666}}{Doctor Eulogio Estigarribia (Campo 9) é o celeiro industrial do
Paraguai. Oferece a melhor combinação de solos férteis e capacidade
instalada de processamento de alimentos do país. Sua viabilidade como
refúgio é sustentada pela organização impecável das cooperativas locais
e pela abundância de recursos alimentares vitais (leite e trigo). A
principal cautela reside na sua importância estratégica nacional, que a
torna um ponto de controle obrigatório em qualquer cenário de
instabilidade sistêmica.

\subsubsection{Por que sim}

\begin{itemize}
\tightlist
\item
  Maior concentração de indústrias de laticínios e moinhos do país.
\item
  Solos de altíssima produtividade e clima excepcional.
\item
  Coesão social e segurança comunitária robusta (Cooperativas).
\item
  Sem restrições da Lei de Fronteira.
\end{itemize}

\subsubsection{Por que não}

\begin{itemize}
\tightlist
\item
  Alvo estratégico alimentar óbvio em crises nacionais.
\item
  Densidade urbana metropolitana em crescimento acelerado.
\item
  Dependência tática do eixo rodoviário PY02.
\end{itemize}

\subsubsection{Recomendação
Técnica}

Altamente indicado para investimentos agroindustriais e estabelecimentos
de autossuficiência de alto desempenho. Requer um plano de segurança que
priorize a autonomia em cenários de bloqueio rodoviário nacional. O GSS
de 6.9 reflete a excepcional riqueza de recursos e a solidez
institucional da localidade. Referencia atual de score: GSS 6.9 em
2026-03-05.

\subsubsection{}

\begin{itemize}
\tightlist
\item
  \begin{enumerate}
  \def\labelenumi{\arabic{enumi})}
  \tightlist
  \item
    Dados censitários validados (INE\footnote{Instituto Nacional de Estadística (Instituto Nacional de Estatística), órgão oficial de dados demográficos e censitários.} 2022).
  \end{enumerate}
\item
  \begin{enumerate}
  \def\labelenumi{\arabic{enumi})}
  \setcounter{enumi}{1}
  \tightlist
  \item
    Relatórios de produção de laticínios e grãos (Lactolanda/\allowbreak CAPAINLAC).
  \end{enumerate}
\item
  \begin{enumerate}
  \def\labelenumi{\arabic{enumi})}
  \setcounter{enumi}{2}
  \tightlist
  \item
    Mapa de solos e aptidão do Centro-Leste (MAG).
  \end{enumerate}
\item
  \begin{enumerate}
  \def\labelenumi{\arabic{enumi})}
  \setcounter{enumi}{3}
  \tightlist
  \item
    Registro de segurança pública e governança cooperativa.
  \end{enumerate}
\end{itemize}}
\begin{table}[ht!]
\small \begin{tabular}{p{3.0cm}p{0.8cm}p{6.0cm}} \toprule
\textbf{Critério} & \textbf{Nota} & \textbf{Justificativa} \\ \midrule
A - Ameaças Estratégicas & 5.0 & Hub agroindustrial vital; alvo estratégico de interesse nacional para controle de suprimentos \\
B - Risco Social & 6.5 & População significativa e densidade controlada; forte coesão social cooperativista \\
C - Riscos Naturais & 7.0 & Geologia muito estável e baixo risco de grandes desastres naturais \\
D - Autossuficiência & 9.0 & Recursos alimentares excepcionais: Lactolanda, moinhos e solos de elite; autossuficiência de nível superior \\
E - Institucional & 7.5 & Forte segurança institucional e social sustentada pelo setor privado cooperativo \\
\midrule \textbf{GSS FINAL} & \textbf{6.9} & \textbf{Moderadamente Seguro (Localidade de elite para suporte logístico e produção de recursos de sobrevivência).} \\ \bottomrule \end{tabular}
\end{table}
\subsubsection{Vulnerabilidades e
Mitigação}\label{vuln-Doctor_Eulogio_Estigarribia-vulnerabilidades-e-mitigauxe7uxe3o}

\begin{itemize}
\tightlist
\item
  \textbf{Alvo Estratégico Alimentar:} O controle de Campo 9 é vital
  para o abastecimento nacional, tornando-o um alvo em crises civis
  (Mitigação: residência rural afastada 10-15 km do eixo urbano
  industrial).
\item
  \textbf{Dependência de Rodovias:} Logística interna vinculada à PY02
  (Mitigação: manutenção de veículos de grande porte e rotas de terra
  mapeadas para o interior do departamento).
\item
  \textbf{Poluição Industrial:} Risco de contaminação pontual em arroios
  urbanos (Mitigação: uso de poços artesianos profundos com filtragem).
\end{itemize}
\subsubsection{Dossiê de Campo}\label{dossie-Doctor_Eulogio_Estigarribia-dossiuxea-de-campo}
\begin{description}[style=nextline,leftmargin=0.5cm,font=\bfseries]
\item[Coordenadas]  25°21'S 55°46'W.
\item[Topografia]  Relevo predominantemente plano a suavemente ondulado. Localizado em zona de solos de alta fertilidade. Altitude \textasciitilde 250m. Área de \textasciitilde 629 km². Geologicamente muito estável, risco sísmico desprezível.
\item[Alvos Estratégicos]  Polo de Segurança Alimentar Nacional (Maior concentração de moinhos de trigo e indústrias de laticínios do Paraguai); Sede da Cooperativa Lactolanda; Grandes terminais de silos; Eixo logístico da Ruta PY02. Infraestrutura crítica absoluta em cenários de escassez de alimentos.
\item[Fallout]  Ventos predominantes do Norte (quentes e úmidos) e Leste. Baixo risco de fallout direto.
\item[População]  38.894 habitantes (Censo 2022 Final). População composta por paraguaios, descendentes de colonos menonitas e imigrantes brasileiros, com forte cultura de cooperação produtiva.
\item[Densidade]  \textasciitilde 61,8 hab/\allowbreak km² (Densidade moderada, com núcleo urbano industrial compacto e colônias rurais altamente produtivas).
\item[Vias de Acesso]  Cortado pela Ruta PY02 duplicada. Conectividade asfáltica de primeiro nível para Assunção (213 km) e Ciudad del Este. Nó logístico vital para o agronegócio regional.
\item[Serviços]  Infraestrutura de serviços superior, com foco em logística pesada, bancos e comércio industrial. Unidades de saúde modernas apoiadas por cooperativas privadas.
\item[Custo de Vida]  Médio; economia dolarizada pelo agronegócio, mas com abundância de produtos alimentícios básicos a preços de origem.
\item[Preço da Terra]  US\$ 6.000-10.000/\allowbreak ha (áreas de elite mecanizadas); lotes industriais com alta valorização nacional.
\item[Hidrologia]  Diversos arroios locais integrados a sistemas de irrigação. Baixo risco de inundações sistêmicas catastróficas devido ao relevo favorável à drenagem.
\item[Clima]  Subtropical Úmido. Chuvas abundantes e bem distribuídas (1.600-1.800 mm/\allowbreak ano). Verões com alta incidência solar, favorecendo a produtividade agrícola.
\item[Energia]  Rede nacional (Itaipu); subestações industriais de alta carga. Muitas unidades produtivas possuem sistemas de geração térmica/\allowbreak biomassa de reserva.
\item[Água]  Abundante via poços artesianos profundos (Aquífero Guarani).
\item[Qualidade do Solo]  Latossolo Vermelho de Elite; solos profundos, estruturados e de altíssima fertilidade natural.
\item[Resources Locais]  Maior bacia leiteira e farinheira do Paraguai. Produção massiva de leite, queijo, farinha de trigo, massas, soja e milho. Elevadíssimo potencial para autossuficiência e suporte logístico de longa duração.
\item[Segurança]  Zona agroindustrial muito estável. Índices de criminalidade violenta baixos. Forte controle comunitário exercido pelas cooperativas menonitas (Sommerfeld e Bergthal). Perfil de segurança tática focado na proteção patrimonial industrial.
\item[Leis Local]  Município com forte influência econômica nacional. \\ \textbf{Livre de Restrição} de Fronteira. \\ \textit{Aquisição de terra rural proibida para estrangeiros limítrofes.}
\end{description}
\textbf{Fonte climática:} NASA POWER Climatology API (período 2001-2020)
\textbf{Fonte luminosa:} estimativa world atlas

\paragraph{Irradiação Solar
(kWh/\allowbreak m²/\allowbreak dia)}\label{irradiauxe7uxe3o-solar-kwhmuxb2dia}

\begin{longtable}[]{@{}
  >{\raggedright\arraybackslash}p{(\linewidth - 24\tabcolsep) * \real{0.0746}}
  >{\raggedright\arraybackslash}p{(\linewidth - 24\tabcolsep) * \real{0.0746}}
  >{\raggedright\arraybackslash}p{(\linewidth - 24\tabcolsep) * \real{0.0746}}
  >{\raggedright\arraybackslash}p{(\linewidth - 24\tabcolsep) * \real{0.0746}}
  >{\raggedright\arraybackslash}p{(\linewidth - 24\tabcolsep) * \real{0.0746}}
  >{\raggedright\arraybackslash}p{(\linewidth - 24\tabcolsep) * \real{0.0746}}
  >{\raggedright\arraybackslash}p{(\linewidth - 24\tabcolsep) * \real{0.0746}}
  >{\raggedright\arraybackslash}p{(\linewidth - 24\tabcolsep) * \real{0.0746}}
  >{\raggedright\arraybackslash}p{(\linewidth - 24\tabcolsep) * \real{0.0746}}
  >{\raggedright\arraybackslash}p{(\linewidth - 24\tabcolsep) * \real{0.0746}}
  >{\raggedright\arraybackslash}p{(\linewidth - 24\tabcolsep) * \real{0.0746}}
  >{\raggedright\arraybackslash}p{(\linewidth - 24\tabcolsep) * \real{0.0746}}
  >{\raggedright\arraybackslash}p{(\linewidth - 24\tabcolsep) * \real{0.1045}}@{}}
\toprule\noalign{}
\begin{minipage}[b]{\linewidth}\raggedright
Jan
\end{minipage} & \begin{minipage}[b]{\linewidth}\raggedright
Fev
\end{minipage} & \begin{minipage}[b]{\linewidth}\raggedright
Mar
\end{minipage} & \begin{minipage}[b]{\linewidth}\raggedright
Abr
\end{minipage} & \begin{minipage}[b]{\linewidth}\raggedright
Mai
\end{minipage} & \begin{minipage}[b]{\linewidth}\raggedright
Jun
\end{minipage} & \begin{minipage}[b]{\linewidth}\raggedright
Jul
\end{minipage} & \begin{minipage}[b]{\linewidth}\raggedright
Ago
\end{minipage} & \begin{minipage}[b]{\linewidth}\raggedright
Set
\end{minipage} & \begin{minipage}[b]{\linewidth}\raggedright
Out
\end{minipage} & \begin{minipage}[b]{\linewidth}\raggedright
Nov
\end{minipage} & \begin{minipage}[b]{\linewidth}\raggedright
Dez
\end{minipage} & \begin{minipage}[b]{\linewidth}\raggedright
Média
\end{minipage} \\
\midrule\noalign{}
\endhead
\bottomrule\noalign{}
\endlastfoot
6.58 & 6.08 & 5.45 & 4.47 & 3.30 & 2.90 & 3.20 & 3.96 & 4.53 & 5.31 &
6.35 & 6.66 & \textbf{4.9} \\
\end{longtable}

\textbf{Inclinação solar recomendada:} 25° N (anual) · 35° N (inverno
jun-ago) · 15° N (verão nov-jan)

\paragraph{Precipitação (mm/\allowbreak mês)}\label{precipitauxe7uxe3o-mmmuxeas}

\begin{longtable}[]{@{}
  >{\raggedright\arraybackslash}p{(\linewidth - 24\tabcolsep) * \real{0.0704}}
  >{\raggedright\arraybackslash}p{(\linewidth - 24\tabcolsep) * \real{0.0704}}
  >{\raggedright\arraybackslash}p{(\linewidth - 24\tabcolsep) * \real{0.0704}}
  >{\raggedright\arraybackslash}p{(\linewidth - 24\tabcolsep) * \real{0.0704}}
  >{\raggedright\arraybackslash}p{(\linewidth - 24\tabcolsep) * \real{0.0704}}
  >{\raggedright\arraybackslash}p{(\linewidth - 24\tabcolsep) * \real{0.0704}}
  >{\raggedright\arraybackslash}p{(\linewidth - 24\tabcolsep) * \real{0.0704}}
  >{\raggedright\arraybackslash}p{(\linewidth - 24\tabcolsep) * \real{0.0704}}
  >{\raggedright\arraybackslash}p{(\linewidth - 24\tabcolsep) * \real{0.0704}}
  >{\raggedright\arraybackslash}p{(\linewidth - 24\tabcolsep) * \real{0.0704}}
  >{\raggedright\arraybackslash}p{(\linewidth - 24\tabcolsep) * \real{0.0704}}
  >{\raggedright\arraybackslash}p{(\linewidth - 24\tabcolsep) * \real{0.0704}}
  >{\raggedright\arraybackslash}p{(\linewidth - 24\tabcolsep) * \real{0.1549}}@{}}
\toprule\noalign{}
\begin{minipage}[b]{\linewidth}\raggedright
Jan
\end{minipage} & \begin{minipage}[b]{\linewidth}\raggedright
Fev
\end{minipage} & \begin{minipage}[b]{\linewidth}\raggedright
Mar
\end{minipage} & \begin{minipage}[b]{\linewidth}\raggedright
Abr
\end{minipage} & \begin{minipage}[b]{\linewidth}\raggedright
Mai
\end{minipage} & \begin{minipage}[b]{\linewidth}\raggedright
Jun
\end{minipage} & \begin{minipage}[b]{\linewidth}\raggedright
Jul
\end{minipage} & \begin{minipage}[b]{\linewidth}\raggedright
Ago
\end{minipage} & \begin{minipage}[b]{\linewidth}\raggedright
Set
\end{minipage} & \begin{minipage}[b]{\linewidth}\raggedright
Out
\end{minipage} & \begin{minipage}[b]{\linewidth}\raggedright
Nov
\end{minipage} & \begin{minipage}[b]{\linewidth}\raggedright
Dez
\end{minipage} & \begin{minipage}[b]{\linewidth}\raggedright
Total/\allowbreak ano
\end{minipage} \\
\midrule\noalign{}
\endhead
\bottomrule\noalign{}
\endlastfoot
142 & 130 & 126 & 152 & 177 & 94 & 79 & 55 & 103 & 189 & 188 & 172 &
\textbf{1607 mm} \\
\end{longtable}

\paragraph{Poluição Luminosa}\label{poluiuxe7uxe3o-luminosa}

\begin{longtable}[]{@{}ll@{}}
\toprule\noalign{}
Parâmetro & Valor \\
\midrule\noalign{}
\endhead
\bottomrule\noalign{}
\endlastfoot
Escala Bortle & 4 --- Céu rural-suburbano \\
Radiância artificial & 3.0 nW/\allowbreak cm²/\allowbreak sr \\
\end{longtable}
\paragraph{Solo (SoilGrids 2.0, média ponderada 0--30
cm)}

\begin{longtable}[]{@{}ll@{}}
\toprule\noalign{}
Parâmetro & Valor \\
\midrule\noalign{}
\endhead
\bottomrule\noalign{}
\endlastfoot
pH (H$_2$O) & 5.27 \\
Carbono orgânico (SOC) & 18.89 g/\allowbreak kg \\
Argila & 25.64 \% \\
Areia & 53.24 \% \\
Silte (calc.) & 21.1 \% \\
Densidade aparente & 1.31 g/\allowbreak cm³ \\
\textbf{Aptidão agrícola} & \textbf{Média} \\
\end{longtable}

Fonte: ISRIC SoilGrids 2.0 via WCS. Coords: -25.35°, -55.7667°. Média
ponderada camadas 0-5, 5-15, 15-30 cm.
\subsubsection{Indicadores Sociais}

\textbf{Fontes:} PNUD 2020, INE\footnote{Instituto Nacional de Estadística (Instituto Nacional de Estatística), órgão oficial de dados demográficos e censitários.} Censo 2022, INE\footnote{Instituto Nacional de Estadística (Instituto Nacional de Estatística), órgão oficial de dados demográficos e censitários.} EPHC 2023, Ministerio
Público 2024\\
\textbf{Nota:} Valores marcados como \emph{(dept.)} referem-se ao
departamento de Caaguazú; valores \emph{(dist.)} são específicos deste
distrito. Dados marcados \emph{(est.)} são estimativas.

\begin{longtable}[]{@{}
  >{\raggedright\arraybackslash}p{(\linewidth - 4\tabcolsep) * \real{0.4231}}
  >{\raggedright\arraybackslash}p{(\linewidth - 4\tabcolsep) * \real{0.2692}}
  >{\raggedright\arraybackslash}p{(\linewidth - 4\tabcolsep) * \real{0.3077}}@{}}
\toprule\noalign{}
\begin{minipage}[b]{\linewidth}\raggedright
Indicador
\end{minipage} & \begin{minipage}[b]{\linewidth}\raggedright
Valor
\end{minipage} & \begin{minipage}[b]{\linewidth}\raggedright
Âmbito
\end{minipage} \\
\midrule\noalign{}
\endhead
\bottomrule\noalign{}
\endlastfoot
IDH (2020) & 0,715 & dept. \\
Ranking IDH nacional & 7/\allowbreak 18 & dept. \\
Esperança de vida & 73,1 anos & dept. \\
Escolaridade média & 8,4 anos & dept. \\
RNB per capita (USD PPA) & 9.000 & dept. \\
Pobreza monetária (\%) & 28,5\% & dept. \\
Pobreza extrema (\%) & 6,8\% & dept. \\
Índice de Gini & 0,447 & dept. \\
Acesso a água potável (\%) & 79,8\% & dept. \\
Acesso a saneamento (\%) & 78,3\% & dept. \\
Taxa de homicídios (est., /\allowbreak 100k hab) & \textasciitilde6--8 (estimativa,
próxima ou acima da média nacional) & dept. \\
Índice de segurança & médio & dept. \\
População (Censo 2022) & 38.894 hab. & dist. \\
Idade mediana & 24.0 anos & dist. \\
Taxa de fecundidade & N/\allowbreak D & dist. \\
\end{longtable}
\subsubsection{Combustível}

\textbf{Referência:} PETROPAR /\allowbreak  postos locais (2024) \textbf{Tipo de
localidade:} Interior

\begin{longtable}[]{@{}lll@{}}
\toprule\noalign{}
Combustível & USD/\allowbreak litro & Gs/\allowbreak litro (aprox.) \\
\midrule\noalign{}
\endhead
\bottomrule\noalign{}
\endlastfoot
Gasolina 93 oct & 0.98 & 7,252 \\
Gasolina 97 oct (premium) & 1.08 & 7,992 \\
Diesel & 0.91 & 6,734 \\
\end{longtable}

\begin{quote}
Preços podem variar ±5\% conforme posto e sazonalidade. Chaco e interior
remoto apresentam maior variação.
\end{quote}

\subsubsection{Cobertura Celular}

\textbf{Fonte:} CONATEL PY /\allowbreak  operadoras (2024)

\begin{longtable}[]{@{}ll@{}}
\toprule\noalign{}
Parâmetro & Valor \\
\midrule\noalign{}
\endhead
\bottomrule\noalign{}
\endlastfoot
Cobertura 4G população (dept.) & 88\% \\
Cobertura 4G área rural & 72\% \\
Melhor operadora & Tigo \\
Qualidade rural & boa \\
\end{longtable}

\begin{quote}
Para áreas rurais fora do núcleo urbano, recomenda-se chip Tigo como
principal e Personal como backup.
\end{quote}

\subsubsection{Internet}

\textbf{Fonte:} CONATEL /\allowbreak  Speedtest Ookla (2024)

\begin{longtable}[]{@{}ll@{}}
\toprule\noalign{}
Parâmetro & Valor \\
\midrule\noalign{}
\endhead
\bottomrule\noalign{}
\endlastfoot
Velocidade média download & 55 Mbps \\
Domicílios com internet (dept.) & 62\% \\
Tecnologia predominante & rádio/\allowbreak fibra \\
Opção rural & Starlink disponível (\textasciitilde USD 44/\allowbreak mês) \\
\end{longtable}

\subsubsection{Mercado Imobiliário e Terra
Rural}

\textbf{Fonte:} INDERT /\allowbreak  Clasificados.com.py (2024)

\begin{longtable}[]{@{}ll@{}}
\toprule\noalign{}
Tipo & Referência \\
\midrule\noalign{}
\endhead
\bottomrule\noalign{}
\endlastfoot
Terra agrícola alta prod. (USD/\allowbreak ha) & 8,800 \\
Imóvel urbano (USD/\allowbreak m²) & 850 \\
Aluguel 2 quartos (USD/\allowbreak mês) & 350 \\
\end{longtable}

\begin{quote}
Valores de referência departamental. Localidades menores podem ter
preços 20--40\% abaixo da capital departamental.
\end{quote}

\subsubsection{Saúde}

\textbf{Fonte:} MSPBS /\allowbreak  IPS Paraguay (2024)

\begin{longtable}[]{@{}ll@{}}
\toprule\noalign{}
Serviço & Disponibilidade \\
\midrule\noalign{}
\endhead
\bottomrule\noalign{}
\endlastfoot
USF /\allowbreak  Posto de Saúde & sim \\
Hospital Regional & não \\
IPS (seguro social) & não \\
Farmácia & sim \\
Distância ao hospital de referência & \textasciitilde80 km (Coronel
Oviedo) \\
\end{longtable}

\textbf{Principais estabelecimentos:} USF local; farmácias distritais.

\textbf{Observação para imigrantes:} Atendimento primário disponível.
Casos de maior complexidade encaminhados a Coronel Oviedo. Cobertura
privada recomendada; IPS não disponível localmente.
\newpage
\secaoDiagnostico{Doctor Juan Manuel Frutos}{\mapaConteudo{-25.3833}{-55.8333}}{Doctor Juan Manuel Frutos (Caaguazu) apresenta perfil operacional
consolidado para comparacao territorial, com leitura conjunta de
infraestrutura, risco hidroclimatico e capacidade de continuidade de
servicos essenciais.

\subsubsection{Por que sim}

\begin{itemize}
\tightlist
\item
  Base institucional de referencia disponivel e rastreavel.
\item
  Estrutura comparativa (A-E/\allowbreak GSS) consistente para decisao relativa.
\item
  Plano de mitigacao acionavel ja definido no dossie.
\end{itemize}

\subsubsection{Por que não}

\begin{itemize}
\tightlist
\item
  Serie oficial distrital granular ainda pode apresentar lacunas
  temporais.
\item
  Sensibilidade do score exige monitoramento continuo de risco e
  conectividade.
\item
  Decisao final depende de validacao de campo e janela temporal recente.
\end{itemize}

\subsubsection{Pre-condicoes Minimas de
Decisao}

\begin{itemize}
\tightlist
\item
  Atualizar indicadores criticos em ciclo trimestral.
\item
  Confirmar trafegabilidade e contingencia hidrica/\allowbreak energetica local.
\item
  Validar checklist de resiliencia domiciliar/\allowbreak logistica antes de decisao
  final.
\end{itemize}

\subsubsection{Recomendação
Técnica}

Uso recomendado como candidato em analise multicriterio, condicionado ao
cumprimento das pre-condicoes acima. Referencia atual de score: GSS 7.7
em 2026-03-05; risco predominante: baixo.

\subsubsection{}

\begin{itemize}
\tightlist
\item
  \begin{enumerate}
  \def\labelenumi{\arabic{enumi})}
  \tightlist
  \item
    Delimitacao territorial validada por georreferencia institucional
    (Fonte: acesso: 2026-03-05).
  \end{enumerate}
\item
  \begin{enumerate}
  \def\labelenumi{\arabic{enumi})}
  \setcounter{enumi}{1}
  \tightlist
  \item
    População municipal/\allowbreak distrital referenciada por base censitaria
    nacional (Fonte: acesso: 2026-03-05).
  \end{enumerate}
\item
  \begin{enumerate}
  \def\labelenumi{\arabic{enumi})}
  \setcounter{enumi}{2}
  \tightlist
  \item
    Comparabilidade intra-departamental apoiada em indicadores
    distritais oficiais (Fonte: acesso: 2026-03-05).
  \end{enumerate}
\item
  \begin{enumerate}
  \def\labelenumi{\arabic{enumi})}
  \setcounter{enumi}{3}
  \tightlist
  \item
    Estado macro de conectividade viaria ancorado em informacoes de
    rodovias (Fonte: acesso: 2026-03-05).
  \end{enumerate}
\item
  \begin{enumerate}
  \def\labelenumi{\arabic{enumi})}
  \setcounter{enumi}{4}
  \tightlist
  \item
    Exposicao hidrometeorologica monitorada por avisos oficiais (Fonte:
    acesso: 2026-03-05).
  \end{enumerate}
\item
  \begin{enumerate}
  \def\labelenumi{\arabic{enumi})}
  \setcounter{enumi}{5}
  \tightlist
  \item
    Historico de contexto meteorologico apoiado em publicacoes tecnicas
    (Fonte: acesso: 2026-03-05).
  \end{enumerate}
\item
  \begin{enumerate}
  \def\labelenumi{\arabic{enumi})}
  \setcounter{enumi}{6}
  \tightlist
  \item
    Capacidade institucional de resposta a eventos apoiada em acoes da
    SEN\footnote{Secretaría de Emergencia Nacional (Secretaria de Emergência Nacional), órgão governamental de resposta a desastres.} (Fonte: acesso: 2026-03-05).
  \end{enumerate}
\item
  \begin{enumerate}
  \def\labelenumi{\arabic{enumi})}
  \setcounter{enumi}{7}
  \tightlist
  \item
    Infraestrutura energetica analisada via operador nacional (Fonte:
    acesso: 2026-03-05).
  \end{enumerate}
\item
  \begin{enumerate}
  \def\labelenumi{\arabic{enumi})}
  \setcounter{enumi}{8}
  \tightlist
  \item
    Referencias de obras e logistica nacional verificadas no MOPC
    (Fonte: acesso: 2026-03-05).
  \end{enumerate}
\item
  \begin{enumerate}
  \def\labelenumi{\arabic{enumi})}
  \setcounter{enumi}{9}
  \tightlist
  \item
    Contexto sociopolitico institucional referenciado por autoridade
    eleitoral (Fonte: acesso: 2026-03-05).
  \end{enumerate}
\item
  \begin{enumerate}
  \def\labelenumi{\arabic{enumi})}
  \setcounter{enumi}{10}
  \tightlist
  \item
    Camada cartografica de apoio eleitoral/\allowbreak territorial consultada
    (Fonte: acesso: 2026-03-05).
  \end{enumerate}
\item
  \begin{enumerate}
  \def\labelenumi{\arabic{enumi})}
  \setcounter{enumi}{11}
  \tightlist
  \item
    Dados publicos complementares para verificacao cruzada (Fonte:
    acesso: 2026-03-05).
  \end{enumerate}
\end{itemize}}
\begin{table}[ht!]
\small \begin{tabular}{p{3.0cm}p{0.8cm}p{6.0cm}} \toprule
\textbf{Critério} & \textbf{Nota} & \textbf{Justificativa} \\ \midrule
A - Ameaças Estratégicas & 8.1 & - \\
B - Risco Social & 8.0 & - \\
C - Riscos Naturais & 7.0 & - \\
D - Autossuficiência & 6.8 & - \\
E - Institucional & 8.1 & - \\
\midrule \textbf{GSS FINAL} & \textbf{7.7} & \textbf{Seguro (bom potencial com preparacao).} \\ \bottomrule \end{tabular}
\end{table}
\subsubsection{Vulnerabilidades e
Mitigação}\label{vuln-Doctor_Juan_Manuel_Frutos-vulnerabilidades-e-mitigauxe7uxe3o}

\begin{itemize}
\tightlist
\item
  Exposicao hidrometeorologica controlada (\textless45/\allowbreak 100).
\item
  Conectividade viaria favoravel (\textgreater=70/\allowbreak 100).
\item
  Estoque de contingencia para 14 dias (agua/\allowbreak alimentos/\allowbreak combustivel),
  priorizado quando risco hidro=38/\allowbreak 100.
\item
  Duplicar rotas/\allowbreak logistica quando conectividade viaria=74/\allowbreak 100 for
  \textless70; manter rota secundaria mapeada e testada trimestralmente.
\item
  Plano de energia resiliente com autonomia minima de 72h
  (geracao/\allowbreak backup), revisao semestral.
\end{itemize}
\subsubsection{Dossiê de Campo}\label{dossie-Doctor_Juan_Manuel_Frutos-dossiuxea-de-campo}
\begin{description}[style=nextline,leftmargin=0.5cm,font=\bfseries]
\item[Coordenadas]  25°23'S 55°50'W (geocodificado via Nominatim).
\end{description}
\textbf{Fonte climática:} NASA POWER Climatology API (período 2001-2020)
\textbf{Fonte luminosa:} estimativa world atlas

\paragraph{Irradiação Solar
(kWh/\allowbreak m²/\allowbreak dia)}\label{irradiauxe7uxe3o-solar-kwhmuxb2dia}

\begin{longtable}[]{@{}
  >{\raggedright\arraybackslash}p{(\linewidth - 24\tabcolsep) * \real{0.0746}}
  >{\raggedright\arraybackslash}p{(\linewidth - 24\tabcolsep) * \real{0.0746}}
  >{\raggedright\arraybackslash}p{(\linewidth - 24\tabcolsep) * \real{0.0746}}
  >{\raggedright\arraybackslash}p{(\linewidth - 24\tabcolsep) * \real{0.0746}}
  >{\raggedright\arraybackslash}p{(\linewidth - 24\tabcolsep) * \real{0.0746}}
  >{\raggedright\arraybackslash}p{(\linewidth - 24\tabcolsep) * \real{0.0746}}
  >{\raggedright\arraybackslash}p{(\linewidth - 24\tabcolsep) * \real{0.0746}}
  >{\raggedright\arraybackslash}p{(\linewidth - 24\tabcolsep) * \real{0.0746}}
  >{\raggedright\arraybackslash}p{(\linewidth - 24\tabcolsep) * \real{0.0746}}
  >{\raggedright\arraybackslash}p{(\linewidth - 24\tabcolsep) * \real{0.0746}}
  >{\raggedright\arraybackslash}p{(\linewidth - 24\tabcolsep) * \real{0.0746}}
  >{\raggedright\arraybackslash}p{(\linewidth - 24\tabcolsep) * \real{0.0746}}
  >{\raggedright\arraybackslash}p{(\linewidth - 24\tabcolsep) * \real{0.1045}}@{}}
\toprule\noalign{}
\begin{minipage}[b]{\linewidth}\raggedright
Jan
\end{minipage} & \begin{minipage}[b]{\linewidth}\raggedright
Fev
\end{minipage} & \begin{minipage}[b]{\linewidth}\raggedright
Mar
\end{minipage} & \begin{minipage}[b]{\linewidth}\raggedright
Abr
\end{minipage} & \begin{minipage}[b]{\linewidth}\raggedright
Mai
\end{minipage} & \begin{minipage}[b]{\linewidth}\raggedright
Jun
\end{minipage} & \begin{minipage}[b]{\linewidth}\raggedright
Jul
\end{minipage} & \begin{minipage}[b]{\linewidth}\raggedright
Ago
\end{minipage} & \begin{minipage}[b]{\linewidth}\raggedright
Set
\end{minipage} & \begin{minipage}[b]{\linewidth}\raggedright
Out
\end{minipage} & \begin{minipage}[b]{\linewidth}\raggedright
Nov
\end{minipage} & \begin{minipage}[b]{\linewidth}\raggedright
Dez
\end{minipage} & \begin{minipage}[b]{\linewidth}\raggedright
Média
\end{minipage} \\
\midrule\noalign{}
\endhead
\bottomrule\noalign{}
\endlastfoot
6.58 & 6.08 & 5.45 & 4.47 & 3.30 & 2.90 & 3.20 & 3.96 & 4.53 & 5.31 &
6.35 & 6.66 & \textbf{4.9} \\
\end{longtable}

\textbf{Inclinação solar recomendada:} 25° N (anual) · 35° N (inverno
jun-ago) · 15° N (verão nov-jan)

\paragraph{Precipitação (mm/\allowbreak mês)}\label{precipitauxe7uxe3o-mmmuxeas}

\begin{longtable}[]{@{}
  >{\raggedright\arraybackslash}p{(\linewidth - 24\tabcolsep) * \real{0.0704}}
  >{\raggedright\arraybackslash}p{(\linewidth - 24\tabcolsep) * \real{0.0704}}
  >{\raggedright\arraybackslash}p{(\linewidth - 24\tabcolsep) * \real{0.0704}}
  >{\raggedright\arraybackslash}p{(\linewidth - 24\tabcolsep) * \real{0.0704}}
  >{\raggedright\arraybackslash}p{(\linewidth - 24\tabcolsep) * \real{0.0704}}
  >{\raggedright\arraybackslash}p{(\linewidth - 24\tabcolsep) * \real{0.0704}}
  >{\raggedright\arraybackslash}p{(\linewidth - 24\tabcolsep) * \real{0.0704}}
  >{\raggedright\arraybackslash}p{(\linewidth - 24\tabcolsep) * \real{0.0704}}
  >{\raggedright\arraybackslash}p{(\linewidth - 24\tabcolsep) * \real{0.0704}}
  >{\raggedright\arraybackslash}p{(\linewidth - 24\tabcolsep) * \real{0.0704}}
  >{\raggedright\arraybackslash}p{(\linewidth - 24\tabcolsep) * \real{0.0704}}
  >{\raggedright\arraybackslash}p{(\linewidth - 24\tabcolsep) * \real{0.0704}}
  >{\raggedright\arraybackslash}p{(\linewidth - 24\tabcolsep) * \real{0.1549}}@{}}
\toprule\noalign{}
\begin{minipage}[b]{\linewidth}\raggedright
Jan
\end{minipage} & \begin{minipage}[b]{\linewidth}\raggedright
Fev
\end{minipage} & \begin{minipage}[b]{\linewidth}\raggedright
Mar
\end{minipage} & \begin{minipage}[b]{\linewidth}\raggedright
Abr
\end{minipage} & \begin{minipage}[b]{\linewidth}\raggedright
Mai
\end{minipage} & \begin{minipage}[b]{\linewidth}\raggedright
Jun
\end{minipage} & \begin{minipage}[b]{\linewidth}\raggedright
Jul
\end{minipage} & \begin{minipage}[b]{\linewidth}\raggedright
Ago
\end{minipage} & \begin{minipage}[b]{\linewidth}\raggedright
Set
\end{minipage} & \begin{minipage}[b]{\linewidth}\raggedright
Out
\end{minipage} & \begin{minipage}[b]{\linewidth}\raggedright
Nov
\end{minipage} & \begin{minipage}[b]{\linewidth}\raggedright
Dez
\end{minipage} & \begin{minipage}[b]{\linewidth}\raggedright
Total/\allowbreak ano
\end{minipage} \\
\midrule\noalign{}
\endhead
\bottomrule\noalign{}
\endlastfoot
142 & 130 & 126 & 152 & 177 & 94 & 79 & 55 & 103 & 189 & 188 & 172 &
\textbf{1607 mm} \\
\end{longtable}

\paragraph{Poluição Luminosa}\label{poluiuxe7uxe3o-luminosa}

\begin{longtable}[]{@{}ll@{}}
\toprule\noalign{}
Parâmetro & Valor \\
\midrule\noalign{}
\endhead
\bottomrule\noalign{}
\endlastfoot
Escala Bortle & 4 --- Céu rural-suburbano \\
Radiância artificial & 3.0 nW/\allowbreak cm²/\allowbreak sr \\
\end{longtable}
\paragraph{Solo (SoilGrids 2.0, média ponderada 0--30
cm)}

\begin{longtable}[]{@{}ll@{}}
\toprule\noalign{}
Parâmetro & Valor \\
\midrule\noalign{}
\endhead
\bottomrule\noalign{}
\endlastfoot
pH (H$_2$O) & 5.28 \\
Carbono orgânico (SOC) & 18.51 g/\allowbreak kg \\
Argila & 27.42 \% \\
Areia & 52.8 \% \\
Silte (calc.) & 19.8 \% \\
Densidade aparente & 1.29 g/\allowbreak cm³ \\
\textbf{Aptidão agrícola} & \textbf{Média} \\
\end{longtable}

Fonte: ISRIC SoilGrids 2.0 via WCS. Coords: -25.3833°, -55.8333°. Média
ponderada camadas 0-5, 5-15, 15-30 cm.
\subsubsection{Indicadores Sociais}

\textbf{Fontes:} PNUD 2020, INE\footnote{Instituto Nacional de Estadística (Instituto Nacional de Estatística), órgão oficial de dados demográficos e censitários.} Censo 2022, INE\footnote{Instituto Nacional de Estadística (Instituto Nacional de Estatística), órgão oficial de dados demográficos e censitários.} EPHC 2023, Ministerio
Público 2024\\
\textbf{Nota:} Valores marcados como \emph{(dept.)} referem-se ao
departamento de Caaguazú; valores \emph{(dist.)} são específicos deste
distrito. Dados marcados \emph{(est.)} são estimativas.

\begin{longtable}[]{@{}
  >{\raggedright\arraybackslash}p{(\linewidth - 4\tabcolsep) * \real{0.4231}}
  >{\raggedright\arraybackslash}p{(\linewidth - 4\tabcolsep) * \real{0.2692}}
  >{\raggedright\arraybackslash}p{(\linewidth - 4\tabcolsep) * \real{0.3077}}@{}}
\toprule\noalign{}
\begin{minipage}[b]{\linewidth}\raggedright
Indicador
\end{minipage} & \begin{minipage}[b]{\linewidth}\raggedright
Valor
\end{minipage} & \begin{minipage}[b]{\linewidth}\raggedright
Âmbito
\end{minipage} \\
\midrule\noalign{}
\endhead
\bottomrule\noalign{}
\endlastfoot
IDH (2020) & 0,715 & dept. \\
Ranking IDH nacional & 7/\allowbreak 18 & dept. \\
Esperança de vida & 73,1 anos & dept. \\
Escolaridade média & 8,4 anos & dept. \\
RNB per capita (USD PPA) & 9.000 & dept. \\
Pobreza monetária (\%) & 28,5\% & dept. \\
Pobreza extrema (\%) & 6,8\% & dept. \\
Índice de Gini & 0,447 & dept. \\
Acesso a água potável (\%) & 79,8\% & dept. \\
Acesso a saneamento (\%) & 78,3\% & dept. \\
Taxa de homicídios (est., /\allowbreak 100k hab) & \textasciitilde6--8 (estimativa,
próxima ou acima da média nacional) & dept. \\
Índice de segurança & médio & dept. \\
População (Censo 2022) & 20.451 hab. & dist. \\
Idade mediana & 28.0 anos & dist. \\
Taxa de fecundidade & N/\allowbreak D & dist. \\
\end{longtable}
\subsubsection{Combustível}

\textbf{Referência:} PETROPAR /\allowbreak  postos locais (2024) \textbf{Tipo de
localidade:} Interior

\begin{longtable}[]{@{}lll@{}}
\toprule\noalign{}
Combustível & USD/\allowbreak litro & Gs/\allowbreak litro (aprox.) \\
\midrule\noalign{}
\endhead
\bottomrule\noalign{}
\endlastfoot
Gasolina 93 oct & 0.98 & 7,252 \\
Gasolina 97 oct (premium) & 1.08 & 7,992 \\
Diesel & 0.91 & 6,734 \\
\end{longtable}

\begin{quote}
Preços podem variar ±5\% conforme posto e sazonalidade. Chaco e interior
remoto apresentam maior variação.
\end{quote}

\subsubsection{Cobertura Celular}

\textbf{Fonte:} CONATEL PY /\allowbreak  operadoras (2024)

\begin{longtable}[]{@{}ll@{}}
\toprule\noalign{}
Parâmetro & Valor \\
\midrule\noalign{}
\endhead
\bottomrule\noalign{}
\endlastfoot
Cobertura 4G população (dept.) & 88\% \\
Cobertura 4G área rural & 72\% \\
Melhor operadora & Tigo \\
Qualidade rural & boa \\
\end{longtable}

\begin{quote}
Para áreas rurais fora do núcleo urbano, recomenda-se chip Tigo como
principal e Personal como backup.
\end{quote}

\subsubsection{Internet}

\textbf{Fonte:} CONATEL /\allowbreak  Speedtest Ookla (2024)

\begin{longtable}[]{@{}ll@{}}
\toprule\noalign{}
Parâmetro & Valor \\
\midrule\noalign{}
\endhead
\bottomrule\noalign{}
\endlastfoot
Velocidade média download & 55 Mbps \\
Domicílios com internet (dept.) & 62\% \\
Tecnologia predominante & rádio/\allowbreak fibra \\
Opção rural & Starlink disponível (\textasciitilde USD 44/\allowbreak mês) \\
\end{longtable}

\subsubsection{Mercado Imobiliário e Terra
Rural}

\textbf{Fonte:} INDERT /\allowbreak  Clasificados.com.py (2024)

\begin{longtable}[]{@{}ll@{}}
\toprule\noalign{}
Tipo & Referência \\
\midrule\noalign{}
\endhead
\bottomrule\noalign{}
\endlastfoot
Terra agrícola alta prod. (USD/\allowbreak ha) & 8,800 \\
Imóvel urbano (USD/\allowbreak m²) & 850 \\
Aluguel 2 quartos (USD/\allowbreak mês) & 350 \\
\end{longtable}

\begin{quote}
Valores de referência departamental. Localidades menores podem ter
preços 20--40\% abaixo da capital departamental.
\end{quote}

\subsubsection{Saúde}

\textbf{Fonte:} MSPBS /\allowbreak  IPS Paraguay (2024)

\begin{longtable}[]{@{}ll@{}}
\toprule\noalign{}
Serviço & Disponibilidade \\
\midrule\noalign{}
\endhead
\bottomrule\noalign{}
\endlastfoot
USF /\allowbreak  Posto de Saúde & sim \\
Hospital Regional & não \\
IPS (seguro social) & não \\
Farmácia & sim \\
Distância ao hospital de referência & \textasciitilde90 km (Coronel
Oviedo) \\
\end{longtable}

\textbf{Principais estabelecimentos:} USF local; farmácias distritais.

\textbf{Observação para imigrantes:} Atendimento primário disponível.
Casos de maior complexidade encaminhados a Coronel Oviedo. Cobertura
privada recomendada; IPS não disponível localmente.
\newpage
\secaoDiagnostico{Jose Domingo Ocampos}{\mapaConteudo{-25.4500}{-55.6833}}{José Domingo Ocampos é uma localidade estratégica que combina a
facilidade logística da Ruta PY02 com o potencial hídrico massivo do
reservatório de Yguazú. Oferece solos férteis para policultivos e um
ambiente social estável e seguro. Sua principal vantagem é o acesso
rápido aos serviços de elite de Campo 9, sem a densidade populacional
daquela cidade industrial. É um local de alto potencial para
autossuficiência hídrica e alimentar.

\subsubsection{Por que sim}

\begin{itemize}
\tightlist
\item
  Abundância hídrica excepcional via reservatório e aquíferos.
\item
  Solo de alta produtividade e topografia favorável à mecanização.
\item
  Conectividade rápida com o polo de serviços de Campo 9 e com o Leste.
\item
  Sem restrições da Lei de Fronteira.
\end{itemize}

\subsubsection{Por que não}

\begin{itemize}
\tightlist
\item
  Alvo logístico rodoviário de alta visibilidade estratégica em crises.
\item
  Infraestrutura de saúde local limitada a atendimentos primários.
\item
  Risco de fluxos populacionais em cenários de colapso urbano nacional.
\end{itemize}

\subsubsection{Recomendação
Técnica}

Altamente indicado para estabelecimentos de autossuficiência familiar ou
comercial que valorizem a segurança hídrica. Requer um planejamento de
recuo físico em relação à rodovia principal e autonomia em saúde. O GSS
de 6.4 reflete a força da base de recursos naturais da localidade.
Referencia atual de score: GSS 6.4 em 2026-03-05.

\subsubsection{}

\begin{itemize}
\tightlist
\item
  \begin{enumerate}
  \def\labelenumi{\arabic{enumi})}
  \tightlist
  \item
    Dados censitários validados (INE\footnote{Instituto Nacional de Estadística (Instituto Nacional de Estatística), órgão oficial de dados demográficos e censitários.} 2022).
  \end{enumerate}
\item
  \begin{enumerate}
  \def\labelenumi{\arabic{enumi})}
  \setcounter{enumi}{1}
  \tightlist
  \item
    Relatórios de produção agrícola e pesca regional (MAG).
  \end{enumerate}
\item
  \begin{enumerate}
  \def\labelenumi{\arabic{enumi})}
  \setcounter{enumi}{2}
  \tightlist
  \item
    Mapa de conectividade do corredor PY02 (MOPC 2026).
  \end{enumerate}
\item
  \begin{enumerate}
  \def\labelenumi{\arabic{enumi})}
  \setcounter{enumi}{3}
  \tightlist
  \item
    Registro de segurança pública departamental.
  \end{enumerate}
\end{itemize}}
\begin{table}[ht!]
\small \begin{tabular}{p{3.0cm}p{0.8cm}p{6.0cm}} \toprule
\textbf{Critério} & \textbf{Nota} & \textbf{Justificativa} \\ \midrule
A - Ameaças Estratégicas & 5.0 & Posicionamento estratégico na PY02 e proximidade com reservatório vital; visibilidade tática moderada \\
B - Risco Social & 6.5 & População pequena e densidade controlada; risco social urbano desprezível \\
C - Riscos Naturais & 7.0 & Geologia muito estável e riscos naturais climáticos típicos manejáveis \\
D - Autossuficiência & 7.5 & Excelentes recursos: solo fértil, abundância hídrica massiva e potencial de pesca \\
E - Institucional & 6.0 & Ambiente institucional estável, mas dependente de serviços de cidades vizinhas \\
\midrule \textbf{GSS FINAL} & \textbf{6.4} & \textbf{Moderadamente Seguro (Excelente opção para quem busca integração agrícola com abundância de recursos hídricos).} \\ \bottomrule \end{tabular}
\end{table}
\subsubsection{Vulnerabilidades e
Mitigação}\label{vuln-Jose_Domingo_Ocampos-vulnerabilidades-e-mitigauxe7uxe3o}

\begin{itemize}
\tightlist
\item
  \textbf{Exposição Rodoviária:} Localização direta na PY02 atrai fluxos
  de trânsito intenso em crises (Mitigação: residência rural afastada
  3-5 km do eixo asfáltico principal).
\item
  \textbf{Dependência de Serviços Externos:} Necessidade de deslocamento
  para Campo 9 para serviços hospitalares e comerciais avançados
  (Mitigação: manutenção de estoque médico básico e autonomia
  logística).
\item
  \textbf{Risco Hidrográfico:} Monitoramento das margens do reservatório
  de Yguazú em anos de chuvas extremas (Mitigação: construção em cotas
  de segurança \textgreater230m).
\end{itemize}
\subsubsection{Dossiê de Campo}\label{dossie-Jose_Domingo_Ocampos-dossiuxea-de-campo}
\begin{description}[style=nextline,leftmargin=0.5cm,font=\bfseries]
\item[Coordenadas]  25°27'S 55°41'W.
\item[Topografia]  Relevo predominantemente plano a suavemente ondulado. Localizado nas imediações do reservatório de Yguazú. Altitude \textasciitilde 220m. Área de \textasciitilde 164 km². Geologicamente muito estável, risco sísmico desprezível.
\item[Alvos Estratégicos]  Ponto de trânsito e serviços na Ruta PY02 duplicada; Proximidade estratégica com a barragem e reservatório de Yguazú. Ausência de alvos militares diretos, mas vulnerável pela exposição rodoviária internacional.
\item[Fallout]  Ventos predominantes do Norte (NE) e Leste. Baixo risco de fallout direto.
\item[População]  7.459 habitantes (Censo 2022 Final). Perfil majoritariamente rural (70\%) com núcleo urbano comercial linear ao longo da rodovia.
\item[Densidade]  \textasciitilde 45,4 hab/\allowbreak km² (Baixa densidade, oferecendo isolamento relativo em áreas de colônias).
\item[Vias de Acesso]  Localização estratégica sobre a Ruta PY02 duplicada. Conectividade asfáltica de elite para Ciudad del Este (aprox. 100 km) e Assunção. Risco elevado de tornar-se rota de fuga massiva e controle logístico em crises nacionais.
\item[Serviços]  Centro de Saúde local para atendimentos primários. Dependência direta de J. Eulogio Estigarribia (Campo 9), a apenas 15 km, para saúde de alta complexidade, bancos e insumos agroindustriais.
\item[Custo de Vida]  Baixo; economia baseada na agricultura familiar e pecuária.
\item[Preço da Terra]  US\$ 4.000-7.000/\allowbreak ha (áreas rurais agrícolas); US\$ 2.000-3.500/\allowbreak ha (áreas de campo/\allowbreak pecuária).
\item[Hidrologia]  Baixo risco de inundações sistêmicas catastróficas. Áreas baixas próximas ao Lago Yguazú podem sofrer variações de nível em eventos climáticos extremos. Drenagem natural eficiente.
\item[Clima]  Subtropical Úmido. Precipitação anual de 1.600-1.800 mm. Verões quentes e chuvosos.
\item[Energia]  Rede nacional (Itaipu); boa estabilidade por estar no eixo principal de transmissão do país.
\item[Água]  Abundante via poços artesianos e proximidade com o reservatório de Yguazú (potencial hídrico massivo).
\item[Qualidade do Solo]  Latossolos e Podzólicos; férteis e profundos, ideais para milho, soja, mandioca e pastagens cultivadas.
\item[Resources Locais]  Produção de gado de corte, grãos e pesca comercial/\allowbreak subsistência no lago. Elevado potencial para autossuficiência alimentar e hídrica em nível de propriedade.
\item[Segurança]  Zona rural pacífica e estável. Baixíssima criminalidade urbana violenta. Estabilidade social sustentada pela cultura agrícola e proximidade com o polo próspero de Campo 9. Exposição moderada a fluxos de trânsito rodoviário.
\item[Leis Local: Município consolidado. Livre de Restrição de Fronteira]  Fora da faixa de 50km da fronteira internacional, permitindo plena titularidade para brasileiros.
\end{description}
\textbf{Fonte climática:} NASA POWER Climatology API (período 2001-2020)
\textbf{Fonte luminosa:} estimativa world atlas

\paragraph{Irradiação Solar
(kWh/\allowbreak m²/\allowbreak dia)}\label{irradiauxe7uxe3o-solar-kwhmuxb2dia}

\begin{longtable}[]{@{}
  >{\raggedright\arraybackslash}p{(\linewidth - 24\tabcolsep) * \real{0.0746}}
  >{\raggedright\arraybackslash}p{(\linewidth - 24\tabcolsep) * \real{0.0746}}
  >{\raggedright\arraybackslash}p{(\linewidth - 24\tabcolsep) * \real{0.0746}}
  >{\raggedright\arraybackslash}p{(\linewidth - 24\tabcolsep) * \real{0.0746}}
  >{\raggedright\arraybackslash}p{(\linewidth - 24\tabcolsep) * \real{0.0746}}
  >{\raggedright\arraybackslash}p{(\linewidth - 24\tabcolsep) * \real{0.0746}}
  >{\raggedright\arraybackslash}p{(\linewidth - 24\tabcolsep) * \real{0.0746}}
  >{\raggedright\arraybackslash}p{(\linewidth - 24\tabcolsep) * \real{0.0746}}
  >{\raggedright\arraybackslash}p{(\linewidth - 24\tabcolsep) * \real{0.0746}}
  >{\raggedright\arraybackslash}p{(\linewidth - 24\tabcolsep) * \real{0.0746}}
  >{\raggedright\arraybackslash}p{(\linewidth - 24\tabcolsep) * \real{0.0746}}
  >{\raggedright\arraybackslash}p{(\linewidth - 24\tabcolsep) * \real{0.0746}}
  >{\raggedright\arraybackslash}p{(\linewidth - 24\tabcolsep) * \real{0.1045}}@{}}
\toprule\noalign{}
\begin{minipage}[b]{\linewidth}\raggedright
Jan
\end{minipage} & \begin{minipage}[b]{\linewidth}\raggedright
Fev
\end{minipage} & \begin{minipage}[b]{\linewidth}\raggedright
Mar
\end{minipage} & \begin{minipage}[b]{\linewidth}\raggedright
Abr
\end{minipage} & \begin{minipage}[b]{\linewidth}\raggedright
Mai
\end{minipage} & \begin{minipage}[b]{\linewidth}\raggedright
Jun
\end{minipage} & \begin{minipage}[b]{\linewidth}\raggedright
Jul
\end{minipage} & \begin{minipage}[b]{\linewidth}\raggedright
Ago
\end{minipage} & \begin{minipage}[b]{\linewidth}\raggedright
Set
\end{minipage} & \begin{minipage}[b]{\linewidth}\raggedright
Out
\end{minipage} & \begin{minipage}[b]{\linewidth}\raggedright
Nov
\end{minipage} & \begin{minipage}[b]{\linewidth}\raggedright
Dez
\end{minipage} & \begin{minipage}[b]{\linewidth}\raggedright
Média
\end{minipage} \\
\midrule\noalign{}
\endhead
\bottomrule\noalign{}
\endlastfoot
6.58 & 6.08 & 5.45 & 4.47 & 3.30 & 2.90 & 3.20 & 3.96 & 4.53 & 5.31 &
6.35 & 6.66 & \textbf{4.9} \\
\end{longtable}

\textbf{Inclinação solar recomendada:} 25° N (anual) · 35° N (inverno
jun-ago) · 15° N (verão nov-jan)

\paragraph{Precipitação (mm/\allowbreak mês)}\label{precipitauxe7uxe3o-mmmuxeas}

\begin{longtable}[]{@{}
  >{\raggedright\arraybackslash}p{(\linewidth - 24\tabcolsep) * \real{0.0704}}
  >{\raggedright\arraybackslash}p{(\linewidth - 24\tabcolsep) * \real{0.0704}}
  >{\raggedright\arraybackslash}p{(\linewidth - 24\tabcolsep) * \real{0.0704}}
  >{\raggedright\arraybackslash}p{(\linewidth - 24\tabcolsep) * \real{0.0704}}
  >{\raggedright\arraybackslash}p{(\linewidth - 24\tabcolsep) * \real{0.0704}}
  >{\raggedright\arraybackslash}p{(\linewidth - 24\tabcolsep) * \real{0.0704}}
  >{\raggedright\arraybackslash}p{(\linewidth - 24\tabcolsep) * \real{0.0704}}
  >{\raggedright\arraybackslash}p{(\linewidth - 24\tabcolsep) * \real{0.0704}}
  >{\raggedright\arraybackslash}p{(\linewidth - 24\tabcolsep) * \real{0.0704}}
  >{\raggedright\arraybackslash}p{(\linewidth - 24\tabcolsep) * \real{0.0704}}
  >{\raggedright\arraybackslash}p{(\linewidth - 24\tabcolsep) * \real{0.0704}}
  >{\raggedright\arraybackslash}p{(\linewidth - 24\tabcolsep) * \real{0.0704}}
  >{\raggedright\arraybackslash}p{(\linewidth - 24\tabcolsep) * \real{0.1549}}@{}}
\toprule\noalign{}
\begin{minipage}[b]{\linewidth}\raggedright
Jan
\end{minipage} & \begin{minipage}[b]{\linewidth}\raggedright
Fev
\end{minipage} & \begin{minipage}[b]{\linewidth}\raggedright
Mar
\end{minipage} & \begin{minipage}[b]{\linewidth}\raggedright
Abr
\end{minipage} & \begin{minipage}[b]{\linewidth}\raggedright
Mai
\end{minipage} & \begin{minipage}[b]{\linewidth}\raggedright
Jun
\end{minipage} & \begin{minipage}[b]{\linewidth}\raggedright
Jul
\end{minipage} & \begin{minipage}[b]{\linewidth}\raggedright
Ago
\end{minipage} & \begin{minipage}[b]{\linewidth}\raggedright
Set
\end{minipage} & \begin{minipage}[b]{\linewidth}\raggedright
Out
\end{minipage} & \begin{minipage}[b]{\linewidth}\raggedright
Nov
\end{minipage} & \begin{minipage}[b]{\linewidth}\raggedright
Dez
\end{minipage} & \begin{minipage}[b]{\linewidth}\raggedright
Total/\allowbreak ano
\end{minipage} \\
\midrule\noalign{}
\endhead
\bottomrule\noalign{}
\endlastfoot
142 & 130 & 126 & 152 & 177 & 94 & 79 & 55 & 103 & 189 & 188 & 172 &
\textbf{1607 mm} \\
\end{longtable}

\paragraph{Poluição Luminosa}\label{poluiuxe7uxe3o-luminosa}

\begin{longtable}[]{@{}ll@{}}
\toprule\noalign{}
Parâmetro & Valor \\
\midrule\noalign{}
\endhead
\bottomrule\noalign{}
\endlastfoot
Escala Bortle & 4 --- Céu rural-suburbano \\
Radiância artificial & 3.0 nW/\allowbreak cm²/\allowbreak sr \\
\end{longtable}
\paragraph{Solo (SoilGrids 2.0, média ponderada 0--30
cm)}

\begin{longtable}[]{@{}ll@{}}
\toprule\noalign{}
Parâmetro & Valor \\
\midrule\noalign{}
\endhead
\bottomrule\noalign{}
\endlastfoot
pH (H$_2$O) & 5.25 \\
Carbono orgânico (SOC) & 19.1 g/\allowbreak kg \\
Argila & 29.44 \% \\
Areia & 44.89 \% \\
Silte (calc.) & 25.7 \% \\
Densidade aparente & 1.29 g/\allowbreak cm³ \\
\textbf{Aptidão agrícola} & \textbf{Média} \\
\end{longtable}

Fonte: ISRIC SoilGrids 2.0 via WCS. Coords: -25.45°, -55.6833°. Média
ponderada camadas 0-5, 5-15, 15-30 cm.
\subsubsection{Indicadores Sociais}

\textbf{Fontes:} PNUD 2020, INE\footnote{Instituto Nacional de Estadística (Instituto Nacional de Estatística), órgão oficial de dados demográficos e censitários.} Censo 2022, INE\footnote{Instituto Nacional de Estadística (Instituto Nacional de Estatística), órgão oficial de dados demográficos e censitários.} EPHC 2023, Ministerio
Público 2024\\
\textbf{Nota:} Valores marcados como \emph{(dept.)} referem-se ao
departamento de Caaguazú; valores \emph{(dist.)} são específicos deste
distrito. Dados marcados \emph{(est.)} são estimativas.

\begin{longtable}[]{@{}
  >{\raggedright\arraybackslash}p{(\linewidth - 4\tabcolsep) * \real{0.4231}}
  >{\raggedright\arraybackslash}p{(\linewidth - 4\tabcolsep) * \real{0.2692}}
  >{\raggedright\arraybackslash}p{(\linewidth - 4\tabcolsep) * \real{0.3077}}@{}}
\toprule\noalign{}
\begin{minipage}[b]{\linewidth}\raggedright
Indicador
\end{minipage} & \begin{minipage}[b]{\linewidth}\raggedright
Valor
\end{minipage} & \begin{minipage}[b]{\linewidth}\raggedright
Âmbito
\end{minipage} \\
\midrule\noalign{}
\endhead
\bottomrule\noalign{}
\endlastfoot
IDH (2020) & 0,715 & dept. \\
Ranking IDH nacional & 7/\allowbreak 18 & dept. \\
Esperança de vida & 73,1 anos & dept. \\
Escolaridade média & 8.3 anos & dist. \\
RNB per capita (USD PPA) & 9.000 & dept. \\
Pobreza monetária (\%) & 28,5\% & dept. \\
Pobreza extrema (\%) & 6,8\% & dept. \\
Índice de Gini & 0,447 & dept. \\
Acesso a água potável (\%) & 79,8\% & dept. \\
Acesso a saneamento (\%) & 78,3\% & dept. \\
Taxa de homicídios (est., /\allowbreak 100k hab) & \textasciitilde6--8 (estimativa,
próxima ou acima da média nacional) & dept. \\
Índice de segurança & médio & dept. \\
População (Censo 2022) & 7.459 hab. & dist. \\
Idade mediana & 28.0 anos & dist. \\
Taxa de fecundidade & N/\allowbreak D & dist. \\
\end{longtable}
\subsubsection{Combustível}

\textbf{Referência:} PETROPAR /\allowbreak  postos locais (2024) \textbf{Tipo de
localidade:} Interior

\begin{longtable}[]{@{}lll@{}}
\toprule\noalign{}
Combustível & USD/\allowbreak litro & Gs/\allowbreak litro (aprox.) \\
\midrule\noalign{}
\endhead
\bottomrule\noalign{}
\endlastfoot
Gasolina 93 oct & 0.98 & 7,252 \\
Gasolina 97 oct (premium) & 1.08 & 7,992 \\
Diesel & 0.91 & 6,734 \\
\end{longtable}

\begin{quote}
Preços podem variar ±5\% conforme posto e sazonalidade. Chaco e interior
remoto apresentam maior variação.
\end{quote}

\subsubsection{Cobertura Celular}

\textbf{Fonte:} CONATEL PY /\allowbreak  operadoras (2024)

\begin{longtable}[]{@{}ll@{}}
\toprule\noalign{}
Parâmetro & Valor \\
\midrule\noalign{}
\endhead
\bottomrule\noalign{}
\endlastfoot
Cobertura 4G população (dept.) & 88\% \\
Cobertura 4G área rural & 72\% \\
Melhor operadora & Tigo \\
Qualidade rural & boa \\
\end{longtable}

\begin{quote}
Para áreas rurais fora do núcleo urbano, recomenda-se chip Tigo como
principal e Personal como backup.
\end{quote}

\subsubsection{Internet}

\textbf{Fonte:} CONATEL /\allowbreak  Speedtest Ookla (2024)

\begin{longtable}[]{@{}ll@{}}
\toprule\noalign{}
Parâmetro & Valor \\
\midrule\noalign{}
\endhead
\bottomrule\noalign{}
\endlastfoot
Velocidade média download & 55 Mbps \\
Domicílios com internet (dept.) & 62\% \\
Tecnologia predominante & rádio/\allowbreak fibra \\
Opção rural & Starlink disponível (\textasciitilde USD 44/\allowbreak mês) \\
\end{longtable}

\subsubsection{Mercado Imobiliário e Terra
Rural}

\textbf{Fonte:} INDERT /\allowbreak  Clasificados.com.py (2024)

\begin{longtable}[]{@{}ll@{}}
\toprule\noalign{}
Tipo & Referência \\
\midrule\noalign{}
\endhead
\bottomrule\noalign{}
\endlastfoot
Terra agrícola alta prod. (USD/\allowbreak ha) & 8,800 \\
Imóvel urbano (USD/\allowbreak m²) & 850 \\
Aluguel 2 quartos (USD/\allowbreak mês) & 350 \\
\end{longtable}

\begin{quote}
Valores de referência departamental. Localidades menores podem ter
preços 20--40\% abaixo da capital departamental.
\end{quote}

\subsubsection{Saúde}

\textbf{Fonte:} MSPBS /\allowbreak  IPS Paraguay (2024)

\begin{longtable}[]{@{}ll@{}}
\toprule\noalign{}
Serviço & Disponibilidade \\
\midrule\noalign{}
\endhead
\bottomrule\noalign{}
\endlastfoot
USF /\allowbreak  Posto de Saúde & sim \\
Hospital Regional & não \\
IPS (seguro social) & não \\
Farmácia & sim \\
Distância ao hospital de referência & \textasciitilde70 km (Coronel
Oviedo) \\
\end{longtable}

\textbf{Principais estabelecimentos:} USF local; farmácias distritais.

\textbf{Observação para imigrantes:} Atendimento primário disponível.
Casos de maior complexidade encaminhados a Coronel Oviedo. Cobertura
privada recomendada; IPS não disponível localmente.
\newpage
\secaoDiagnostico{La Pastora}{\mapaConteudo{-25.2500}{-56.5333}}{La Pastora (Caaguazu) apresenta perfil operacional consolidado para
comparacao territorial, com leitura conjunta de infraestrutura, risco
hidroclimatico e capacidade de continuidade de servicos essenciais.

\subsubsection{Por que sim}

\begin{itemize}
\tightlist
\item
  Base institucional de referencia disponivel e rastreavel.
\item
  Estrutura comparativa (A-E/\allowbreak GSS) consistente para decisao relativa.
\item
  Plano de mitigacao acionavel ja definido no dossie.
\end{itemize}

\subsubsection{Por que não}

\begin{itemize}
\tightlist
\item
  Serie oficial distrital granular ainda pode apresentar lacunas
  temporais.
\item
  Sensibilidade do score exige monitoramento continuo de risco e
  conectividade.
\item
  Decisao final depende de validacao de campo e janela temporal recente.
\end{itemize}

\subsubsection{Pre-condicoes Minimas de
Decisao}

\begin{itemize}
\tightlist
\item
  Atualizar indicadores criticos em ciclo trimestral.
\item
  Confirmar trafegabilidade e contingencia hidrica/\allowbreak energetica local.
\item
  Validar checklist de resiliencia domiciliar/\allowbreak logistica antes de decisao
  final.
\end{itemize}

\subsubsection{Recomendação
Técnica}

Uso recomendado como candidato em analise multicriterio, condicionado ao
cumprimento das pre-condicoes acima. Referencia atual de score: GSS 7.0
em 2026-03-05; risco predominante: baixo.

\subsubsection{}

\begin{itemize}
\tightlist
\item
  \begin{enumerate}
  \def\labelenumi{\arabic{enumi})}
  \tightlist
  \item
    Delimitacao territorial validada por georreferencia institucional
    (Fonte: acesso: 2026-03-05).
  \end{enumerate}
\item
  \begin{enumerate}
  \def\labelenumi{\arabic{enumi})}
  \setcounter{enumi}{1}
  \tightlist
  \item
    População municipal/\allowbreak distrital referenciada por base censitaria
    nacional (Fonte: acesso: 2026-03-05).
  \end{enumerate}
\item
  \begin{enumerate}
  \def\labelenumi{\arabic{enumi})}
  \setcounter{enumi}{2}
  \tightlist
  \item
    Comparabilidade intra-departamental apoiada em indicadores
    distritais oficiais (Fonte: acesso: 2026-03-05).
  \end{enumerate}
\item
  \begin{enumerate}
  \def\labelenumi{\arabic{enumi})}
  \setcounter{enumi}{3}
  \tightlist
  \item
    Estado macro de conectividade viaria ancorado em informacoes de
    rodovias (Fonte: acesso: 2026-03-05).
  \end{enumerate}
\item
  \begin{enumerate}
  \def\labelenumi{\arabic{enumi})}
  \setcounter{enumi}{4}
  \tightlist
  \item
    Exposicao hidrometeorologica monitorada por avisos oficiais (Fonte:
    acesso: 2026-03-05).
  \end{enumerate}
\item
  \begin{enumerate}
  \def\labelenumi{\arabic{enumi})}
  \setcounter{enumi}{5}
  \tightlist
  \item
    Historico de contexto meteorologico apoiado em publicacoes tecnicas
    (Fonte: acesso: 2026-03-05).
  \end{enumerate}
\item
  \begin{enumerate}
  \def\labelenumi{\arabic{enumi})}
  \setcounter{enumi}{6}
  \tightlist
  \item
    Capacidade institucional de resposta a eventos apoiada em acoes da
    SEN\footnote{Secretaría de Emergencia Nacional (Secretaria de Emergência Nacional), órgão governamental de resposta a desastres.} (Fonte: acesso: 2026-03-05).
  \end{enumerate}
\item
  \begin{enumerate}
  \def\labelenumi{\arabic{enumi})}
  \setcounter{enumi}{7}
  \tightlist
  \item
    Infraestrutura energetica analisada via operador nacional (Fonte:
    acesso: 2026-03-05).
  \end{enumerate}
\item
  \begin{enumerate}
  \def\labelenumi{\arabic{enumi})}
  \setcounter{enumi}{8}
  \tightlist
  \item
    Referencias de obras e logistica nacional verificadas no MOPC
    (Fonte: acesso: 2026-03-05).
  \end{enumerate}
\item
  \begin{enumerate}
  \def\labelenumi{\arabic{enumi})}
  \setcounter{enumi}{9}
  \tightlist
  \item
    Contexto sociopolitico institucional referenciado por autoridade
    eleitoral (Fonte: acesso: 2026-03-05).
  \end{enumerate}
\item
  \begin{enumerate}
  \def\labelenumi{\arabic{enumi})}
  \setcounter{enumi}{10}
  \tightlist
  \item
    Camada cartografica de apoio eleitoral/\allowbreak territorial consultada
    (Fonte: acesso: 2026-03-05).
  \end{enumerate}
\item
  \begin{enumerate}
  \def\labelenumi{\arabic{enumi})}
  \setcounter{enumi}{11}
  \tightlist
  \item
    Dados publicos complementares para verificacao cruzada (Fonte:
    acesso: 2026-03-05).
  \end{enumerate}
\end{itemize}}
\begin{table}[ht!]
\small \begin{tabular}{p{3.0cm}p{0.8cm}p{6.0cm}} \toprule
\textbf{Critério} & \textbf{Nota} & \textbf{Justificativa} \\ \midrule
A - Ameaças Estratégicas & 7.7 & - \\
B - Risco Social & 5.2 & - \\
C - Riscos Naturais & 7.7 & - \\
D - Autossuficiência & 7.7 & - \\
E - Institucional & 6.5 & - \\
\midrule \textbf{GSS FINAL} & \textbf{7.0} & \textbf{Seguro (bom potencial com preparacao).} \\ \bottomrule \end{tabular}
\end{table}
\subsubsection{Vulnerabilidades e
Mitigação}\label{vuln-La_Pastora-vulnerabilidades-e-mitigauxe7uxe3o}

\begin{itemize}
\tightlist
\item
  Exposicao hidrometeorologica controlada (\textless45/\allowbreak 100).
\item
  Conectividade viaria intermediaria (50-69/\allowbreak 100).
\item
  Estoque de contingencia para 14 dias (agua/\allowbreak alimentos/\allowbreak combustivel),
  priorizado quando risco hidro=29/\allowbreak 100.
\item
  Duplicar rotas/\allowbreak logistica quando conectividade viaria=52/\allowbreak 100 for
  \textless70; manter rota secundaria mapeada e testada trimestralmente.
\item
  Plano de energia resiliente com autonomia minima de 72h
  (geracao/\allowbreak backup), revisao semestral.
\end{itemize}
\subsubsection{Dossiê de Campo}\label{dossie-La_Pastora-dossiuxea-de-campo}
\begin{description}[style=nextline,leftmargin=0.5cm,font=\bfseries]
\item[Coordenadas]  25°15'S 56°32'W (geocodificado via Nominatim).
\end{description}
\textbf{Fonte climática:} NASA POWER Climatology API (período 2001-2020)
\textbf{Fonte luminosa:} estimativa world atlas

\paragraph{Irradiação Solar
(kWh/\allowbreak m²/\allowbreak dia)}\label{irradiauxe7uxe3o-solar-kwhmuxb2dia}

\begin{longtable}[]{@{}
  >{\raggedright\arraybackslash}p{(\linewidth - 24\tabcolsep) * \real{0.0746}}
  >{\raggedright\arraybackslash}p{(\linewidth - 24\tabcolsep) * \real{0.0746}}
  >{\raggedright\arraybackslash}p{(\linewidth - 24\tabcolsep) * \real{0.0746}}
  >{\raggedright\arraybackslash}p{(\linewidth - 24\tabcolsep) * \real{0.0746}}
  >{\raggedright\arraybackslash}p{(\linewidth - 24\tabcolsep) * \real{0.0746}}
  >{\raggedright\arraybackslash}p{(\linewidth - 24\tabcolsep) * \real{0.0746}}
  >{\raggedright\arraybackslash}p{(\linewidth - 24\tabcolsep) * \real{0.0746}}
  >{\raggedright\arraybackslash}p{(\linewidth - 24\tabcolsep) * \real{0.0746}}
  >{\raggedright\arraybackslash}p{(\linewidth - 24\tabcolsep) * \real{0.0746}}
  >{\raggedright\arraybackslash}p{(\linewidth - 24\tabcolsep) * \real{0.0746}}
  >{\raggedright\arraybackslash}p{(\linewidth - 24\tabcolsep) * \real{0.0746}}
  >{\raggedright\arraybackslash}p{(\linewidth - 24\tabcolsep) * \real{0.0746}}
  >{\raggedright\arraybackslash}p{(\linewidth - 24\tabcolsep) * \real{0.1045}}@{}}
\toprule\noalign{}
\begin{minipage}[b]{\linewidth}\raggedright
Jan
\end{minipage} & \begin{minipage}[b]{\linewidth}\raggedright
Fev
\end{minipage} & \begin{minipage}[b]{\linewidth}\raggedright
Mar
\end{minipage} & \begin{minipage}[b]{\linewidth}\raggedright
Abr
\end{minipage} & \begin{minipage}[b]{\linewidth}\raggedright
Mai
\end{minipage} & \begin{minipage}[b]{\linewidth}\raggedright
Jun
\end{minipage} & \begin{minipage}[b]{\linewidth}\raggedright
Jul
\end{minipage} & \begin{minipage}[b]{\linewidth}\raggedright
Ago
\end{minipage} & \begin{minipage}[b]{\linewidth}\raggedright
Set
\end{minipage} & \begin{minipage}[b]{\linewidth}\raggedright
Out
\end{minipage} & \begin{minipage}[b]{\linewidth}\raggedright
Nov
\end{minipage} & \begin{minipage}[b]{\linewidth}\raggedright
Dez
\end{minipage} & \begin{minipage}[b]{\linewidth}\raggedright
Média
\end{minipage} \\
\midrule\noalign{}
\endhead
\bottomrule\noalign{}
\endlastfoot
6.74 & 6.20 & 5.50 & 4.48 & 3.35 & 2.88 & 3.21 & 3.96 & 4.59 & 5.43 &
6.43 & 6.81 & \textbf{4.96} \\
\end{longtable}

\textbf{Inclinação solar recomendada:} 25° N (anual) · 35° N (inverno
jun-ago) · 15° N (verão nov-jan)

\paragraph{Precipitação (mm/\allowbreak mês)}\label{precipitauxe7uxe3o-mmmuxeas}

\begin{longtable}[]{@{}
  >{\raggedright\arraybackslash}p{(\linewidth - 24\tabcolsep) * \real{0.0704}}
  >{\raggedright\arraybackslash}p{(\linewidth - 24\tabcolsep) * \real{0.0704}}
  >{\raggedright\arraybackslash}p{(\linewidth - 24\tabcolsep) * \real{0.0704}}
  >{\raggedright\arraybackslash}p{(\linewidth - 24\tabcolsep) * \real{0.0704}}
  >{\raggedright\arraybackslash}p{(\linewidth - 24\tabcolsep) * \real{0.0704}}
  >{\raggedright\arraybackslash}p{(\linewidth - 24\tabcolsep) * \real{0.0704}}
  >{\raggedright\arraybackslash}p{(\linewidth - 24\tabcolsep) * \real{0.0704}}
  >{\raggedright\arraybackslash}p{(\linewidth - 24\tabcolsep) * \real{0.0704}}
  >{\raggedright\arraybackslash}p{(\linewidth - 24\tabcolsep) * \real{0.0704}}
  >{\raggedright\arraybackslash}p{(\linewidth - 24\tabcolsep) * \real{0.0704}}
  >{\raggedright\arraybackslash}p{(\linewidth - 24\tabcolsep) * \real{0.0704}}
  >{\raggedright\arraybackslash}p{(\linewidth - 24\tabcolsep) * \real{0.0704}}
  >{\raggedright\arraybackslash}p{(\linewidth - 24\tabcolsep) * \real{0.1549}}@{}}
\toprule\noalign{}
\begin{minipage}[b]{\linewidth}\raggedright
Jan
\end{minipage} & \begin{minipage}[b]{\linewidth}\raggedright
Fev
\end{minipage} & \begin{minipage}[b]{\linewidth}\raggedright
Mar
\end{minipage} & \begin{minipage}[b]{\linewidth}\raggedright
Abr
\end{minipage} & \begin{minipage}[b]{\linewidth}\raggedright
Mai
\end{minipage} & \begin{minipage}[b]{\linewidth}\raggedright
Jun
\end{minipage} & \begin{minipage}[b]{\linewidth}\raggedright
Jul
\end{minipage} & \begin{minipage}[b]{\linewidth}\raggedright
Ago
\end{minipage} & \begin{minipage}[b]{\linewidth}\raggedright
Set
\end{minipage} & \begin{minipage}[b]{\linewidth}\raggedright
Out
\end{minipage} & \begin{minipage}[b]{\linewidth}\raggedright
Nov
\end{minipage} & \begin{minipage}[b]{\linewidth}\raggedright
Dez
\end{minipage} & \begin{minipage}[b]{\linewidth}\raggedright
Total/\allowbreak ano
\end{minipage} \\
\midrule\noalign{}
\endhead
\bottomrule\noalign{}
\endlastfoot
124 & 143 & 125 & 147 & 173 & 90 & 70 & 46 & 94 & 184 & 190 & 168 &
\textbf{1557 mm} \\
\end{longtable}

\paragraph{Poluição Luminosa}\label{poluiuxe7uxe3o-luminosa}

\begin{longtable}[]{@{}ll@{}}
\toprule\noalign{}
Parâmetro & Valor \\
\midrule\noalign{}
\endhead
\bottomrule\noalign{}
\endlastfoot
Escala Bortle & 4 --- Céu rural-suburbano \\
Radiância artificial & 3.0 nW/\allowbreak cm²/\allowbreak sr \\
\end{longtable}
\paragraph{Solo (SoilGrids 2.0, média ponderada 0--30
cm)}

\begin{longtable}[]{@{}ll@{}}
\toprule\noalign{}
Parâmetro & Valor \\
\midrule\noalign{}
\endhead
\bottomrule\noalign{}
\endlastfoot
pH (H$_2$O) & 5.3 \\
Carbono orgânico (SOC) & 17.29 g/\allowbreak kg \\
Argila & 25.77 \% \\
Areia & 53.32 \% \\
Silte (calc.) & 20.9 \% \\
Densidade aparente & 1.33 g/\allowbreak cm³ \\
\textbf{Aptidão agrícola} & \textbf{Média} \\
\end{longtable}

Fonte: ISRIC SoilGrids 2.0 via WCS. Coords: -25.25°, -56.5333°. Média
ponderada camadas 0-5, 5-15, 15-30 cm.
\subsubsection{Indicadores Sociais}

\textbf{Fontes:} PNUD 2020, INE\footnote{Instituto Nacional de Estadística (Instituto Nacional de Estatística), órgão oficial de dados demográficos e censitários.} Censo 2022, INE\footnote{Instituto Nacional de Estadística (Instituto Nacional de Estatística), órgão oficial de dados demográficos e censitários.} EPHC 2023, Ministerio
Público 2024\\
\textbf{Nota:} Valores marcados como \emph{(dept.)} referem-se ao
departamento de Caaguazú; valores \emph{(dist.)} são específicos deste
distrito. Dados marcados \emph{(est.)} são estimativas.

\begin{longtable}[]{@{}
  >{\raggedright\arraybackslash}p{(\linewidth - 4\tabcolsep) * \real{0.4231}}
  >{\raggedright\arraybackslash}p{(\linewidth - 4\tabcolsep) * \real{0.2692}}
  >{\raggedright\arraybackslash}p{(\linewidth - 4\tabcolsep) * \real{0.3077}}@{}}
\toprule\noalign{}
\begin{minipage}[b]{\linewidth}\raggedright
Indicador
\end{minipage} & \begin{minipage}[b]{\linewidth}\raggedright
Valor
\end{minipage} & \begin{minipage}[b]{\linewidth}\raggedright
Âmbito
\end{minipage} \\
\midrule\noalign{}
\endhead
\bottomrule\noalign{}
\endlastfoot
IDH (2020) & 0,715 & dept. \\
Ranking IDH nacional & 7/\allowbreak 18 & dept. \\
Esperança de vida & 73,1 anos & dept. \\
Escolaridade média & 8.5 anos & dist. \\
RNB per capita (USD PPA) & 9.000 & dept. \\
Pobreza monetária (\%) & 28,5\% & dept. \\
Pobreza extrema (\%) & 6,8\% & dept. \\
Índice de Gini & 0,447 & dept. \\
Acesso a água potável (\%) & 79,8\% & dept. \\
Acesso a saneamento (\%) & 78,3\% & dept. \\
Taxa de homicídios (est., /\allowbreak 100k hab) & \textasciitilde6--8 (estimativa,
próxima ou acima da média nacional) & dept. \\
Índice de segurança & médio & dept. \\
População (Censo 2022) & 3.294 hab. & dist. \\
Idade mediana & 35.0 anos & dist. \\
Taxa de fecundidade & N/\allowbreak D & dist. \\
\end{longtable}
\subsubsection{Combustível}

\textbf{Referência:} PETROPAR /\allowbreak  postos locais (2024) \textbf{Tipo de
localidade:} Interior

\begin{longtable}[]{@{}lll@{}}
\toprule\noalign{}
Combustível & USD/\allowbreak litro & Gs/\allowbreak litro (aprox.) \\
\midrule\noalign{}
\endhead
\bottomrule\noalign{}
\endlastfoot
Gasolina 93 oct & 0.98 & 7,252 \\
Gasolina 97 oct (premium) & 1.08 & 7,992 \\
Diesel & 0.91 & 6,734 \\
\end{longtable}

\begin{quote}
Preços podem variar ±5\% conforme posto e sazonalidade. Chaco e interior
remoto apresentam maior variação.
\end{quote}

\subsubsection{Cobertura Celular}

\textbf{Fonte:} CONATEL PY /\allowbreak  operadoras (2024)

\begin{longtable}[]{@{}ll@{}}
\toprule\noalign{}
Parâmetro & Valor \\
\midrule\noalign{}
\endhead
\bottomrule\noalign{}
\endlastfoot
Cobertura 4G população (dept.) & 88\% \\
Cobertura 4G área rural & 72\% \\
Melhor operadora & Tigo \\
Qualidade rural & boa \\
\end{longtable}

\begin{quote}
Para áreas rurais fora do núcleo urbano, recomenda-se chip Tigo como
principal e Personal como backup.
\end{quote}

\subsubsection{Internet}

\textbf{Fonte:} CONATEL /\allowbreak  Speedtest Ookla (2024)

\begin{longtable}[]{@{}ll@{}}
\toprule\noalign{}
Parâmetro & Valor \\
\midrule\noalign{}
\endhead
\bottomrule\noalign{}
\endlastfoot
Velocidade média download & 55 Mbps \\
Domicílios com internet (dept.) & 62\% \\
Tecnologia predominante & rádio/\allowbreak fibra \\
Opção rural & Starlink disponível (\textasciitilde USD 44/\allowbreak mês) \\
\end{longtable}

\subsubsection{Mercado Imobiliário e Terra
Rural}

\textbf{Fonte:} INDERT /\allowbreak  Clasificados.com.py (2024)

\begin{longtable}[]{@{}ll@{}}
\toprule\noalign{}
Tipo & Referência \\
\midrule\noalign{}
\endhead
\bottomrule\noalign{}
\endlastfoot
Terra agrícola alta prod. (USD/\allowbreak ha) & 8,800 \\
Imóvel urbano (USD/\allowbreak m²) & 850 \\
Aluguel 2 quartos (USD/\allowbreak mês) & 350 \\
\end{longtable}

\begin{quote}
Valores de referência departamental. Localidades menores podem ter
preços 20--40\% abaixo da capital departamental.
\end{quote}

\subsubsection{Saúde}

\textbf{Fonte:} MSPBS /\allowbreak  IPS Paraguay (2024)

\begin{longtable}[]{@{}ll@{}}
\toprule\noalign{}
Serviço & Disponibilidade \\
\midrule\noalign{}
\endhead
\bottomrule\noalign{}
\endlastfoot
USF /\allowbreak  Posto de Saúde & sim \\
Hospital Regional & não \\
IPS (seguro social) & não \\
Farmácia & sim \\
Distância ao hospital de referência & \textasciitilde80 km (Coronel
Oviedo) \\
\end{longtable}

\textbf{Principais estabelecimentos:} USF local; farmácias distritais.

\textbf{Observação para imigrantes:} Atendimento primário disponível.
Casos de maior complexidade encaminhados a Coronel Oviedo. Cobertura
privada recomendada; IPS não disponível localmente.
\newpage
\secaoDiagnostico{Mariscal Francisco Solano Lopez}{\mapaConteudo{-25.4666}{-56.0166}}{Mariscal Francisco Solano Lopez (Caaguazu) apresenta perfil operacional
consolidado para comparacao territorial, com leitura conjunta de
infraestrutura, risco hidroclimatico e capacidade de continuidade de
servicos essenciais.

\subsubsection{Por que sim}

\begin{itemize}
\tightlist
\item
  Base institucional de referencia disponivel e rastreavel.
\item
  Estrutura comparativa (A-E/\allowbreak GSS) consistente para decisao relativa.
\item
  Plano de mitigacao acionavel ja definido no dossie.
\end{itemize}

\subsubsection{Por que não}

\begin{itemize}
\tightlist
\item
  Serie oficial distrital granular ainda pode apresentar lacunas
  temporais.
\item
  Sensibilidade do score exige monitoramento continuo de risco e
  conectividade.
\item
  Decisao final depende de validacao de campo e janela temporal recente.
\end{itemize}

\subsubsection{Pre-condicoes Minimas de
Decisao}

\begin{itemize}
\tightlist
\item
  Atualizar indicadores criticos em ciclo trimestral.
\item
  Confirmar trafegabilidade e contingencia hidrica/\allowbreak energetica local.
\item
  Validar checklist de resiliencia domiciliar/\allowbreak logistica antes de decisao
  final.
\end{itemize}

\subsubsection{Recomendação
Técnica}

Uso recomendado como candidato em analise multicriterio, condicionado ao
cumprimento das pre-condicoes acima. Referencia atual de score: GSS 7.7
em 2026-03-05; risco predominante: baixo.

\subsubsection{}

\begin{itemize}
\tightlist
\item
  \begin{enumerate}
  \def\labelenumi{\arabic{enumi})}
  \tightlist
  \item
    Delimitacao territorial validada por georreferencia institucional
    (Fonte: acesso: 2026-03-05).
  \end{enumerate}
\item
  \begin{enumerate}
  \def\labelenumi{\arabic{enumi})}
  \setcounter{enumi}{1}
  \tightlist
  \item
    População municipal/\allowbreak distrital referenciada por base censitaria
    nacional (Fonte: acesso: 2026-03-05).
  \end{enumerate}
\item
  \begin{enumerate}
  \def\labelenumi{\arabic{enumi})}
  \setcounter{enumi}{2}
  \tightlist
  \item
    Comparabilidade intra-departamental apoiada em indicadores
    distritais oficiais (Fonte: acesso: 2026-03-05).
  \end{enumerate}
\item
  \begin{enumerate}
  \def\labelenumi{\arabic{enumi})}
  \setcounter{enumi}{3}
  \tightlist
  \item
    Estado macro de conectividade viaria ancorado em informacoes de
    rodovias (Fonte: acesso: 2026-03-05).
  \end{enumerate}
\item
  \begin{enumerate}
  \def\labelenumi{\arabic{enumi})}
  \setcounter{enumi}{4}
  \tightlist
  \item
    Exposicao hidrometeorologica monitorada por avisos oficiais (Fonte:
    acesso: 2026-03-05).
  \end{enumerate}
\item
  \begin{enumerate}
  \def\labelenumi{\arabic{enumi})}
  \setcounter{enumi}{5}
  \tightlist
  \item
    Historico de contexto meteorologico apoiado em publicacoes tecnicas
    (Fonte: acesso: 2026-03-05).
  \end{enumerate}
\item
  \begin{enumerate}
  \def\labelenumi{\arabic{enumi})}
  \setcounter{enumi}{6}
  \tightlist
  \item
    Capacidade institucional de resposta a eventos apoiada em acoes da
    SEN\footnote{Secretaría de Emergencia Nacional (Secretaria de Emergência Nacional), órgão governamental de resposta a desastres.} (Fonte: acesso: 2026-03-05).
  \end{enumerate}
\item
  \begin{enumerate}
  \def\labelenumi{\arabic{enumi})}
  \setcounter{enumi}{7}
  \tightlist
  \item
    Infraestrutura energetica analisada via operador nacional (Fonte:
    acesso: 2026-03-05).
  \end{enumerate}
\item
  \begin{enumerate}
  \def\labelenumi{\arabic{enumi})}
  \setcounter{enumi}{8}
  \tightlist
  \item
    Referencias de obras e logistica nacional verificadas no MOPC
    (Fonte: acesso: 2026-03-05).
  \end{enumerate}
\item
  \begin{enumerate}
  \def\labelenumi{\arabic{enumi})}
  \setcounter{enumi}{9}
  \tightlist
  \item
    Contexto sociopolitico institucional referenciado por autoridade
    eleitoral (Fonte: acesso: 2026-03-05).
  \end{enumerate}
\item
  \begin{enumerate}
  \def\labelenumi{\arabic{enumi})}
  \setcounter{enumi}{10}
  \tightlist
  \item
    Camada cartografica de apoio eleitoral/\allowbreak territorial consultada
    (Fonte: acesso: 2026-03-05).
  \end{enumerate}
\item
  \begin{enumerate}
  \def\labelenumi{\arabic{enumi})}
  \setcounter{enumi}{11}
  \tightlist
  \item
    Dados publicos complementares para verificacao cruzada (Fonte:
    acesso: 2026-03-05).
  \end{enumerate}
\end{itemize}}
\begin{table}[ht!]
\small \begin{tabular}{p{3.0cm}p{0.8cm}p{6.0cm}} \toprule
\textbf{Critério} & \textbf{Nota} & \textbf{Justificativa} \\ \midrule
A - Ameaças Estratégicas & 8.7 & - \\
B - Risco Social & 6.9 & - \\
C - Riscos Naturais & 8.2 & - \\
D - Autossuficiência & 7.0 & - \\
E - Institucional & 7.6 & - \\
\midrule \textbf{GSS FINAL} & \textbf{7.7} & \textbf{Seguro (bom potencial com preparacao).} \\ \bottomrule \end{tabular}
\end{table}
\subsubsection{Vulnerabilidades e
Mitigação}\label{vuln-Mariscal_Francisco_Solano_Lopez-vulnerabilidades-e-mitigauxe7uxe3o}

\begin{itemize}
\tightlist
\item
  Exposicao hidrometeorologica controlada (\textless45/\allowbreak 100).
\item
  Conectividade viaria favoravel (\textgreater=70/\allowbreak 100).
\item
  Estoque de contingencia para 14 dias (agua/\allowbreak alimentos/\allowbreak combustivel),
  priorizado quando risco hidro=23/\allowbreak 100.
\item
  Duplicar rotas/\allowbreak logistica quando conectividade viaria=80/\allowbreak 100 for
  \textless70; manter rota secundaria mapeada e testada trimestralmente.
\item
  Plano de energia resiliente com autonomia minima de 72h
  (geracao/\allowbreak backup), revisao semestral.
\end{itemize}
\subsubsection{Dossiê de Campo}\label{dossie-Mariscal_Francisco_Solano_Lopez-dossiuxea-de-campo}
\begin{description}[style=nextline,leftmargin=0.5cm,font=\bfseries]
\item[Coordenadas]  25°28'S 56°01'W (geocodificado via Nominatim).
\end{description}
\textbf{Fonte climática:} NASA POWER Climatology API (período 2001-2020)
\textbf{Fonte luminosa:} estimativa world atlas

\paragraph{Irradiação Solar
(kWh/\allowbreak m²/\allowbreak dia)}\label{irradiauxe7uxe3o-solar-kwhmuxb2dia}

\begin{longtable}[]{@{}
  >{\raggedright\arraybackslash}p{(\linewidth - 24\tabcolsep) * \real{0.0746}}
  >{\raggedright\arraybackslash}p{(\linewidth - 24\tabcolsep) * \real{0.0746}}
  >{\raggedright\arraybackslash}p{(\linewidth - 24\tabcolsep) * \real{0.0746}}
  >{\raggedright\arraybackslash}p{(\linewidth - 24\tabcolsep) * \real{0.0746}}
  >{\raggedright\arraybackslash}p{(\linewidth - 24\tabcolsep) * \real{0.0746}}
  >{\raggedright\arraybackslash}p{(\linewidth - 24\tabcolsep) * \real{0.0746}}
  >{\raggedright\arraybackslash}p{(\linewidth - 24\tabcolsep) * \real{0.0746}}
  >{\raggedright\arraybackslash}p{(\linewidth - 24\tabcolsep) * \real{0.0746}}
  >{\raggedright\arraybackslash}p{(\linewidth - 24\tabcolsep) * \real{0.0746}}
  >{\raggedright\arraybackslash}p{(\linewidth - 24\tabcolsep) * \real{0.0746}}
  >{\raggedright\arraybackslash}p{(\linewidth - 24\tabcolsep) * \real{0.0746}}
  >{\raggedright\arraybackslash}p{(\linewidth - 24\tabcolsep) * \real{0.0746}}
  >{\raggedright\arraybackslash}p{(\linewidth - 24\tabcolsep) * \real{0.1045}}@{}}
\toprule\noalign{}
\begin{minipage}[b]{\linewidth}\raggedright
Jan
\end{minipage} & \begin{minipage}[b]{\linewidth}\raggedright
Fev
\end{minipage} & \begin{minipage}[b]{\linewidth}\raggedright
Mar
\end{minipage} & \begin{minipage}[b]{\linewidth}\raggedright
Abr
\end{minipage} & \begin{minipage}[b]{\linewidth}\raggedright
Mai
\end{minipage} & \begin{minipage}[b]{\linewidth}\raggedright
Jun
\end{minipage} & \begin{minipage}[b]{\linewidth}\raggedright
Jul
\end{minipage} & \begin{minipage}[b]{\linewidth}\raggedright
Ago
\end{minipage} & \begin{minipage}[b]{\linewidth}\raggedright
Set
\end{minipage} & \begin{minipage}[b]{\linewidth}\raggedright
Out
\end{minipage} & \begin{minipage}[b]{\linewidth}\raggedright
Nov
\end{minipage} & \begin{minipage}[b]{\linewidth}\raggedright
Dez
\end{minipage} & \begin{minipage}[b]{\linewidth}\raggedright
Média
\end{minipage} \\
\midrule\noalign{}
\endhead
\bottomrule\noalign{}
\endlastfoot
6.74 & 6.20 & 5.50 & 4.48 & 3.35 & 2.88 & 3.21 & 3.96 & 4.59 & 5.43 &
6.43 & 6.81 & \textbf{4.96} \\
\end{longtable}

\textbf{Inclinação solar recomendada:} 25° N (anual) · 35° N (inverno
jun-ago) · 15° N (verão nov-jan)

\paragraph{Precipitação (mm/\allowbreak mês)}\label{precipitauxe7uxe3o-mmmuxeas}

\begin{longtable}[]{@{}
  >{\raggedright\arraybackslash}p{(\linewidth - 24\tabcolsep) * \real{0.0704}}
  >{\raggedright\arraybackslash}p{(\linewidth - 24\tabcolsep) * \real{0.0704}}
  >{\raggedright\arraybackslash}p{(\linewidth - 24\tabcolsep) * \real{0.0704}}
  >{\raggedright\arraybackslash}p{(\linewidth - 24\tabcolsep) * \real{0.0704}}
  >{\raggedright\arraybackslash}p{(\linewidth - 24\tabcolsep) * \real{0.0704}}
  >{\raggedright\arraybackslash}p{(\linewidth - 24\tabcolsep) * \real{0.0704}}
  >{\raggedright\arraybackslash}p{(\linewidth - 24\tabcolsep) * \real{0.0704}}
  >{\raggedright\arraybackslash}p{(\linewidth - 24\tabcolsep) * \real{0.0704}}
  >{\raggedright\arraybackslash}p{(\linewidth - 24\tabcolsep) * \real{0.0704}}
  >{\raggedright\arraybackslash}p{(\linewidth - 24\tabcolsep) * \real{0.0704}}
  >{\raggedright\arraybackslash}p{(\linewidth - 24\tabcolsep) * \real{0.0704}}
  >{\raggedright\arraybackslash}p{(\linewidth - 24\tabcolsep) * \real{0.0704}}
  >{\raggedright\arraybackslash}p{(\linewidth - 24\tabcolsep) * \real{0.1549}}@{}}
\toprule\noalign{}
\begin{minipage}[b]{\linewidth}\raggedright
Jan
\end{minipage} & \begin{minipage}[b]{\linewidth}\raggedright
Fev
\end{minipage} & \begin{minipage}[b]{\linewidth}\raggedright
Mar
\end{minipage} & \begin{minipage}[b]{\linewidth}\raggedright
Abr
\end{minipage} & \begin{minipage}[b]{\linewidth}\raggedright
Mai
\end{minipage} & \begin{minipage}[b]{\linewidth}\raggedright
Jun
\end{minipage} & \begin{minipage}[b]{\linewidth}\raggedright
Jul
\end{minipage} & \begin{minipage}[b]{\linewidth}\raggedright
Ago
\end{minipage} & \begin{minipage}[b]{\linewidth}\raggedright
Set
\end{minipage} & \begin{minipage}[b]{\linewidth}\raggedright
Out
\end{minipage} & \begin{minipage}[b]{\linewidth}\raggedright
Nov
\end{minipage} & \begin{minipage}[b]{\linewidth}\raggedright
Dez
\end{minipage} & \begin{minipage}[b]{\linewidth}\raggedright
Total/\allowbreak ano
\end{minipage} \\
\midrule\noalign{}
\endhead
\bottomrule\noalign{}
\endlastfoot
131 & 136 & 131 & 152 & 172 & 90 & 75 & 52 & 92 & 187 & 190 & 166 &
\textbf{1577 mm} \\
\end{longtable}

\paragraph{Poluição Luminosa}\label{poluiuxe7uxe3o-luminosa}

\begin{longtable}[]{@{}ll@{}}
\toprule\noalign{}
Parâmetro & Valor \\
\midrule\noalign{}
\endhead
\bottomrule\noalign{}
\endlastfoot
Escala Bortle & 4 --- Céu rural-suburbano \\
Radiância artificial & 3.0 nW/\allowbreak cm²/\allowbreak sr \\
\end{longtable}
\paragraph{Solo (SoilGrids 2.0, média ponderada 0--30
cm)}

\begin{longtable}[]{@{}ll@{}}
\toprule\noalign{}
Parâmetro & Valor \\
\midrule\noalign{}
\endhead
\bottomrule\noalign{}
\endlastfoot
pH (H$_2$O) & 5.23 \\
Carbono orgânico (SOC) & 17.38 g/\allowbreak kg \\
Argila & 27.91 \% \\
Areia & 49.93 \% \\
Silte (calc.) & 22.2 \% \\
Densidade aparente & 1.3 g/\allowbreak cm³ \\
\textbf{Aptidão agrícola} & \textbf{Média} \\
\end{longtable}

Fonte: ISRIC SoilGrids 2.0 via WCS. Coords: -25.4667°, -56.0167°. Média
ponderada camadas 0-5, 5-15, 15-30 cm.
\subsubsection{Indicadores Sociais}

\textbf{Fontes:} PNUD 2020, INE\footnote{Instituto Nacional de Estadística (Instituto Nacional de Estatística), órgão oficial de dados demográficos e censitários.} Censo 2022, INE\footnote{Instituto Nacional de Estadística (Instituto Nacional de Estatística), órgão oficial de dados demográficos e censitários.} EPHC 2023, Ministerio
Público 2024\\
\textbf{Nota:} Valores marcados como \emph{(dept.)} referem-se ao
departamento de Caaguazú; valores \emph{(dist.)} são específicos deste
distrito. Dados marcados \emph{(est.)} são estimativas.

\begin{longtable}[]{@{}
  >{\raggedright\arraybackslash}p{(\linewidth - 4\tabcolsep) * \real{0.4231}}
  >{\raggedright\arraybackslash}p{(\linewidth - 4\tabcolsep) * \real{0.2692}}
  >{\raggedright\arraybackslash}p{(\linewidth - 4\tabcolsep) * \real{0.3077}}@{}}
\toprule\noalign{}
\begin{minipage}[b]{\linewidth}\raggedright
Indicador
\end{minipage} & \begin{minipage}[b]{\linewidth}\raggedright
Valor
\end{minipage} & \begin{minipage}[b]{\linewidth}\raggedright
Âmbito
\end{minipage} \\
\midrule\noalign{}
\endhead
\bottomrule\noalign{}
\endlastfoot
IDH (2020) & 0,715 & dept. \\
Ranking IDH nacional & 7/\allowbreak 18 & dept. \\
Esperança de vida & 73,1 anos & dept. \\
Escolaridade média & 6.9 anos & dist. \\
RNB per capita (USD PPA) & 9.000 & dept. \\
Pobreza monetária (\%) & 28,5\% & dept. \\
Pobreza extrema (\%) & 6,8\% & dept. \\
Índice de Gini & 0,447 & dept. \\
Acesso a água potável (\%) & 79,8\% & dept. \\
Acesso a saneamento (\%) & 78,3\% & dept. \\
Taxa de homicídios (est., /\allowbreak 100k hab) & \textasciitilde6--8 (estimativa,
próxima ou acima da média nacional) & dept. \\
Índice de segurança & médio & dept. \\
População (Censo 2022) & 5.338 hab. & dist. \\
Idade mediana & 25.0 anos & dist. \\
Taxa de fecundidade & N/\allowbreak D & dist. \\
\end{longtable}
\subsubsection{Combustível}

\textbf{Referência:} PETROPAR /\allowbreak  postos locais (2024) \textbf{Tipo de
localidade:} Interior

\begin{longtable}[]{@{}lll@{}}
\toprule\noalign{}
Combustível & USD/\allowbreak litro & Gs/\allowbreak litro (aprox.) \\
\midrule\noalign{}
\endhead
\bottomrule\noalign{}
\endlastfoot
Gasolina 93 oct & 0.98 & 7,252 \\
Gasolina 97 oct (premium) & 1.08 & 7,992 \\
Diesel & 0.91 & 6,734 \\
\end{longtable}

\begin{quote}
Preços podem variar ±5\% conforme posto e sazonalidade. Chaco e interior
remoto apresentam maior variação.
\end{quote}

\subsubsection{Cobertura Celular}

\textbf{Fonte:} CONATEL PY /\allowbreak  operadoras (2024)

\begin{longtable}[]{@{}ll@{}}
\toprule\noalign{}
Parâmetro & Valor \\
\midrule\noalign{}
\endhead
\bottomrule\noalign{}
\endlastfoot
Cobertura 4G população (dept.) & 88\% \\
Cobertura 4G área rural & 72\% \\
Melhor operadora & Tigo \\
Qualidade rural & boa \\
\end{longtable}

\begin{quote}
Para áreas rurais fora do núcleo urbano, recomenda-se chip Tigo como
principal e Personal como backup.
\end{quote}

\subsubsection{Internet}

\textbf{Fonte:} CONATEL /\allowbreak  Speedtest Ookla (2024)

\begin{longtable}[]{@{}ll@{}}
\toprule\noalign{}
Parâmetro & Valor \\
\midrule\noalign{}
\endhead
\bottomrule\noalign{}
\endlastfoot
Velocidade média download & 55 Mbps \\
Domicílios com internet (dept.) & 62\% \\
Tecnologia predominante & rádio/\allowbreak fibra \\
Opção rural & Starlink disponível (\textasciitilde USD 44/\allowbreak mês) \\
\end{longtable}

\subsubsection{Mercado Imobiliário e Terra
Rural}

\textbf{Fonte:} INDERT /\allowbreak  Clasificados.com.py (2024)

\begin{longtable}[]{@{}ll@{}}
\toprule\noalign{}
Tipo & Referência \\
\midrule\noalign{}
\endhead
\bottomrule\noalign{}
\endlastfoot
Terra agrícola alta prod. (USD/\allowbreak ha) & 8,800 \\
Imóvel urbano (USD/\allowbreak m²) & 850 \\
Aluguel 2 quartos (USD/\allowbreak mês) & 350 \\
\end{longtable}

\begin{quote}
Valores de referência departamental. Localidades menores podem ter
preços 20--40\% abaixo da capital departamental.
\end{quote}

\subsubsection{Saúde}

\textbf{Fonte:} MSPBS /\allowbreak  IPS Paraguay (2024)

\begin{longtable}[]{@{}ll@{}}
\toprule\noalign{}
Serviço & Disponibilidade \\
\midrule\noalign{}
\endhead
\bottomrule\noalign{}
\endlastfoot
USF /\allowbreak  Posto de Saúde & sim \\
Hospital Regional & não \\
IPS (seguro social) & não \\
Farmácia & sim \\
Distância ao hospital de referência & \textasciitilde100 km (Coronel
Oviedo) \\
\end{longtable}

\textbf{Principais estabelecimentos:} USF local; farmácias distritais.

\textbf{Observação para imigrantes:} Atendimento primário disponível.
Casos de maior complexidade encaminhados a Coronel Oviedo. Cobertura
privada recomendada; IPS não disponível localmente.
\newpage
\secaoDiagnostico{Nueva Londres}{\mapaConteudo{-25.4000}{-56.5500}}{Nueva Londres (Caaguazu) apresenta perfil operacional consolidado para
comparacao territorial, com leitura conjunta de infraestrutura, risco
hidroclimatico e capacidade de continuidade de servicos essenciais.

\subsubsection{Por que sim}

\begin{itemize}
\tightlist
\item
  Base institucional de referencia disponivel e rastreavel.
\item
  Estrutura comparativa (A-E/\allowbreak GSS) consistente para decisao relativa.
\item
  Plano de mitigacao acionavel ja definido no dossie.
\end{itemize}

\subsubsection{Por que não}

\begin{itemize}
\tightlist
\item
  Serie oficial distrital granular ainda pode apresentar lacunas
  temporais.
\item
  Sensibilidade do score exige monitoramento continuo de risco e
  conectividade.
\item
  Decisao final depende de validacao de campo e janela temporal recente.
\end{itemize}

\subsubsection{Pre-condicoes Minimas de
Decisao}

\begin{itemize}
\tightlist
\item
  Atualizar indicadores criticos em ciclo trimestral.
\item
  Confirmar trafegabilidade e contingencia hidrica/\allowbreak energetica local.
\item
  Validar checklist de resiliencia domiciliar/\allowbreak logistica antes de decisao
  final.
\end{itemize}

\subsubsection{Recomendação
Técnica}

Uso recomendado como candidato em analise multicriterio, condicionado ao
cumprimento das pre-condicoes acima. Referencia atual de score: GSS 7.2
em 2026-03-05; risco predominante: baixo.

\subsubsection{}

\begin{itemize}
\tightlist
\item
  \begin{enumerate}
  \def\labelenumi{\arabic{enumi})}
  \tightlist
  \item
    Delimitacao territorial validada por georreferencia institucional
    (Fonte: acesso: 2026-03-05).
  \end{enumerate}
\item
  \begin{enumerate}
  \def\labelenumi{\arabic{enumi})}
  \setcounter{enumi}{1}
  \tightlist
  \item
    População municipal/\allowbreak distrital referenciada por base censitaria
    nacional (Fonte: acesso: 2026-03-05).
  \end{enumerate}
\item
  \begin{enumerate}
  \def\labelenumi{\arabic{enumi})}
  \setcounter{enumi}{2}
  \tightlist
  \item
    Comparabilidade intra-departamental apoiada em indicadores
    distritais oficiais (Fonte: acesso: 2026-03-05).
  \end{enumerate}
\item
  \begin{enumerate}
  \def\labelenumi{\arabic{enumi})}
  \setcounter{enumi}{3}
  \tightlist
  \item
    Estado macro de conectividade viaria ancorado em informacoes de
    rodovias (Fonte: acesso: 2026-03-05).
  \end{enumerate}
\item
  \begin{enumerate}
  \def\labelenumi{\arabic{enumi})}
  \setcounter{enumi}{4}
  \tightlist
  \item
    Exposicao hidrometeorologica monitorada por avisos oficiais (Fonte:
    acesso: 2026-03-05).
  \end{enumerate}
\item
  \begin{enumerate}
  \def\labelenumi{\arabic{enumi})}
  \setcounter{enumi}{5}
  \tightlist
  \item
    Historico de contexto meteorologico apoiado em publicacoes tecnicas
    (Fonte: acesso: 2026-03-05).
  \end{enumerate}
\item
  \begin{enumerate}
  \def\labelenumi{\arabic{enumi})}
  \setcounter{enumi}{6}
  \tightlist
  \item
    Capacidade institucional de resposta a eventos apoiada em acoes da
    SEN\footnote{Secretaría de Emergencia Nacional (Secretaria de Emergência Nacional), órgão governamental de resposta a desastres.} (Fonte: acesso: 2026-03-05).
  \end{enumerate}
\item
  \begin{enumerate}
  \def\labelenumi{\arabic{enumi})}
  \setcounter{enumi}{7}
  \tightlist
  \item
    Infraestrutura energetica analisada via operador nacional (Fonte:
    acesso: 2026-03-05).
  \end{enumerate}
\item
  \begin{enumerate}
  \def\labelenumi{\arabic{enumi})}
  \setcounter{enumi}{8}
  \tightlist
  \item
    Referencias de obras e logistica nacional verificadas no MOPC
    (Fonte: acesso: 2026-03-05).
  \end{enumerate}
\item
  \begin{enumerate}
  \def\labelenumi{\arabic{enumi})}
  \setcounter{enumi}{9}
  \tightlist
  \item
    Contexto sociopolitico institucional referenciado por autoridade
    eleitoral (Fonte: acesso: 2026-03-05).
  \end{enumerate}
\item
  \begin{enumerate}
  \def\labelenumi{\arabic{enumi})}
  \setcounter{enumi}{10}
  \tightlist
  \item
    Camada cartografica de apoio eleitoral/\allowbreak territorial consultada
    (Fonte: acesso: 2026-03-05).
  \end{enumerate}
\item
  \begin{enumerate}
  \def\labelenumi{\arabic{enumi})}
  \setcounter{enumi}{11}
  \tightlist
  \item
    Dados publicos complementares para verificacao cruzada (Fonte:
    acesso: 2026-03-05).
  \end{enumerate}
\end{itemize}}
\begin{table}[ht!]
\small \begin{tabular}{p{3.0cm}p{0.8cm}p{6.0cm}} \toprule
\textbf{Critério} & \textbf{Nota} & \textbf{Justificativa} \\ \midrule
A - Ameaças Estratégicas & 7.8 & - \\
B - Risco Social & 6.6 & - \\
C - Riscos Naturais & 7.8 & - \\
D - Autossuficiência & 6.2 & - \\
E - Institucional & 7.5 & - \\
\midrule \textbf{GSS FINAL} & \textbf{7.2} & \textbf{Seguro (bom potencial com preparacao).} \\ \bottomrule \end{tabular}
\end{table}
\subsubsection{Vulnerabilidades e
Mitigação}\label{vuln-Nueva_Londres-vulnerabilidades-e-mitigauxe7uxe3o}

\begin{itemize}
\tightlist
\item
  Exposicao hidrometeorologica controlada (\textless45/\allowbreak 100).
\item
  Conectividade viaria intermediaria (50-69/\allowbreak 100).
\item
  Estoque de contingencia para 14 dias (agua/\allowbreak alimentos/\allowbreak combustivel),
  priorizado quando risco hidro=28/\allowbreak 100.
\item
  Duplicar rotas/\allowbreak logistica quando conectividade viaria=54/\allowbreak 100 for
  \textless70; manter rota secundaria mapeada e testada trimestralmente.
\item
  Plano de energia resiliente com autonomia minima de 72h
  (geracao/\allowbreak backup), revisao semestral.
\end{itemize}
\subsubsection{Dossiê de Campo}\label{dossie-Nueva_Londres-dossiuxea-de-campo}
\begin{description}[style=nextline,leftmargin=0.5cm,font=\bfseries]
\item[Coordenadas]  25°24'S 56°33'W (geocodificado via Nominatim).
\end{description}
\textbf{Fonte climática:} NASA POWER Climatology API (período 2001-2020)
\textbf{Fonte luminosa:} estimativa world atlas

\paragraph{Irradiação Solar
(kWh/\allowbreak m²/\allowbreak dia)}\label{irradiauxe7uxe3o-solar-kwhmuxb2dia}

\begin{longtable}[]{@{}
  >{\raggedright\arraybackslash}p{(\linewidth - 24\tabcolsep) * \real{0.0746}}
  >{\raggedright\arraybackslash}p{(\linewidth - 24\tabcolsep) * \real{0.0746}}
  >{\raggedright\arraybackslash}p{(\linewidth - 24\tabcolsep) * \real{0.0746}}
  >{\raggedright\arraybackslash}p{(\linewidth - 24\tabcolsep) * \real{0.0746}}
  >{\raggedright\arraybackslash}p{(\linewidth - 24\tabcolsep) * \real{0.0746}}
  >{\raggedright\arraybackslash}p{(\linewidth - 24\tabcolsep) * \real{0.0746}}
  >{\raggedright\arraybackslash}p{(\linewidth - 24\tabcolsep) * \real{0.0746}}
  >{\raggedright\arraybackslash}p{(\linewidth - 24\tabcolsep) * \real{0.0746}}
  >{\raggedright\arraybackslash}p{(\linewidth - 24\tabcolsep) * \real{0.0746}}
  >{\raggedright\arraybackslash}p{(\linewidth - 24\tabcolsep) * \real{0.0746}}
  >{\raggedright\arraybackslash}p{(\linewidth - 24\tabcolsep) * \real{0.0746}}
  >{\raggedright\arraybackslash}p{(\linewidth - 24\tabcolsep) * \real{0.0746}}
  >{\raggedright\arraybackslash}p{(\linewidth - 24\tabcolsep) * \real{0.1045}}@{}}
\toprule\noalign{}
\begin{minipage}[b]{\linewidth}\raggedright
Jan
\end{minipage} & \begin{minipage}[b]{\linewidth}\raggedright
Fev
\end{minipage} & \begin{minipage}[b]{\linewidth}\raggedright
Mar
\end{minipage} & \begin{minipage}[b]{\linewidth}\raggedright
Abr
\end{minipage} & \begin{minipage}[b]{\linewidth}\raggedright
Mai
\end{minipage} & \begin{minipage}[b]{\linewidth}\raggedright
Jun
\end{minipage} & \begin{minipage}[b]{\linewidth}\raggedright
Jul
\end{minipage} & \begin{minipage}[b]{\linewidth}\raggedright
Ago
\end{minipage} & \begin{minipage}[b]{\linewidth}\raggedright
Set
\end{minipage} & \begin{minipage}[b]{\linewidth}\raggedright
Out
\end{minipage} & \begin{minipage}[b]{\linewidth}\raggedright
Nov
\end{minipage} & \begin{minipage}[b]{\linewidth}\raggedright
Dez
\end{minipage} & \begin{minipage}[b]{\linewidth}\raggedright
Média
\end{minipage} \\
\midrule\noalign{}
\endhead
\bottomrule\noalign{}
\endlastfoot
6.74 & 6.20 & 5.50 & 4.48 & 3.35 & 2.88 & 3.21 & 3.96 & 4.59 & 5.43 &
6.43 & 6.81 & \textbf{4.96} \\
\end{longtable}

\textbf{Inclinação solar recomendada:} 25° N (anual) · 35° N (inverno
jun-ago) · 15° N (verão nov-jan)

\paragraph{Precipitação (mm/\allowbreak mês)}\label{precipitauxe7uxe3o-mmmuxeas}

\begin{longtable}[]{@{}
  >{\raggedright\arraybackslash}p{(\linewidth - 24\tabcolsep) * \real{0.0704}}
  >{\raggedright\arraybackslash}p{(\linewidth - 24\tabcolsep) * \real{0.0704}}
  >{\raggedright\arraybackslash}p{(\linewidth - 24\tabcolsep) * \real{0.0704}}
  >{\raggedright\arraybackslash}p{(\linewidth - 24\tabcolsep) * \real{0.0704}}
  >{\raggedright\arraybackslash}p{(\linewidth - 24\tabcolsep) * \real{0.0704}}
  >{\raggedright\arraybackslash}p{(\linewidth - 24\tabcolsep) * \real{0.0704}}
  >{\raggedright\arraybackslash}p{(\linewidth - 24\tabcolsep) * \real{0.0704}}
  >{\raggedright\arraybackslash}p{(\linewidth - 24\tabcolsep) * \real{0.0704}}
  >{\raggedright\arraybackslash}p{(\linewidth - 24\tabcolsep) * \real{0.0704}}
  >{\raggedright\arraybackslash}p{(\linewidth - 24\tabcolsep) * \real{0.0704}}
  >{\raggedright\arraybackslash}p{(\linewidth - 24\tabcolsep) * \real{0.0704}}
  >{\raggedright\arraybackslash}p{(\linewidth - 24\tabcolsep) * \real{0.0704}}
  >{\raggedright\arraybackslash}p{(\linewidth - 24\tabcolsep) * \real{0.1549}}@{}}
\toprule\noalign{}
\begin{minipage}[b]{\linewidth}\raggedright
Jan
\end{minipage} & \begin{minipage}[b]{\linewidth}\raggedright
Fev
\end{minipage} & \begin{minipage}[b]{\linewidth}\raggedright
Mar
\end{minipage} & \begin{minipage}[b]{\linewidth}\raggedright
Abr
\end{minipage} & \begin{minipage}[b]{\linewidth}\raggedright
Mai
\end{minipage} & \begin{minipage}[b]{\linewidth}\raggedright
Jun
\end{minipage} & \begin{minipage}[b]{\linewidth}\raggedright
Jul
\end{minipage} & \begin{minipage}[b]{\linewidth}\raggedright
Ago
\end{minipage} & \begin{minipage}[b]{\linewidth}\raggedright
Set
\end{minipage} & \begin{minipage}[b]{\linewidth}\raggedright
Out
\end{minipage} & \begin{minipage}[b]{\linewidth}\raggedright
Nov
\end{minipage} & \begin{minipage}[b]{\linewidth}\raggedright
Dez
\end{minipage} & \begin{minipage}[b]{\linewidth}\raggedright
Total/\allowbreak ano
\end{minipage} \\
\midrule\noalign{}
\endhead
\bottomrule\noalign{}
\endlastfoot
131 & 136 & 131 & 152 & 172 & 90 & 75 & 52 & 92 & 187 & 190 & 166 &
\textbf{1577 mm} \\
\end{longtable}

\paragraph{Poluição Luminosa}\label{poluiuxe7uxe3o-luminosa}

\begin{longtable}[]{@{}ll@{}}
\toprule\noalign{}
Parâmetro & Valor \\
\midrule\noalign{}
\endhead
\bottomrule\noalign{}
\endlastfoot
Escala Bortle & 4 --- Céu rural-suburbano \\
Radiância artificial & 3.0 nW/\allowbreak cm²/\allowbreak sr \\
\end{longtable}
\paragraph{Solo (SoilGrids 2.0, média ponderada 0--30
cm)}

\begin{longtable}[]{@{}ll@{}}
\toprule\noalign{}
Parâmetro & Valor \\
\midrule\noalign{}
\endhead
\bottomrule\noalign{}
\endlastfoot
pH (H$_2$O) & 5.4 \\
Carbono orgânico (SOC) & 17.2 g/\allowbreak kg \\
Argila & 26.56 \% \\
Areia & 53.42 \% \\
Silte (calc.) & 20.0 \% \\
Densidade aparente & 1.35 g/\allowbreak cm³ \\
\textbf{Aptidão agrícola} & \textbf{Média} \\
\end{longtable}

Fonte: ISRIC SoilGrids 2.0 via WCS. Coords: -25.4°, -56.55°. Média
ponderada camadas 0-5, 5-15, 15-30 cm.
\subsubsection{Indicadores Sociais}

\textbf{Fontes:} PNUD 2020, INE\footnote{Instituto Nacional de Estadística (Instituto Nacional de Estatística), órgão oficial de dados demográficos e censitários.} Censo 2022, INE\footnote{Instituto Nacional de Estadística (Instituto Nacional de Estatística), órgão oficial de dados demográficos e censitários.} EPHC 2023, Ministerio
Público 2024\\
\textbf{Nota:} Valores marcados como \emph{(dept.)} referem-se ao
departamento de Caaguazú; valores \emph{(dist.)} são específicos deste
distrito. Dados marcados \emph{(est.)} são estimativas.

\begin{longtable}[]{@{}
  >{\raggedright\arraybackslash}p{(\linewidth - 4\tabcolsep) * \real{0.4231}}
  >{\raggedright\arraybackslash}p{(\linewidth - 4\tabcolsep) * \real{0.2692}}
  >{\raggedright\arraybackslash}p{(\linewidth - 4\tabcolsep) * \real{0.3077}}@{}}
\toprule\noalign{}
\begin{minipage}[b]{\linewidth}\raggedright
Indicador
\end{minipage} & \begin{minipage}[b]{\linewidth}\raggedright
Valor
\end{minipage} & \begin{minipage}[b]{\linewidth}\raggedright
Âmbito
\end{minipage} \\
\midrule\noalign{}
\endhead
\bottomrule\noalign{}
\endlastfoot
IDH (2020) & 0,715 & dept. \\
Ranking IDH nacional & 7/\allowbreak 18 & dept. \\
Esperança de vida & 73,1 anos & dept. \\
Escolaridade média & 9.1 anos & dist. \\
RNB per capita (USD PPA) & 9.000 & dept. \\
Pobreza monetária (\%) & 28,5\% & dept. \\
Pobreza extrema (\%) & 6,8\% & dept. \\
Índice de Gini & 0,447 & dept. \\
Acesso a água potável (\%) & 79,8\% & dept. \\
Acesso a saneamento (\%) & 78,3\% & dept. \\
Taxa de homicídios (est., /\allowbreak 100k hab) & \textasciitilde6--8 (estimativa,
próxima ou acima da média nacional) & dept. \\
Índice de segurança & médio & dept. \\
População (Censo 2022) & 3.594 hab. & dist. \\
Idade mediana & 34.0 anos & dist. \\
Taxa de fecundidade & N/\allowbreak D & dist. \\
\end{longtable}
\subsubsection{Combustível}

\textbf{Referência:} PETROPAR /\allowbreak  postos locais (2024) \textbf{Tipo de
localidade:} Interior

\begin{longtable}[]{@{}lll@{}}
\toprule\noalign{}
Combustível & USD/\allowbreak litro & Gs/\allowbreak litro (aprox.) \\
\midrule\noalign{}
\endhead
\bottomrule\noalign{}
\endlastfoot
Gasolina 93 oct & 0.98 & 7,252 \\
Gasolina 97 oct (premium) & 1.08 & 7,992 \\
Diesel & 0.91 & 6,734 \\
\end{longtable}

\begin{quote}
Preços podem variar ±5\% conforme posto e sazonalidade. Chaco e interior
remoto apresentam maior variação.
\end{quote}

\subsubsection{Cobertura Celular}

\textbf{Fonte:} CONATEL PY /\allowbreak  operadoras (2024)

\begin{longtable}[]{@{}ll@{}}
\toprule\noalign{}
Parâmetro & Valor \\
\midrule\noalign{}
\endhead
\bottomrule\noalign{}
\endlastfoot
Cobertura 4G população (dept.) & 88\% \\
Cobertura 4G área rural & 72\% \\
Melhor operadora & Tigo \\
Qualidade rural & boa \\
\end{longtable}

\begin{quote}
Para áreas rurais fora do núcleo urbano, recomenda-se chip Tigo como
principal e Personal como backup.
\end{quote}

\subsubsection{Internet}

\textbf{Fonte:} CONATEL /\allowbreak  Speedtest Ookla (2024)

\begin{longtable}[]{@{}ll@{}}
\toprule\noalign{}
Parâmetro & Valor \\
\midrule\noalign{}
\endhead
\bottomrule\noalign{}
\endlastfoot
Velocidade média download & 55 Mbps \\
Domicílios com internet (dept.) & 62\% \\
Tecnologia predominante & rádio/\allowbreak fibra \\
Opção rural & Starlink disponível (\textasciitilde USD 44/\allowbreak mês) \\
\end{longtable}

\subsubsection{Mercado Imobiliário e Terra
Rural}

\textbf{Fonte:} INDERT /\allowbreak  Clasificados.com.py (2024)

\begin{longtable}[]{@{}ll@{}}
\toprule\noalign{}
Tipo & Referência \\
\midrule\noalign{}
\endhead
\bottomrule\noalign{}
\endlastfoot
Terra agrícola alta prod. (USD/\allowbreak ha) & 8,800 \\
Imóvel urbano (USD/\allowbreak m²) & 850 \\
Aluguel 2 quartos (USD/\allowbreak mês) & 350 \\
\end{longtable}

\begin{quote}
Valores de referência departamental. Localidades menores podem ter
preços 20--40\% abaixo da capital departamental.
\end{quote}

\subsubsection{Saúde}

\textbf{Fonte:} MSPBS /\allowbreak  IPS Paraguay (2024)

\begin{longtable}[]{@{}ll@{}}
\toprule\noalign{}
Serviço & Disponibilidade \\
\midrule\noalign{}
\endhead
\bottomrule\noalign{}
\endlastfoot
USF /\allowbreak  Posto de Saúde & sim \\
Hospital Regional & não \\
IPS (seguro social) & não \\
Farmácia & sim \\
Distância ao hospital de referência & \textasciitilde25 km (Coronel
Oviedo) \\
\end{longtable}

\textbf{Principais estabelecimentos:} USF local; farmácias distritais.

\textbf{Observação para imigrantes:} Atendimento primário disponível.
Proximidade com Coronel Oviedo garante rápido acesso ao hospital
regional. Cobertura privada recomendada; IPS não disponível localmente.
\newpage
\secaoDiagnostico{Nueva Toledo}{\mapaConteudo{-24.8666}{-55.5666}}{Nueva Toledo é uma localidade que personifica o isolamento demográfico
em uma das regiões agrícolas mais produtivas do país. Sua maior força
reside na baixíssima carga populacional e na invisibilidade estratégica,
sendo ideal para quem prioriza a privacidade e a segurança social.
Oferece solo de alta qualidade para grãos e um ambiente cultural
familiar e trabalhador, exigindo do ocupante autonomia logística e de
saúde devido à infraestrutura local limitada.

\subsubsection{Por que sim}

\begin{itemize}
\tightlist
\item
  Densidade populacional entre as menores do departamento (privacidade
  total).
\item
  Solo de excelente qualidade para agricultura mecanizada e pecuária.
\item
  Nova rota asfáltica que melhora o acesso sem atrair fluxos migratórios
  urbanos.
\item
  Sem restrições da Lei de Fronteira para estrangeiros.
\end{itemize}

\subsubsection{Por que não}

\begin{itemize}
\tightlist
\item
  Infraestrutura de saúde local limitada a atendimentos primários.
\item
  Dependência de estradas vicinais de terra para acesso às colônias.
\item
  Distância de centros comerciais e financeiros diversificados.
\end{itemize}

\subsubsection{Recomendação
Técnica}

Altamente indicado para estabelecimentos de autossuficiência de médio
porte ou exploração agrícola mecanizada que busquem isolamento tático.
Requer autonomia em termos de logística básica e saúde. O GSS de 5.7
reflete a segurança do isolamento em contraste com a limitação
institucional local. Referencia atual de score: GSS 5.7 em 2026-03-05.

\subsubsection{}

\begin{itemize}
\tightlist
\item
  \begin{enumerate}
  \def\labelenumi{\arabic{enumi})}
  \tightlist
  \item
    Dados censitários validados (INE\footnote{Instituto Nacional de Estadística (Instituto Nacional de Estatística), órgão oficial de dados demográficos e censitários.} 2022).
  \end{enumerate}
\item
  \begin{enumerate}
  \def\labelenumi{\arabic{enumi})}
  \setcounter{enumi}{1}
  \tightlist
  \item
    Relatórios de pavimentação asfáltica recente (MOPC 2024-2026).
  \end{enumerate}
\item
  \begin{enumerate}
  \def\labelenumi{\arabic{enumi})}
  \setcounter{enumi}{2}
  \tightlist
  \item
    Mapa de solos e aptidão agrícola de Caaguazú (MAG).
  \end{enumerate}
\item
  \begin{enumerate}
  \def\labelenumi{\arabic{enumi})}
  \setcounter{enumi}{3}
  \tightlist
  \item
    Registro de segurança pública departamental.
  \end{enumerate}
\end{itemize}}
\begin{table}[ht!]
\small \begin{tabular}{p{3.0cm}p{0.8cm}p{6.0cm}} \toprule
\textbf{Critério} & \textbf{Nota} & \textbf{Justificativa} \\ \midrule
A - Ameaças Estratégicas & 6.0 & Localização discreta e fora dos eixos de conflito primário; excelente invisibilidade tática \\
B - Risco Social & 7.0 & População pequena e baixíssima densidade; risco social urbano desprezível \\
C - Riscos Naturais & 6.5 & Geologia muito estável; penalizado pelo risco moderado de isolamento pluvial \\
D - Autossuficiência & 6.5 & Bons recursos: soja, milho e gado, porém limitado pela infraestrutura logística local \\
E - Institucional & 4.5 & Município pequeno com infraestrutura institucional e de saúde mínima \\
\midrule \textbf{GSS FINAL} & \textbf{5.7} & \textbf{Moderadamente Seguro (Ideal para quem busca o máximo de isolamento demográfico em uma região agrícola próspera).} \\ \bottomrule \end{tabular}
\end{table}
\subsubsection{Vulnerabilidades e
Mitigação}\label{vuln-Nueva_Toledo-vulnerabilidades-e-mitigauxe7uxe3o}

\begin{itemize}
\tightlist
\item
  \textbf{Isolamento de Saúde:} Distância significativa de hospitais
  especializados (Mitigação: manutenção de estoque médico básico e
  treinamento de socorrista local).
\item
  \textbf{Dependência Logística:} Necessidade de deslocamento para
  cidades maiores para insumos técnicos e comerciais (Mitigação: estoque
  de suprimentos para 60 dias).
\item
  \textbf{Instabilidade Elétrica Rural:} Quedas ocasionais em temporais
  (Mitigação: instalação de sistemas solares fotovoltaicos redundantes).
\end{itemize}
\subsubsection{Dossiê de Campo}\label{dossie-Nueva_Toledo-dossiuxea-de-campo}
\begin{description}[style=nextline,leftmargin=0.5cm,font=\bfseries]
\item[Coordenadas]  24°52'S 55°34'W.
\item[Topografia]  Relevo predominantemente plano a suavemente ondulado. Localizado no nordeste do departamento. Altitude \textasciitilde 240m. Área de \textasciitilde 600 km². Geologicamente muito estável, risco sísmico desprezível.
\item[Alvos Estratégicos]  Áreas de produção agropecuária mecanizada de alta produtividade; Ponto de transição secundário entre Caaguazú e Canindeyú. Ausência de alvos militares ou industriais primários, oferecendo excelente invisibilidade tática e discrição estratégica.
\item[Fallout]  Ventos predominantes do Norte (quentes e úmidos) e Leste. Baixo risco de fallout direto.
\item[População]  4.826 habitantes (Censo 2022 Final). Perfil essencialmente rural, com forte presença de imigrantes e descendentes de origem brasileira (brasiguaios).
\item[Densidade]  \textasciitilde 8,0 hab/\allowbreak km² (Densidade baixíssima, oferecendo isolamento tático superior e privacidade estratégica).
\item[Vias de Acesso]  Acesso via Rota D024, recentemente beneficiada pela pavimentação asfáltica do trecho Nueva Toledo -- Vaquería (20 km). Conectividade interna dependente de estradas vicinais de terra que podem ser críticas em períodos chuvosos intensos (média 1.600 mm/\allowbreak ano).
\item[Serviços]  Unidade de Saúde da Família (USF) local. Dependência de cidades maiores como Caaguazú (cidade) ou Coronel Oviedo para serviços hospitalares de alta complexidade.
\item[Custo de Vida]  Baixo; economia estritamente vinculada ao agronegócio e agricultura familiar.
\item[Preço da Terra]  US\$ 5.000-8.500/\allowbreak ha (áreas rurais mecanizadas para soja); US\$ 2.500-4.000/\allowbreak ha (áreas de pecuária ou campo natural).
\item[Hidrologia]  Baixo risco de inundações sistêmicas catastróficas. Presença de arroios locais que exigem atenção apenas em eventos de precipitação extrema. Drenagem natural favorável.
\item[Clima]  Subtropical Úmido. Verões intensos e invernos com entradas ocasionais de ar polar.
\item[Energia]  Rede nacional (Itaipu); instável em áreas rurais remotas devido ao isolamento da rede.
\item[Água]  Captação via poços artesianos e sistemas de junta de saneamento local; boa disponibilidade hídrica subterrânea.
\item[Qualidade do Solo]  Latossolo Vermelho (Terra Vermelha); fértil, profundo e bem drenado. Excelente aptidão para soja, milho, trigo e pastagens cultivadas.
\item[Resources Locais]  Grande produção de grãos (soja e milho) e gado de corte. Elevado potencial para autossuficiência alimentar básica em nível de propriedade.
\item[Segurança]  Zona rural pacífica e estável. Baixíssima criminalidade urbana violenta. Vigilância necessária contra furtos de gado e questões agrárias pontuais ligadas a comunidades indígenas e fronteira agrícola.
\item[Leis Local: Município consolidado. Livre de Restrição de Fronteira]  Fora da faixa de 50km da fronteira internacional, permitindo plena titularidade para brasileiros.
\end{description}
\textbf{Fonte climática:} NASA POWER Climatology API (período 2001-2020)
\textbf{Fonte luminosa:} estimativa world atlas

\paragraph{Irradiação Solar
(kWh/\allowbreak m²/\allowbreak dia)}\label{irradiauxe7uxe3o-solar-kwhmuxb2dia}

\begin{longtable}[]{@{}
  >{\raggedright\arraybackslash}p{(\linewidth - 24\tabcolsep) * \real{0.0746}}
  >{\raggedright\arraybackslash}p{(\linewidth - 24\tabcolsep) * \real{0.0746}}
  >{\raggedright\arraybackslash}p{(\linewidth - 24\tabcolsep) * \real{0.0746}}
  >{\raggedright\arraybackslash}p{(\linewidth - 24\tabcolsep) * \real{0.0746}}
  >{\raggedright\arraybackslash}p{(\linewidth - 24\tabcolsep) * \real{0.0746}}
  >{\raggedright\arraybackslash}p{(\linewidth - 24\tabcolsep) * \real{0.0746}}
  >{\raggedright\arraybackslash}p{(\linewidth - 24\tabcolsep) * \real{0.0746}}
  >{\raggedright\arraybackslash}p{(\linewidth - 24\tabcolsep) * \real{0.0746}}
  >{\raggedright\arraybackslash}p{(\linewidth - 24\tabcolsep) * \real{0.0746}}
  >{\raggedright\arraybackslash}p{(\linewidth - 24\tabcolsep) * \real{0.0746}}
  >{\raggedright\arraybackslash}p{(\linewidth - 24\tabcolsep) * \real{0.0746}}
  >{\raggedright\arraybackslash}p{(\linewidth - 24\tabcolsep) * \real{0.0746}}
  >{\raggedright\arraybackslash}p{(\linewidth - 24\tabcolsep) * \real{0.1045}}@{}}
\toprule\noalign{}
\begin{minipage}[b]{\linewidth}\raggedright
Jan
\end{minipage} & \begin{minipage}[b]{\linewidth}\raggedright
Fev
\end{minipage} & \begin{minipage}[b]{\linewidth}\raggedright
Mar
\end{minipage} & \begin{minipage}[b]{\linewidth}\raggedright
Abr
\end{minipage} & \begin{minipage}[b]{\linewidth}\raggedright
Mai
\end{minipage} & \begin{minipage}[b]{\linewidth}\raggedright
Jun
\end{minipage} & \begin{minipage}[b]{\linewidth}\raggedright
Jul
\end{minipage} & \begin{minipage}[b]{\linewidth}\raggedright
Ago
\end{minipage} & \begin{minipage}[b]{\linewidth}\raggedright
Set
\end{minipage} & \begin{minipage}[b]{\linewidth}\raggedright
Out
\end{minipage} & \begin{minipage}[b]{\linewidth}\raggedright
Nov
\end{minipage} & \begin{minipage}[b]{\linewidth}\raggedright
Dez
\end{minipage} & \begin{minipage}[b]{\linewidth}\raggedright
Média
\end{minipage} \\
\midrule\noalign{}
\endhead
\bottomrule\noalign{}
\endlastfoot
6.56 & 6.04 & 5.59 & 4.57 & 3.37 & 2.97 & 3.27 & 4.04 & 4.58 & 5.29 &
6.28 & 6.54 & \textbf{4.92} \\
\end{longtable}

\textbf{Inclinação solar recomendada:} 25° N (anual) · 35° N (inverno
jun-ago) · 15° N (verão nov-jan)

\paragraph{Precipitação (mm/\allowbreak mês)}\label{precipitauxe7uxe3o-mmmuxeas}

\begin{longtable}[]{@{}
  >{\raggedright\arraybackslash}p{(\linewidth - 24\tabcolsep) * \real{0.0704}}
  >{\raggedright\arraybackslash}p{(\linewidth - 24\tabcolsep) * \real{0.0704}}
  >{\raggedright\arraybackslash}p{(\linewidth - 24\tabcolsep) * \real{0.0704}}
  >{\raggedright\arraybackslash}p{(\linewidth - 24\tabcolsep) * \real{0.0704}}
  >{\raggedright\arraybackslash}p{(\linewidth - 24\tabcolsep) * \real{0.0704}}
  >{\raggedright\arraybackslash}p{(\linewidth - 24\tabcolsep) * \real{0.0704}}
  >{\raggedright\arraybackslash}p{(\linewidth - 24\tabcolsep) * \real{0.0704}}
  >{\raggedright\arraybackslash}p{(\linewidth - 24\tabcolsep) * \real{0.0704}}
  >{\raggedright\arraybackslash}p{(\linewidth - 24\tabcolsep) * \real{0.0704}}
  >{\raggedright\arraybackslash}p{(\linewidth - 24\tabcolsep) * \real{0.0704}}
  >{\raggedright\arraybackslash}p{(\linewidth - 24\tabcolsep) * \real{0.0704}}
  >{\raggedright\arraybackslash}p{(\linewidth - 24\tabcolsep) * \real{0.0704}}
  >{\raggedright\arraybackslash}p{(\linewidth - 24\tabcolsep) * \real{0.1549}}@{}}
\toprule\noalign{}
\begin{minipage}[b]{\linewidth}\raggedright
Jan
\end{minipage} & \begin{minipage}[b]{\linewidth}\raggedright
Fev
\end{minipage} & \begin{minipage}[b]{\linewidth}\raggedright
Mar
\end{minipage} & \begin{minipage}[b]{\linewidth}\raggedright
Abr
\end{minipage} & \begin{minipage}[b]{\linewidth}\raggedright
Mai
\end{minipage} & \begin{minipage}[b]{\linewidth}\raggedright
Jun
\end{minipage} & \begin{minipage}[b]{\linewidth}\raggedright
Jul
\end{minipage} & \begin{minipage}[b]{\linewidth}\raggedright
Ago
\end{minipage} & \begin{minipage}[b]{\linewidth}\raggedright
Set
\end{minipage} & \begin{minipage}[b]{\linewidth}\raggedright
Out
\end{minipage} & \begin{minipage}[b]{\linewidth}\raggedright
Nov
\end{minipage} & \begin{minipage}[b]{\linewidth}\raggedright
Dez
\end{minipage} & \begin{minipage}[b]{\linewidth}\raggedright
Total/\allowbreak ano
\end{minipage} \\
\midrule\noalign{}
\endhead
\bottomrule\noalign{}
\endlastfoot
134 & 138 & 121 & 148 & 174 & 93 & 73 & 51 & 104 & 186 & 185 & 172 &
\textbf{1581 mm} \\
\end{longtable}

\paragraph{Poluição Luminosa}\label{poluiuxe7uxe3o-luminosa}

\begin{longtable}[]{@{}ll@{}}
\toprule\noalign{}
Parâmetro & Valor \\
\midrule\noalign{}
\endhead
\bottomrule\noalign{}
\endlastfoot
Escala Bortle & 4 --- Céu rural-suburbano \\
Radiância artificial & 3.0 nW/\allowbreak cm²/\allowbreak sr \\
\end{longtable}
\paragraph{Solo (SoilGrids 2.0, média ponderada 0--30
cm)}

\begin{longtable}[]{@{}ll@{}}
\toprule\noalign{}
Parâmetro & Valor \\
\midrule\noalign{}
\endhead
\bottomrule\noalign{}
\endlastfoot
pH (H$_2$O) & 5.37 \\
Carbono orgânico (SOC) & 19.09 g/\allowbreak kg \\
Argila & 28.0 \% \\
Areia & 52.17 \% \\
Silte (calc.) & 19.8 \% \\
Densidade aparente & 1.33 g/\allowbreak cm³ \\
\textbf{Aptidão agrícola} & \textbf{Média} \\
\end{longtable}

Fonte: ISRIC SoilGrids 2.0 via WCS. Coords: -24.8667°, -55.5667°. Média
ponderada camadas 0-5, 5-15, 15-30 cm.
\subsubsection{Indicadores Sociais}

\textbf{Fontes:} PNUD 2020, INE\footnote{Instituto Nacional de Estadística (Instituto Nacional de Estatística), órgão oficial de dados demográficos e censitários.} Censo 2022, INE\footnote{Instituto Nacional de Estadística (Instituto Nacional de Estatística), órgão oficial de dados demográficos e censitários.} EPHC 2023, Ministerio
Público 2024\\
\textbf{Nota:} Valores marcados como \emph{(dept.)} referem-se ao
departamento de Caaguazú; valores \emph{(dist.)} são específicos deste
distrito. Dados marcados \emph{(est.)} são estimativas.

\begin{longtable}[]{@{}
  >{\raggedright\arraybackslash}p{(\linewidth - 4\tabcolsep) * \real{0.4231}}
  >{\raggedright\arraybackslash}p{(\linewidth - 4\tabcolsep) * \real{0.2692}}
  >{\raggedright\arraybackslash}p{(\linewidth - 4\tabcolsep) * \real{0.3077}}@{}}
\toprule\noalign{}
\begin{minipage}[b]{\linewidth}\raggedright
Indicador
\end{minipage} & \begin{minipage}[b]{\linewidth}\raggedright
Valor
\end{minipage} & \begin{minipage}[b]{\linewidth}\raggedright
Âmbito
\end{minipage} \\
\midrule\noalign{}
\endhead
\bottomrule\noalign{}
\endlastfoot
IDH (2020) & 0,715 & dept. \\
Ranking IDH nacional & 7/\allowbreak 18 & dept. \\
Esperança de vida & 73,1 anos & dept. \\
Escolaridade média & 6.9 anos & dist. \\
RNB per capita (USD PPA) & 9.000 & dept. \\
Pobreza monetária (\%) & 28,5\% & dept. \\
Pobreza extrema (\%) & 6,8\% & dept. \\
Índice de Gini & 0,447 & dept. \\
Acesso a água potável (\%) & 79,8\% & dept. \\
Acesso a saneamento (\%) & 78,3\% & dept. \\
Taxa de homicídios (est., /\allowbreak 100k hab) & \textasciitilde6--8 (estimativa,
próxima ou acima da média nacional) & dept. \\
Índice de segurança & médio & dept. \\
População (Censo 2022) & 4.826 hab. & dist. \\
Idade mediana & 23.0 anos & dist. \\
Taxa de fecundidade & N/\allowbreak D & dist. \\
\end{longtable}
\subsubsection{Combustível}

\textbf{Referência:} PETROPAR /\allowbreak  postos locais (2024) \textbf{Tipo de
localidade:} Interior

\begin{longtable}[]{@{}lll@{}}
\toprule\noalign{}
Combustível & USD/\allowbreak litro & Gs/\allowbreak litro (aprox.) \\
\midrule\noalign{}
\endhead
\bottomrule\noalign{}
\endlastfoot
Gasolina 93 oct & 0.98 & 7,252 \\
Gasolina 97 oct (premium) & 1.08 & 7,992 \\
Diesel & 0.91 & 6,734 \\
\end{longtable}

\begin{quote}
Preços podem variar ±5\% conforme posto e sazonalidade. Chaco e interior
remoto apresentam maior variação.
\end{quote}

\subsubsection{Cobertura Celular}

\textbf{Fonte:} CONATEL PY /\allowbreak  operadoras (2024)

\begin{longtable}[]{@{}ll@{}}
\toprule\noalign{}
Parâmetro & Valor \\
\midrule\noalign{}
\endhead
\bottomrule\noalign{}
\endlastfoot
Cobertura 4G população (dept.) & 88\% \\
Cobertura 4G área rural & 72\% \\
Melhor operadora & Tigo \\
Qualidade rural & boa \\
\end{longtable}

\begin{quote}
Para áreas rurais fora do núcleo urbano, recomenda-se chip Tigo como
principal e Personal como backup.
\end{quote}

\subsubsection{Internet}

\textbf{Fonte:} CONATEL /\allowbreak  Speedtest Ookla (2024)

\begin{longtable}[]{@{}ll@{}}
\toprule\noalign{}
Parâmetro & Valor \\
\midrule\noalign{}
\endhead
\bottomrule\noalign{}
\endlastfoot
Velocidade média download & 55 Mbps \\
Domicílios com internet (dept.) & 62\% \\
Tecnologia predominante & rádio/\allowbreak fibra \\
Opção rural & Starlink disponível (\textasciitilde USD 44/\allowbreak mês) \\
\end{longtable}

\subsubsection{Mercado Imobiliário e Terra
Rural}

\textbf{Fonte:} INDERT /\allowbreak  Clasificados.com.py (2024)

\begin{longtable}[]{@{}ll@{}}
\toprule\noalign{}
Tipo & Referência \\
\midrule\noalign{}
\endhead
\bottomrule\noalign{}
\endlastfoot
Terra agrícola alta prod. (USD/\allowbreak ha) & 8,800 \\
Imóvel urbano (USD/\allowbreak m²) & 850 \\
Aluguel 2 quartos (USD/\allowbreak mês) & 350 \\
\end{longtable}

\begin{quote}
Valores de referência departamental. Localidades menores podem ter
preços 20--40\% abaixo da capital departamental.
\end{quote}

\subsubsection{Saúde}

\textbf{Fonte:} MSPBS /\allowbreak  IPS Paraguay (2024)

\begin{longtable}[]{@{}ll@{}}
\toprule\noalign{}
Serviço & Disponibilidade \\
\midrule\noalign{}
\endhead
\bottomrule\noalign{}
\endlastfoot
USF /\allowbreak  Posto de Saúde & sim \\
Hospital Regional & não \\
IPS (seguro social) & não \\
Farmácia & sim \\
Distância ao hospital de referência & \textasciitilde80 km (Coronel
Oviedo) \\
\end{longtable}

\textbf{Principais estabelecimentos:} USF local; farmácias distritais.

\textbf{Observação para imigrantes:} Atendimento primário disponível.
Casos de maior complexidade encaminhados a Coronel Oviedo. Cobertura
privada recomendada; IPS não disponível localmente.
\newpage
\secaoDiagnostico{Raul Arsenio Oviedo}{\mapaConteudo{-25.1666}{-55.6333}}{Raúl Arsenio Oviedo é um refúgio de produtividade no centro-leste de
Caaguazú. Sua maior vantagem é a combinação de baixíssima densidade
populacional com solos de elite (``terra vermelha'') e proximidade com o
reservatório de Yguazú. Oferece um santuário de paz social e recursos
naturais abundantes, exigindo do ocupante autonomia logística para
serviços de saúde e comércio avançado, concentrados na vizinha Campo 9.

\subsubsection{Por que sim}

\begin{itemize}
\tightlist
\item
  Solo de altíssima produtividade e clima favorável para grãos e
  pecuária.
\item
  Baixíssima densidade demográfica e ambiente social extremamente
  pacífico.
\item
  Segurança hídrica regional garantida pela proximidade com o Lago
  Yguazú.
\item
  Sem restrições da Lei de Fronteira.
\end{itemize}

\subsubsection{Por que não}

\begin{itemize}
\tightlist
\item
  Infraestrutura de saúde local limitada a atendimentos primários.
\item
  Dependência de estradas vicinais de terra para acesso às propriedades.
\item
  Valor da terra agrícola acima da média nacional devido à
  produtividade.
\end{itemize}

\subsubsection{Recomendação
Técnica}

Altamente indicado para estabelecimentos de autossuficiência de médio
porte ou exploração agrícola que busque o máximo de isolamento
demográfico em solo de alta performance. Requer autonomia em logística
básica e saúde. O GSS de 6.5 reflete a segurança do isolamento em uma
região próspera. Referencia atual de score: GSS 6.5 em 2026-03-05.

\subsubsection{}

\begin{itemize}
\tightlist
\item
  \begin{enumerate}
  \def\labelenumi{\arabic{enumi})}
  \tightlist
  \item
    Dados censitários validados (INE\footnote{Instituto Nacional de Estadística (Instituto Nacional de Estatística), órgão oficial de dados demográficos e censitários.} 2022).
  \end{enumerate}
\item
  \begin{enumerate}
  \def\labelenumi{\arabic{enumi})}
  \setcounter{enumi}{1}
  \tightlist
  \item
    Registro de produtores de soja e trigo (MAG).
  \end{enumerate}
\item
  \begin{enumerate}
  \def\labelenumi{\arabic{enumi})}
  \setcounter{enumi}{2}
  \tightlist
  \item
    Mapa de solos e aptidão agrícola de Caaguazú (MAG).
  \end{enumerate}
\item
  \begin{enumerate}
  \def\labelenumi{\arabic{enumi})}
  \setcounter{enumi}{3}
  \tightlist
  \item
    Registro de segurança pública departamental.
  \end{enumerate}
\end{itemize}}
\begin{table}[ht!]
\small \begin{tabular}{p{3.0cm}p{0.8cm}p{6.0cm}} \toprule
\textbf{Critério} & \textbf{Nota} & \textbf{Justificativa} \\ \midrule
A - Ameaças Estratégicas & 6.0 & Localização discreta e fora dos eixos de conflito primário; boa invisibilidade estratégica \\
B - Risco Social & 7.0 & População pequena e baixa densidade demográfica; risco social urbano desprezível \\
C - Riscos Naturais & 7.0 & Geologia muito estável e riscos climáticos típicos manejáveis \\
D - Autossuficiência & 7.5 & Riqueza em recursos: solo de elite, gado e proximidade com o reservatório de Yguazú \\
E - Institucional & 5.0 & Município consolidado com segurança social, mas infraestrutura de serviços dependente de vizinhos \\
\midrule \textbf{GSS FINAL} & \textbf{6.5} & \textbf{Moderadamente Seguro (Ideal para quem busca isolamento demográfico em terras de altíssima produtividade).} \\ \bottomrule \end{tabular}
\end{table}
\subsubsection{Vulnerabilidades e
Mitigação}\label{vuln-Raul_Arsenio_Oviedo-vulnerabilidades-e-mitigauxe7uxe3o}

\begin{itemize}
\tightlist
\item
  \textbf{Isolamento de Saúde:} Distância de centros hospitalares
  cirúrgicos (Mitigação: manutenção de estoque médico básico e
  treinamento de primeiros socorros).
\item
  \textbf{Dependência Logística:} Necessidade de deslocamento para Campo
  9 para insumos técnicos (Mitigação: estoque de suprimentos para 60
  dias).
\item
  \textbf{Instabilidade Elétrica Rural:} Quedas ocasionais em temporais
  (Mitigação: instalação de sistemas solares fotovoltaicos redundantes).
\end{itemize}
\subsubsection{Dossiê de Campo}\label{dossie-Raul_Arsenio_Oviedo-dossiuxea-de-campo}
\begin{description}[style=nextline,leftmargin=0.5cm,font=\bfseries]
\item[Coordenadas]  25°10'S 55°38'W.
\item[Topografia]  Relevo predominantemente plano a suavemente ondulado. Localizado na bacia do Rio Monday, com proximidade estratégica com o reservatório de Yguazú. Altitude \textasciitilde 230m. Área de \textasciitilde 625 km². Geologicamente muito estável, risco sísmico desprezível.
\item[Alvos Estratégicos]  Áreas de produção agrícola mecanizada de larga escala; Terminais de armazenamento de grãos. Ausência de alvos militares ou industriais pesados primários, conferindo excelente invisibilidade tática e discrição.
\item[Fallout]  Ventos predominantes do Norte (NE) e Leste. Baixo risco de fallout direto.
\item[População]  12.543 habitantes (Censo 2022 Final). Perfil demográfico essencialmente rural (85\%) com forte cultura de trabalho agrícola mecanizado.
\item[Densidade]  \textasciitilde 20,1 hab/\allowbreak km² (Baixa densidade, oferecendo excelente isolamento tático e privacidade estratégica).
\item[Vias de Acesso]  Acesso via ramais da Ruta PY13. Conectividade interna por estradas de terra batida que exigem manutenção e veículos robustos em períodos de chuva intensa (média 1.600 mm/\allowbreak ano).
\item[Serviços]  Unidade de Saúde da Família (USF) local. Dependência direta do polo de Campo 9 (Dr. J. Eulogio Estigarribia) a 25 km para saúde de alta complexidade e logística comercial avançada.
\item[Custo de Vida]  Baixo; economia baseada na exportação de grãos e pecuária.
\item[Preço da Terra]  US\$ 5.000-8.000/\allowbreak ha (terras mecanizadas de elite); US\$ 2.500-4.000/\allowbreak ha (áreas de campo natural/\allowbreak pecuária).
\item[Hidrologia]  Banhado por arroios tributários do Rio Monday. Baixo risco de inundações sistêmicas. Gestão hídrica favorável à agricultura.
\item[Clima]  Subtropical Úmido. Verões quentes; invernos com geadas ocasionais.
\item[Energia]  Rede nacional (Itaipu); instável em áreas rurais remotas durante picos de verão.
\item[Água]  Abundante via poços artesianos e lençol freático produtivo; proximidade com o reservatório de Yguazú garante segurança hídrica regional.
\item[Qualidade do Solo]  Latossolo Vermelho (Terra Vermelha); profundo, fértil e bem drenado. Excelente aptidão para soja, milho, trigo e pastagens de alta performance.
\item[Resources Locais]  Grande produção de soja, milho, trigo e pecuária de corte. Elevadíssimo potencial para autossuficiência alimentar básica e geração de excedentes.
\item[Segurança]  Zona rural extremamente pacífica e estável. Baixíssima criminalidade urbana. Coesão social baseada na cultura produtora brasiguaia e paraguaia tradicional. Baixa visibilidade para agitações sociais de massa.
\item[Leis Local: Município consolidado. Livre de Restrição de Fronteira]  Fora da faixa de 50km da fronteira, permitindo plena titularidade para brasileiros.
\end{description}
\textbf{Fonte climática:} NASA POWER Climatology API (período 2001-2020)
\textbf{Fonte luminosa:} estimativa world atlas

\paragraph{Irradiação Solar
(kWh/\allowbreak m²/\allowbreak dia)}\label{irradiauxe7uxe3o-solar-kwhmuxb2dia}

\begin{longtable}[]{@{}
  >{\raggedright\arraybackslash}p{(\linewidth - 24\tabcolsep) * \real{0.0746}}
  >{\raggedright\arraybackslash}p{(\linewidth - 24\tabcolsep) * \real{0.0746}}
  >{\raggedright\arraybackslash}p{(\linewidth - 24\tabcolsep) * \real{0.0746}}
  >{\raggedright\arraybackslash}p{(\linewidth - 24\tabcolsep) * \real{0.0746}}
  >{\raggedright\arraybackslash}p{(\linewidth - 24\tabcolsep) * \real{0.0746}}
  >{\raggedright\arraybackslash}p{(\linewidth - 24\tabcolsep) * \real{0.0746}}
  >{\raggedright\arraybackslash}p{(\linewidth - 24\tabcolsep) * \real{0.0746}}
  >{\raggedright\arraybackslash}p{(\linewidth - 24\tabcolsep) * \real{0.0746}}
  >{\raggedright\arraybackslash}p{(\linewidth - 24\tabcolsep) * \real{0.0746}}
  >{\raggedright\arraybackslash}p{(\linewidth - 24\tabcolsep) * \real{0.0746}}
  >{\raggedright\arraybackslash}p{(\linewidth - 24\tabcolsep) * \real{0.0746}}
  >{\raggedright\arraybackslash}p{(\linewidth - 24\tabcolsep) * \real{0.0746}}
  >{\raggedright\arraybackslash}p{(\linewidth - 24\tabcolsep) * \real{0.1045}}@{}}
\toprule\noalign{}
\begin{minipage}[b]{\linewidth}\raggedright
Jan
\end{minipage} & \begin{minipage}[b]{\linewidth}\raggedright
Fev
\end{minipage} & \begin{minipage}[b]{\linewidth}\raggedright
Mar
\end{minipage} & \begin{minipage}[b]{\linewidth}\raggedright
Abr
\end{minipage} & \begin{minipage}[b]{\linewidth}\raggedright
Mai
\end{minipage} & \begin{minipage}[b]{\linewidth}\raggedright
Jun
\end{minipage} & \begin{minipage}[b]{\linewidth}\raggedright
Jul
\end{minipage} & \begin{minipage}[b]{\linewidth}\raggedright
Ago
\end{minipage} & \begin{minipage}[b]{\linewidth}\raggedright
Set
\end{minipage} & \begin{minipage}[b]{\linewidth}\raggedright
Out
\end{minipage} & \begin{minipage}[b]{\linewidth}\raggedright
Nov
\end{minipage} & \begin{minipage}[b]{\linewidth}\raggedright
Dez
\end{minipage} & \begin{minipage}[b]{\linewidth}\raggedright
Média
\end{minipage} \\
\midrule\noalign{}
\endhead
\bottomrule\noalign{}
\endlastfoot
6.58 & 6.08 & 5.45 & 4.47 & 3.30 & 2.90 & 3.20 & 3.96 & 4.53 & 5.31 &
6.35 & 6.66 & \textbf{4.9} \\
\end{longtable}

\textbf{Inclinação solar recomendada:} 25° N (anual) · 35° N (inverno
jun-ago) · 15° N (verão nov-jan)

\paragraph{Precipitação (mm/\allowbreak mês)}\label{precipitauxe7uxe3o-mmmuxeas}

\begin{longtable}[]{@{}
  >{\raggedright\arraybackslash}p{(\linewidth - 24\tabcolsep) * \real{0.0704}}
  >{\raggedright\arraybackslash}p{(\linewidth - 24\tabcolsep) * \real{0.0704}}
  >{\raggedright\arraybackslash}p{(\linewidth - 24\tabcolsep) * \real{0.0704}}
  >{\raggedright\arraybackslash}p{(\linewidth - 24\tabcolsep) * \real{0.0704}}
  >{\raggedright\arraybackslash}p{(\linewidth - 24\tabcolsep) * \real{0.0704}}
  >{\raggedright\arraybackslash}p{(\linewidth - 24\tabcolsep) * \real{0.0704}}
  >{\raggedright\arraybackslash}p{(\linewidth - 24\tabcolsep) * \real{0.0704}}
  >{\raggedright\arraybackslash}p{(\linewidth - 24\tabcolsep) * \real{0.0704}}
  >{\raggedright\arraybackslash}p{(\linewidth - 24\tabcolsep) * \real{0.0704}}
  >{\raggedright\arraybackslash}p{(\linewidth - 24\tabcolsep) * \real{0.0704}}
  >{\raggedright\arraybackslash}p{(\linewidth - 24\tabcolsep) * \real{0.0704}}
  >{\raggedright\arraybackslash}p{(\linewidth - 24\tabcolsep) * \real{0.0704}}
  >{\raggedright\arraybackslash}p{(\linewidth - 24\tabcolsep) * \real{0.1549}}@{}}
\toprule\noalign{}
\begin{minipage}[b]{\linewidth}\raggedright
Jan
\end{minipage} & \begin{minipage}[b]{\linewidth}\raggedright
Fev
\end{minipage} & \begin{minipage}[b]{\linewidth}\raggedright
Mar
\end{minipage} & \begin{minipage}[b]{\linewidth}\raggedright
Abr
\end{minipage} & \begin{minipage}[b]{\linewidth}\raggedright
Mai
\end{minipage} & \begin{minipage}[b]{\linewidth}\raggedright
Jun
\end{minipage} & \begin{minipage}[b]{\linewidth}\raggedright
Jul
\end{minipage} & \begin{minipage}[b]{\linewidth}\raggedright
Ago
\end{minipage} & \begin{minipage}[b]{\linewidth}\raggedright
Set
\end{minipage} & \begin{minipage}[b]{\linewidth}\raggedright
Out
\end{minipage} & \begin{minipage}[b]{\linewidth}\raggedright
Nov
\end{minipage} & \begin{minipage}[b]{\linewidth}\raggedright
Dez
\end{minipage} & \begin{minipage}[b]{\linewidth}\raggedright
Total/\allowbreak ano
\end{minipage} \\
\midrule\noalign{}
\endhead
\bottomrule\noalign{}
\endlastfoot
134 & 138 & 121 & 148 & 174 & 93 & 73 & 51 & 104 & 186 & 185 & 172 &
\textbf{1581 mm} \\
\end{longtable}

\paragraph{Poluição Luminosa}\label{poluiuxe7uxe3o-luminosa}

\begin{longtable}[]{@{}ll@{}}
\toprule\noalign{}
Parâmetro & Valor \\
\midrule\noalign{}
\endhead
\bottomrule\noalign{}
\endlastfoot
Escala Bortle & 4 --- Céu rural-suburbano \\
Radiância artificial & 3.0 nW/\allowbreak cm²/\allowbreak sr \\
\end{longtable}
\paragraph{Solo (SoilGrids 2.0, média ponderada 0--30
cm)}

\begin{longtable}[]{@{}ll@{}}
\toprule\noalign{}
Parâmetro & Valor \\
\midrule\noalign{}
\endhead
\bottomrule\noalign{}
\endlastfoot
pH (H$_2$O) & 5.3 \\
Carbono orgânico (SOC) & 18.47 g/\allowbreak kg \\
Argila & 22.88 \% \\
Areia & 57.92 \% \\
Silte (calc.) & 19.2 \% \\
Densidade aparente & 1.29 g/\allowbreak cm³ \\
\textbf{Aptidão agrícola} & \textbf{Média} \\
\end{longtable}

Fonte: ISRIC SoilGrids 2.0 via WCS. Coords: -25.1667°, -55.6333°. Média
ponderada camadas 0-5, 5-15, 15-30 cm.

\begin{itemize}
\tightlist
\item
  \textbf{Homicidios/\allowbreak 100k:} \textasciitilde8,5 (Regional/\allowbreak Caaguazú).
\end{itemize}

\subsection{10. Analise de Riscos}

\begin{itemize}
\tightlist
\item
  \textbf{Cenario curto prazo:} Manutenção da rentabilidade das culturas
  de grãos e valorização das terras mecanizadas.
\item
  \textbf{Cenario medio prazo:} Impacto positivo da melhoria de acessos
  vicinais e da rede elétrica rural.
\end{itemize}

\subsection{11. Redacao Final}

\subsubsection{Diagnostico Integrado}

Raúl Arsenio Oviedo é um refúgio de produtividade no centro-leste de
Caaguazú. Sua maior vantagem é a combinação de baixíssima densidade
populacional com solos de elite (``terra vermelha'') e proximidade com o
reservatório de Yguazú. Oferece um santuário de paz social e recursos
naturais abundantes, exigindo do ocupante autonomia logística para
serviços de saúde e comércio avançado, concentrados na vizinha Campo 9.

\subsubsection{Por que sim}

\begin{itemize}
\tightlist
\item
  Solo de altíssima produtividade e clima favorável para grãos e
  pecuária.
\item
  Baixíssima densidade demográfica e ambiente social extremamente
  pacífico.
\item
  Segurança hídrica regional garantida pela proximidade com o Lago
  Yguazú.
\item
  Sem restrições da Lei de Fronteira.
\end{itemize}

\subsubsection{Por que nao}

\begin{itemize}
\tightlist
\item
  Infraestrutura de saúde local limitada a atendimentos primários.
\item
  Dependência de estradas vicinais de terra para acesso às propriedades.
\item
  Valor da terra agrícola acima da média nacional devido à
  produtividade.
\end{itemize}

\subsubsection{Recomendacao Tecnica}

Altamente indicado para estabelecimentos de autossuficiência de médio
porte ou exploração agrícola que busque o máximo de isolamento
demográfico em solo de alta performance. Requer autonomia em logística
básica e saúde. O GSS de 6.5 reflete a segurança do isolamento em uma
região próspera. Referencia atual de score: GSS 6.5 em 2026-03-05.

\subsection{13.}

\begin{itemize}
\tightlist
\item
  \begin{enumerate}
  \def\labelenumi{\arabic{enumi})}
  \tightlist
  \item
    Dados censitários validados (INE\footnote{Instituto Nacional de Estadística (Instituto Nacional de Estatística), órgão oficial de dados demográficos e censitários.} 2022).
  \end{enumerate}
\item
  \begin{enumerate}
  \def\labelenumi{\arabic{enumi})}
  \setcounter{enumi}{1}
  \tightlist
  \item
    Registro de produtores de soja e trigo (MAG).
  \end{enumerate}
\item
  \begin{enumerate}
  \def\labelenumi{\arabic{enumi})}
  \setcounter{enumi}{2}
  \tightlist
  \item
    Mapa de solos e aptidão agrícola de Caaguazú (MAG).
  \end{enumerate}
\item
  \begin{enumerate}
  \def\labelenumi{\arabic{enumi})}
  \setcounter{enumi}{3}
  \tightlist
  \item
    Registro de segurança pública departamental.
  \end{enumerate}
\end{itemize}

\subsubsection{Combustível}

\textbf{Referência:} PETROPAR /\allowbreak  postos locais (2024) \textbf{Tipo de
localidade:} Interior

\begin{longtable}[]{@{}lll@{}}
\toprule\noalign{}
Combustível & USD/\allowbreak litro & Gs/\allowbreak litro (aprox.) \\
\midrule\noalign{}
\endhead
\bottomrule\noalign{}
\endlastfoot
Gasolina 93 oct & 0.98 & 7,252 \\
Gasolina 97 oct (premium) & 1.08 & 7,992 \\
Diesel & 0.91 & 6,734 \\
\end{longtable}

\begin{quote}
Preços podem variar ±5\% conforme posto e sazonalidade. Chaco e interior
remoto apresentam maior variação.
\end{quote}

\subsubsection{Cobertura Celular}

\textbf{Fonte:} CONATEL PY /\allowbreak  operadoras (2024)

\begin{longtable}[]{@{}ll@{}}
\toprule\noalign{}
Parâmetro & Valor \\
\midrule\noalign{}
\endhead
\bottomrule\noalign{}
\endlastfoot
Cobertura 4G população (dept.) & 88\% \\
Cobertura 4G área rural & 72\% \\
Melhor operadora & Tigo \\
Qualidade rural & boa \\
\end{longtable}

\begin{quote}
Para áreas rurais fora do núcleo urbano, recomenda-se chip Tigo como
principal e Personal como backup.
\end{quote}

\subsubsection{Internet}

\textbf{Fonte:} CONATEL /\allowbreak  Speedtest Ookla (2024)

\begin{longtable}[]{@{}ll@{}}
\toprule\noalign{}
Parâmetro & Valor \\
\midrule\noalign{}
\endhead
\bottomrule\noalign{}
\endlastfoot
Velocidade média download & 55 Mbps \\
Domicílios com internet (dept.) & 62\% \\
Tecnologia predominante & rádio/\allowbreak fibra \\
Opção rural & Starlink disponível (\textasciitilde USD 44/\allowbreak mês) \\
\end{longtable}

\subsubsection{Mercado Imobiliário e Terra
Rural}

\textbf{Fonte:} INDERT /\allowbreak  Clasificados.com.py (2024)

\begin{longtable}[]{@{}ll@{}}
\toprule\noalign{}
Tipo & Referência \\
\midrule\noalign{}
\endhead
\bottomrule\noalign{}
\endlastfoot
Terra agrícola alta prod. (USD/\allowbreak ha) & 8,800 \\
Imóvel urbano (USD/\allowbreak m²) & 850 \\
Aluguel 2 quartos (USD/\allowbreak mês) & 350 \\
\end{longtable}

\begin{quote}
Valores de referência departamental. Localidades menores podem ter
preços 20--40\% abaixo da capital departamental. \#\#\# Saúde
\end{quote}

\textbf{Fonte:} MSPBS /\allowbreak  IPS Paraguay (2024-2026), consolidação
departamental e proxy local

\begin{longtable}[]{@{}ll@{}}
\toprule\noalign{}
Serviço & Disponibilidade \\
\midrule\noalign{}
\endhead
\bottomrule\noalign{}
\endlastfoot
USF /\allowbreak  Posto de Saúde & sim \\
Hospital Regional & não \\
IPS (seguro social) & sim \\
Farmácia & sim \\
Distância ao hospital de referência & \textasciitilde111 km (Coronel
Oviedo) \\
\end{longtable}

\textbf{Principais estabelecimentos:} USF local; referencia hospitalar
em Coronel Oviedo

\textbf{Observação para imigrantes:} Atendimento primario local ou em
raio curto; casos de maior complexidade seguem para Coronel Oviedo.
Cobertura privada continua recomendada para especialidades e urgências
de maior complexidade.
\subsubsection{Indicadores Sociais}

\textbf{Fontes:} PNUD 2020, INE\footnote{Instituto Nacional de Estadística (Instituto Nacional de Estatística), órgão oficial de dados demográficos e censitários.} Censo 2022, INE\footnote{Instituto Nacional de Estadística (Instituto Nacional de Estatística), órgão oficial de dados demográficos e censitários.} EPHC 2023, Ministerio
Público 2024\\
\textbf{Nota:} Valores marcados como \emph{(dept.)} referem-se ao
departamento de Caaguazú; valores \emph{(dist.)} são específicos deste
distrito. Dados marcados \emph{(est.)} são estimativas.

\begin{longtable}[]{@{}
  >{\raggedright\arraybackslash}p{(\linewidth - 4\tabcolsep) * \real{0.4231}}
  >{\raggedright\arraybackslash}p{(\linewidth - 4\tabcolsep) * \real{0.2692}}
  >{\raggedright\arraybackslash}p{(\linewidth - 4\tabcolsep) * \real{0.3077}}@{}}
\toprule\noalign{}
\begin{minipage}[b]{\linewidth}\raggedright
Indicador
\end{minipage} & \begin{minipage}[b]{\linewidth}\raggedright
Valor
\end{minipage} & \begin{minipage}[b]{\linewidth}\raggedright
Âmbito
\end{minipage} \\
\midrule\noalign{}
\endhead
\bottomrule\noalign{}
\endlastfoot
IDH (2020) & 0,715 & dept. \\
Ranking IDH nacional & 7/\allowbreak 18 & dept. \\
Esperança de vida & 73,1 anos & dept. \\
Escolaridade média & 7.8 anos & dist. \\
RNB per capita (USD PPA) & 9.000 & dept. \\
Pobreza monetária (\%) & 28,5\% & dept. \\
Pobreza extrema (\%) & 6,8\% & dept. \\
Índice de Gini & 0,447 & dept. \\
Acesso a água potável (\%) & 79,8\% & dept. \\
Acesso a saneamento (\%) & 78,3\% & dept. \\
Taxa de homicídios (est., /\allowbreak 100k hab) & \textasciitilde6--8 (estimativa,
próxima ou acima da média nacional) & dept. \\
Índice de segurança & médio & dept. \\
População (Censo 2022) & 12.543 hab. & dist. \\
Idade mediana & 26.0 anos & dist. \\
Taxa de fecundidade & N/\allowbreak D & dist. \\
\end{longtable}
\subsubsection{Combustível}

\textbf{Referência:} PETROPAR /\allowbreak  postos locais (2024) \textbf{Tipo de
localidade:} Interior

\begin{longtable}[]{@{}lll@{}}
\toprule\noalign{}
Combustível & USD/\allowbreak litro & Gs/\allowbreak litro (aprox.) \\
\midrule\noalign{}
\endhead
\bottomrule\noalign{}
\endlastfoot
Gasolina 93 oct & 0.98 & 7,252 \\
Gasolina 97 oct (premium) & 1.08 & 7,992 \\
Diesel & 0.91 & 6,734 \\
\end{longtable}

\begin{quote}
Preços podem variar ±5\% conforme posto e sazonalidade. Chaco e interior
remoto apresentam maior variação.
\end{quote}

\subsubsection{Cobertura Celular}

\textbf{Fonte:} CONATEL PY /\allowbreak  operadoras (2024)

\begin{longtable}[]{@{}ll@{}}
\toprule\noalign{}
Parâmetro & Valor \\
\midrule\noalign{}
\endhead
\bottomrule\noalign{}
\endlastfoot
Cobertura 4G população (dept.) & 88\% \\
Cobertura 4G área rural & 72\% \\
Melhor operadora & Tigo \\
Qualidade rural & boa \\
\end{longtable}

\begin{quote}
Para áreas rurais fora do núcleo urbano, recomenda-se chip Tigo como
principal e Personal como backup.
\end{quote}

\subsubsection{Internet}

\textbf{Fonte:} CONATEL /\allowbreak  Speedtest Ookla (2024)

\begin{longtable}[]{@{}ll@{}}
\toprule\noalign{}
Parâmetro & Valor \\
\midrule\noalign{}
\endhead
\bottomrule\noalign{}
\endlastfoot
Velocidade média download & 55 Mbps \\
Domicílios com internet (dept.) & 62\% \\
Tecnologia predominante & rádio/\allowbreak fibra \\
Opção rural & Starlink disponível (\textasciitilde USD 44/\allowbreak mês) \\
\end{longtable}

\subsubsection{Mercado Imobiliário e Terra
Rural}

\textbf{Fonte:} INDERT /\allowbreak  Clasificados.com.py (2024)

\begin{longtable}[]{@{}ll@{}}
\toprule\noalign{}
Tipo & Referência \\
\midrule\noalign{}
\endhead
\bottomrule\noalign{}
\endlastfoot
Terra agrícola alta prod. (USD/\allowbreak ha) & 8,800 \\
Imóvel urbano (USD/\allowbreak m²) & 850 \\
Aluguel 2 quartos (USD/\allowbreak mês) & 350 \\
\end{longtable}

\begin{quote}
Valores de referência departamental. Localidades menores podem ter
preços 20--40\% abaixo da capital departamental. \#\#\# Saúde
\end{quote}

\textbf{Fonte:} MSPBS /\allowbreak  IPS Paraguay (2024-2026), consolidação
departamental e proxy local

\begin{longtable}[]{@{}ll@{}}
\toprule\noalign{}
Serviço & Disponibilidade \\
\midrule\noalign{}
\endhead
\bottomrule\noalign{}
\endlastfoot
USF /\allowbreak  Posto de Saúde & sim \\
Hospital Regional & não \\
IPS (seguro social) & sim \\
Farmácia & sim \\
Distância ao hospital de referência & \textasciitilde111 km (Coronel
Oviedo) \\
\end{longtable}

\textbf{Principais estabelecimentos:} USF local; referencia hospitalar
em Coronel Oviedo

\textbf{Observação para imigrantes:} Atendimento primario local ou em
raio curto; casos de maior complexidade seguem para Coronel Oviedo.
Cobertura privada continua recomendada para especialidades e urgências
de maior complexidade.
\newpage
\secaoDiagnostico{Repatriacion}{\mapaConteudo{-25.5166}{-55.9500}}{Repatriación é uma localidade que oferece um dos maiores potenciais de
privacidade demográfica e isolamento no departamento de Caaguazú. Sua
força reside na vasta extensão territorial com baixa densidade
populacional, combinada com uma segurança social superior e proximidade
estratégica com o polo de Caaguazú (cidade). É o local ideal para
projetos de autossuficiência de médio/\allowbreak largo porte que priorizam a
discrição e a abundância de recursos hídricos e de solo.

\subsubsection{Por que sim}

\begin{itemize}
\tightlist
\item
  Vasta área territorial com baixíssima densidade populacional
  (privacidade total).
\item
  Ambiente social extremamente seguro e coeso.
\item
  Proximidade com serviços em Caaguazú sem a exposição urbana da
  capital.
\item
  Sem restrições da Lei de Fronteira.
\end{itemize}

\subsubsection{Por que não}

\begin{itemize}
\tightlist
\item
  Infraestrutura de saúde local limitada a atendimentos primários.
\item
  Dependência de estradas vicinais que podem sofrer em chuvas severas.
\item
  Solo exige manejo técnico para evitar erosão em áreas mecanizadas.
\end{itemize}

\subsubsection{Recomendação
Técnica}

Altamente indicado para estabelecimentos de autossuficiência de
médio/\allowbreak grande porte que busquem isolamento demográfico em uma região
institucionalmente estável. Requer autonomia em termos de logística
básica e saúde. O GSS de 6.7 reflete a segurança do isolamento em
contraste com a limitação institucional local. Referencia atual de
score: GSS 6.7 em 2026-03-05.

\subsubsection{}

\begin{itemize}
\tightlist
\item
  \begin{enumerate}
  \def\labelenumi{\arabic{enumi})}
  \tightlist
  \item
    Dados censitários validados (INE\footnote{Instituto Nacional de Estadística (Instituto Nacional de Estatística), órgão oficial de dados demográficos e censitários.} 2022).
  \end{enumerate}
\item
  \begin{enumerate}
  \def\labelenumi{\arabic{enumi})}
  \setcounter{enumi}{1}
  \tightlist
  \item
    Levantamento territorial oficial do distrito (Portal
    Geoestadístico).
  \end{enumerate}
\item
  \begin{enumerate}
  \def\labelenumi{\arabic{enumi})}
  \setcounter{enumi}{2}
  \tightlist
  \item
    Mapa de solos e hidrografia de Caaguazú (MAG).
  \end{enumerate}
\item
  \begin{enumerate}
  \def\labelenumi{\arabic{enumi})}
  \setcounter{enumi}{3}
  \tightlist
  \item
    Registro de segurança pública departamental.
  \end{enumerate}
\end{itemize}}
\begin{table}[ht!]
\small \begin{tabular}{p{3.0cm}p{0.8cm}p{6.0cm}} \toprule
\textbf{Critério} & \textbf{Nota} & \textbf{Justificativa} \\ \midrule
A - Ameaças Estratégicas & 6.5 & Localização discreta e protegida de rotas principais; excelente invisibilidade tática \\
B - Risco Social & 7.5 & Baixíssima densidade demográfica em território vasto; isolamento social superior \\
C - Riscos Naturais & 7.0 & Geologia muito estável e riscos naturais climáticos típicos manejáveis \\
D - Autossuficiência & 7.0 & Bons recursos: abundância hídrica, gado e solos aptos para policultivos \\
E - Institucional & 5.5 & Ambiente institucional seguro, mas com infraestrutura pública local limitada \\
\midrule \textbf{GSS FINAL} & \textbf{6.7} & \textbf{Moderadamente Seguro (Ideal para quem busca o máximo de isolamento demográfico em uma região segura e produtiva).} \\ \bottomrule \end{tabular}
\end{table}
\subsubsection{Vulnerabilidades e
Mitigação}\label{vuln-Repatriacion-vulnerabilidades-e-mitigauxe7uxe3o}

\begin{itemize}
\tightlist
\item
  \textbf{Dependência de Serviços Externos:} Saúde especializada em
  Caaguazú (Mitigação: manutenção de estoque médico básico e kit trauma
  local).
\item
  \textbf{Logística Interna:} Acesso terrestre precário em chuvas
  intensas (Mitigação: estoque de suprimentos para 60 dias e posse de
  veículo 4x4 robusto).
\item
  \textbf{Instabilidade Elétrica Rural:} Quedas ocasionais em temporais
  (Mitigação: instalação de sistemas solares fotovoltaicos redundantes).
\end{itemize}
\subsubsection{Dossiê de Campo}\label{dossie-Repatriacion-dossiuxea-de-campo}
\begin{description}[style=nextline,leftmargin=0.5cm,font=\bfseries]
\item[Coordenadas]  25°31'S 55°57'W.
\item[Topografia]  Relevo predominantemente plano a suavemente ondulado. Localizado ao sul da cidade de Caaguazú. Altitude \textasciitilde 200m. Área vasta de \textasciitilde 1.012 km². Geologicamente muito estável, risco sísmico desprezível.
\item[Alvos Estratégicos]  Áreas de produção agropecuária e reflorestamento; Ponto de transição logística para o sul do departamento (bacia do Tebicuary). Ausência de alvos militares ou industriais primários, oferecendo excelente invisibilidade tática e discrição.
\item[Fallout]  Ventos predominantes do Norte (quentes e úmidos) e Leste. Baixo risco de fallout direto.
\item[População]  24.459 habitantes (Censo 2022 Final). População majoritariamente rural dispersa em diversas "companhias" e colônias.
\item[Densidade]  \textasciitilde 24,2 hab/\allowbreak km² (Baixa densidade, considerando a grande extensão territorial, oferecendo excelente isolamento tático).
\item[Vias de Acesso]  Acesso asfáltico conectando à cidade de Caaguazú e à Ruta PY02. Malha interna de estradas de terra batida que exigem manutenção contra erosão e veículos robustos em períodos chuvosos (média 1.600 mm/\allowbreak ano).
\item[Serviços]  Unidade de Saúde da Família (USF) local. Forte dependência da cidade de Caaguazú (10-15 km) para saúde de alta complexidade, bancos e logística avançada.
\item[Custo de Vida]  Muito baixo; economia baseada na agricultura familiar, pecuária e extração sustentável de madeira.
\item[Preço da Terra]  US\$ 2.500-5.000/\allowbreak ha (áreas rurais agrícolas/\allowbreak mistas); US\$ 1.200-2.000/\allowbreak ha (áreas de pastagem natural ou campo baixo).
\item[Hidrologia]  Diversos arroios locais (ex: Arroyo Capi'ibary). Baixo risco de inundações sistêmicas catastróficas. Risco de isolamento temporário de colônias por danos em pontes vicinais durante tempestades de verão.
\item[Clima]  Subtropical Úmido. Verões intensos e invernos com entradas ocasionais de ar polar.
\item[Energia]  Rede nacional (Itaipu); instável em áreas rurais remotas.
\item[Água]  Abundante via poços artesianos e lençol freático produtivo; sistemas de junta de saneamento em núcleos povoados.
\item[Qualidade do Solo]  Solo franco-arenoso a argiloso; fértil e profundo, ideal para mandioca, milho, gado e projetos de reflorestamento.
\item[Resources Locais]  Grande produção de mandioca, milho, gado de corte e madeira. Elevadíssimo potencial para autossuficiência alimentar básica em nível de propriedade.
\item[Segurança]  Zona rural tradicionalmente estável e pacífica. Baixíssima criminalidade urbana. Coesão social baseada na cultura agrícola e na identidade de "repatriados" (retornados de países vizinhos). Baixa visibilidade para conflitos de massa.
\item[Leis Local: Município consolidado. Livre de Restrição de Fronteira]  Localizado no interior do país, permitindo plena titularidade para investidores estrangeiros.
\end{description}
\textbf{Fonte climática:} NASA POWER Climatology API (período 2001-2020)
\textbf{Fonte luminosa:} estimativa world atlas

\paragraph{Irradiação Solar
(kWh/\allowbreak m²/\allowbreak dia)}\label{irradiauxe7uxe3o-solar-kwhmuxb2dia}

\begin{longtable}[]{@{}
  >{\raggedright\arraybackslash}p{(\linewidth - 24\tabcolsep) * \real{0.0746}}
  >{\raggedright\arraybackslash}p{(\linewidth - 24\tabcolsep) * \real{0.0746}}
  >{\raggedright\arraybackslash}p{(\linewidth - 24\tabcolsep) * \real{0.0746}}
  >{\raggedright\arraybackslash}p{(\linewidth - 24\tabcolsep) * \real{0.0746}}
  >{\raggedright\arraybackslash}p{(\linewidth - 24\tabcolsep) * \real{0.0746}}
  >{\raggedright\arraybackslash}p{(\linewidth - 24\tabcolsep) * \real{0.0746}}
  >{\raggedright\arraybackslash}p{(\linewidth - 24\tabcolsep) * \real{0.0746}}
  >{\raggedright\arraybackslash}p{(\linewidth - 24\tabcolsep) * \real{0.0746}}
  >{\raggedright\arraybackslash}p{(\linewidth - 24\tabcolsep) * \real{0.0746}}
  >{\raggedright\arraybackslash}p{(\linewidth - 24\tabcolsep) * \real{0.0746}}
  >{\raggedright\arraybackslash}p{(\linewidth - 24\tabcolsep) * \real{0.0746}}
  >{\raggedright\arraybackslash}p{(\linewidth - 24\tabcolsep) * \real{0.0746}}
  >{\raggedright\arraybackslash}p{(\linewidth - 24\tabcolsep) * \real{0.1045}}@{}}
\toprule\noalign{}
\begin{minipage}[b]{\linewidth}\raggedright
Jan
\end{minipage} & \begin{minipage}[b]{\linewidth}\raggedright
Fev
\end{minipage} & \begin{minipage}[b]{\linewidth}\raggedright
Mar
\end{minipage} & \begin{minipage}[b]{\linewidth}\raggedright
Abr
\end{minipage} & \begin{minipage}[b]{\linewidth}\raggedright
Mai
\end{minipage} & \begin{minipage}[b]{\linewidth}\raggedright
Jun
\end{minipage} & \begin{minipage}[b]{\linewidth}\raggedright
Jul
\end{minipage} & \begin{minipage}[b]{\linewidth}\raggedright
Ago
\end{minipage} & \begin{minipage}[b]{\linewidth}\raggedright
Set
\end{minipage} & \begin{minipage}[b]{\linewidth}\raggedright
Out
\end{minipage} & \begin{minipage}[b]{\linewidth}\raggedright
Nov
\end{minipage} & \begin{minipage}[b]{\linewidth}\raggedright
Dez
\end{minipage} & \begin{minipage}[b]{\linewidth}\raggedright
Média
\end{minipage} \\
\midrule\noalign{}
\endhead
\bottomrule\noalign{}
\endlastfoot
6.58 & 6.08 & 5.45 & 4.47 & 3.30 & 2.90 & 3.20 & 3.96 & 4.53 & 5.31 &
6.35 & 6.66 & \textbf{4.9} \\
\end{longtable}

\textbf{Inclinação solar recomendada:} 26° N (anual) · 36° N (inverno
jun-ago) · 16° N (verão nov-jan)

\paragraph{Precipitação (mm/\allowbreak mês)}\label{precipitauxe7uxe3o-mmmuxeas}

\begin{longtable}[]{@{}
  >{\raggedright\arraybackslash}p{(\linewidth - 24\tabcolsep) * \real{0.0704}}
  >{\raggedright\arraybackslash}p{(\linewidth - 24\tabcolsep) * \real{0.0704}}
  >{\raggedright\arraybackslash}p{(\linewidth - 24\tabcolsep) * \real{0.0704}}
  >{\raggedright\arraybackslash}p{(\linewidth - 24\tabcolsep) * \real{0.0704}}
  >{\raggedright\arraybackslash}p{(\linewidth - 24\tabcolsep) * \real{0.0704}}
  >{\raggedright\arraybackslash}p{(\linewidth - 24\tabcolsep) * \real{0.0704}}
  >{\raggedright\arraybackslash}p{(\linewidth - 24\tabcolsep) * \real{0.0704}}
  >{\raggedright\arraybackslash}p{(\linewidth - 24\tabcolsep) * \real{0.0704}}
  >{\raggedright\arraybackslash}p{(\linewidth - 24\tabcolsep) * \real{0.0704}}
  >{\raggedright\arraybackslash}p{(\linewidth - 24\tabcolsep) * \real{0.0704}}
  >{\raggedright\arraybackslash}p{(\linewidth - 24\tabcolsep) * \real{0.0704}}
  >{\raggedright\arraybackslash}p{(\linewidth - 24\tabcolsep) * \real{0.0704}}
  >{\raggedright\arraybackslash}p{(\linewidth - 24\tabcolsep) * \real{0.1549}}@{}}
\toprule\noalign{}
\begin{minipage}[b]{\linewidth}\raggedright
Jan
\end{minipage} & \begin{minipage}[b]{\linewidth}\raggedright
Fev
\end{minipage} & \begin{minipage}[b]{\linewidth}\raggedright
Mar
\end{minipage} & \begin{minipage}[b]{\linewidth}\raggedright
Abr
\end{minipage} & \begin{minipage}[b]{\linewidth}\raggedright
Mai
\end{minipage} & \begin{minipage}[b]{\linewidth}\raggedright
Jun
\end{minipage} & \begin{minipage}[b]{\linewidth}\raggedright
Jul
\end{minipage} & \begin{minipage}[b]{\linewidth}\raggedright
Ago
\end{minipage} & \begin{minipage}[b]{\linewidth}\raggedright
Set
\end{minipage} & \begin{minipage}[b]{\linewidth}\raggedright
Out
\end{minipage} & \begin{minipage}[b]{\linewidth}\raggedright
Nov
\end{minipage} & \begin{minipage}[b]{\linewidth}\raggedright
Dez
\end{minipage} & \begin{minipage}[b]{\linewidth}\raggedright
Total/\allowbreak ano
\end{minipage} \\
\midrule\noalign{}
\endhead
\bottomrule\noalign{}
\endlastfoot
131 & 136 & 131 & 152 & 172 & 90 & 75 & 52 & 92 & 187 & 190 & 166 &
\textbf{1577 mm} \\
\end{longtable}

\paragraph{Poluição Luminosa}\label{poluiuxe7uxe3o-luminosa}

\begin{longtable}[]{@{}ll@{}}
\toprule\noalign{}
Parâmetro & Valor \\
\midrule\noalign{}
\endhead
\bottomrule\noalign{}
\endlastfoot
Escala Bortle & 4 --- Céu rural-suburbano \\
Radiância artificial & 3.0 nW/\allowbreak cm²/\allowbreak sr \\
\end{longtable}
\paragraph{Solo (SoilGrids 2.0, média ponderada 0--30
cm)}

\begin{longtable}[]{@{}ll@{}}
\toprule\noalign{}
Parâmetro & Valor \\
\midrule\noalign{}
\endhead
\bottomrule\noalign{}
\endlastfoot
pH (H$_2$O) & 5.2 \\
Carbono orgânico (SOC) & 19.25 g/\allowbreak kg \\
Argila & 27.1 \% \\
Areia & 50.77 \% \\
Silte (calc.) & 22.1 \% \\
Densidade aparente & 1.27 g/\allowbreak cm³ \\
\textbf{Aptidão agrícola} & \textbf{Média} \\
\end{longtable}

Fonte: ISRIC SoilGrids 2.0 via WCS. Coords: -25.5167°, -55.95°. Média
ponderada camadas 0-5, 5-15, 15-30 cm.
\subsubsection{Indicadores Sociais}

\textbf{Fontes:} PNUD 2020, INE\footnote{Instituto Nacional de Estadística (Instituto Nacional de Estatística), órgão oficial de dados demográficos e censitários.} Censo 2022, INE\footnote{Instituto Nacional de Estadística (Instituto Nacional de Estatística), órgão oficial de dados demográficos e censitários.} EPHC 2023, Ministerio
Público 2024\\
\textbf{Nota:} Valores marcados como \emph{(dept.)} referem-se ao
departamento de Caaguazú; valores \emph{(dist.)} são específicos deste
distrito. Dados marcados \emph{(est.)} são estimativas.

\begin{longtable}[]{@{}
  >{\raggedright\arraybackslash}p{(\linewidth - 4\tabcolsep) * \real{0.4231}}
  >{\raggedright\arraybackslash}p{(\linewidth - 4\tabcolsep) * \real{0.2692}}
  >{\raggedright\arraybackslash}p{(\linewidth - 4\tabcolsep) * \real{0.3077}}@{}}
\toprule\noalign{}
\begin{minipage}[b]{\linewidth}\raggedright
Indicador
\end{minipage} & \begin{minipage}[b]{\linewidth}\raggedright
Valor
\end{minipage} & \begin{minipage}[b]{\linewidth}\raggedright
Âmbito
\end{minipage} \\
\midrule\noalign{}
\endhead
\bottomrule\noalign{}
\endlastfoot
IDH (2020) & 0,715 & dept. \\
Ranking IDH nacional & 7/\allowbreak 18 & dept. \\
Esperança de vida & 73,1 anos & dept. \\
Escolaridade média & 7.5 anos & dist. \\
RNB per capita (USD PPA) & 9.000 & dept. \\
Pobreza monetária (\%) & 28,5\% & dept. \\
Pobreza extrema (\%) & 6,8\% & dept. \\
Índice de Gini & 0,447 & dept. \\
Acesso a água potável (\%) & 79,8\% & dept. \\
Acesso a saneamento (\%) & 78,3\% & dept. \\
Taxa de homicídios (est., /\allowbreak 100k hab) & \textasciitilde6--8 (estimativa,
próxima ou acima da média nacional) & dept. \\
Índice de segurança & médio & dept. \\
População (Censo 2022) & 24.459 hab. & dist. \\
Idade mediana & 28.0 anos & dist. \\
Taxa de fecundidade & N/\allowbreak D & dist. \\
\end{longtable}
\subsubsection{Combustível}

\textbf{Referência:} PETROPAR /\allowbreak  postos locais (2024) \textbf{Tipo de
localidade:} Interior

\begin{longtable}[]{@{}lll@{}}
\toprule\noalign{}
Combustível & USD/\allowbreak litro & Gs/\allowbreak litro (aprox.) \\
\midrule\noalign{}
\endhead
\bottomrule\noalign{}
\endlastfoot
Gasolina 93 oct & 0.98 & 7,252 \\
Gasolina 97 oct (premium) & 1.08 & 7,992 \\
Diesel & 0.91 & 6,734 \\
\end{longtable}

\begin{quote}
Preços podem variar ±5\% conforme posto e sazonalidade. Chaco e interior
remoto apresentam maior variação.
\end{quote}

\subsubsection{Cobertura Celular}

\textbf{Fonte:} CONATEL PY /\allowbreak  operadoras (2024)

\begin{longtable}[]{@{}ll@{}}
\toprule\noalign{}
Parâmetro & Valor \\
\midrule\noalign{}
\endhead
\bottomrule\noalign{}
\endlastfoot
Cobertura 4G população (dept.) & 88\% \\
Cobertura 4G área rural & 72\% \\
Melhor operadora & Tigo \\
Qualidade rural & boa \\
\end{longtable}

\begin{quote}
Para áreas rurais fora do núcleo urbano, recomenda-se chip Tigo como
principal e Personal como backup.
\end{quote}

\subsubsection{Internet}

\textbf{Fonte:} CONATEL /\allowbreak  Speedtest Ookla (2024)

\begin{longtable}[]{@{}ll@{}}
\toprule\noalign{}
Parâmetro & Valor \\
\midrule\noalign{}
\endhead
\bottomrule\noalign{}
\endlastfoot
Velocidade média download & 55 Mbps \\
Domicílios com internet (dept.) & 62\% \\
Tecnologia predominante & rádio/\allowbreak fibra \\
Opção rural & Starlink disponível (\textasciitilde USD 44/\allowbreak mês) \\
\end{longtable}

\subsubsection{Mercado Imobiliário e Terra
Rural}

\textbf{Fonte:} INDERT /\allowbreak  Clasificados.com.py (2024)

\begin{longtable}[]{@{}ll@{}}
\toprule\noalign{}
Tipo & Referência \\
\midrule\noalign{}
\endhead
\bottomrule\noalign{}
\endlastfoot
Terra agrícola alta prod. (USD/\allowbreak ha) & 8,800 \\
Imóvel urbano (USD/\allowbreak m²) & 850 \\
Aluguel 2 quartos (USD/\allowbreak mês) & 350 \\
\end{longtable}

\begin{quote}
Valores de referência departamental. Localidades menores podem ter
preços 20--40\% abaixo da capital departamental. \#\#\# Saúde
\end{quote}

\textbf{Fonte:} MSPBS /\allowbreak  IPS Paraguay (2024-2026), consolidação
departamental e proxy local

\begin{longtable}[]{@{}ll@{}}
\toprule\noalign{}
Serviço & Disponibilidade \\
\midrule\noalign{}
\endhead
\bottomrule\noalign{}
\endlastfoot
USF /\allowbreak  Posto de Saúde & sim \\
Hospital Regional & não \\
IPS (seguro social) & sim \\
Farmácia & sim \\
Distância ao hospital de referência & \textasciitilde57 km (Coronel
Oviedo) \\
\end{longtable}

\textbf{Principais estabelecimentos:} USF local; referencia hospitalar
em Coronel Oviedo

\textbf{Observação para imigrantes:} Atendimento primario local ou em
raio curto; casos de maior complexidade seguem para Coronel Oviedo.
Cobertura privada continua recomendada para especialidades e urgências
de maior complexidade.
\newpage
\secaoDiagnostico{RI Tres Corrales}{\mapaConteudo{-25.2833}{-56.3333}}{RI Tres Corrales (Caaguazu) apresenta perfil operacional consolidado
para comparacao territorial, com leitura conjunta de infraestrutura,
risco hidroclimatico e capacidade de continuidade de servicos
essenciais.

\subsubsection{Por que sim}

\begin{itemize}
\tightlist
\item
  Base institucional de referencia disponivel e rastreavel.
\item
  Estrutura comparativa (A-E/\allowbreak GSS) consistente para decisao relativa.
\item
  Plano de mitigacao acionavel ja definido no dossie.
\end{itemize}

\subsubsection{Por que não}

\begin{itemize}
\tightlist
\item
  Serie oficial distrital granular ainda pode apresentar lacunas
  temporais.
\item
  Sensibilidade do score exige monitoramento continuo de risco e
  conectividade.
\item
  Decisao final depende de validacao de campo e janela temporal recente.
\end{itemize}

\subsubsection{Pre-condicoes Minimas de
Decisao}

\begin{itemize}
\tightlist
\item
  Atualizar indicadores criticos em ciclo trimestral.
\item
  Confirmar trafegabilidade e contingencia hidrica/\allowbreak energetica local.
\item
  Validar checklist de resiliencia domiciliar/\allowbreak logistica antes de decisao
  final.
\end{itemize}

\subsubsection{Recomendação
Técnica}

Uso recomendado como candidato em analise multicriterio, condicionado ao
cumprimento das pre-condicoes acima. Referencia atual de score: GSS 7.1
em 2026-03-05; risco predominante: baixo.

\subsubsection{}

\begin{itemize}
\tightlist
\item
  \begin{enumerate}
  \def\labelenumi{\arabic{enumi})}
  \tightlist
  \item
    Delimitacao territorial validada por georreferencia institucional
    (Fonte: acesso: 2026-03-05).
  \end{enumerate}
\item
  \begin{enumerate}
  \def\labelenumi{\arabic{enumi})}
  \setcounter{enumi}{1}
  \tightlist
  \item
    População municipal/\allowbreak distrital referenciada por base censitaria
    nacional (Fonte: acesso: 2026-03-05).
  \end{enumerate}
\item
  \begin{enumerate}
  \def\labelenumi{\arabic{enumi})}
  \setcounter{enumi}{2}
  \tightlist
  \item
    Comparabilidade intra-departamental apoiada em indicadores
    distritais oficiais (Fonte: acesso: 2026-03-05).
  \end{enumerate}
\item
  \begin{enumerate}
  \def\labelenumi{\arabic{enumi})}
  \setcounter{enumi}{3}
  \tightlist
  \item
    Estado macro de conectividade viaria ancorado em informacoes de
    rodovias (Fonte: acesso: 2026-03-05).
  \end{enumerate}
\item
  \begin{enumerate}
  \def\labelenumi{\arabic{enumi})}
  \setcounter{enumi}{4}
  \tightlist
  \item
    Exposicao hidrometeorologica monitorada por avisos oficiais (Fonte:
    acesso: 2026-03-05).
  \end{enumerate}
\item
  \begin{enumerate}
  \def\labelenumi{\arabic{enumi})}
  \setcounter{enumi}{5}
  \tightlist
  \item
    Historico de contexto meteorologico apoiado em publicacoes tecnicas
    (Fonte: acesso: 2026-03-05).
  \end{enumerate}
\item
  \begin{enumerate}
  \def\labelenumi{\arabic{enumi})}
  \setcounter{enumi}{6}
  \tightlist
  \item
    Capacidade institucional de resposta a eventos apoiada em acoes da
    SEN\footnote{Secretaría de Emergencia Nacional (Secretaria de Emergência Nacional), órgão governamental de resposta a desastres.} (Fonte: acesso: 2026-03-05).
  \end{enumerate}
\item
  \begin{enumerate}
  \def\labelenumi{\arabic{enumi})}
  \setcounter{enumi}{7}
  \tightlist
  \item
    Infraestrutura energetica analisada via operador nacional (Fonte:
    acesso: 2026-03-05).
  \end{enumerate}
\item
  \begin{enumerate}
  \def\labelenumi{\arabic{enumi})}
  \setcounter{enumi}{8}
  \tightlist
  \item
    Referencias de obras e logistica nacional verificadas no MOPC
    (Fonte: acesso: 2026-03-05).
  \end{enumerate}
\item
  \begin{enumerate}
  \def\labelenumi{\arabic{enumi})}
  \setcounter{enumi}{9}
  \tightlist
  \item
    Contexto sociopolitico institucional referenciado por autoridade
    eleitoral (Fonte: acesso: 2026-03-05).
  \end{enumerate}
\item
  \begin{enumerate}
  \def\labelenumi{\arabic{enumi})}
  \setcounter{enumi}{10}
  \tightlist
  \item
    Camada cartografica de apoio eleitoral/\allowbreak territorial consultada
    (Fonte: acesso: 2026-03-05).
  \end{enumerate}
\item
  \begin{enumerate}
  \def\labelenumi{\arabic{enumi})}
  \setcounter{enumi}{11}
  \tightlist
  \item
    Dados publicos complementares para verificacao cruzada (Fonte:
    acesso: 2026-03-05).
  \end{enumerate}
\end{itemize}}
\begin{table}[ht!]
\small \begin{tabular}{p{3.0cm}p{0.8cm}p{6.0cm}} \toprule
\textbf{Critério} & \textbf{Nota} & \textbf{Justificativa} \\ \midrule
A - Ameaças Estratégicas & 7.9 & - \\
B - Risco Social & 6.0 & - \\
C - Riscos Naturais & 8.4 & - \\
D - Autossuficiência & 6.2 & - \\
E - Institucional & 7.2 & - \\
\midrule \textbf{GSS FINAL} & \textbf{7.1} & \textbf{Seguro (bom potencial com preparacao).} \\ \bottomrule \end{tabular}
\end{table}
\subsubsection{Vulnerabilidades e
Mitigação}\label{vuln-RI_Tres_Corrales-vulnerabilidades-e-mitigauxe7uxe3o}

\begin{itemize}
\tightlist
\item
  Exposicao hidrometeorologica controlada (\textless45/\allowbreak 100).
\item
  Conectividade viaria intermediaria (50-69/\allowbreak 100).
\item
  Estoque de contingencia para 14 dias (agua/\allowbreak alimentos/\allowbreak combustivel),
  priorizado quando risco hidro=20/\allowbreak 100.
\item
  Duplicar rotas/\allowbreak logistica quando conectividade viaria=51/\allowbreak 100 for
  \textless70; manter rota secundaria mapeada e testada trimestralmente.
\item
  Plano de energia resiliente com autonomia minima de 72h
  (geracao/\allowbreak backup), revisao semestral.
\end{itemize}
\subsubsection{Dossiê de Campo}\label{dossie-RI_Tres_Corrales-dossiuxea-de-campo}
\begin{description}[style=nextline,leftmargin=0.5cm,font=\bfseries]
\item[Coordenadas]  25°17'S 56°20'W (geocodificado via Nominatim).
\end{description}
\textbf{Fonte climática:} NASA POWER Climatology API (período 2001-2020)
\textbf{Fonte luminosa:} estimativa world atlas

\paragraph{Irradiação Solar
(kWh/\allowbreak m²/\allowbreak dia)}\label{irradiauxe7uxe3o-solar-kwhmuxb2dia}

\begin{longtable}[]{@{}
  >{\raggedright\arraybackslash}p{(\linewidth - 24\tabcolsep) * \real{0.0746}}
  >{\raggedright\arraybackslash}p{(\linewidth - 24\tabcolsep) * \real{0.0746}}
  >{\raggedright\arraybackslash}p{(\linewidth - 24\tabcolsep) * \real{0.0746}}
  >{\raggedright\arraybackslash}p{(\linewidth - 24\tabcolsep) * \real{0.0746}}
  >{\raggedright\arraybackslash}p{(\linewidth - 24\tabcolsep) * \real{0.0746}}
  >{\raggedright\arraybackslash}p{(\linewidth - 24\tabcolsep) * \real{0.0746}}
  >{\raggedright\arraybackslash}p{(\linewidth - 24\tabcolsep) * \real{0.0746}}
  >{\raggedright\arraybackslash}p{(\linewidth - 24\tabcolsep) * \real{0.0746}}
  >{\raggedright\arraybackslash}p{(\linewidth - 24\tabcolsep) * \real{0.0746}}
  >{\raggedright\arraybackslash}p{(\linewidth - 24\tabcolsep) * \real{0.0746}}
  >{\raggedright\arraybackslash}p{(\linewidth - 24\tabcolsep) * \real{0.0746}}
  >{\raggedright\arraybackslash}p{(\linewidth - 24\tabcolsep) * \real{0.0746}}
  >{\raggedright\arraybackslash}p{(\linewidth - 24\tabcolsep) * \real{0.1045}}@{}}
\toprule\noalign{}
\begin{minipage}[b]{\linewidth}\raggedright
Jan
\end{minipage} & \begin{minipage}[b]{\linewidth}\raggedright
Fev
\end{minipage} & \begin{minipage}[b]{\linewidth}\raggedright
Mar
\end{minipage} & \begin{minipage}[b]{\linewidth}\raggedright
Abr
\end{minipage} & \begin{minipage}[b]{\linewidth}\raggedright
Mai
\end{minipage} & \begin{minipage}[b]{\linewidth}\raggedright
Jun
\end{minipage} & \begin{minipage}[b]{\linewidth}\raggedright
Jul
\end{minipage} & \begin{minipage}[b]{\linewidth}\raggedright
Ago
\end{minipage} & \begin{minipage}[b]{\linewidth}\raggedright
Set
\end{minipage} & \begin{minipage}[b]{\linewidth}\raggedright
Out
\end{minipage} & \begin{minipage}[b]{\linewidth}\raggedright
Nov
\end{minipage} & \begin{minipage}[b]{\linewidth}\raggedright
Dez
\end{minipage} & \begin{minipage}[b]{\linewidth}\raggedright
Média
\end{minipage} \\
\midrule\noalign{}
\endhead
\bottomrule\noalign{}
\endlastfoot
6.58 & 6.08 & 5.45 & 4.47 & 3.30 & 2.90 & 3.20 & 3.96 & 4.53 & 5.31 &
6.35 & 6.66 & \textbf{4.9} \\
\end{longtable}

\textbf{Inclinação solar recomendada:} 25° N (anual) · 35° N (inverno
jun-ago) · 15° N (verão nov-jan)

\paragraph{Precipitação (mm/\allowbreak mês)}\label{precipitauxe7uxe3o-mmmuxeas}

\begin{longtable}[]{@{}
  >{\raggedright\arraybackslash}p{(\linewidth - 24\tabcolsep) * \real{0.0704}}
  >{\raggedright\arraybackslash}p{(\linewidth - 24\tabcolsep) * \real{0.0704}}
  >{\raggedright\arraybackslash}p{(\linewidth - 24\tabcolsep) * \real{0.0704}}
  >{\raggedright\arraybackslash}p{(\linewidth - 24\tabcolsep) * \real{0.0704}}
  >{\raggedright\arraybackslash}p{(\linewidth - 24\tabcolsep) * \real{0.0704}}
  >{\raggedright\arraybackslash}p{(\linewidth - 24\tabcolsep) * \real{0.0704}}
  >{\raggedright\arraybackslash}p{(\linewidth - 24\tabcolsep) * \real{0.0704}}
  >{\raggedright\arraybackslash}p{(\linewidth - 24\tabcolsep) * \real{0.0704}}
  >{\raggedright\arraybackslash}p{(\linewidth - 24\tabcolsep) * \real{0.0704}}
  >{\raggedright\arraybackslash}p{(\linewidth - 24\tabcolsep) * \real{0.0704}}
  >{\raggedright\arraybackslash}p{(\linewidth - 24\tabcolsep) * \real{0.0704}}
  >{\raggedright\arraybackslash}p{(\linewidth - 24\tabcolsep) * \real{0.0704}}
  >{\raggedright\arraybackslash}p{(\linewidth - 24\tabcolsep) * \real{0.1549}}@{}}
\toprule\noalign{}
\begin{minipage}[b]{\linewidth}\raggedright
Jan
\end{minipage} & \begin{minipage}[b]{\linewidth}\raggedright
Fev
\end{minipage} & \begin{minipage}[b]{\linewidth}\raggedright
Mar
\end{minipage} & \begin{minipage}[b]{\linewidth}\raggedright
Abr
\end{minipage} & \begin{minipage}[b]{\linewidth}\raggedright
Mai
\end{minipage} & \begin{minipage}[b]{\linewidth}\raggedright
Jun
\end{minipage} & \begin{minipage}[b]{\linewidth}\raggedright
Jul
\end{minipage} & \begin{minipage}[b]{\linewidth}\raggedright
Ago
\end{minipage} & \begin{minipage}[b]{\linewidth}\raggedright
Set
\end{minipage} & \begin{minipage}[b]{\linewidth}\raggedright
Out
\end{minipage} & \begin{minipage}[b]{\linewidth}\raggedright
Nov
\end{minipage} & \begin{minipage}[b]{\linewidth}\raggedright
Dez
\end{minipage} & \begin{minipage}[b]{\linewidth}\raggedright
Total/\allowbreak ano
\end{minipage} \\
\midrule\noalign{}
\endhead
\bottomrule\noalign{}
\endlastfoot
142 & 130 & 126 & 152 & 177 & 94 & 79 & 55 & 103 & 189 & 188 & 172 &
\textbf{1607 mm} \\
\end{longtable}

\paragraph{Poluição Luminosa}\label{poluiuxe7uxe3o-luminosa}

\begin{longtable}[]{@{}ll@{}}
\toprule\noalign{}
Parâmetro & Valor \\
\midrule\noalign{}
\endhead
\bottomrule\noalign{}
\endlastfoot
Escala Bortle & 4 --- Céu rural-suburbano \\
Radiância artificial & 3.0 nW/\allowbreak cm²/\allowbreak sr \\
\end{longtable}
\paragraph{Solo (SoilGrids 2.0, média ponderada 0--30
cm)}

\begin{longtable}[]{@{}ll@{}}
\toprule\noalign{}
Parâmetro & Valor \\
\midrule\noalign{}
\endhead
\bottomrule\noalign{}
\endlastfoot
pH (H$_2$O) & 5.35 \\
Carbono orgânico (SOC) & 19.22 g/\allowbreak kg \\
Argila & 23.43 \% \\
Areia & 56.45 \% \\
Silte (calc.) & 20.1 \% \\
Densidade aparente & 1.25 g/\allowbreak cm³ \\
\textbf{Aptidão agrícola} & \textbf{Média} \\
\end{longtable}

Fonte: ISRIC SoilGrids 2.0 via WCS. Coords: -25.2833°, -56.3333°. Média
ponderada camadas 0-5, 5-15, 15-30 cm.
\subsubsection{Indicadores Sociais}

\textbf{Fontes:} PNUD 2020, INE\footnote{Instituto Nacional de Estadística (Instituto Nacional de Estatística), órgão oficial de dados demográficos e censitários.} Censo 2022, INE\footnote{Instituto Nacional de Estadística (Instituto Nacional de Estatística), órgão oficial de dados demográficos e censitários.} EPHC 2023, Ministerio
Público 2024\\
\textbf{Nota:} Valores marcados como \emph{(dept.)} referem-se ao
departamento de Caaguazú; valores \emph{(dist.)} são específicos deste
distrito. Dados marcados \emph{(est.)} são estimativas.

\begin{longtable}[]{@{}
  >{\raggedright\arraybackslash}p{(\linewidth - 4\tabcolsep) * \real{0.4231}}
  >{\raggedright\arraybackslash}p{(\linewidth - 4\tabcolsep) * \real{0.2692}}
  >{\raggedright\arraybackslash}p{(\linewidth - 4\tabcolsep) * \real{0.3077}}@{}}
\toprule\noalign{}
\begin{minipage}[b]{\linewidth}\raggedright
Indicador
\end{minipage} & \begin{minipage}[b]{\linewidth}\raggedright
Valor
\end{minipage} & \begin{minipage}[b]{\linewidth}\raggedright
Âmbito
\end{minipage} \\
\midrule\noalign{}
\endhead
\bottomrule\noalign{}
\endlastfoot
IDH (2020) & 0,715 & dept. \\
Ranking IDH nacional & 7/\allowbreak 18 & dept. \\
Esperança de vida & 73,1 anos & dept. \\
Escolaridade média & 8,4 anos & dept. \\
RNB per capita (USD PPA) & 9.000 & dept. \\
Pobreza monetária (\%) & 28,5\% & dept. \\
Pobreza extrema (\%) & 6,8\% & dept. \\
Índice de Gini & 0,447 & dept. \\
Acesso a água potável (\%) & 79,8\% & dept. \\
Acesso a saneamento (\%) & 78,3\% & dept. \\
Taxa de homicídios (est., /\allowbreak 100k hab) & \textasciitilde6--8 (estimativa,
próxima ou acima da média nacional) & dept. \\
Índice de segurança & médio & dept. \\
População (Censo 2022) & 6.021 hab. & dist. \\
Idade mediana & 30.0 anos & dist. \\
Taxa de fecundidade & N/\allowbreak D & dist. \\
\end{longtable}
\subsubsection{Combustível}

\textbf{Referência:} PETROPAR /\allowbreak  postos locais (2024) \textbf{Tipo de
localidade:} Interior

\begin{longtable}[]{@{}lll@{}}
\toprule\noalign{}
Combustível & USD/\allowbreak litro & Gs/\allowbreak litro (aprox.) \\
\midrule\noalign{}
\endhead
\bottomrule\noalign{}
\endlastfoot
Gasolina 93 oct & 0.98 & 7,252 \\
Gasolina 97 oct (premium) & 1.08 & 7,992 \\
Diesel & 0.91 & 6,734 \\
\end{longtable}

\begin{quote}
Preços podem variar ±5\% conforme posto e sazonalidade. Chaco e interior
remoto apresentam maior variação.
\end{quote}

\subsubsection{Cobertura Celular}

\textbf{Fonte:} CONATEL PY /\allowbreak  operadoras (2024)

\begin{longtable}[]{@{}ll@{}}
\toprule\noalign{}
Parâmetro & Valor \\
\midrule\noalign{}
\endhead
\bottomrule\noalign{}
\endlastfoot
Cobertura 4G população (dept.) & 88\% \\
Cobertura 4G área rural & 72\% \\
Melhor operadora & Tigo \\
Qualidade rural & boa \\
\end{longtable}

\begin{quote}
Para áreas rurais fora do núcleo urbano, recomenda-se chip Tigo como
principal e Personal como backup.
\end{quote}

\subsubsection{Internet}

\textbf{Fonte:} CONATEL /\allowbreak  Speedtest Ookla (2024)

\begin{longtable}[]{@{}ll@{}}
\toprule\noalign{}
Parâmetro & Valor \\
\midrule\noalign{}
\endhead
\bottomrule\noalign{}
\endlastfoot
Velocidade média download & 55 Mbps \\
Domicílios com internet (dept.) & 62\% \\
Tecnologia predominante & rádio/\allowbreak fibra \\
Opção rural & Starlink disponível (\textasciitilde USD 44/\allowbreak mês) \\
\end{longtable}

\subsubsection{Mercado Imobiliário e Terra
Rural}

\textbf{Fonte:} INDERT /\allowbreak  Clasificados.com.py (2024)

\begin{longtable}[]{@{}ll@{}}
\toprule\noalign{}
Tipo & Referência \\
\midrule\noalign{}
\endhead
\bottomrule\noalign{}
\endlastfoot
Terra agrícola alta prod. (USD/\allowbreak ha) & 8,800 \\
Imóvel urbano (USD/\allowbreak m²) & 850 \\
Aluguel 2 quartos (USD/\allowbreak mês) & 350 \\
\end{longtable}

\begin{quote}
Valores de referência departamental. Localidades menores podem ter
preços 20--40\% abaixo da capital departamental.
\end{quote}

\subsubsection{Saúde}

\textbf{Fonte:} MSPBS /\allowbreak  IPS Paraguay (2024)

\begin{longtable}[]{@{}ll@{}}
\toprule\noalign{}
Serviço & Disponibilidade \\
\midrule\noalign{}
\endhead
\bottomrule\noalign{}
\endlastfoot
USF /\allowbreak  Posto de Saúde & sim \\
Hospital Regional & não \\
IPS (seguro social) & não \\
Farmácia & sim \\
Distância ao hospital de referência & \textasciitilde60 km (Coronel
Oviedo) \\
\end{longtable}

\textbf{Principais estabelecimentos:} USF local; farmácias distritais.

\textbf{Observação para imigrantes:} Atendimento primário disponível.
Casos de maior complexidade encaminhados a Coronel Oviedo. Cobertura
privada recomendada; IPS não disponível localmente.
\newpage
\secaoDiagnostico{San Joaquin}{\mapaConteudo{-24.9666}{-56.1166}}{San Joaquín é um santuário de estabilidade e tradição no coração de
Caaguazú. Sua maior vantagem é a combinação de baixa visibilidade tática
com uma segurança social baseada em raízes históricas profundas e
coesas. Oferece abundância de água e terras aptas para pecuária e
policultivos, sendo um local ideal para quem busca um refúgio seguro e
discreto, aceitando a necessidade de autonomia em serviços de saúde e
suprimentos técnicos especializados.

\subsubsection{Por que sim}

\begin{itemize}
\tightlist
\item
  Ambiente social extremamente seguro e coeso baseado na cultura local.
\item
  Localização ``fora do radar'' das grandes crises demográficas da Ruta
  PY02.
\item
  Abundância de água superficial e solo fértil para pecuária e
  subsistência.
\item
  Sem restrições da Lei de Fronteira.
\end{itemize}

\subsubsection{Por que não}

\begin{itemize}
\tightlist
\item
  Infraestrutura de saúde local limitada a atendimentos básicos.
\item
  Dependência de estradas vicinais de terra para mobilidade interna.
\item
  Pequena escala comercial local exige deslocamento para centros
  maiores.
\end{itemize}

\subsubsection{Recomendação
Técnica}

Altamente indicado para estabelecimentos de autossuficiência familiar de
pequeno/\allowbreak médio porte que valorizem a segurança social e a paz rural.
Requer autonomia em logística básica e saúde. O GSS de 6.2 reflete a
segurança do isolamento geográfico e histórico em contraste com a
limitação institucional local. Referencia atual de score: GSS 6.2 em
2026-03-05.

\subsubsection{}

\begin{itemize}
\tightlist
\item
  \begin{enumerate}
  \def\labelenumi{\arabic{enumi})}
  \tightlist
  \item
    Dados censitários validados (INE\footnote{Instituto Nacional de Estadística (Instituto Nacional de Estatística), órgão oficial de dados demográficos e censitários.} 2022).
  \end{enumerate}
\item
  \begin{enumerate}
  \def\labelenumi{\arabic{enumi})}
  \setcounter{enumi}{1}
  \tightlist
  \item
    Relatórios de patrimônio histórico nacional (Secretaría Nacional de
    Cultura).
  \end{enumerate}
\item
  \begin{enumerate}
  \def\labelenumi{\arabic{enumi})}
  \setcounter{enumi}{2}
  \tightlist
  \item
    Mapa de solos e aptidão agropecuária de Caaguazú (MAG).
  \end{enumerate}
\item
  \begin{enumerate}
  \def\labelenumi{\arabic{enumi})}
  \setcounter{enumi}{3}
  \tightlist
  \item
    Registro de segurança pública departamental.
  \end{enumerate}
\end{itemize}}
\begin{table}[ht!]
\small \begin{tabular}{p{3.0cm}p{0.8cm}p{6.0cm}} \toprule
\textbf{Critério} & \textbf{Nota} & \textbf{Justificativa} \\ \midrule
A - Ameaças Estratégicas & 6.0 & Posicionamento estratégico discreto e protegido de fluxos rodoviários massivos; alta invisibilidade tática \\
B - Risco Social & 6.5 & População pequena e densidade controlada; baixíssimo risco social urbano \\
C - Riscos Naturais & 7.0 & Geologia muito estável e riscos naturais climáticos manejáveis \\
D - Autossuficiência & 6.5 & Bons recursos naturais - gado e água - mas limitado pela infraestrutura logística de acesso \\
E - Institucional & 5.0 & Município histórico seguro, mas com infraestrutura institucional e de saúde básica limitada \\
\midrule \textbf{GSS FINAL} & \textbf{6.2} & \textbf{Moderadamente Seguro (Ideal para quem busca isolamento rural estável em região historicamente segura).} \\ \bottomrule \end{tabular}
\end{table}
\subsubsection{Vulnerabilidades e
Mitigação}\label{vuln-San_Joaquin-vulnerabilidades-e-mitigauxe7uxe3o}

\begin{itemize}
\tightlist
\item
  \textbf{Isolamento de Saúde:} Distância significativa de hospitais
  cirúrgicos (Mitigação: manutenção de estoque médico básico e kit
  trauma local).
\item
  \textbf{Logística Vicinal:} Dependência de estradas de terra para
  acesso interno (Mitigação: estoque de suprimentos para 60 dias e posse
  de veículo 4x4 robusto).
\item
  \textbf{Vulnerabilidade Energética:} Rede ANDE\footnote{Administración Nacional de Electricidade (Administração Nacional de Eletricidade), companhia estatal de energia.} rural instável em
  temporais (Mitigação: instalação de sistemas solares fotovoltaicos
  redundantes).
\end{itemize}
\subsubsection{Dossiê de Campo}\label{dossie-San_Joaquin-dossiuxea-de-campo}
\begin{description}[style=nextline,leftmargin=0.5cm,font=\bfseries]
\item[Coordenadas]  24°58'S 56°07'W.
\item[Topografia]  Relevo ondulado ("lomadas") com vales férteis e arroios perenes. Altitude \textasciitilde 180m. Área de \textasciitilde 226 km². Geologicamente muito estável, risco sísmico desprezível.
\item[Alvos Estratégicos]  Igreja Jesuítica de San Joaquín (Único remanescente jesuítico habitado no Paraguai, alto valor histórico/\allowbreak patrimonial); Ponto de transição logística na Ruta PY13 (Ruta de la Mentira); Áreas de pecuária extensiva. Ausência de alvos militares de grande porte, oferecendo excelente invisibilidade estratégica.
\item[Fallout]  Ventos predominantes do Norte (quentes e úmidos) e Leste. Baixo risco de fallout direto de alvos continentais primários.
\item[População]  11.949 habitantes (Censo 2022 Final). Perfil demográfico maduro e essencialmente rural (85\%).
\item[Densidade]  \textasciitilde 52,9 hab/\allowbreak km² (Densidade moderada-baixa, oferecendo isolamento relativo em áreas de propriedades tradicionais).
\item[Vias de Acesso]  Acesso centralizado na Ruta PY13, que conecta o distrito a Vaquería e Caaguazú (cidade). Conectividade asfáltica recente que quebrou o isolamento histórico, mas mantém a localidade fora do radar dos grandes fluxos de carga internacional da PY02.
\item[Serviços]  Unidade de Saúde da Família (USF) local. Dependência de Caaguazú (50 km) ou Coronel Oviedo para atendimentos hospitalares de alta complexidade e serviços financeiros diversificados.
\item[Custo de Vida]  Muito baixo; economia baseada na pecuária de corte e agricultura familiar.
\item[Preço da Terra]  US\$ 3.000-5.500/\allowbreak ha (áreas rurais agrícolas); US\$ 1.500-2.500/\allowbreak ha (áreas de pastagem natural).
\item[Hidrologia]  Baixo risco de inundações sistêmicas catastróficas. Presença de arroios locais que favorecem a pequena pecuária. Risco moderado de isolamento de caminhos vicinais de terra em tempestades severas (média pluvial 1.600 mm/\allowbreak ano).
\item[Clima]  Subtropical Úmido. Verões quentes; invernos secos com entradas ocasionais de ar polar do Sul.
\item[Energia]  Rede nacional (Itaipu); instabilidade na rede rural devido ao isolamento geográfico.
\item[Água]  Abundante via poços artesianos e lençol freático superficial produtivo; juntas de saneamento em núcleos povoados.
\item[Qualidade do Solo]  Solo franco-arenoso a argiloso (avermelhado); fértil e profundo, ideal para pastagens, mandioca, milho e pequena produção de erva-mate.
\item[Resources Locais]  Grande produção de gado de corte, leite e alimentos básicos de subsistência. Elevado potencial para autossuficiência alimentar básica em nível comunitário.
\item[Segurança]  Zona extremamente pacífica e estável. Taxa de homicídios regional de \textasciitilde 8,5/\allowbreak 100k. Coesão social baseada na tradição histórica e religiosa (missão jesuítica). Baixíssima incidência de crimes violentos graves.
\item[Leis Local: Município histórico consolidado. Livre de Restrição de Fronteira]  Fora da faixa de 50km da fronteira, permitindo plena titularidade para investidores brasileiros.
\end{description}
\textbf{Fonte climática:} NASA POWER Climatology API (período 2001-2020)
\textbf{Fonte luminosa:} estimativa world atlas

\paragraph{Irradiação Solar
(kWh/\allowbreak m²/\allowbreak dia)}\label{irradiauxe7uxe3o-solar-kwhmuxb2dia}

\begin{longtable}[]{@{}
  >{\raggedright\arraybackslash}p{(\linewidth - 24\tabcolsep) * \real{0.0746}}
  >{\raggedright\arraybackslash}p{(\linewidth - 24\tabcolsep) * \real{0.0746}}
  >{\raggedright\arraybackslash}p{(\linewidth - 24\tabcolsep) * \real{0.0746}}
  >{\raggedright\arraybackslash}p{(\linewidth - 24\tabcolsep) * \real{0.0746}}
  >{\raggedright\arraybackslash}p{(\linewidth - 24\tabcolsep) * \real{0.0746}}
  >{\raggedright\arraybackslash}p{(\linewidth - 24\tabcolsep) * \real{0.0746}}
  >{\raggedright\arraybackslash}p{(\linewidth - 24\tabcolsep) * \real{0.0746}}
  >{\raggedright\arraybackslash}p{(\linewidth - 24\tabcolsep) * \real{0.0746}}
  >{\raggedright\arraybackslash}p{(\linewidth - 24\tabcolsep) * \real{0.0746}}
  >{\raggedright\arraybackslash}p{(\linewidth - 24\tabcolsep) * \real{0.0746}}
  >{\raggedright\arraybackslash}p{(\linewidth - 24\tabcolsep) * \real{0.0746}}
  >{\raggedright\arraybackslash}p{(\linewidth - 24\tabcolsep) * \real{0.0746}}
  >{\raggedright\arraybackslash}p{(\linewidth - 24\tabcolsep) * \real{0.1045}}@{}}
\toprule\noalign{}
\begin{minipage}[b]{\linewidth}\raggedright
Jan
\end{minipage} & \begin{minipage}[b]{\linewidth}\raggedright
Fev
\end{minipage} & \begin{minipage}[b]{\linewidth}\raggedright
Mar
\end{minipage} & \begin{minipage}[b]{\linewidth}\raggedright
Abr
\end{minipage} & \begin{minipage}[b]{\linewidth}\raggedright
Mai
\end{minipage} & \begin{minipage}[b]{\linewidth}\raggedright
Jun
\end{minipage} & \begin{minipage}[b]{\linewidth}\raggedright
Jul
\end{minipage} & \begin{minipage}[b]{\linewidth}\raggedright
Ago
\end{minipage} & \begin{minipage}[b]{\linewidth}\raggedright
Set
\end{minipage} & \begin{minipage}[b]{\linewidth}\raggedright
Out
\end{minipage} & \begin{minipage}[b]{\linewidth}\raggedright
Nov
\end{minipage} & \begin{minipage}[b]{\linewidth}\raggedright
Dez
\end{minipage} & \begin{minipage}[b]{\linewidth}\raggedright
Média
\end{minipage} \\
\midrule\noalign{}
\endhead
\bottomrule\noalign{}
\endlastfoot
6.74 & 6.20 & 5.50 & 4.48 & 3.35 & 2.88 & 3.21 & 3.96 & 4.59 & 5.43 &
6.43 & 6.81 & \textbf{4.96} \\
\end{longtable}

\textbf{Inclinação solar recomendada:} 25° N (anual) · 35° N (inverno
jun-ago) · 15° N (verão nov-jan)

\paragraph{Precipitação (mm/\allowbreak mês)}\label{precipitauxe7uxe3o-mmmuxeas}

\begin{longtable}[]{@{}
  >{\raggedright\arraybackslash}p{(\linewidth - 24\tabcolsep) * \real{0.0704}}
  >{\raggedright\arraybackslash}p{(\linewidth - 24\tabcolsep) * \real{0.0704}}
  >{\raggedright\arraybackslash}p{(\linewidth - 24\tabcolsep) * \real{0.0704}}
  >{\raggedright\arraybackslash}p{(\linewidth - 24\tabcolsep) * \real{0.0704}}
  >{\raggedright\arraybackslash}p{(\linewidth - 24\tabcolsep) * \real{0.0704}}
  >{\raggedright\arraybackslash}p{(\linewidth - 24\tabcolsep) * \real{0.0704}}
  >{\raggedright\arraybackslash}p{(\linewidth - 24\tabcolsep) * \real{0.0704}}
  >{\raggedright\arraybackslash}p{(\linewidth - 24\tabcolsep) * \real{0.0704}}
  >{\raggedright\arraybackslash}p{(\linewidth - 24\tabcolsep) * \real{0.0704}}
  >{\raggedright\arraybackslash}p{(\linewidth - 24\tabcolsep) * \real{0.0704}}
  >{\raggedright\arraybackslash}p{(\linewidth - 24\tabcolsep) * \real{0.0704}}
  >{\raggedright\arraybackslash}p{(\linewidth - 24\tabcolsep) * \real{0.0704}}
  >{\raggedright\arraybackslash}p{(\linewidth - 24\tabcolsep) * \real{0.1549}}@{}}
\toprule\noalign{}
\begin{minipage}[b]{\linewidth}\raggedright
Jan
\end{minipage} & \begin{minipage}[b]{\linewidth}\raggedright
Fev
\end{minipage} & \begin{minipage}[b]{\linewidth}\raggedright
Mar
\end{minipage} & \begin{minipage}[b]{\linewidth}\raggedright
Abr
\end{minipage} & \begin{minipage}[b]{\linewidth}\raggedright
Mai
\end{minipage} & \begin{minipage}[b]{\linewidth}\raggedright
Jun
\end{minipage} & \begin{minipage}[b]{\linewidth}\raggedright
Jul
\end{minipage} & \begin{minipage}[b]{\linewidth}\raggedright
Ago
\end{minipage} & \begin{minipage}[b]{\linewidth}\raggedright
Set
\end{minipage} & \begin{minipage}[b]{\linewidth}\raggedright
Out
\end{minipage} & \begin{minipage}[b]{\linewidth}\raggedright
Nov
\end{minipage} & \begin{minipage}[b]{\linewidth}\raggedright
Dez
\end{minipage} & \begin{minipage}[b]{\linewidth}\raggedright
Total/\allowbreak ano
\end{minipage} \\
\midrule\noalign{}
\endhead
\bottomrule\noalign{}
\endlastfoot
124 & 143 & 125 & 147 & 173 & 90 & 70 & 46 & 94 & 184 & 190 & 168 &
\textbf{1557 mm} \\
\end{longtable}

\paragraph{Poluição Luminosa}\label{poluiuxe7uxe3o-luminosa}

\begin{longtable}[]{@{}ll@{}}
\toprule\noalign{}
Parâmetro & Valor \\
\midrule\noalign{}
\endhead
\bottomrule\noalign{}
\endlastfoot
Escala Bortle & 4 --- Céu rural-suburbano \\
Radiância artificial & 3.0 nW/\allowbreak cm²/\allowbreak sr \\
\end{longtable}
\paragraph{Solo (SoilGrids 2.0, média ponderada 0--30
cm)}

\begin{longtable}[]{@{}ll@{}}
\toprule\noalign{}
Parâmetro & Valor \\
\midrule\noalign{}
\endhead
\bottomrule\noalign{}
\endlastfoot
pH (H$_2$O) & 5.4 \\
Carbono orgânico (SOC) & 19.54 g/\allowbreak kg \\
Argila & 23.3 \% \\
Areia & 57.87 \% \\
Silte (calc.) & 18.8 \% \\
Densidade aparente & 1.23 g/\allowbreak cm³ \\
\textbf{Aptidão agrícola} & \textbf{Média} \\
\end{longtable}

Fonte: ISRIC SoilGrids 2.0 via WCS. Coords: -24.9667°, -56.1167°. Média
ponderada camadas 0-5, 5-15, 15-30 cm.
\subsubsection{Indicadores Sociais}

\textbf{Fontes:} PNUD 2020, INE\footnote{Instituto Nacional de Estadística (Instituto Nacional de Estatística), órgão oficial de dados demográficos e censitários.} Censo 2022, INE\footnote{Instituto Nacional de Estadística (Instituto Nacional de Estatística), órgão oficial de dados demográficos e censitários.} EPHC 2023, Ministerio
Público 2024\\
\textbf{Nota:} Valores marcados como \emph{(dept.)} referem-se ao
departamento de Caaguazú; valores \emph{(dist.)} são específicos deste
distrito. Dados marcados \emph{(est.)} são estimativas.

\begin{longtable}[]{@{}
  >{\raggedright\arraybackslash}p{(\linewidth - 4\tabcolsep) * \real{0.4231}}
  >{\raggedright\arraybackslash}p{(\linewidth - 4\tabcolsep) * \real{0.2692}}
  >{\raggedright\arraybackslash}p{(\linewidth - 4\tabcolsep) * \real{0.3077}}@{}}
\toprule\noalign{}
\begin{minipage}[b]{\linewidth}\raggedright
Indicador
\end{minipage} & \begin{minipage}[b]{\linewidth}\raggedright
Valor
\end{minipage} & \begin{minipage}[b]{\linewidth}\raggedright
Âmbito
\end{minipage} \\
\midrule\noalign{}
\endhead
\bottomrule\noalign{}
\endlastfoot
IDH (2020) & 0,715 & dept. \\
Ranking IDH nacional & 7/\allowbreak 18 & dept. \\
Esperança de vida & 73,1 anos & dept. \\
Escolaridade média & 7.3 anos & dist. \\
RNB per capita (USD PPA) & 9.000 & dept. \\
Pobreza monetária (\%) & 28,5\% & dept. \\
Pobreza extrema (\%) & 6,8\% & dept. \\
Índice de Gini & 0,447 & dept. \\
Acesso a água potável (\%) & 79,8\% & dept. \\
Acesso a saneamento (\%) & 78,3\% & dept. \\
Taxa de homicídios (est., /\allowbreak 100k hab) & \textasciitilde6--8 (estimativa,
próxima ou acima da média nacional) & dept. \\
Índice de segurança & médio & dept. \\
População (Censo 2022) & 11.949 hab. & dist. \\
Idade mediana & 28.0 anos & dist. \\
Taxa de fecundidade & N/\allowbreak D & dist. \\
\end{longtable}
\subsubsection{Combustível}

\textbf{Referência:} PETROPAR /\allowbreak  postos locais (2024) \textbf{Tipo de
localidade:} Interior

\begin{longtable}[]{@{}lll@{}}
\toprule\noalign{}
Combustível & USD/\allowbreak litro & Gs/\allowbreak litro (aprox.) \\
\midrule\noalign{}
\endhead
\bottomrule\noalign{}
\endlastfoot
Gasolina 93 oct & 0.98 & 7,252 \\
Gasolina 97 oct (premium) & 1.08 & 7,992 \\
Diesel & 0.91 & 6,734 \\
\end{longtable}

\begin{quote}
Preços podem variar ±5\% conforme posto e sazonalidade. Chaco e interior
remoto apresentam maior variação.
\end{quote}

\subsubsection{Cobertura Celular}

\textbf{Fonte:} CONATEL PY /\allowbreak  operadoras (2024)

\begin{longtable}[]{@{}ll@{}}
\toprule\noalign{}
Parâmetro & Valor \\
\midrule\noalign{}
\endhead
\bottomrule\noalign{}
\endlastfoot
Cobertura 4G população (dept.) & 88\% \\
Cobertura 4G área rural & 72\% \\
Melhor operadora & Tigo \\
Qualidade rural & boa \\
\end{longtable}

\begin{quote}
Para áreas rurais fora do núcleo urbano, recomenda-se chip Tigo como
principal e Personal como backup.
\end{quote}

\subsubsection{Internet}

\textbf{Fonte:} CONATEL /\allowbreak  Speedtest Ookla (2024)

\begin{longtable}[]{@{}ll@{}}
\toprule\noalign{}
Parâmetro & Valor \\
\midrule\noalign{}
\endhead
\bottomrule\noalign{}
\endlastfoot
Velocidade média download & 55 Mbps \\
Domicílios com internet (dept.) & 62\% \\
Tecnologia predominante & rádio/\allowbreak fibra \\
Opção rural & Starlink disponível (\textasciitilde USD 44/\allowbreak mês) \\
\end{longtable}

\subsubsection{Mercado Imobiliário e Terra
Rural}

\textbf{Fonte:} INDERT /\allowbreak  Clasificados.com.py (2024)

\begin{longtable}[]{@{}ll@{}}
\toprule\noalign{}
Tipo & Referência \\
\midrule\noalign{}
\endhead
\bottomrule\noalign{}
\endlastfoot
Terra agrícola alta prod. (USD/\allowbreak ha) & 8,800 \\
Imóvel urbano (USD/\allowbreak m²) & 850 \\
Aluguel 2 quartos (USD/\allowbreak mês) & 350 \\
\end{longtable}

\begin{quote}
Valores de referência departamental. Localidades menores podem ter
preços 20--40\% abaixo da capital departamental. \#\#\# Saúde
\end{quote}

\textbf{Fonte:} MSPBS /\allowbreak  IPS Paraguay (2024-2026), consolidação
departamental e proxy local

\begin{longtable}[]{@{}ll@{}}
\toprule\noalign{}
Serviço & Disponibilidade \\
\midrule\noalign{}
\endhead
\bottomrule\noalign{}
\endlastfoot
USF /\allowbreak  Posto de Saúde & sim \\
Hospital Regional & sim \\
IPS (seguro social) & sim \\
Farmácia & sim \\
Distância ao hospital de referência & \textasciitilde77 km (Coronel
Oviedo) \\
\end{longtable}

\textbf{Principais estabelecimentos:} USF local; referencia hospitalar
em Coronel Oviedo

\textbf{Observação para imigrantes:} Atendimento primario local ou em
raio curto; casos de maior complexidade seguem para Coronel Oviedo.
Cobertura privada continua recomendada para especialidades e urgências
de maior complexidade.
\newpage
\secaoDiagnostico{San Jose de los Arroyos}{\mapaConteudo{-25.3833}{-56.7333}}{San José de los Arroyos é uma localidade estratégica devido à sua
posição na principal artéria logística do país. Possui baixa densidade
populacional, o que favorece o isolamento em cenários de crise, mas sua
proximidade com a capital exige vigilância sobre fluxos migratórios.

\subsubsection{Por que sim}

\begin{itemize}
\tightlist
\item
  Conectividade logística excepcional.
\item
  Baixa densidade populacional (segurança social).
\item
  Solo fértil para subsistência.
\end{itemize}

\subsubsection{Por que não}

\begin{itemize}
\tightlist
\item
  Exposição direta ao tráfego nacional intenso.
\item
  Dependência de serviços de saúde externos para casos graves.
\end{itemize}

\subsubsection{Recomendação
Técnica}

Candidato viável para quem prioriza logística e acesso fácil, desde que
possua infraestrutura própria de defesa e estoque para mitigar a
exposição à rodovia principal.

\subsubsection{Analise de Sensibilidade}

\begin{itemize}
\tightlist
\item
  Cenario otimista: GSS 6.7 (Melhoria em serviços locais).
\item
  Cenario conservador: GSS 5.8 (Instabilidade na rota PY02).
\end{itemize}

\subsubsection{}

\begin{itemize}
\tightlist
\item
  Dados populacionais validados pelo Censo 2022 (INE\footnote{Instituto Nacional de Estadística (Instituto Nacional de Estatística), órgão oficial de dados demográficos e censitários.}).
\item
  Infraestrutura rodoviária confirmada pelo MOPC.
\item
  Segurança departamental referenciada pelo Ministério do Interior.
\end{itemize}}
\begin{table}[ht!]
\small \begin{tabular}{p{3.0cm}p{0.8cm}p{6.0cm}} \toprule
\textbf{Critério} & \textbf{Nota} & \textbf{Justificativa} \\ \midrule
A - Ameaças Estratégicas & 6.2 & - \\
B - Risco Social & 6.0 & - \\
C - Riscos Naturais & 7.0 & - \\
D - Autossuficiência & 6.0 & - \\
E - Institucional & 6.5 & - \\
\midrule \textbf{GSS FINAL} & \textbf{6.3} & \textbf{Seguro (Potencial moderado de resiliência).} \\ \bottomrule \end{tabular}
\end{table}
\subsubsection{Vulnerabilidades e
Mitigação}\label{vuln-San_Jose_de_los_Arroyos-vulnerabilidades-e-mitigauxe7uxe3o}

\begin{itemize}
\item
  \textbf{Fluxo Rodoviário:} A Ruta PY02 é um ativo e um risco
  (bloqueios/\allowbreak refugiados). Mitigação: Mapear rotas secundárias por
  caminhos de terra.
\item
  \textbf{Abastecimento:} Dependência de logística externa para itens
  processados. Mitigação: Estoque de 30 dias.
\item
  \textbf{Saúde:} Distância de centros cirúrgicos. Mitigação: Kit de
  primeiros socorros avançado e telemedicina.
\item
  INE\footnote{Instituto Nacional de Estadística (Instituto Nacional de Estatística), órgão oficial de dados demográficos e censitários.} Censo 2022: 13.926 habitantes.
\item
  Preço da Terra: Gs 30M a 60M por hectare (Agro/\allowbreak Pecuária).
\item
  Homicídios: 3,5/\allowbreak 100k (Caaguazú).
\item
  População distrital: 13.926 habitantes (Censo 2022).
\item
  Densidade demografica: 15,7 hab/\allowbreak km2.
\item
  Conectividade viaria: 90/\allowbreak 100 (Acesso direto à PY02).
\item
  Exposicao hidrometeorologica: 40/\allowbreak 100.
\item
  Curto prazo: Saturação da Ruta PY02 em caso de instabilidade nacional.
\item
  Médio prazo: Mudanças climáticas afetando a produção de
  cana-de-açúcar.
\end{itemize}

\subsubsection{Diagnostico Integrado}\label{vuln-San_Jose_de_los_Arroyos-diagnostico-integrado}

San José de los Arroyos é uma localidade estratégica devido à sua
posição na principal artéria logística do país. Possui baixa densidade
populacional, o que favorece o isolamento em cenários de crise, mas sua
proximidade com a capital exige vigilância sobre fluxos migratórios.

\subsubsection{Por que sim}\label{vuln-San_Jose_de_los_Arroyos-por-que-sim}

\begin{itemize}
\tightlist
\item
  Conectividade logística excepcional.
\item
  Baixa densidade populacional (segurança social).
\item
  Solo fértil para subsistência.
\end{itemize}

\subsubsection{Por que não}\label{vuln-San_Jose_de_los_Arroyos-por-que-nuxe3o}

\begin{itemize}
\tightlist
\item
  Exposição direta ao tráfego nacional intenso.
\item
  Dependência de serviços de saúde externos para casos graves.
\end{itemize}

\subsubsection{Recomendação
Técnica}\label{vuln-San_Jose_de_los_Arroyos-recomendauxe7uxe3o-tuxe9cnica}

Candidato viável para quem prioriza logística e acesso fácil, desde que
possua infraestrutura própria de defesa e estoque para mitigar a
exposição à rodovia principal.

\begin{itemize}
\item
  Cenario otimista: GSS 6.7 (Melhoria em serviços locais).
\item
  Cenario conservador: GSS 5.8 (Instabilidade na rota PY02).
\item
  Dados populacionais validados pelo Censo 2022 (INE\footnote{Instituto Nacional de Estadística (Instituto Nacional de Estatística), órgão oficial de dados demográficos e censitários.}).
\item
  Infraestrutura rodoviária confirmada pelo MOPC.
\item
  Segurança departamental referenciada pelo Ministério do Interior.
\end{itemize}

\subsubsection{Combustível}\label{vuln-San_Jose_de_los_Arroyos-combustuxedvel}

\textbf{Referência:} PETROPAR /\allowbreak  postos locais (2024) \textbf{Tipo de
localidade:} Interior

\begin{longtable}[]{@{}lll@{}}
\toprule\noalign{}
Combustível & USD/\allowbreak litro & Gs/\allowbreak litro (aprox.) \\
\midrule\noalign{}
\endhead
\bottomrule\noalign{}
\endlastfoot
Gasolina 93 oct & 0.98 & 7,252 \\
Gasolina 97 oct (premium) & 1.08 & 7,992 \\
Diesel & 0.91 & 6,734 \\
\end{longtable}

\begin{quote}
Preços podem variar ±5\% conforme posto e sazonalidade. Chaco e interior
remoto apresentam maior variação.
\end{quote}

\subsubsection{Cobertura Celular}\label{vuln-San_Jose_de_los_Arroyos-cobertura-celular}

\textbf{Fonte:} CONATEL PY /\allowbreak  operadoras (2024)

\begin{longtable}[]{@{}ll@{}}
\toprule\noalign{}
Parâmetro & Valor \\
\midrule\noalign{}
\endhead
\bottomrule\noalign{}
\endlastfoot
Cobertura 4G população (dept.) & 88\% \\
Cobertura 4G área rural & 72\% \\
Melhor operadora & Tigo \\
Qualidade rural & boa \\
\end{longtable}

\begin{quote}
Para áreas rurais fora do núcleo urbano, recomenda-se chip Tigo como
principal e Personal como backup.
\end{quote}

\subsubsection{Internet}\label{vuln-San_Jose_de_los_Arroyos-internet}

\textbf{Fonte:} CONATEL /\allowbreak  Speedtest Ookla (2024)

\begin{longtable}[]{@{}ll@{}}
\toprule\noalign{}
Parâmetro & Valor \\
\midrule\noalign{}
\endhead
\bottomrule\noalign{}
\endlastfoot
Velocidade média download & 55 Mbps \\
Domicílios com internet (dept.) & 62\% \\
Tecnologia predominante & rádio/\allowbreak fibra \\
Opção rural & Starlink disponível (\textasciitilde USD 44/\allowbreak mês) \\
\end{longtable}

\subsubsection{Mercado Imobiliário e Terra
Rural}\label{vuln-San_Jose_de_los_Arroyos-mercado-imobiliuxe1rio-e-terra-rural}

\textbf{Fonte:} INDERT /\allowbreak  Clasificados.com.py (2024)

\begin{longtable}[]{@{}ll@{}}
\toprule\noalign{}
Tipo & Referência \\
\midrule\noalign{}
\endhead
\bottomrule\noalign{}
\endlastfoot
Terra agrícola alta prod. (USD/\allowbreak ha) & 8,800 \\
Imóvel urbano (USD/\allowbreak m²) & 850 \\
Aluguel 2 quartos (USD/\allowbreak mês) & 350 \\
\end{longtable}

\begin{quote}
Valores de referência departamental. Localidades menores podem ter
preços 20--40\% abaixo da capital departamental. \#\#\# Saúde
\end{quote}

\textbf{Fonte:} MSPBS /\allowbreak  IPS Paraguay (2024-2026), consolidação
departamental e proxy local

\begin{longtable}[]{@{}ll@{}}
\toprule\noalign{}
Serviço & Disponibilidade \\
\midrule\noalign{}
\endhead
\bottomrule\noalign{}
\endlastfoot
USF /\allowbreak  Posto de Saúde & sim \\
Hospital Regional & não \\
IPS (seguro social) & sim \\
Farmácia & sim \\
Distância ao hospital de referência & \textasciitilde38 km (Coronel
Oviedo) \\
\end{longtable}

\textbf{Principais estabelecimentos:} USF local; referencia hospitalar
em Coronel Oviedo

\textbf{Observação para imigrantes:} Atendimento primario local ou em
raio curto; casos de maior complexidade seguem para Coronel Oviedo.
Cobertura privada continua recomendada para especialidades e urgências
de maior complexidade.
\subsubsection{Dossiê de Campo}\label{dossie-San_Jose_de_los_Arroyos-dossiuxea-de-campo}
\begin{description}[style=nextline,leftmargin=0.5cm,font=\bfseries]
\item[Coordenadas]  25°23'S 56°44'W.
\item[Topografia]  Relevo predominantemente plano a suavemente ondulado. Área de \textasciitilde 887 km².
\item[Alvos Estrategicos]  Ponto de transito e logistica na Ruta PY02. Areas de producao intensiva de cana-de-acucar. Ausencia de alvos militares diretos.
\item[Fallout]  Ventos predominantes do Norte/\allowbreak Nordeste.
\item[Populacao]  13.926 (Censo 2022 Final). População significativamente menor que as projeções estatísticas anteriores.
\item[Densidade]  \textasciitilde 15,70 hab/\allowbreak km² (Baixa densidade, excelente para segurança social e isolamento relativo).
\item[Vias de Acesso]  Localização estratégica sobre a Ruta PY02 (Eixo principal Assunção - Ciudad del Este). Conexão direta com o departamento de Cordillera.
\item[Servicos]  Infraestrutura básica de saúde e educação. Dependência de Coronel Oviedo ou da zona metropolitana de Assunção (\textasciitilde 100 km) para serviços complexos.
\item[Hidrologia]  Cruzado por diversos arroios. Baixo risco de inundações graves na zona urbana. Risco moderado de danos em acessos rurais em eventos de chuva extrema.
\item[Sismicidade]  Baixa sismicidade (intraplaca), sem registros significativos locais.
\item[Solo]  Alfisois e Ultisois predominantes, fertilidade media a alta.
\item[Agua]  Poços artesianos e sistemas de distribuição local (Juntas de Saneamiento).
\item[Energia]  Rede nacional (ANDE\footnote{Administración Nacional de Electricidade (Administração Nacional de Eletricidade), companhia estatal de energia.} - Itaipu).
\item[Producao]  Cana-de-açúcar, gado e agricultura de subsistência.
\item[Seguranca]  Taxa de homicídios departamental de 3,5/\allowbreak 100k hab (2022). Baixa criminalidade urbana.
\item[Clima Social]  Comunidade agrícola tradicional, estabilidade social elevada.
\end{description}
\textbf{Fonte climática:} NASA POWER Climatology API (período 2001-2020)
\textbf{Fonte luminosa:} estimativa world atlas

\paragraph{Irradiação Solar
(kWh/\allowbreak m²/\allowbreak dia)}\label{irradiauxe7uxe3o-solar-kwhmuxb2dia}

\begin{longtable}[]{@{}
  >{\raggedright\arraybackslash}p{(\linewidth - 24\tabcolsep) * \real{0.0746}}
  >{\raggedright\arraybackslash}p{(\linewidth - 24\tabcolsep) * \real{0.0746}}
  >{\raggedright\arraybackslash}p{(\linewidth - 24\tabcolsep) * \real{0.0746}}
  >{\raggedright\arraybackslash}p{(\linewidth - 24\tabcolsep) * \real{0.0746}}
  >{\raggedright\arraybackslash}p{(\linewidth - 24\tabcolsep) * \real{0.0746}}
  >{\raggedright\arraybackslash}p{(\linewidth - 24\tabcolsep) * \real{0.0746}}
  >{\raggedright\arraybackslash}p{(\linewidth - 24\tabcolsep) * \real{0.0746}}
  >{\raggedright\arraybackslash}p{(\linewidth - 24\tabcolsep) * \real{0.0746}}
  >{\raggedright\arraybackslash}p{(\linewidth - 24\tabcolsep) * \real{0.0746}}
  >{\raggedright\arraybackslash}p{(\linewidth - 24\tabcolsep) * \real{0.0746}}
  >{\raggedright\arraybackslash}p{(\linewidth - 24\tabcolsep) * \real{0.0746}}
  >{\raggedright\arraybackslash}p{(\linewidth - 24\tabcolsep) * \real{0.0746}}
  >{\raggedright\arraybackslash}p{(\linewidth - 24\tabcolsep) * \real{0.1045}}@{}}
\toprule\noalign{}
\begin{minipage}[b]{\linewidth}\raggedright
Jan
\end{minipage} & \begin{minipage}[b]{\linewidth}\raggedright
Fev
\end{minipage} & \begin{minipage}[b]{\linewidth}\raggedright
Mar
\end{minipage} & \begin{minipage}[b]{\linewidth}\raggedright
Abr
\end{minipage} & \begin{minipage}[b]{\linewidth}\raggedright
Mai
\end{minipage} & \begin{minipage}[b]{\linewidth}\raggedright
Jun
\end{minipage} & \begin{minipage}[b]{\linewidth}\raggedright
Jul
\end{minipage} & \begin{minipage}[b]{\linewidth}\raggedright
Ago
\end{minipage} & \begin{minipage}[b]{\linewidth}\raggedright
Set
\end{minipage} & \begin{minipage}[b]{\linewidth}\raggedright
Out
\end{minipage} & \begin{minipage}[b]{\linewidth}\raggedright
Nov
\end{minipage} & \begin{minipage}[b]{\linewidth}\raggedright
Dez
\end{minipage} & \begin{minipage}[b]{\linewidth}\raggedright
Média
\end{minipage} \\
\midrule\noalign{}
\endhead
\bottomrule\noalign{}
\endlastfoot
6.74 & 6.20 & 5.50 & 4.48 & 3.35 & 2.88 & 3.21 & 3.96 & 4.59 & 5.43 &
6.43 & 6.81 & \textbf{4.96} \\
\end{longtable}

\textbf{Inclinação solar recomendada:} 26° N (anual) · 36° N (inverno
jun-ago) · 16° N (verão nov-jan)

\paragraph{Precipitação (mm/\allowbreak mês)}\label{precipitauxe7uxe3o-mmmuxeas}

\begin{longtable}[]{@{}
  >{\raggedright\arraybackslash}p{(\linewidth - 24\tabcolsep) * \real{0.0704}}
  >{\raggedright\arraybackslash}p{(\linewidth - 24\tabcolsep) * \real{0.0704}}
  >{\raggedright\arraybackslash}p{(\linewidth - 24\tabcolsep) * \real{0.0704}}
  >{\raggedright\arraybackslash}p{(\linewidth - 24\tabcolsep) * \real{0.0704}}
  >{\raggedright\arraybackslash}p{(\linewidth - 24\tabcolsep) * \real{0.0704}}
  >{\raggedright\arraybackslash}p{(\linewidth - 24\tabcolsep) * \real{0.0704}}
  >{\raggedright\arraybackslash}p{(\linewidth - 24\tabcolsep) * \real{0.0704}}
  >{\raggedright\arraybackslash}p{(\linewidth - 24\tabcolsep) * \real{0.0704}}
  >{\raggedright\arraybackslash}p{(\linewidth - 24\tabcolsep) * \real{0.0704}}
  >{\raggedright\arraybackslash}p{(\linewidth - 24\tabcolsep) * \real{0.0704}}
  >{\raggedright\arraybackslash}p{(\linewidth - 24\tabcolsep) * \real{0.0704}}
  >{\raggedright\arraybackslash}p{(\linewidth - 24\tabcolsep) * \real{0.0704}}
  >{\raggedright\arraybackslash}p{(\linewidth - 24\tabcolsep) * \real{0.1549}}@{}}
\toprule\noalign{}
\begin{minipage}[b]{\linewidth}\raggedright
Jan
\end{minipage} & \begin{minipage}[b]{\linewidth}\raggedright
Fev
\end{minipage} & \begin{minipage}[b]{\linewidth}\raggedright
Mar
\end{minipage} & \begin{minipage}[b]{\linewidth}\raggedright
Abr
\end{minipage} & \begin{minipage}[b]{\linewidth}\raggedright
Mai
\end{minipage} & \begin{minipage}[b]{\linewidth}\raggedright
Jun
\end{minipage} & \begin{minipage}[b]{\linewidth}\raggedright
Jul
\end{minipage} & \begin{minipage}[b]{\linewidth}\raggedright
Ago
\end{minipage} & \begin{minipage}[b]{\linewidth}\raggedright
Set
\end{minipage} & \begin{minipage}[b]{\linewidth}\raggedright
Out
\end{minipage} & \begin{minipage}[b]{\linewidth}\raggedright
Nov
\end{minipage} & \begin{minipage}[b]{\linewidth}\raggedright
Dez
\end{minipage} & \begin{minipage}[b]{\linewidth}\raggedright
Total/\allowbreak ano
\end{minipage} \\
\midrule\noalign{}
\endhead
\bottomrule\noalign{}
\endlastfoot
127 & 140 & 133 & 148 & 152 & 78 & 66 & 41 & 77 & 171 & 193 & 175 &
\textbf{1506 mm} \\
\end{longtable}

\paragraph{Poluição Luminosa}\label{poluiuxe7uxe3o-luminosa}

\begin{longtable}[]{@{}ll@{}}
\toprule\noalign{}
Parâmetro & Valor \\
\midrule\noalign{}
\endhead
\bottomrule\noalign{}
\endlastfoot
Escala Bortle & 4 --- Céu rural-suburbano \\
Radiância artificial & 3.0 nW/\allowbreak cm²/\allowbreak sr \\
\end{longtable}
\paragraph{Solo (SoilGrids 2.0, média ponderada 0--30
cm)}

\begin{longtable}[]{@{}ll@{}}
\toprule\noalign{}
Parâmetro & Valor \\
\midrule\noalign{}
\endhead
\bottomrule\noalign{}
\endlastfoot
pH (H$_2$O) & 5.4 \\
Carbono orgânico (SOC) & 17.12 g/\allowbreak kg \\
Argila & 30.97 \% \\
Areia & 48.71 \% \\
Silte (calc.) & 20.3 \% \\
Densidade aparente & 1.35 g/\allowbreak cm³ \\
\textbf{Aptidão agrícola} & \textbf{Média} \\
\end{longtable}

Fonte: ISRIC SoilGrids 2.0 via WCS. Coords: -25.3833°, -56.7333°. Média
ponderada camadas 0-5, 5-15, 15-30 cm.

\begin{itemize}
\tightlist
\item
  \textbf{Agua:} Poços artesianos e sistemas de distribuição local
  (Juntas de Saneamiento).
\item
  \textbf{Energia:} Rede nacional (ANDE\footnote{Administración Nacional de Electricidade (Administração Nacional de Eletricidade), companhia estatal de energia.} - Itaipu).
\item
  \textbf{Producao:} Cana-de-açúcar, gado e agricultura de subsistência.
\end{itemize}
\subsubsection{Indicadores Sociais}

\textbf{Fontes:} PNUD 2020, INE\footnote{Instituto Nacional de Estadística (Instituto Nacional de Estatística), órgão oficial de dados demográficos e censitários.} Censo 2022, INE\footnote{Instituto Nacional de Estadística (Instituto Nacional de Estatística), órgão oficial de dados demográficos e censitários.} EPHC 2023, Ministerio
Público 2024\\
\textbf{Nota:} Valores marcados como \emph{(dept.)} referem-se ao
departamento de Caaguazú; valores \emph{(dist.)} são específicos deste
distrito. Dados marcados \emph{(est.)} são estimativas.

\begin{longtable}[]{@{}
  >{\raggedright\arraybackslash}p{(\linewidth - 4\tabcolsep) * \real{0.4231}}
  >{\raggedright\arraybackslash}p{(\linewidth - 4\tabcolsep) * \real{0.2692}}
  >{\raggedright\arraybackslash}p{(\linewidth - 4\tabcolsep) * \real{0.3077}}@{}}
\toprule\noalign{}
\begin{minipage}[b]{\linewidth}\raggedright
Indicador
\end{minipage} & \begin{minipage}[b]{\linewidth}\raggedright
Valor
\end{minipage} & \begin{minipage}[b]{\linewidth}\raggedright
Âmbito
\end{minipage} \\
\midrule\noalign{}
\endhead
\bottomrule\noalign{}
\endlastfoot
IDH (2020) & 0,715 & dept. \\
Ranking IDH nacional & 7/\allowbreak 18 & dept. \\
Esperança de vida & 73,1 anos & dept. \\
Escolaridade média & 9.2 anos & dist. \\
RNB per capita (USD PPA) & 9.000 & dept. \\
Pobreza monetária (\%) & 28,5\% & dept. \\
Pobreza extrema (\%) & 6,8\% & dept. \\
Índice de Gini & 0,447 & dept. \\
Acesso a água potável (\%) & 79,8\% & dept. \\
Acesso a saneamento (\%) & 78,3\% & dept. \\
Taxa de homicídios (est., /\allowbreak 100k hab) & \textasciitilde6--8 (estimativa,
próxima ou acima da média nacional) & dept. \\
Índice de segurança & médio & dept. \\
População (Censo 2022) & 13.926 hab. & dist. \\
Idade mediana & 33.0 anos & dist. \\
Taxa de fecundidade & N/\allowbreak D & dist. \\
\end{longtable}
\subsubsection{Combustível}

\textbf{Referência:} PETROPAR /\allowbreak  postos locais (2024) \textbf{Tipo de
localidade:} Interior

\begin{longtable}[]{@{}lll@{}}
\toprule\noalign{}
Combustível & USD/\allowbreak litro & Gs/\allowbreak litro (aprox.) \\
\midrule\noalign{}
\endhead
\bottomrule\noalign{}
\endlastfoot
Gasolina 93 oct & 0.98 & 7,252 \\
Gasolina 97 oct (premium) & 1.08 & 7,992 \\
Diesel & 0.91 & 6,734 \\
\end{longtable}

\begin{quote}
Preços podem variar ±5\% conforme posto e sazonalidade. Chaco e interior
remoto apresentam maior variação.
\end{quote}

\subsubsection{Cobertura Celular}

\textbf{Fonte:} CONATEL PY /\allowbreak  operadoras (2024)

\begin{longtable}[]{@{}ll@{}}
\toprule\noalign{}
Parâmetro & Valor \\
\midrule\noalign{}
\endhead
\bottomrule\noalign{}
\endlastfoot
Cobertura 4G população (dept.) & 88\% \\
Cobertura 4G área rural & 72\% \\
Melhor operadora & Tigo \\
Qualidade rural & boa \\
\end{longtable}

\begin{quote}
Para áreas rurais fora do núcleo urbano, recomenda-se chip Tigo como
principal e Personal como backup.
\end{quote}

\subsubsection{Internet}

\textbf{Fonte:} CONATEL /\allowbreak  Speedtest Ookla (2024)

\begin{longtable}[]{@{}ll@{}}
\toprule\noalign{}
Parâmetro & Valor \\
\midrule\noalign{}
\endhead
\bottomrule\noalign{}
\endlastfoot
Velocidade média download & 55 Mbps \\
Domicílios com internet (dept.) & 62\% \\
Tecnologia predominante & rádio/\allowbreak fibra \\
Opção rural & Starlink disponível (\textasciitilde USD 44/\allowbreak mês) \\
\end{longtable}

\subsubsection{Mercado Imobiliário e Terra
Rural}

\textbf{Fonte:} INDERT /\allowbreak  Clasificados.com.py (2024)

\begin{longtable}[]{@{}ll@{}}
\toprule\noalign{}
Tipo & Referência \\
\midrule\noalign{}
\endhead
\bottomrule\noalign{}
\endlastfoot
Terra agrícola alta prod. (USD/\allowbreak ha) & 8,800 \\
Imóvel urbano (USD/\allowbreak m²) & 850 \\
Aluguel 2 quartos (USD/\allowbreak mês) & 350 \\
\end{longtable}

\begin{quote}
Valores de referência departamental. Localidades menores podem ter
preços 20--40\% abaixo da capital departamental. \#\#\# Saúde
\end{quote}

\textbf{Fonte:} MSPBS /\allowbreak  IPS Paraguay (2024-2026), consolidação
departamental e proxy local

\begin{longtable}[]{@{}ll@{}}
\toprule\noalign{}
Serviço & Disponibilidade \\
\midrule\noalign{}
\endhead
\bottomrule\noalign{}
\endlastfoot
USF /\allowbreak  Posto de Saúde & sim \\
Hospital Regional & não \\
IPS (seguro social) & sim \\
Farmácia & sim \\
Distância ao hospital de referência & \textasciitilde38 km (Coronel
Oviedo) \\
\end{longtable}

\textbf{Principais estabelecimentos:} USF local; referencia hospitalar
em Coronel Oviedo

\textbf{Observação para imigrantes:} Atendimento primario local ou em
raio curto; casos de maior complexidade seguem para Coronel Oviedo.
Cobertura privada continua recomendada para especialidades e urgências
de maior complexidade.
\newpage
\secaoDiagnostico{Santa Rosa del Mbutuy}{\mapaConteudo{-25.0666}{-56.3000}}{Santa Rosa del Mbutuy é um distrito equilibrado, com densidade
demográfica baixa e perfil agrícola resiliente. Sua posição na Ruta PY08
facilita a logística mas requer atenção em cenários de instabilidade
regional.

\subsubsection{Por que sim}

\begin{itemize}
\tightlist
\item
  Perfil comunitário sólido.
\item
  Baixo risco de inundações graves.
\end{itemize}

\subsubsection{Por que não}

\begin{itemize}
\tightlist
\item
  Dependência excessiva de centros urbanos maiores para infraestrutura
  crítica.
\end{itemize}

\subsubsection{Recomendação
Técnica}

Recomendado para pequenos grupos ou famílias que buscam autossuficiência
agrícola em uma região tranquila e acessível.

\subsubsection{Analise de Sensibilidade}

\begin{itemize}
\tightlist
\item
  Cenario otimista: GSS 6.2 (Pavimentação de rotas secundárias).
\item
  Cenario conservador: GSS 5.3 (Degradação da infraestrutura
  rodoviária).
\end{itemize}

\subsubsection{}

\begin{itemize}
\tightlist
\item
  Dados populacionais INE\footnote{Instituto Nacional de Estadística (Instituto Nacional de Estatística), órgão oficial de dados demográficos e censitários.} 2022.
\item
  Mapeamento viário MOPC.
\item
  Relatórios de segurança regional.
\end{itemize}}
\begin{table}[ht!]
\small \begin{tabular}{p{3.0cm}p{0.8cm}p{6.0cm}} \toprule
\textbf{Critério} & \textbf{Nota} & \textbf{Justificativa} \\ \midrule
A - Ameaças Estratégicas & 5.0 & - \\
B - Risco Social & 5.5 & - \\
C - Riscos Naturais & 7.0 & - \\
D - Autossuficiência & 6.0 & - \\
E - Institucional & 6.0 & - \\
\midrule \textbf{GSS FINAL} & \textbf{5.8} & \textbf{Seguro (Potencial moderado de resiliência).} \\ \bottomrule \end{tabular}
\end{table}
\subsubsection{Vulnerabilidades e
Mitigação}\label{vuln-Santa_Rosa_del_Mbutuy-vulnerabilidades-e-mitigauxe7uxe3o}

\begin{itemize}
\item
  \textbf{Isolamento Rural:} Estradas de terra vulneráveis a chuvas.
  Mitigação: Veículos 4x4 e manutenção de acessos.
\item
  \textbf{Saúde:} Carência de UTI local. Mitigação: Convênio com centros
  em Coronel Oviedo.
\item
  INE\footnote{Instituto Nacional de Estadística (Instituto Nacional de Estatística), órgão oficial de dados demográficos e censitários.} Censo 2022: 8.467 habitantes.
\item
  Solo: Alfisois/\allowbreak Ultisois.
\item
  Homicídios: 3,5/\allowbreak 100k.
\item
  População distrital: 8.467 habitantes.
\item
  Densidade demografica: 28,4 hab/\allowbreak km2.
\item
  Conectividade viaria: 75/\allowbreak 100.
\item
  Exposicao hidrometeorologica: 35/\allowbreak 100.
\item
  Curto prazo: Impacto de secas sazonais na agricultura de subsistência.
\item
  Médio prazo: Expansão de monoculturas mecanizadas pressionando
  pequenas propriedades.
\end{itemize}

\subsubsection{Diagnostico Integrado}\label{vuln-Santa_Rosa_del_Mbutuy-diagnostico-integrado}

Santa Rosa del Mbutuy é um distrito equilibrado, com densidade
demográfica baixa e perfil agrícola resiliente. Sua posição na Ruta PY08
facilita a logística mas requer atenção em cenários de instabilidade
regional.

\subsubsection{Por que sim}\label{vuln-Santa_Rosa_del_Mbutuy-por-que-sim}

\begin{itemize}
\tightlist
\item
  Perfil comunitário sólido.
\item
  Baixo risco de inundações graves.
\end{itemize}

\subsubsection{Por que não}\label{vuln-Santa_Rosa_del_Mbutuy-por-que-nuxe3o}

\begin{itemize}
\tightlist
\item
  Dependência excessiva de centros urbanos maiores para infraestrutura
  crítica.
\end{itemize}

\subsubsection{Recomendação
Técnica}\label{vuln-Santa_Rosa_del_Mbutuy-recomendauxe7uxe3o-tuxe9cnica}

Recomendado para pequenos grupos ou famílias que buscam autossuficiência
agrícola em uma região tranquila e acessível.

\begin{itemize}
\item
  Cenario otimista: GSS 6.2 (Pavimentação de rotas secundárias).
\item
  Cenario conservador: GSS 5.3 (Degradação da infraestrutura
  rodoviária).
\item
  Dados populacionais INE\footnote{Instituto Nacional de Estadística (Instituto Nacional de Estatística), órgão oficial de dados demográficos e censitários.} 2022.
\item
  Mapeamento viário MOPC.
\item
  Relatórios de segurança regional.
\end{itemize}

\subsubsection{Combustível}\label{vuln-Santa_Rosa_del_Mbutuy-combustuxedvel}

\textbf{Referência:} PETROPAR /\allowbreak  postos locais (2024) \textbf{Tipo de
localidade:} Interior

\begin{longtable}[]{@{}lll@{}}
\toprule\noalign{}
Combustível & USD/\allowbreak litro & Gs/\allowbreak litro (aprox.) \\
\midrule\noalign{}
\endhead
\bottomrule\noalign{}
\endlastfoot
Gasolina 93 oct & 0.98 & 7,252 \\
Gasolina 97 oct (premium) & 1.08 & 7,992 \\
Diesel & 0.91 & 6,734 \\
\end{longtable}

\begin{quote}
Preços podem variar ±5\% conforme posto e sazonalidade. Chaco e interior
remoto apresentam maior variação.
\end{quote}

\subsubsection{Cobertura Celular}\label{vuln-Santa_Rosa_del_Mbutuy-cobertura-celular}

\textbf{Fonte:} CONATEL PY /\allowbreak  operadoras (2024)

\begin{longtable}[]{@{}ll@{}}
\toprule\noalign{}
Parâmetro & Valor \\
\midrule\noalign{}
\endhead
\bottomrule\noalign{}
\endlastfoot
Cobertura 4G população (dept.) & 88\% \\
Cobertura 4G área rural & 72\% \\
Melhor operadora & Tigo \\
Qualidade rural & boa \\
\end{longtable}

\begin{quote}
Para áreas rurais fora do núcleo urbano, recomenda-se chip Tigo como
principal e Personal como backup.
\end{quote}

\subsubsection{Internet}\label{vuln-Santa_Rosa_del_Mbutuy-internet}

\textbf{Fonte:} CONATEL /\allowbreak  Speedtest Ookla (2024)

\begin{longtable}[]{@{}ll@{}}
\toprule\noalign{}
Parâmetro & Valor \\
\midrule\noalign{}
\endhead
\bottomrule\noalign{}
\endlastfoot
Velocidade média download & 55 Mbps \\
Domicílios com internet (dept.) & 62\% \\
Tecnologia predominante & rádio/\allowbreak fibra \\
Opção rural & Starlink disponível (\textasciitilde USD 44/\allowbreak mês) \\
\end{longtable}

\subsubsection{Mercado Imobiliário e Terra
Rural}\label{vuln-Santa_Rosa_del_Mbutuy-mercado-imobiliuxe1rio-e-terra-rural}

\textbf{Fonte:} INDERT /\allowbreak  Clasificados.com.py (2024)

\begin{longtable}[]{@{}ll@{}}
\toprule\noalign{}
Tipo & Referência \\
\midrule\noalign{}
\endhead
\bottomrule\noalign{}
\endlastfoot
Terra agrícola alta prod. (USD/\allowbreak ha) & 8,800 \\
Imóvel urbano (USD/\allowbreak m²) & 850 \\
Aluguel 2 quartos (USD/\allowbreak mês) & 350 \\
\end{longtable}

\begin{quote}
Valores de referência departamental. Localidades menores podem ter
preços 20--40\% abaixo da capital departamental. \#\#\# Saúde
\end{quote}

\textbf{Fonte:} MSPBS /\allowbreak  IPS Paraguay (2024-2026), consolidação
departamental e proxy local

\begin{longtable}[]{@{}ll@{}}
\toprule\noalign{}
Serviço & Disponibilidade \\
\midrule\noalign{}
\endhead
\bottomrule\noalign{}
\endlastfoot
USF /\allowbreak  Posto de Saúde & sim \\
Hospital Regional & não \\
IPS (seguro social) & sim \\
Farmácia & sim \\
Distância ao hospital de referência & \textasciitilde69 km (Coronel
Oviedo) \\
\end{longtable}

\textbf{Principais estabelecimentos:} USF local; referencia hospitalar
em Coronel Oviedo

\textbf{Observação para imigrantes:} Atendimento primario local ou em
raio curto; casos de maior complexidade seguem para Coronel Oviedo.
Cobertura privada continua recomendada para especialidades e urgências
de maior complexidade.
\subsubsection{Dossiê de Campo}\label{dossie-Santa_Rosa_del_Mbutuy-dossiuxea-de-campo}
\begin{description}[style=nextline,leftmargin=0.5cm,font=\bfseries]
\item[Coordenadas]  25°04'S 56°18'W.
\item[Topografia]  Terreno ondulado com áreas de vale. Área de \textasciitilde 298 km².
\item[Alvos Estrategicos]  Situado na Ruta PY08, importante eixo norte-sul. Ausência de alvos militares.
\item[Fallout]  Ventos predominantes do Norte.
\item[Populacao]  8.467 (Censo 2022 Final).
\item[Densidade]  \textasciitilde 28,4 hab/\allowbreak km².
\item[Vias de Acesso]  Cruzamento da Ruta PY08 com rotas departamentais. Boa conectividade com San Estanislao e Coronel Oviedo.
\item[Servicos]  Infraestrutura de saúde básica. Dependência de Coronel Oviedo para serviços especializados.
\item[Hidrologia]  Presença de arroios tributários do Rio Tebicuary-mí. Risco de isolamento de pontes em chuvas extremas.
\item[Sismicidade]  Baixa.
\item[Solo]  Alfisois, boa aptidão para agricultura diversificada.
\item[Agua]  Poços comunitários e nascentes protegidas.
\item[Energia]  Rede ANDE\footnote{Administración Nacional de Electricidade (Administração Nacional de Eletricidade), companhia estatal de energia.}.
\item[Producao]  Mandioca, milho, gado leiteiro e de corte.
\item[Seguranca]  Baixa taxa de criminalidade, característica de zona rural estável.
\item[Clima Social]  Forte coesão comunitária baseada em pequenas propriedades.
\end{description}
\textbf{Fonte climática:} NASA POWER Climatology API (período 2001-2020)
\textbf{Fonte luminosa:} estimativa world atlas

\paragraph{Irradiação Solar
(kWh/\allowbreak m²/\allowbreak dia)}\label{irradiauxe7uxe3o-solar-kwhmuxb2dia}

\begin{longtable}[]{@{}
  >{\raggedright\arraybackslash}p{(\linewidth - 24\tabcolsep) * \real{0.0746}}
  >{\raggedright\arraybackslash}p{(\linewidth - 24\tabcolsep) * \real{0.0746}}
  >{\raggedright\arraybackslash}p{(\linewidth - 24\tabcolsep) * \real{0.0746}}
  >{\raggedright\arraybackslash}p{(\linewidth - 24\tabcolsep) * \real{0.0746}}
  >{\raggedright\arraybackslash}p{(\linewidth - 24\tabcolsep) * \real{0.0746}}
  >{\raggedright\arraybackslash}p{(\linewidth - 24\tabcolsep) * \real{0.0746}}
  >{\raggedright\arraybackslash}p{(\linewidth - 24\tabcolsep) * \real{0.0746}}
  >{\raggedright\arraybackslash}p{(\linewidth - 24\tabcolsep) * \real{0.0746}}
  >{\raggedright\arraybackslash}p{(\linewidth - 24\tabcolsep) * \real{0.0746}}
  >{\raggedright\arraybackslash}p{(\linewidth - 24\tabcolsep) * \real{0.0746}}
  >{\raggedright\arraybackslash}p{(\linewidth - 24\tabcolsep) * \real{0.0746}}
  >{\raggedright\arraybackslash}p{(\linewidth - 24\tabcolsep) * \real{0.0746}}
  >{\raggedright\arraybackslash}p{(\linewidth - 24\tabcolsep) * \real{0.1045}}@{}}
\toprule\noalign{}
\begin{minipage}[b]{\linewidth}\raggedright
Jan
\end{minipage} & \begin{minipage}[b]{\linewidth}\raggedright
Fev
\end{minipage} & \begin{minipage}[b]{\linewidth}\raggedright
Mar
\end{minipage} & \begin{minipage}[b]{\linewidth}\raggedright
Abr
\end{minipage} & \begin{minipage}[b]{\linewidth}\raggedright
Mai
\end{minipage} & \begin{minipage}[b]{\linewidth}\raggedright
Jun
\end{minipage} & \begin{minipage}[b]{\linewidth}\raggedright
Jul
\end{minipage} & \begin{minipage}[b]{\linewidth}\raggedright
Ago
\end{minipage} & \begin{minipage}[b]{\linewidth}\raggedright
Set
\end{minipage} & \begin{minipage}[b]{\linewidth}\raggedright
Out
\end{minipage} & \begin{minipage}[b]{\linewidth}\raggedright
Nov
\end{minipage} & \begin{minipage}[b]{\linewidth}\raggedright
Dez
\end{minipage} & \begin{minipage}[b]{\linewidth}\raggedright
Média
\end{minipage} \\
\midrule\noalign{}
\endhead
\bottomrule\noalign{}
\endlastfoot
6.73 & 6.14 & 5.58 & 4.61 & 3.40 & 2.96 & 3.24 & 4.04 & 4.62 & 5.38 &
6.37 & 6.67 & \textbf{4.98} \\
\end{longtable}

\textbf{Inclinação solar recomendada:} 25° N (anual) · 35° N (inverno
jun-ago) · 15° N (verão nov-jan)

\paragraph{Precipitação (mm/\allowbreak mês)}\label{precipitauxe7uxe3o-mmmuxeas}

\begin{longtable}[]{@{}
  >{\raggedright\arraybackslash}p{(\linewidth - 24\tabcolsep) * \real{0.0704}}
  >{\raggedright\arraybackslash}p{(\linewidth - 24\tabcolsep) * \real{0.0704}}
  >{\raggedright\arraybackslash}p{(\linewidth - 24\tabcolsep) * \real{0.0704}}
  >{\raggedright\arraybackslash}p{(\linewidth - 24\tabcolsep) * \real{0.0704}}
  >{\raggedright\arraybackslash}p{(\linewidth - 24\tabcolsep) * \real{0.0704}}
  >{\raggedright\arraybackslash}p{(\linewidth - 24\tabcolsep) * \real{0.0704}}
  >{\raggedright\arraybackslash}p{(\linewidth - 24\tabcolsep) * \real{0.0704}}
  >{\raggedright\arraybackslash}p{(\linewidth - 24\tabcolsep) * \real{0.0704}}
  >{\raggedright\arraybackslash}p{(\linewidth - 24\tabcolsep) * \real{0.0704}}
  >{\raggedright\arraybackslash}p{(\linewidth - 24\tabcolsep) * \real{0.0704}}
  >{\raggedright\arraybackslash}p{(\linewidth - 24\tabcolsep) * \real{0.0704}}
  >{\raggedright\arraybackslash}p{(\linewidth - 24\tabcolsep) * \real{0.0704}}
  >{\raggedright\arraybackslash}p{(\linewidth - 24\tabcolsep) * \real{0.1549}}@{}}
\toprule\noalign{}
\begin{minipage}[b]{\linewidth}\raggedright
Jan
\end{minipage} & \begin{minipage}[b]{\linewidth}\raggedright
Fev
\end{minipage} & \begin{minipage}[b]{\linewidth}\raggedright
Mar
\end{minipage} & \begin{minipage}[b]{\linewidth}\raggedright
Abr
\end{minipage} & \begin{minipage}[b]{\linewidth}\raggedright
Mai
\end{minipage} & \begin{minipage}[b]{\linewidth}\raggedright
Jun
\end{minipage} & \begin{minipage}[b]{\linewidth}\raggedright
Jul
\end{minipage} & \begin{minipage}[b]{\linewidth}\raggedright
Ago
\end{minipage} & \begin{minipage}[b]{\linewidth}\raggedright
Set
\end{minipage} & \begin{minipage}[b]{\linewidth}\raggedright
Out
\end{minipage} & \begin{minipage}[b]{\linewidth}\raggedright
Nov
\end{minipage} & \begin{minipage}[b]{\linewidth}\raggedright
Dez
\end{minipage} & \begin{minipage}[b]{\linewidth}\raggedright
Total/\allowbreak ano
\end{minipage} \\
\midrule\noalign{}
\endhead
\bottomrule\noalign{}
\endlastfoot
124 & 143 & 125 & 147 & 173 & 90 & 70 & 46 & 94 & 184 & 190 & 168 &
\textbf{1557 mm} \\
\end{longtable}

\paragraph{Poluição Luminosa}\label{poluiuxe7uxe3o-luminosa}

\begin{longtable}[]{@{}ll@{}}
\toprule\noalign{}
Parâmetro & Valor \\
\midrule\noalign{}
\endhead
\bottomrule\noalign{}
\endlastfoot
Escala Bortle & 4 --- Céu rural-suburbano \\
Radiância artificial & 3.0 nW/\allowbreak cm²/\allowbreak sr \\
\end{longtable}
\paragraph{Solo (SoilGrids 2.0, média ponderada 0--30
cm)}

\begin{longtable}[]{@{}ll@{}}
\toprule\noalign{}
Parâmetro & Valor \\
\midrule\noalign{}
\endhead
\bottomrule\noalign{}
\endlastfoot
pH (H$_2$O) & 5.3 \\
Carbono orgânico (SOC) & 18.23 g/\allowbreak kg \\
Argila & 23.89 \% \\
Areia & 51.8 \% \\
Silte (calc.) & 24.3 \% \\
Densidade aparente & 1.31 g/\allowbreak cm³ \\
\textbf{Aptidão agrícola} & \textbf{Média} \\
\end{longtable}

Fonte: ISRIC SoilGrids 2.0 via WCS. Coords: -25.0667°, -56.3°. Média
ponderada camadas 0-5, 5-15, 15-30 cm.

\begin{itemize}
\tightlist
\item
  \textbf{Agua:} Poços comunitários e nascentes protegidas.
\item
  \textbf{Energia:} Rede ANDE\footnote{Administración Nacional de Electricidade (Administração Nacional de Eletricidade), companhia estatal de energia.}.
\item
  \textbf{Producao:} Mandioca, milho, gado leiteiro e de corte.
\end{itemize}
\subsubsection{Indicadores Sociais}

\textbf{Fontes:} PNUD 2020, INE\footnote{Instituto Nacional de Estadística (Instituto Nacional de Estatística), órgão oficial de dados demográficos e censitários.} Censo 2022, INE\footnote{Instituto Nacional de Estadística (Instituto Nacional de Estatística), órgão oficial de dados demográficos e censitários.} EPHC 2023, Ministerio
Público 2024\\
\textbf{Nota:} Valores marcados como \emph{(dept.)} referem-se ao
departamento de Caaguazú; valores \emph{(dist.)} são específicos deste
distrito. Dados marcados \emph{(est.)} são estimativas.

\begin{longtable}[]{@{}
  >{\raggedright\arraybackslash}p{(\linewidth - 4\tabcolsep) * \real{0.4231}}
  >{\raggedright\arraybackslash}p{(\linewidth - 4\tabcolsep) * \real{0.2692}}
  >{\raggedright\arraybackslash}p{(\linewidth - 4\tabcolsep) * \real{0.3077}}@{}}
\toprule\noalign{}
\begin{minipage}[b]{\linewidth}\raggedright
Indicador
\end{minipage} & \begin{minipage}[b]{\linewidth}\raggedright
Valor
\end{minipage} & \begin{minipage}[b]{\linewidth}\raggedright
Âmbito
\end{minipage} \\
\midrule\noalign{}
\endhead
\bottomrule\noalign{}
\endlastfoot
IDH (2020) & 0,715 & dept. \\
Ranking IDH nacional & 7/\allowbreak 18 & dept. \\
Esperança de vida & 73,1 anos & dept. \\
Escolaridade média & 8.3 anos & dist. \\
RNB per capita (USD PPA) & 9.000 & dept. \\
Pobreza monetária (\%) & 28,5\% & dept. \\
Pobreza extrema (\%) & 6,8\% & dept. \\
Índice de Gini & 0,447 & dept. \\
Acesso a água potável (\%) & 79,8\% & dept. \\
Acesso a saneamento (\%) & 78,3\% & dept. \\
Taxa de homicídios (est., /\allowbreak 100k hab) & \textasciitilde6--8 (estimativa,
próxima ou acima da média nacional) & dept. \\
Índice de segurança & médio & dept. \\
População (Censo 2022) & 8.467 hab. & dist. \\
Idade mediana & 30.0 anos & dist. \\
Taxa de fecundidade & N/\allowbreak D & dist. \\
\end{longtable}
\subsubsection{Combustível}

\textbf{Referência:} PETROPAR /\allowbreak  postos locais (2024) \textbf{Tipo de
localidade:} Interior

\begin{longtable}[]{@{}lll@{}}
\toprule\noalign{}
Combustível & USD/\allowbreak litro & Gs/\allowbreak litro (aprox.) \\
\midrule\noalign{}
\endhead
\bottomrule\noalign{}
\endlastfoot
Gasolina 93 oct & 0.98 & 7,252 \\
Gasolina 97 oct (premium) & 1.08 & 7,992 \\
Diesel & 0.91 & 6,734 \\
\end{longtable}

\begin{quote}
Preços podem variar ±5\% conforme posto e sazonalidade. Chaco e interior
remoto apresentam maior variação.
\end{quote}

\subsubsection{Cobertura Celular}

\textbf{Fonte:} CONATEL PY /\allowbreak  operadoras (2024)

\begin{longtable}[]{@{}ll@{}}
\toprule\noalign{}
Parâmetro & Valor \\
\midrule\noalign{}
\endhead
\bottomrule\noalign{}
\endlastfoot
Cobertura 4G população (dept.) & 88\% \\
Cobertura 4G área rural & 72\% \\
Melhor operadora & Tigo \\
Qualidade rural & boa \\
\end{longtable}

\begin{quote}
Para áreas rurais fora do núcleo urbano, recomenda-se chip Tigo como
principal e Personal como backup.
\end{quote}

\subsubsection{Internet}

\textbf{Fonte:} CONATEL /\allowbreak  Speedtest Ookla (2024)

\begin{longtable}[]{@{}ll@{}}
\toprule\noalign{}
Parâmetro & Valor \\
\midrule\noalign{}
\endhead
\bottomrule\noalign{}
\endlastfoot
Velocidade média download & 55 Mbps \\
Domicílios com internet (dept.) & 62\% \\
Tecnologia predominante & rádio/\allowbreak fibra \\
Opção rural & Starlink disponível (\textasciitilde USD 44/\allowbreak mês) \\
\end{longtable}

\subsubsection{Mercado Imobiliário e Terra
Rural}

\textbf{Fonte:} INDERT /\allowbreak  Clasificados.com.py (2024)

\begin{longtable}[]{@{}ll@{}}
\toprule\noalign{}
Tipo & Referência \\
\midrule\noalign{}
\endhead
\bottomrule\noalign{}
\endlastfoot
Terra agrícola alta prod. (USD/\allowbreak ha) & 8,800 \\
Imóvel urbano (USD/\allowbreak m²) & 850 \\
Aluguel 2 quartos (USD/\allowbreak mês) & 350 \\
\end{longtable}

\begin{quote}
Valores de referência departamental. Localidades menores podem ter
preços 20--40\% abaixo da capital departamental. \#\#\# Saúde
\end{quote}

\textbf{Fonte:} MSPBS /\allowbreak  IPS Paraguay (2024-2026), consolidação
departamental e proxy local

\begin{longtable}[]{@{}ll@{}}
\toprule\noalign{}
Serviço & Disponibilidade \\
\midrule\noalign{}
\endhead
\bottomrule\noalign{}
\endlastfoot
USF /\allowbreak  Posto de Saúde & sim \\
Hospital Regional & não \\
IPS (seguro social) & sim \\
Farmácia & sim \\
Distância ao hospital de referência & \textasciitilde69 km (Coronel
Oviedo) \\
\end{longtable}

\textbf{Principais estabelecimentos:} USF local; referencia hospitalar
em Coronel Oviedo

\textbf{Observação para imigrantes:} Atendimento primario local ou em
raio curto; casos de maior complexidade seguem para Coronel Oviedo.
Cobertura privada continua recomendada para especialidades e urgências
de maior complexidade.
\newpage
\secaoDiagnostico{Simon Bolivar}{\mapaConteudo{-25.0500}{-56.3333}}{Simon Bolivar (Caaguazu) apresenta perfil operacional consolidado para
comparacao territorial, com leitura conjunta de infraestrutura, risco
hidroclimatico e capacidade de continuidade de servicos essenciais.

\subsubsection{Por que sim}

\begin{itemize}
\tightlist
\item
  Base institucional de referencia disponivel e rastreavel.
\item
  Estrutura comparativa (A-E/\allowbreak GSS) consistente para decisao relativa.
\item
  Plano de mitigacao acionavel ja definido no dossie.
\end{itemize}

\subsubsection{Por que não}

\begin{itemize}
\tightlist
\item
  Serie oficial distrital granular ainda pode apresentar lacunas
  temporais.
\item
  Sensibilidade do score exige monitoramento continuo de risco e
  conectividade.
\item
  Decisao final depende de validacao de campo e janela temporal recente.
\end{itemize}

\subsubsection{Pre-condicoes Minimas de
Decisao}

\begin{itemize}
\tightlist
\item
  Atualizar indicadores criticos em ciclo trimestral.
\item
  Confirmar trafegabilidade e contingencia hidrica/\allowbreak energetica local.
\item
  Validar checklist de resiliencia domiciliar/\allowbreak logistica antes de decisao
  final.
\end{itemize}

\subsubsection{Recomendação
Técnica}

Uso recomendado como candidato em analise multicriterio, condicionado ao
cumprimento das pre-condicoes acima. Referencia atual de score: GSS 7.3
em 2026-03-05; risco predominante: baixo.

\subsubsection{}

\begin{itemize}
\tightlist
\item
  \begin{enumerate}
  \def\labelenumi{\arabic{enumi})}
  \tightlist
  \item
    Delimitacao territorial validada por georreferencia institucional
    (Fonte: acesso: 2026-03-05).
  \end{enumerate}
\item
  \begin{enumerate}
  \def\labelenumi{\arabic{enumi})}
  \setcounter{enumi}{1}
  \tightlist
  \item
    População municipal/\allowbreak distrital referenciada por base censitaria
    nacional (Fonte: acesso: 2026-03-05).
  \end{enumerate}
\item
  \begin{enumerate}
  \def\labelenumi{\arabic{enumi})}
  \setcounter{enumi}{2}
  \tightlist
  \item
    Comparabilidade intra-departamental apoiada em indicadores
    distritais oficiais (Fonte: acesso: 2026-03-05).
  \end{enumerate}
\item
  \begin{enumerate}
  \def\labelenumi{\arabic{enumi})}
  \setcounter{enumi}{3}
  \tightlist
  \item
    Estado macro de conectividade viaria ancorado em informacoes de
    rodovias (Fonte: acesso: 2026-03-05).
  \end{enumerate}
\item
  \begin{enumerate}
  \def\labelenumi{\arabic{enumi})}
  \setcounter{enumi}{4}
  \tightlist
  \item
    Exposicao hidrometeorologica monitorada por avisos oficiais (Fonte:
    acesso: 2026-03-05).
  \end{enumerate}
\item
  \begin{enumerate}
  \def\labelenumi{\arabic{enumi})}
  \setcounter{enumi}{5}
  \tightlist
  \item
    Historico de contexto meteorologico apoiado em publicacoes tecnicas
    (Fonte: acesso: 2026-03-05).
  \end{enumerate}
\item
  \begin{enumerate}
  \def\labelenumi{\arabic{enumi})}
  \setcounter{enumi}{6}
  \tightlist
  \item
    Capacidade institucional de resposta a eventos apoiada em acoes da
    SEN\footnote{Secretaría de Emergencia Nacional (Secretaria de Emergência Nacional), órgão governamental de resposta a desastres.} (Fonte: acesso: 2026-03-05).
  \end{enumerate}
\item
  \begin{enumerate}
  \def\labelenumi{\arabic{enumi})}
  \setcounter{enumi}{7}
  \tightlist
  \item
    Infraestrutura energetica analisada via operador nacional (Fonte:
    acesso: 2026-03-05).
  \end{enumerate}
\item
  \begin{enumerate}
  \def\labelenumi{\arabic{enumi})}
  \setcounter{enumi}{8}
  \tightlist
  \item
    Referencias de obras e logistica nacional verificadas no MOPC
    (Fonte: acesso: 2026-03-05).
  \end{enumerate}
\item
  \begin{enumerate}
  \def\labelenumi{\arabic{enumi})}
  \setcounter{enumi}{9}
  \tightlist
  \item
    Contexto sociopolitico institucional referenciado por autoridade
    eleitoral (Fonte: acesso: 2026-03-05).
  \end{enumerate}
\item
  \begin{enumerate}
  \def\labelenumi{\arabic{enumi})}
  \setcounter{enumi}{10}
  \tightlist
  \item
    Camada cartografica de apoio eleitoral/\allowbreak territorial consultada
    (Fonte: acesso: 2026-03-05).
  \end{enumerate}
\item
  \begin{enumerate}
  \def\labelenumi{\arabic{enumi})}
  \setcounter{enumi}{11}
  \tightlist
  \item
    Dados publicos complementares para verificacao cruzada (Fonte:
    acesso: 2026-03-05).
  \end{enumerate}
\end{itemize}}
\begin{table}[ht!]
\small \begin{tabular}{p{3.0cm}p{0.8cm}p{6.0cm}} \toprule
\textbf{Critério} & \textbf{Nota} & \textbf{Justificativa} \\ \midrule
A - Ameaças Estratégicas & 7.8 & - \\
B - Risco Social & 6.1 & - \\
C - Riscos Naturais & 8.1 & - \\
D - Autossuficiência & 7.2 & - \\
E - Institucional & 7.3 & - \\
\midrule \textbf{GSS FINAL} & \textbf{7.3} & \textbf{Seguro (bom potencial com preparacao).} \\ \bottomrule \end{tabular}
\end{table}
\subsubsection{Vulnerabilidades e
Mitigação}\label{vuln-Simon_Bolivar-vulnerabilidades-e-mitigauxe7uxe3o}

\begin{itemize}
\tightlist
\item
  Exposicao hidrometeorologica controlada (\textless45/\allowbreak 100).
\item
  Conectividade viaria intermediaria (50-69/\allowbreak 100).
\item
  Estoque de contingencia para 14 dias (agua/\allowbreak alimentos/\allowbreak combustivel),
  priorizado quando risco hidro=24/\allowbreak 100.
\item
  Duplicar rotas/\allowbreak logistica quando conectividade viaria=50/\allowbreak 100 for
  \textless70; manter rota secundaria mapeada e testada trimestralmente.
\item
  Plano de energia resiliente com autonomia minima de 72h
  (geracao/\allowbreak backup), revisao semestral.
\end{itemize}
\subsubsection{Dossiê de Campo}\label{dossie-Simon_Bolivar-dossiuxea-de-campo}
\begin{description}[style=nextline,leftmargin=0.5cm,font=\bfseries]
\item[Coordenadas]  25°03'S 56°20'W (geocodificado via Nominatim).
\end{description}
\textbf{Fonte climática:} NASA POWER Climatology API (período 2001-2020)
\textbf{Fonte luminosa:} estimativa world atlas

\paragraph{Irradiação Solar
(kWh/\allowbreak m²/\allowbreak dia)}\label{irradiauxe7uxe3o-solar-kwhmuxb2dia}

\begin{longtable}[]{@{}
  >{\raggedright\arraybackslash}p{(\linewidth - 24\tabcolsep) * \real{0.0746}}
  >{\raggedright\arraybackslash}p{(\linewidth - 24\tabcolsep) * \real{0.0746}}
  >{\raggedright\arraybackslash}p{(\linewidth - 24\tabcolsep) * \real{0.0746}}
  >{\raggedright\arraybackslash}p{(\linewidth - 24\tabcolsep) * \real{0.0746}}
  >{\raggedright\arraybackslash}p{(\linewidth - 24\tabcolsep) * \real{0.0746}}
  >{\raggedright\arraybackslash}p{(\linewidth - 24\tabcolsep) * \real{0.0746}}
  >{\raggedright\arraybackslash}p{(\linewidth - 24\tabcolsep) * \real{0.0746}}
  >{\raggedright\arraybackslash}p{(\linewidth - 24\tabcolsep) * \real{0.0746}}
  >{\raggedright\arraybackslash}p{(\linewidth - 24\tabcolsep) * \real{0.0746}}
  >{\raggedright\arraybackslash}p{(\linewidth - 24\tabcolsep) * \real{0.0746}}
  >{\raggedright\arraybackslash}p{(\linewidth - 24\tabcolsep) * \real{0.0746}}
  >{\raggedright\arraybackslash}p{(\linewidth - 24\tabcolsep) * \real{0.0746}}
  >{\raggedright\arraybackslash}p{(\linewidth - 24\tabcolsep) * \real{0.1045}}@{}}
\toprule\noalign{}
\begin{minipage}[b]{\linewidth}\raggedright
Jan
\end{minipage} & \begin{minipage}[b]{\linewidth}\raggedright
Fev
\end{minipage} & \begin{minipage}[b]{\linewidth}\raggedright
Mar
\end{minipage} & \begin{minipage}[b]{\linewidth}\raggedright
Abr
\end{minipage} & \begin{minipage}[b]{\linewidth}\raggedright
Mai
\end{minipage} & \begin{minipage}[b]{\linewidth}\raggedright
Jun
\end{minipage} & \begin{minipage}[b]{\linewidth}\raggedright
Jul
\end{minipage} & \begin{minipage}[b]{\linewidth}\raggedright
Ago
\end{minipage} & \begin{minipage}[b]{\linewidth}\raggedright
Set
\end{minipage} & \begin{minipage}[b]{\linewidth}\raggedright
Out
\end{minipage} & \begin{minipage}[b]{\linewidth}\raggedright
Nov
\end{minipage} & \begin{minipage}[b]{\linewidth}\raggedright
Dez
\end{minipage} & \begin{minipage}[b]{\linewidth}\raggedright
Média
\end{minipage} \\
\midrule\noalign{}
\endhead
\bottomrule\noalign{}
\endlastfoot
6.74 & 6.20 & 5.50 & 4.48 & 3.35 & 2.88 & 3.21 & 3.96 & 4.59 & 5.43 &
6.43 & 6.81 & \textbf{4.96} \\
\end{longtable}

\textbf{Inclinação solar recomendada:} 25° N (anual) · 35° N (inverno
jun-ago) · 15° N (verão nov-jan)

\paragraph{Precipitação (mm/\allowbreak mês)}\label{precipitauxe7uxe3o-mmmuxeas}

\begin{longtable}[]{@{}
  >{\raggedright\arraybackslash}p{(\linewidth - 24\tabcolsep) * \real{0.0704}}
  >{\raggedright\arraybackslash}p{(\linewidth - 24\tabcolsep) * \real{0.0704}}
  >{\raggedright\arraybackslash}p{(\linewidth - 24\tabcolsep) * \real{0.0704}}
  >{\raggedright\arraybackslash}p{(\linewidth - 24\tabcolsep) * \real{0.0704}}
  >{\raggedright\arraybackslash}p{(\linewidth - 24\tabcolsep) * \real{0.0704}}
  >{\raggedright\arraybackslash}p{(\linewidth - 24\tabcolsep) * \real{0.0704}}
  >{\raggedright\arraybackslash}p{(\linewidth - 24\tabcolsep) * \real{0.0704}}
  >{\raggedright\arraybackslash}p{(\linewidth - 24\tabcolsep) * \real{0.0704}}
  >{\raggedright\arraybackslash}p{(\linewidth - 24\tabcolsep) * \real{0.0704}}
  >{\raggedright\arraybackslash}p{(\linewidth - 24\tabcolsep) * \real{0.0704}}
  >{\raggedright\arraybackslash}p{(\linewidth - 24\tabcolsep) * \real{0.0704}}
  >{\raggedright\arraybackslash}p{(\linewidth - 24\tabcolsep) * \real{0.0704}}
  >{\raggedright\arraybackslash}p{(\linewidth - 24\tabcolsep) * \real{0.1549}}@{}}
\toprule\noalign{}
\begin{minipage}[b]{\linewidth}\raggedright
Jan
\end{minipage} & \begin{minipage}[b]{\linewidth}\raggedright
Fev
\end{minipage} & \begin{minipage}[b]{\linewidth}\raggedright
Mar
\end{minipage} & \begin{minipage}[b]{\linewidth}\raggedright
Abr
\end{minipage} & \begin{minipage}[b]{\linewidth}\raggedright
Mai
\end{minipage} & \begin{minipage}[b]{\linewidth}\raggedright
Jun
\end{minipage} & \begin{minipage}[b]{\linewidth}\raggedright
Jul
\end{minipage} & \begin{minipage}[b]{\linewidth}\raggedright
Ago
\end{minipage} & \begin{minipage}[b]{\linewidth}\raggedright
Set
\end{minipage} & \begin{minipage}[b]{\linewidth}\raggedright
Out
\end{minipage} & \begin{minipage}[b]{\linewidth}\raggedright
Nov
\end{minipage} & \begin{minipage}[b]{\linewidth}\raggedright
Dez
\end{minipage} & \begin{minipage}[b]{\linewidth}\raggedright
Total/\allowbreak ano
\end{minipage} \\
\midrule\noalign{}
\endhead
\bottomrule\noalign{}
\endlastfoot
124 & 143 & 125 & 147 & 173 & 90 & 70 & 46 & 94 & 184 & 190 & 168 &
\textbf{1557 mm} \\
\end{longtable}

\paragraph{Poluição Luminosa}\label{poluiuxe7uxe3o-luminosa}

\begin{longtable}[]{@{}ll@{}}
\toprule\noalign{}
Parâmetro & Valor \\
\midrule\noalign{}
\endhead
\bottomrule\noalign{}
\endlastfoot
Escala Bortle & 4 --- Céu rural-suburbano \\
Radiância artificial & 3.0 nW/\allowbreak cm²/\allowbreak sr \\
\end{longtable}
\paragraph{Solo (SoilGrids 2.0, média ponderada 0--30
cm)}

\begin{longtable}[]{@{}ll@{}}
\toprule\noalign{}
Parâmetro & Valor \\
\midrule\noalign{}
\endhead
\bottomrule\noalign{}
\endlastfoot
pH (H$_2$O) & 5.3 \\
Carbono orgânico (SOC) & 18.43 g/\allowbreak kg \\
Argila & 23.73 \% \\
Areia & 51.33 \% \\
Silte (calc.) & 24.9 \% \\
Densidade aparente & 1.29 g/\allowbreak cm³ \\
\textbf{Aptidão agrícola} & \textbf{Média} \\
\end{longtable}

Fonte: ISRIC SoilGrids 2.0 via WCS. Coords: -25.05°, -56.3333°. Média
ponderada camadas 0-5, 5-15, 15-30 cm.
\subsubsection{Indicadores Sociais}

\textbf{Fontes:} PNUD 2020, INE\footnote{Instituto Nacional de Estadística (Instituto Nacional de Estatística), órgão oficial de dados demográficos e censitários.} Censo 2022, INE\footnote{Instituto Nacional de Estadística (Instituto Nacional de Estatística), órgão oficial de dados demográficos e censitários.} EPHC 2023, Ministerio
Público 2024\\
\textbf{Nota:} Valores marcados como \emph{(dept.)} referem-se ao
departamento de Caaguazú; valores \emph{(dist.)} são específicos deste
distrito. Dados marcados \emph{(est.)} são estimativas.

\begin{longtable}[]{@{}
  >{\raggedright\arraybackslash}p{(\linewidth - 4\tabcolsep) * \real{0.4231}}
  >{\raggedright\arraybackslash}p{(\linewidth - 4\tabcolsep) * \real{0.2692}}
  >{\raggedright\arraybackslash}p{(\linewidth - 4\tabcolsep) * \real{0.3077}}@{}}
\toprule\noalign{}
\begin{minipage}[b]{\linewidth}\raggedright
Indicador
\end{minipage} & \begin{minipage}[b]{\linewidth}\raggedright
Valor
\end{minipage} & \begin{minipage}[b]{\linewidth}\raggedright
Âmbito
\end{minipage} \\
\midrule\noalign{}
\endhead
\bottomrule\noalign{}
\endlastfoot
IDH (2020) & 0,715 & dept. \\
Ranking IDH nacional & 7/\allowbreak 18 & dept. \\
Esperança de vida & 73,1 anos & dept. \\
Escolaridade média & 8.4 anos & dist. \\
RNB per capita (USD PPA) & 9.000 & dept. \\
Pobreza monetária (\%) & 28,5\% & dept. \\
Pobreza extrema (\%) & 6,8\% & dept. \\
Índice de Gini & 0,447 & dept. \\
Acesso a água potável (\%) & 79,8\% & dept. \\
Acesso a saneamento (\%) & 78,3\% & dept. \\
Taxa de homicídios (est., /\allowbreak 100k hab) & \textasciitilde6--8 (estimativa,
próxima ou acima da média nacional) & dept. \\
Índice de segurança & médio & dept. \\
População (Censo 2022) & 474 hab. & dist. \\
Idade mediana & 30.0 anos & dist. \\
Taxa de fecundidade & N/\allowbreak D & dist. \\
\end{longtable}
\subsubsection{Combustível}

\textbf{Referência:} PETROPAR /\allowbreak  postos locais (2024) \textbf{Tipo de
localidade:} Interior

\begin{longtable}[]{@{}lll@{}}
\toprule\noalign{}
Combustível & USD/\allowbreak litro & Gs/\allowbreak litro (aprox.) \\
\midrule\noalign{}
\endhead
\bottomrule\noalign{}
\endlastfoot
Gasolina 93 oct & 0.98 & 7,252 \\
Gasolina 97 oct (premium) & 1.08 & 7,992 \\
Diesel & 0.91 & 6,734 \\
\end{longtable}

\begin{quote}
Preços podem variar ±5\% conforme posto e sazonalidade. Chaco e interior
remoto apresentam maior variação.
\end{quote}

\subsubsection{Cobertura Celular}

\textbf{Fonte:} CONATEL PY /\allowbreak  operadoras (2024)

\begin{longtable}[]{@{}ll@{}}
\toprule\noalign{}
Parâmetro & Valor \\
\midrule\noalign{}
\endhead
\bottomrule\noalign{}
\endlastfoot
Cobertura 4G população (dept.) & 88\% \\
Cobertura 4G área rural & 72\% \\
Melhor operadora & Tigo \\
Qualidade rural & boa \\
\end{longtable}

\begin{quote}
Para áreas rurais fora do núcleo urbano, recomenda-se chip Tigo como
principal e Personal como backup.
\end{quote}

\subsubsection{Internet}

\textbf{Fonte:} CONATEL /\allowbreak  Speedtest Ookla (2024)

\begin{longtable}[]{@{}ll@{}}
\toprule\noalign{}
Parâmetro & Valor \\
\midrule\noalign{}
\endhead
\bottomrule\noalign{}
\endlastfoot
Velocidade média download & 55 Mbps \\
Domicílios com internet (dept.) & 62\% \\
Tecnologia predominante & rádio/\allowbreak fibra \\
Opção rural & Starlink disponível (\textasciitilde USD 44/\allowbreak mês) \\
\end{longtable}

\subsubsection{Mercado Imobiliário e Terra
Rural}

\textbf{Fonte:} INDERT /\allowbreak  Clasificados.com.py (2024)

\begin{longtable}[]{@{}ll@{}}
\toprule\noalign{}
Tipo & Referência \\
\midrule\noalign{}
\endhead
\bottomrule\noalign{}
\endlastfoot
Terra agrícola alta prod. (USD/\allowbreak ha) & 8,800 \\
Imóvel urbano (USD/\allowbreak m²) & 850 \\
Aluguel 2 quartos (USD/\allowbreak mês) & 350 \\
\end{longtable}

\begin{quote}
Valores de referência departamental. Localidades menores podem ter
preços 20--40\% abaixo da capital departamental. \#\#\# Saúde
\end{quote}

\textbf{Fonte:} MSPBS /\allowbreak  IPS Paraguay (2024-2026), consolidação
departamental e proxy local

\begin{longtable}[]{@{}ll@{}}
\toprule\noalign{}
Serviço & Disponibilidade \\
\midrule\noalign{}
\endhead
\bottomrule\noalign{}
\endlastfoot
USF /\allowbreak  Posto de Saúde & sim \\
Hospital Regional & não \\
IPS (seguro social) & sim \\
Farmácia & sim \\
Distância ao hospital de referência & \textasciitilde57 km (Coronel
Oviedo) \\
\end{longtable}

\textbf{Principais estabelecimentos:} USF local; referencia hospitalar
em Coronel Oviedo

\textbf{Observação para imigrantes:} Atendimento primario local ou em
raio curto; casos de maior complexidade seguem para Coronel Oviedo.
Cobertura privada continua recomendada para especialidades e urgências
de maior complexidade.
\newpage
\secaoDiagnostico{Tembiapora}{\mapaConteudo{-25.3166}{-55.5500}}{Tembiaporá é a ``península da abundância'' no departamento de Caaguazú.
Sua geografia única, quase cercada pelas águas do Lago Yguazú, oferece
uma camada natural de isolamento e segurança física contra grandes
fluxos demográficos. É um polo de segurança alimentar insuperável,
garantindo recursos calóricos e hídricos massivos para sua população
dispersa. Sua viabilidade como refúgio de nível superior depende da
capacidade do ocupante de gerir a autonomia logística e energética.

\subsubsection{Por que sim}

\begin{itemize}
\tightlist
\item
  Abundância hídrica e alimentar excepcional (Capital da Banana).
\item
  Geografia favorável ao controle de acessos e invisibilidade tática.
\item
  Baixíssima criminalidade e ambiente social rural estável.
\item
  Sem restrições da Lei de Fronteira.
\end{itemize}

\subsubsection{Por que não}

\begin{itemize}
\tightlist
\item
  Infraestrutura de saúde local limitada (dependência de Campo 9).
\item
  Risco de isolamento logístico terrestre em cenários climáticos
  severos.
\item
  Valor da terra em ascensão pela alta produtividade frutícola.
\end{itemize}

\subsubsection{Recomendação
Técnica}

Altamente indicado para estabelecimentos de autossuficiência de médio
porte que priorizem a segurança alimentar e hídrica. Requer posse de
meios de transporte redundantes (terrestre e aquático) e autonomia em
saúde básica. O GSS de 6.5 reflete a riqueza excepcional de recursos em
um ambiente de isolamento natural. Referencia atual de score: GSS 6.5 em
2026-03-05.

\subsubsection{}

\begin{itemize}
\tightlist
\item
  \begin{enumerate}
  \def\labelenumi{\arabic{enumi})}
  \tightlist
  \item
    Dados censitários validados (INE\footnote{Instituto Nacional de Estadística (Instituto Nacional de Estatística), órgão oficial de dados demográficos e censitários.} 2022).
  \end{enumerate}
\item
  \begin{enumerate}
  \def\labelenumi{\arabic{enumi})}
  \setcounter{enumi}{1}
  \tightlist
  \item
    Relatórios de exportação de banana e frutas tropicais (MAG).
  \end{enumerate}
\item
  \begin{enumerate}
  \def\labelenumi{\arabic{enumi})}
  \setcounter{enumi}{2}
  \tightlist
  \item
    Mapa hidrológico do reservatório de Yguazú (ANDE\footnote{Administración Nacional de Electricidade (Administração Nacional de Eletricidade), companhia estatal de energia.}/\allowbreak SEN\footnote{Secretaría de Emergencia Nacional (Secretaria de Emergência Nacional), órgão governamental de resposta a desastres.}).
  \end{enumerate}
\item
  \begin{enumerate}
  \def\labelenumi{\arabic{enumi})}
  \setcounter{enumi}{3}
  \tightlist
  \item
    Registro de segurança pública departamental.
  \end{enumerate}
\end{itemize}}
\begin{table}[ht!]
\small \begin{tabular}{p{3.0cm}p{0.8cm}p{6.0cm}} \toprule
\textbf{Critério} & \textbf{Nota} & \textbf{Justificativa} \\ \midrule
A - Ameaças Estratégicas & 6.0 & Posicionamento peninsular discreto e protegido; excelente invisibilidade tática \\
B - Risco Social & 7.0 & População pequena e baixa densidade demográfica; risco social urbano nulo \\
C - Riscos Naturais & 6.5 & Geologia muito estável; penalizado pelo risco moderado de isolamento hídrico e tempestades \\
D - Autossuficiência & 8.0 & Recursos alimentares e hídricos excepcionais; capital nacional da banana \\
E - Institucional & 5.0 & Município recente com infraestrutura institucional e de saúde limitada \\
\midrule \textbf{GSS FINAL} & \textbf{6.5} & \textbf{Moderadamente Seguro (Ideal para quem busca isolamento demográfico em região de alta produtividade de alimentos).} \\ \bottomrule \end{tabular}
\end{table}
\subsubsection{Vulnerabilidades e
Mitigação}\label{vuln-Tembiapora-vulnerabilidades-e-mitigauxe7uxe3o}

\begin{itemize}
\tightlist
\item
  \textbf{Isolamento Geográfico:} Risco de desabastecimento de insumos
  técnicos por dependência de acesso terrestre único ou balsa
  (Mitigação: estoque de insumos para 90 dias e posse de embarcação a
  motor).
\item
  \textbf{Instabilidade Elétrica:} Dependência da rede ANDE\footnote{Administración Nacional de Electricidade (Administração Nacional de Eletricidade), companhia estatal de energia.} rural
  (Mitigação: instalação de sistemas solares fotovoltaicos potentes para
  refrigeração de alimentos).
\item
  \textbf{Dependência de Monocultura:} Economia muito vinculada à banana
  (Mitigação: diversificação produtiva com foco em grãos e proteína
  animal de pequeno porte).
\end{itemize}
\subsubsection{Dossiê de Campo}\label{dossie-Tembiapora-dossiuxea-de-campo}
\begin{description}[style=nextline,leftmargin=0.5cm,font=\bfseries]
\item[Coordenadas]  25°19'S 55°33'W.
\item[Topografia]  Relevo plano a suavemente ondulado, margeado em três quadrantes pelo reservatório de Yguazú. Altitude \textasciitilde 220m. Área de \textasciitilde 444 km². Geologicamente muito estável, risco sísmico desprezível.
\item[Alvos Estratégicos]  Capital Nacional da Banana (Maior polo exportador de frutas tropicais do Paraguai); Portos fluviais no Lago Yguazú; Áreas de cultivo intensivo protegidas pela massa de água do reservatório. Ausência de alvos militares diretos, oferecendo boa invisibilidade tática.
\item[Fallout]  Ventos predominantes do Norte (NE) e Leste. A grande massa de água do lago pode influenciar a dispersão de partículas em baixa altitude.
\item[População]  12.877 habitantes (Censo 2022 Final). População muito jovem (idade mediana de 23 anos) e essencialmente rural (90\%).
\item[Densidade]  \textasciitilde 29,0 hab/\allowbreak km² (Baixa densidade, oferecendo isolamento tático superior e privacidade estratégica facilitada pela geografia peninsular).
\item[Vias de Acesso]  Acesso via Rota D025 a partir da Ruta PY02. Localização geográfica isolada por ser uma "península" no Lago Yguazú, o que dificulta o acesso de grandes massas, mas exige logística fluvial ou terrestre redundante.
\item[Serviços]  Unidade de Saúde da Família (USF) local. Dependência de J. Eulogio Estigarribia (Campo 9) para saúde de alta complexidade e logística comercial diversificada.
\item[Custo de Vida]  Baixo; economia baseada na exportação frutícola e agricultura de subsistência.
\item[Preço da Terra]  US\$ 4.000-7.000/\allowbreak ha (áreas mecanizáveis ou com plantios de banana consolidados); áreas próximas ao lago com valorização turística e residencial.
\item[Hidrologia]  Risco Moderado de Isolamento e Cheias. A oscilação do nível do reservatório de Yguazú pode afetar áreas baixas e acessos internos. O isolamento por terra pode ocorrer em eventos climáticos extremos.
\item[Clima]  Subtropical Úmido com forte influência termorreguladora da massa de água do lago. Verões quentes e chuvosos (média 1.700 mm/\allowbreak ano).
\item[Energia]  Rede nacional (Itaipu); instável em áreas rurais remotas.
\item[Água]  Abundância hídrica massiva via reservatório de Yguazú e poços artesianos de alta vazão.
\item[Qualidade do Solo]  Solo franco-arenoso a argiloso; fértil e profundo, ideal para fruticultura (banana, abacaxi), soja e milho.
\item[Resources Locais]  Maior produtor nacional de bananas. Elevadíssimo potencial para autossuficiência calórica e hídrica. Potencial para pesca de subsistência e comercial de água doce.
\item[Segurança]  Zona rural pacífica e estável. Baixíssima criminalidade urbana violenta. Estabilidade social sustentada pelo sucesso do agronegócio familiar de exportação. Geografia peninsular favorece o controle de acessos.
\item[Leis Local: Município de criação recente (2008), emancipado de Raúl A. Oviedo. Livre de Restrição de Fronteira]  Fora da faixa de 50km da fronteira, permitindo plena titularidade para brasileiros.
\end{description}
\textbf{Fonte climática:} NASA POWER Climatology API (período 2001-2020)
\textbf{Fonte luminosa:} estimativa world atlas

\paragraph{Irradiação Solar
(kWh/\allowbreak m²/\allowbreak dia)}\label{irradiauxe7uxe3o-solar-kwhmuxb2dia}

\begin{longtable}[]{@{}
  >{\raggedright\arraybackslash}p{(\linewidth - 24\tabcolsep) * \real{0.0746}}
  >{\raggedright\arraybackslash}p{(\linewidth - 24\tabcolsep) * \real{0.0746}}
  >{\raggedright\arraybackslash}p{(\linewidth - 24\tabcolsep) * \real{0.0746}}
  >{\raggedright\arraybackslash}p{(\linewidth - 24\tabcolsep) * \real{0.0746}}
  >{\raggedright\arraybackslash}p{(\linewidth - 24\tabcolsep) * \real{0.0746}}
  >{\raggedright\arraybackslash}p{(\linewidth - 24\tabcolsep) * \real{0.0746}}
  >{\raggedright\arraybackslash}p{(\linewidth - 24\tabcolsep) * \real{0.0746}}
  >{\raggedright\arraybackslash}p{(\linewidth - 24\tabcolsep) * \real{0.0746}}
  >{\raggedright\arraybackslash}p{(\linewidth - 24\tabcolsep) * \real{0.0746}}
  >{\raggedright\arraybackslash}p{(\linewidth - 24\tabcolsep) * \real{0.0746}}
  >{\raggedright\arraybackslash}p{(\linewidth - 24\tabcolsep) * \real{0.0746}}
  >{\raggedright\arraybackslash}p{(\linewidth - 24\tabcolsep) * \real{0.0746}}
  >{\raggedright\arraybackslash}p{(\linewidth - 24\tabcolsep) * \real{0.1045}}@{}}
\toprule\noalign{}
\begin{minipage}[b]{\linewidth}\raggedright
Jan
\end{minipage} & \begin{minipage}[b]{\linewidth}\raggedright
Fev
\end{minipage} & \begin{minipage}[b]{\linewidth}\raggedright
Mar
\end{minipage} & \begin{minipage}[b]{\linewidth}\raggedright
Abr
\end{minipage} & \begin{minipage}[b]{\linewidth}\raggedright
Mai
\end{minipage} & \begin{minipage}[b]{\linewidth}\raggedright
Jun
\end{minipage} & \begin{minipage}[b]{\linewidth}\raggedright
Jul
\end{minipage} & \begin{minipage}[b]{\linewidth}\raggedright
Ago
\end{minipage} & \begin{minipage}[b]{\linewidth}\raggedright
Set
\end{minipage} & \begin{minipage}[b]{\linewidth}\raggedright
Out
\end{minipage} & \begin{minipage}[b]{\linewidth}\raggedright
Nov
\end{minipage} & \begin{minipage}[b]{\linewidth}\raggedright
Dez
\end{minipage} & \begin{minipage}[b]{\linewidth}\raggedright
Média
\end{minipage} \\
\midrule\noalign{}
\endhead
\bottomrule\noalign{}
\endlastfoot
6.58 & 6.08 & 5.45 & 4.47 & 3.30 & 2.90 & 3.20 & 3.96 & 4.53 & 5.31 &
6.35 & 6.66 & \textbf{4.9} \\
\end{longtable}

\textbf{Inclinação solar recomendada:} 25° N (anual) · 35° N (inverno
jun-ago) · 15° N (verão nov-jan)

\paragraph{Precipitação (mm/\allowbreak mês)}\label{precipitauxe7uxe3o-mmmuxeas}

\begin{longtable}[]{@{}
  >{\raggedright\arraybackslash}p{(\linewidth - 24\tabcolsep) * \real{0.0704}}
  >{\raggedright\arraybackslash}p{(\linewidth - 24\tabcolsep) * \real{0.0704}}
  >{\raggedright\arraybackslash}p{(\linewidth - 24\tabcolsep) * \real{0.0704}}
  >{\raggedright\arraybackslash}p{(\linewidth - 24\tabcolsep) * \real{0.0704}}
  >{\raggedright\arraybackslash}p{(\linewidth - 24\tabcolsep) * \real{0.0704}}
  >{\raggedright\arraybackslash}p{(\linewidth - 24\tabcolsep) * \real{0.0704}}
  >{\raggedright\arraybackslash}p{(\linewidth - 24\tabcolsep) * \real{0.0704}}
  >{\raggedright\arraybackslash}p{(\linewidth - 24\tabcolsep) * \real{0.0704}}
  >{\raggedright\arraybackslash}p{(\linewidth - 24\tabcolsep) * \real{0.0704}}
  >{\raggedright\arraybackslash}p{(\linewidth - 24\tabcolsep) * \real{0.0704}}
  >{\raggedright\arraybackslash}p{(\linewidth - 24\tabcolsep) * \real{0.0704}}
  >{\raggedright\arraybackslash}p{(\linewidth - 24\tabcolsep) * \real{0.0704}}
  >{\raggedright\arraybackslash}p{(\linewidth - 24\tabcolsep) * \real{0.1549}}@{}}
\toprule\noalign{}
\begin{minipage}[b]{\linewidth}\raggedright
Jan
\end{minipage} & \begin{minipage}[b]{\linewidth}\raggedright
Fev
\end{minipage} & \begin{minipage}[b]{\linewidth}\raggedright
Mar
\end{minipage} & \begin{minipage}[b]{\linewidth}\raggedright
Abr
\end{minipage} & \begin{minipage}[b]{\linewidth}\raggedright
Mai
\end{minipage} & \begin{minipage}[b]{\linewidth}\raggedright
Jun
\end{minipage} & \begin{minipage}[b]{\linewidth}\raggedright
Jul
\end{minipage} & \begin{minipage}[b]{\linewidth}\raggedright
Ago
\end{minipage} & \begin{minipage}[b]{\linewidth}\raggedright
Set
\end{minipage} & \begin{minipage}[b]{\linewidth}\raggedright
Out
\end{minipage} & \begin{minipage}[b]{\linewidth}\raggedright
Nov
\end{minipage} & \begin{minipage}[b]{\linewidth}\raggedright
Dez
\end{minipage} & \begin{minipage}[b]{\linewidth}\raggedright
Total/\allowbreak ano
\end{minipage} \\
\midrule\noalign{}
\endhead
\bottomrule\noalign{}
\endlastfoot
142 & 130 & 126 & 152 & 177 & 94 & 79 & 55 & 103 & 189 & 188 & 172 &
\textbf{1607 mm} \\
\end{longtable}

\paragraph{Poluição Luminosa}\label{poluiuxe7uxe3o-luminosa}

\begin{longtable}[]{@{}ll@{}}
\toprule\noalign{}
Parâmetro & Valor \\
\midrule\noalign{}
\endhead
\bottomrule\noalign{}
\endlastfoot
Escala Bortle & 4 --- Céu rural-suburbano \\
Radiância artificial & 3.0 nW/\allowbreak cm²/\allowbreak sr \\
\end{longtable}
\paragraph{Solo (SoilGrids 2.0, média ponderada 0--30
cm)}

\begin{longtable}[]{@{}ll@{}}
\toprule\noalign{}
Parâmetro & Valor \\
\midrule\noalign{}
\endhead
\bottomrule\noalign{}
\endlastfoot
pH (H$_2$O) & 5.25 \\
Carbono orgânico (SOC) & 18.57 g/\allowbreak kg \\
Argila & 28.72 \% \\
Areia & 47.71 \% \\
Silte (calc.) & 23.6 \% \\
Densidade aparente & 1.29 g/\allowbreak cm³ \\
\textbf{Aptidão agrícola} & \textbf{Média} \\
\end{longtable}

Fonte: ISRIC SoilGrids 2.0 via WCS. Coords: -25.3167°, -55.55°. Média
ponderada camadas 0-5, 5-15, 15-30 cm.
\subsubsection{Indicadores Sociais}

\textbf{Fontes:} PNUD 2020, INE\footnote{Instituto Nacional de Estadística (Instituto Nacional de Estatística), órgão oficial de dados demográficos e censitários.} Censo 2022, INE\footnote{Instituto Nacional de Estadística (Instituto Nacional de Estatística), órgão oficial de dados demográficos e censitários.} EPHC 2023, Ministerio
Público 2024\\
\textbf{Nota:} Valores marcados como \emph{(dept.)} referem-se ao
departamento de Caaguazú; valores \emph{(dist.)} são específicos deste
distrito. Dados marcados \emph{(est.)} são estimativas.

\begin{longtable}[]{@{}
  >{\raggedright\arraybackslash}p{(\linewidth - 4\tabcolsep) * \real{0.4231}}
  >{\raggedright\arraybackslash}p{(\linewidth - 4\tabcolsep) * \real{0.2692}}
  >{\raggedright\arraybackslash}p{(\linewidth - 4\tabcolsep) * \real{0.3077}}@{}}
\toprule\noalign{}
\begin{minipage}[b]{\linewidth}\raggedright
Indicador
\end{minipage} & \begin{minipage}[b]{\linewidth}\raggedright
Valor
\end{minipage} & \begin{minipage}[b]{\linewidth}\raggedright
Âmbito
\end{minipage} \\
\midrule\noalign{}
\endhead
\bottomrule\noalign{}
\endlastfoot
IDH (2020) & 0,715 & dept. \\
Ranking IDH nacional & 7/\allowbreak 18 & dept. \\
Esperança de vida & 73,1 anos & dept. \\
Escolaridade média & 7.6 anos & dist. \\
RNB per capita (USD PPA) & 9.000 & dept. \\
Pobreza monetária (\%) & 28,5\% & dept. \\
Pobreza extrema (\%) & 6,8\% & dept. \\
Índice de Gini & 0,447 & dept. \\
Acesso a água potável (\%) & 79,8\% & dept. \\
Acesso a saneamento (\%) & 78,3\% & dept. \\
Taxa de homicídios (est., /\allowbreak 100k hab) & \textasciitilde6--8 (estimativa,
próxima ou acima da média nacional) & dept. \\
Índice de segurança & médio & dept. \\
População (Censo 2022) & 12.877 hab. & dist. \\
Idade mediana & 23.0 anos & dist. \\
Taxa de fecundidade & N/\allowbreak D & dist. \\
\end{longtable}
\subsubsection{Combustível}

\textbf{Referência:} PETROPAR /\allowbreak  postos locais (2024) \textbf{Tipo de
localidade:} Interior

\begin{longtable}[]{@{}lll@{}}
\toprule\noalign{}
Combustível & USD/\allowbreak litro & Gs/\allowbreak litro (aprox.) \\
\midrule\noalign{}
\endhead
\bottomrule\noalign{}
\endlastfoot
Gasolina 93 oct & 0.98 & 7,252 \\
Gasolina 97 oct (premium) & 1.08 & 7,992 \\
Diesel & 0.91 & 6,734 \\
\end{longtable}

\begin{quote}
Preços podem variar ±5\% conforme posto e sazonalidade. Chaco e interior
remoto apresentam maior variação.
\end{quote}

\subsubsection{Cobertura Celular}

\textbf{Fonte:} CONATEL PY /\allowbreak  operadoras (2024)

\begin{longtable}[]{@{}ll@{}}
\toprule\noalign{}
Parâmetro & Valor \\
\midrule\noalign{}
\endhead
\bottomrule\noalign{}
\endlastfoot
Cobertura 4G população (dept.) & 88\% \\
Cobertura 4G área rural & 72\% \\
Melhor operadora & Tigo \\
Qualidade rural & boa \\
\end{longtable}

\begin{quote}
Para áreas rurais fora do núcleo urbano, recomenda-se chip Tigo como
principal e Personal como backup.
\end{quote}

\subsubsection{Internet}

\textbf{Fonte:} CONATEL /\allowbreak  Speedtest Ookla (2024)

\begin{longtable}[]{@{}ll@{}}
\toprule\noalign{}
Parâmetro & Valor \\
\midrule\noalign{}
\endhead
\bottomrule\noalign{}
\endlastfoot
Velocidade média download & 55 Mbps \\
Domicílios com internet (dept.) & 62\% \\
Tecnologia predominante & rádio/\allowbreak fibra \\
Opção rural & Starlink disponível (\textasciitilde USD 44/\allowbreak mês) \\
\end{longtable}

\subsubsection{Mercado Imobiliário e Terra
Rural}

\textbf{Fonte:} INDERT /\allowbreak  Clasificados.com.py (2024)

\begin{longtable}[]{@{}ll@{}}
\toprule\noalign{}
Tipo & Referência \\
\midrule\noalign{}
\endhead
\bottomrule\noalign{}
\endlastfoot
Terra agrícola alta prod. (USD/\allowbreak ha) & 8,800 \\
Imóvel urbano (USD/\allowbreak m²) & 850 \\
Aluguel 2 quartos (USD/\allowbreak mês) & 350 \\
\end{longtable}

\begin{quote}
Valores de referência departamental. Localidades menores podem ter
preços 20--40\% abaixo da capital departamental. \#\#\# Saúde
\end{quote}

\textbf{Fonte:} MSPBS /\allowbreak  IPS Paraguay (2024-2026), consolidação
departamental e proxy local

\begin{longtable}[]{@{}ll@{}}
\toprule\noalign{}
Serviço & Disponibilidade \\
\midrule\noalign{}
\endhead
\bottomrule\noalign{}
\endlastfoot
USF /\allowbreak  Posto de Saúde & sim \\
Hospital Regional & não \\
IPS (seguro social) & sim \\
Farmácia & sim \\
Distância ao hospital de referência & \textasciitilde134 km (Coronel
Oviedo) \\
\end{longtable}

\textbf{Principais estabelecimentos:} USF local; referencia hospitalar
em Coronel Oviedo

\textbf{Observação para imigrantes:} Atendimento primario local ou em
raio curto; casos de maior complexidade seguem para Coronel Oviedo.
Cobertura privada continua recomendada para especialidades e urgências
de maior complexidade.
\newpage
\secaoDiagnostico{Tres de Febrero}{\mapaConteudo{-25.3000}{-55.8333}}{Tres de Febrero é uma excelente opção para quem busca discrição e solo
de alta qualidade em uma região estável do departamento de Caaguazú. Sua
escala reduzida facilita o controle social e a gestão de recursos em
crises.

\subsubsection{Por que sim}

\begin{itemize}
\tightlist
\item
  Qualidade superior do solo.
\item
  Baixa visibilidade estratégica (segurança por obscuridade).
\end{itemize}

\subsubsection{Por que não}

\begin{itemize}
\tightlist
\item
  Acessibilidade limitada em comparação com distritos na PY02.
\end{itemize}

\subsubsection{Recomendação
Técnica}

Indicado para projetos de autossuficiência agrícola intensiva e retiro
estratégico de longo prazo.

\subsubsection{Analise de Sensibilidade}

\begin{itemize}
\tightlist
\item
  Cenario otimista: GSS 6.3 (Melhoria na rede de saúde local).
\item
  Cenario conservador: GSS 5.5 (Isolamento logístico prolongado).
\end{itemize}

\subsubsection{}

\begin{itemize}
\tightlist
\item
  Censo 2022 INE\footnote{Instituto Nacional de Estadística (Instituto Nacional de Estatística), órgão oficial de dados demográficos e censitários.}.
\item
  Boletins meteorológicos regionais.
\item
  Dados de produção agropecuária departamental.
\end{itemize}}
\begin{table}[ht!]
\small \begin{tabular}{p{3.0cm}p{0.8cm}p{6.0cm}} \toprule
\textbf{Critério} & \textbf{Nota} & \textbf{Justificativa} \\ \midrule
A - Ameaças Estratégicas & 4.6 & - \\
B - Risco Social & 6.0 & - \\
C - Riscos Naturais & 7.0 & - \\
D - Autossuficiência & 6.5 & - \\
E - Institucional & 6.0 & - \\
\midrule \textbf{GSS FINAL} & \textbf{5.9} & \textbf{Seguro (Potencial moderado de resiliência).} \\ \bottomrule \end{tabular}
\end{table}
\subsubsection{Vulnerabilidades e
Mitigação}\label{vuln-Tres_de_Febrero-vulnerabilidades-e-mitigauxe7uxe3o}

\begin{itemize}
\item
  \textbf{Conectividade:} Dependência de acessos que podem ser
  bloqueados em crises. Mitigação: Mapear rotas rurais alternativas.
\item
  \textbf{Dependência Técnica:} Insumos agrícolas externos. Mitigação:
  Estoque de sementes crioulas e técnicas de agroecologia.
\item
  INE\footnote{Instituto Nacional de Estadística (Instituto Nacional de Estatística), órgão oficial de dados demográficos e censitários.} Censo 2022: 7.492 habitantes.
\item
  Preço da Terra: Gs 40M a 65M (solo de alta qualidade).
\item
  Homicídios: 3,5/\allowbreak 100k.
\item
  População distrital: 7.492 habitantes.
\item
  Densidade demografica: 35,0 hab/\allowbreak km2.
\item
  Conectividade viaria: 70/\allowbreak 100.
\item
  Exposicao hidrometeorologica: 30/\allowbreak 100.
\item
  Curto prazo: Flutuação de preços de commodities afetando a economia
  local.
\item
  Médio prazo: Pressão por desmatamento em áreas remanescentes.
\end{itemize}

\subsubsection{Diagnostico Integrado}\label{vuln-Tres_de_Febrero-diagnostico-integrado}

Tres de Febrero é uma excelente opção para quem busca discrição e solo
de alta qualidade em uma região estável do departamento de Caaguazú. Sua
escala reduzida facilita o controle social e a gestão de recursos em
crises.

\subsubsection{Por que sim}\label{vuln-Tres_de_Febrero-por-que-sim}

\begin{itemize}
\tightlist
\item
  Qualidade superior do solo.
\item
  Baixa visibilidade estratégica (segurança por obscuridade).
\end{itemize}

\subsubsection{Por que não}\label{vuln-Tres_de_Febrero-por-que-nuxe3o}

\begin{itemize}
\tightlist
\item
  Acessibilidade limitada em comparação com distritos na PY02.
\end{itemize}

\subsubsection{Recomendação
Técnica}\label{vuln-Tres_de_Febrero-recomendauxe7uxe3o-tuxe9cnica}

Indicado para projetos de autossuficiência agrícola intensiva e retiro
estratégico de longo prazo.

\begin{itemize}
\item
  Cenario otimista: GSS 6.3 (Melhoria na rede de saúde local).
\item
  Cenario conservador: GSS 5.5 (Isolamento logístico prolongado).
\item
  Censo 2022 INE\footnote{Instituto Nacional de Estadística (Instituto Nacional de Estatística), órgão oficial de dados demográficos e censitários.}.
\item
  Boletins meteorológicos regionais.
\item
  Dados de produção agropecuária departamental.
\end{itemize}

\subsubsection{Combustível}\label{vuln-Tres_de_Febrero-combustuxedvel}

\textbf{Referência:} PETROPAR /\allowbreak  postos locais (2024) \textbf{Tipo de
localidade:} Interior

\begin{longtable}[]{@{}lll@{}}
\toprule\noalign{}
Combustível & USD/\allowbreak litro & Gs/\allowbreak litro (aprox.) \\
\midrule\noalign{}
\endhead
\bottomrule\noalign{}
\endlastfoot
Gasolina 93 oct & 0.98 & 7,252 \\
Gasolina 97 oct (premium) & 1.08 & 7,992 \\
Diesel & 0.91 & 6,734 \\
\end{longtable}

\begin{quote}
Preços podem variar ±5\% conforme posto e sazonalidade. Chaco e interior
remoto apresentam maior variação.
\end{quote}

\subsubsection{Cobertura Celular}\label{vuln-Tres_de_Febrero-cobertura-celular}

\textbf{Fonte:} CONATEL PY /\allowbreak  operadoras (2024)

\begin{longtable}[]{@{}ll@{}}
\toprule\noalign{}
Parâmetro & Valor \\
\midrule\noalign{}
\endhead
\bottomrule\noalign{}
\endlastfoot
Cobertura 4G população (dept.) & 88\% \\
Cobertura 4G área rural & 72\% \\
Melhor operadora & Tigo \\
Qualidade rural & boa \\
\end{longtable}

\begin{quote}
Para áreas rurais fora do núcleo urbano, recomenda-se chip Tigo como
principal e Personal como backup.
\end{quote}

\subsubsection{Internet}\label{vuln-Tres_de_Febrero-internet}

\textbf{Fonte:} CONATEL /\allowbreak  Speedtest Ookla (2024)

\begin{longtable}[]{@{}ll@{}}
\toprule\noalign{}
Parâmetro & Valor \\
\midrule\noalign{}
\endhead
\bottomrule\noalign{}
\endlastfoot
Velocidade média download & 55 Mbps \\
Domicílios com internet (dept.) & 62\% \\
Tecnologia predominante & rádio/\allowbreak fibra \\
Opção rural & Starlink disponível (\textasciitilde USD 44/\allowbreak mês) \\
\end{longtable}

\subsubsection{Mercado Imobiliário e Terra
Rural}\label{vuln-Tres_de_Febrero-mercado-imobiliuxe1rio-e-terra-rural}

\textbf{Fonte:} INDERT /\allowbreak  Clasificados.com.py (2024)

\begin{longtable}[]{@{}ll@{}}
\toprule\noalign{}
Tipo & Referência \\
\midrule\noalign{}
\endhead
\bottomrule\noalign{}
\endlastfoot
Terra agrícola alta prod. (USD/\allowbreak ha) & 8,800 \\
Imóvel urbano (USD/\allowbreak m²) & 850 \\
Aluguel 2 quartos (USD/\allowbreak mês) & 350 \\
\end{longtable}

\begin{quote}
Valores de referência departamental. Localidades menores podem ter
preços 20--40\% abaixo da capital departamental. \#\#\# Saúde
\end{quote}

\textbf{Fonte:} MSPBS /\allowbreak  IPS Paraguay (2024-2026), consolidação
departamental e proxy local

\begin{longtable}[]{@{}ll@{}}
\toprule\noalign{}
Serviço & Disponibilidade \\
\midrule\noalign{}
\endhead
\bottomrule\noalign{}
\endlastfoot
USF /\allowbreak  Posto de Saúde & sim \\
Hospital Regional & não \\
IPS (seguro social) & sim \\
Farmácia & sim \\
Distância ao hospital de referência & \textasciitilde88 km (Coronel
Oviedo) \\
\end{longtable}

\textbf{Principais estabelecimentos:} USF local; referencia hospitalar
em Coronel Oviedo

\textbf{Observação para imigrantes:} Atendimento primario local ou em
raio curto; casos de maior complexidade seguem para Coronel Oviedo.
Cobertura privada continua recomendada para especialidades e urgências
de maior complexidade.
\subsubsection{Dossiê de Campo}\label{dossie-Tres_de_Febrero-dossiuxea-de-campo}
\begin{description}[style=nextline,leftmargin=0.5cm,font=\bfseries]
\item[Coordenadas]  25°18'S 55°50'W.
\item[Topografia]  Relevo predominantemente plano. Área de \textasciitilde 214 km².
\item[Alvos Estrategicos]  Pequena área de produção agrícola sem alvos militares. Distante de rotas internacionais.
\item[Fallout]  Ventos predominantes do Norte/\allowbreak Nordeste.
\item[Populacao]  7.492 (Censo 2022 Final).
\item[Densidade]  \textasciitilde 35,0 hab/\allowbreak km².
\item[Vias de Acesso]  Acesso por rotas secundárias ligando à Ruta PY02 e PY13.
\item[Servicos]  Saúde básica local. Dependência de Caaguazú ou Coronel Oviedo para complexidade.
\item[Hidrologia]  Baixo risco de inundação urbana. Áreas rurais com drenagem natural eficiente.
\item[Sismicidade]  Nula/\allowbreak Baixa.
\item[Solo]  Oxisois/\allowbreak Alfisois, altamente produtivos.
\item[Agua]  Poços artesianos privados e comunitários.
\item[Energia]  Rede ANDE\footnote{Administración Nacional de Electricidade (Administração Nacional de Eletricidade), companhia estatal de energia.}.
\item[Producao]  Soja, milho e pequena pecuária.
\item[Seguranca]  Ambiente rural pacato, baixo índice de criminalidade.
\item[Clima Social]  Estabilidade social elevada, população predominantemente rural.
\end{description}
\textbf{Fonte climática:} NASA POWER Climatology API (período 2001-2020)
\textbf{Fonte luminosa:} estimativa world atlas

\paragraph{Irradiação Solar
(kWh/\allowbreak m²/\allowbreak dia)}\label{irradiauxe7uxe3o-solar-kwhmuxb2dia}

\begin{longtable}[]{@{}
  >{\raggedright\arraybackslash}p{(\linewidth - 24\tabcolsep) * \real{0.0746}}
  >{\raggedright\arraybackslash}p{(\linewidth - 24\tabcolsep) * \real{0.0746}}
  >{\raggedright\arraybackslash}p{(\linewidth - 24\tabcolsep) * \real{0.0746}}
  >{\raggedright\arraybackslash}p{(\linewidth - 24\tabcolsep) * \real{0.0746}}
  >{\raggedright\arraybackslash}p{(\linewidth - 24\tabcolsep) * \real{0.0746}}
  >{\raggedright\arraybackslash}p{(\linewidth - 24\tabcolsep) * \real{0.0746}}
  >{\raggedright\arraybackslash}p{(\linewidth - 24\tabcolsep) * \real{0.0746}}
  >{\raggedright\arraybackslash}p{(\linewidth - 24\tabcolsep) * \real{0.0746}}
  >{\raggedright\arraybackslash}p{(\linewidth - 24\tabcolsep) * \real{0.0746}}
  >{\raggedright\arraybackslash}p{(\linewidth - 24\tabcolsep) * \real{0.0746}}
  >{\raggedright\arraybackslash}p{(\linewidth - 24\tabcolsep) * \real{0.0746}}
  >{\raggedright\arraybackslash}p{(\linewidth - 24\tabcolsep) * \real{0.0746}}
  >{\raggedright\arraybackslash}p{(\linewidth - 24\tabcolsep) * \real{0.1045}}@{}}
\toprule\noalign{}
\begin{minipage}[b]{\linewidth}\raggedright
Jan
\end{minipage} & \begin{minipage}[b]{\linewidth}\raggedright
Fev
\end{minipage} & \begin{minipage}[b]{\linewidth}\raggedright
Mar
\end{minipage} & \begin{minipage}[b]{\linewidth}\raggedright
Abr
\end{minipage} & \begin{minipage}[b]{\linewidth}\raggedright
Mai
\end{minipage} & \begin{minipage}[b]{\linewidth}\raggedright
Jun
\end{minipage} & \begin{minipage}[b]{\linewidth}\raggedright
Jul
\end{minipage} & \begin{minipage}[b]{\linewidth}\raggedright
Ago
\end{minipage} & \begin{minipage}[b]{\linewidth}\raggedright
Set
\end{minipage} & \begin{minipage}[b]{\linewidth}\raggedright
Out
\end{minipage} & \begin{minipage}[b]{\linewidth}\raggedright
Nov
\end{minipage} & \begin{minipage}[b]{\linewidth}\raggedright
Dez
\end{minipage} & \begin{minipage}[b]{\linewidth}\raggedright
Média
\end{minipage} \\
\midrule\noalign{}
\endhead
\bottomrule\noalign{}
\endlastfoot
6.58 & 6.08 & 5.45 & 4.47 & 3.30 & 2.90 & 3.20 & 3.96 & 4.53 & 5.31 &
6.35 & 6.66 & \textbf{4.9} \\
\end{longtable}

\textbf{Inclinação solar recomendada:} 25° N (anual) · 35° N (inverno
jun-ago) · 15° N (verão nov-jan)

\paragraph{Precipitação (mm/\allowbreak mês)}\label{precipitauxe7uxe3o-mmmuxeas}

\begin{longtable}[]{@{}
  >{\raggedright\arraybackslash}p{(\linewidth - 24\tabcolsep) * \real{0.0704}}
  >{\raggedright\arraybackslash}p{(\linewidth - 24\tabcolsep) * \real{0.0704}}
  >{\raggedright\arraybackslash}p{(\linewidth - 24\tabcolsep) * \real{0.0704}}
  >{\raggedright\arraybackslash}p{(\linewidth - 24\tabcolsep) * \real{0.0704}}
  >{\raggedright\arraybackslash}p{(\linewidth - 24\tabcolsep) * \real{0.0704}}
  >{\raggedright\arraybackslash}p{(\linewidth - 24\tabcolsep) * \real{0.0704}}
  >{\raggedright\arraybackslash}p{(\linewidth - 24\tabcolsep) * \real{0.0704}}
  >{\raggedright\arraybackslash}p{(\linewidth - 24\tabcolsep) * \real{0.0704}}
  >{\raggedright\arraybackslash}p{(\linewidth - 24\tabcolsep) * \real{0.0704}}
  >{\raggedright\arraybackslash}p{(\linewidth - 24\tabcolsep) * \real{0.0704}}
  >{\raggedright\arraybackslash}p{(\linewidth - 24\tabcolsep) * \real{0.0704}}
  >{\raggedright\arraybackslash}p{(\linewidth - 24\tabcolsep) * \real{0.0704}}
  >{\raggedright\arraybackslash}p{(\linewidth - 24\tabcolsep) * \real{0.1549}}@{}}
\toprule\noalign{}
\begin{minipage}[b]{\linewidth}\raggedright
Jan
\end{minipage} & \begin{minipage}[b]{\linewidth}\raggedright
Fev
\end{minipage} & \begin{minipage}[b]{\linewidth}\raggedright
Mar
\end{minipage} & \begin{minipage}[b]{\linewidth}\raggedright
Abr
\end{minipage} & \begin{minipage}[b]{\linewidth}\raggedright
Mai
\end{minipage} & \begin{minipage}[b]{\linewidth}\raggedright
Jun
\end{minipage} & \begin{minipage}[b]{\linewidth}\raggedright
Jul
\end{minipage} & \begin{minipage}[b]{\linewidth}\raggedright
Ago
\end{minipage} & \begin{minipage}[b]{\linewidth}\raggedright
Set
\end{minipage} & \begin{minipage}[b]{\linewidth}\raggedright
Out
\end{minipage} & \begin{minipage}[b]{\linewidth}\raggedright
Nov
\end{minipage} & \begin{minipage}[b]{\linewidth}\raggedright
Dez
\end{minipage} & \begin{minipage}[b]{\linewidth}\raggedright
Total/\allowbreak ano
\end{minipage} \\
\midrule\noalign{}
\endhead
\bottomrule\noalign{}
\endlastfoot
134 & 138 & 121 & 148 & 174 & 93 & 73 & 51 & 104 & 186 & 185 & 172 &
\textbf{1581 mm} \\
\end{longtable}

\paragraph{Poluição Luminosa}\label{poluiuxe7uxe3o-luminosa}

\begin{longtable}[]{@{}ll@{}}
\toprule\noalign{}
Parâmetro & Valor \\
\midrule\noalign{}
\endhead
\bottomrule\noalign{}
\endlastfoot
Escala Bortle & 4 --- Céu rural-suburbano \\
Radiância artificial & 3.0 nW/\allowbreak cm²/\allowbreak sr \\
\end{longtable}
\paragraph{Solo (SoilGrids 2.0, média ponderada 0--30
cm)}

\begin{longtable}[]{@{}ll@{}}
\toprule\noalign{}
Parâmetro & Valor \\
\midrule\noalign{}
\endhead
\bottomrule\noalign{}
\endlastfoot
pH (H$_2$O) & 5.25 \\
Carbono orgânico (SOC) & 18.53 g/\allowbreak kg \\
Argila & 23.95 \% \\
Areia & 58.36 \% \\
Silte (calc.) & 17.7 \% \\
Densidade aparente & 1.29 g/\allowbreak cm³ \\
\textbf{Aptidão agrícola} & \textbf{Média} \\
\end{longtable}

Fonte: ISRIC SoilGrids 2.0 via WCS. Coords: -25.3°, -55.8333°. Média
ponderada camadas 0-5, 5-15, 15-30 cm.

\begin{itemize}
\tightlist
\item
  \textbf{Agua:} Poços artesianos privados e comunitários.
\item
  \textbf{Energia:} Rede ANDE\footnote{Administración Nacional de Electricidade (Administração Nacional de Eletricidade), companhia estatal de energia.}.
\item
  \textbf{Producao:} Soja, milho e pequena pecuária.
\end{itemize}
\subsubsection{Indicadores Sociais}

\textbf{Fontes:} PNUD 2020, INE\footnote{Instituto Nacional de Estadística (Instituto Nacional de Estatística), órgão oficial de dados demográficos e censitários.} Censo 2022, INE\footnote{Instituto Nacional de Estadística (Instituto Nacional de Estatística), órgão oficial de dados demográficos e censitários.} EPHC 2023, Ministerio
Público 2024\\
\textbf{Nota:} Valores marcados como \emph{(dept.)} referem-se ao
departamento de Caaguazú; valores \emph{(dist.)} são específicos deste
distrito. Dados marcados \emph{(est.)} são estimativas.

\begin{longtable}[]{@{}
  >{\raggedright\arraybackslash}p{(\linewidth - 4\tabcolsep) * \real{0.4231}}
  >{\raggedright\arraybackslash}p{(\linewidth - 4\tabcolsep) * \real{0.2692}}
  >{\raggedright\arraybackslash}p{(\linewidth - 4\tabcolsep) * \real{0.3077}}@{}}
\toprule\noalign{}
\begin{minipage}[b]{\linewidth}\raggedright
Indicador
\end{minipage} & \begin{minipage}[b]{\linewidth}\raggedright
Valor
\end{minipage} & \begin{minipage}[b]{\linewidth}\raggedright
Âmbito
\end{minipage} \\
\midrule\noalign{}
\endhead
\bottomrule\noalign{}
\endlastfoot
IDH (2020) & 0,715 & dept. \\
Ranking IDH nacional & 7/\allowbreak 18 & dept. \\
Esperança de vida & 73,1 anos & dept. \\
Escolaridade média & 8,4 anos & dept. \\
RNB per capita (USD PPA) & 9.000 & dept. \\
Pobreza monetária (\%) & 28,5\% & dept. \\
Pobreza extrema (\%) & 6,8\% & dept. \\
Índice de Gini & 0,447 & dept. \\
Acesso a água potável (\%) & 79,8\% & dept. \\
Acesso a saneamento (\%) & 78,3\% & dept. \\
Taxa de homicídios (est., /\allowbreak 100k hab) & \textasciitilde6--8 (estimativa,
próxima ou acima da média nacional) & dept. \\
Índice de segurança & médio & dept. \\
População (Censo 2022) & 7.492 hab. & dist. \\
Idade mediana & 32.0 anos & dist. \\
Taxa de fecundidade & N/\allowbreak D & dist. \\
\end{longtable}
\subsubsection{Combustível}

\textbf{Referência:} PETROPAR /\allowbreak  postos locais (2024) \textbf{Tipo de
localidade:} Interior

\begin{longtable}[]{@{}lll@{}}
\toprule\noalign{}
Combustível & USD/\allowbreak litro & Gs/\allowbreak litro (aprox.) \\
\midrule\noalign{}
\endhead
\bottomrule\noalign{}
\endlastfoot
Gasolina 93 oct & 0.98 & 7,252 \\
Gasolina 97 oct (premium) & 1.08 & 7,992 \\
Diesel & 0.91 & 6,734 \\
\end{longtable}

\begin{quote}
Preços podem variar ±5\% conforme posto e sazonalidade. Chaco e interior
remoto apresentam maior variação.
\end{quote}

\subsubsection{Cobertura Celular}

\textbf{Fonte:} CONATEL PY /\allowbreak  operadoras (2024)

\begin{longtable}[]{@{}ll@{}}
\toprule\noalign{}
Parâmetro & Valor \\
\midrule\noalign{}
\endhead
\bottomrule\noalign{}
\endlastfoot
Cobertura 4G população (dept.) & 88\% \\
Cobertura 4G área rural & 72\% \\
Melhor operadora & Tigo \\
Qualidade rural & boa \\
\end{longtable}

\begin{quote}
Para áreas rurais fora do núcleo urbano, recomenda-se chip Tigo como
principal e Personal como backup.
\end{quote}

\subsubsection{Internet}

\textbf{Fonte:} CONATEL /\allowbreak  Speedtest Ookla (2024)

\begin{longtable}[]{@{}ll@{}}
\toprule\noalign{}
Parâmetro & Valor \\
\midrule\noalign{}
\endhead
\bottomrule\noalign{}
\endlastfoot
Velocidade média download & 55 Mbps \\
Domicílios com internet (dept.) & 62\% \\
Tecnologia predominante & rádio/\allowbreak fibra \\
Opção rural & Starlink disponível (\textasciitilde USD 44/\allowbreak mês) \\
\end{longtable}

\subsubsection{Mercado Imobiliário e Terra
Rural}

\textbf{Fonte:} INDERT /\allowbreak  Clasificados.com.py (2024)

\begin{longtable}[]{@{}ll@{}}
\toprule\noalign{}
Tipo & Referência \\
\midrule\noalign{}
\endhead
\bottomrule\noalign{}
\endlastfoot
Terra agrícola alta prod. (USD/\allowbreak ha) & 8,800 \\
Imóvel urbano (USD/\allowbreak m²) & 850 \\
Aluguel 2 quartos (USD/\allowbreak mês) & 350 \\
\end{longtable}

\begin{quote}
Valores de referência departamental. Localidades menores podem ter
preços 20--40\% abaixo da capital departamental. \#\#\# Saúde
\end{quote}

\textbf{Fonte:} MSPBS /\allowbreak  IPS Paraguay (2024-2026), consolidação
departamental e proxy local

\begin{longtable}[]{@{}ll@{}}
\toprule\noalign{}
Serviço & Disponibilidade \\
\midrule\noalign{}
\endhead
\bottomrule\noalign{}
\endlastfoot
USF /\allowbreak  Posto de Saúde & sim \\
Hospital Regional & não \\
IPS (seguro social) & sim \\
Farmácia & sim \\
Distância ao hospital de referência & \textasciitilde88 km (Coronel
Oviedo) \\
\end{longtable}

\textbf{Principais estabelecimentos:} USF local; referencia hospitalar
em Coronel Oviedo

\textbf{Observação para imigrantes:} Atendimento primario local ou em
raio curto; casos de maior complexidade seguem para Coronel Oviedo.
Cobertura privada continua recomendada para especialidades e urgências
de maior complexidade.
\newpage
\secaoDiagnostico{Vaqueria}{\mapaConteudo{-24.9166}{-55.7500}}{Vaquería destaca-se pela sua baixíssima densidade demográfica e grande
extensão territorial, oferecendo um ambiente de alta privacidade e
segurança contra ameaças externas. É um refúgio ideal para quem busca
isolamento com potencial produtivo.

\subsubsection{Por que sim}

\begin{itemize}
\tightlist
\item
  Maior disponibilidade de terra per capita na região.
\item
  Ambiente rural extremamente tranquilo.
\end{itemize}

\subsubsection{Por que não}

\begin{itemize}
\tightlist
\item
  Infraestrutura rodoviária e de serviços menos desenvolvida que
  vizinhos.
\end{itemize}

\subsubsection{Recomendação
Técnica}

Altamente recomendado para grandes propriedades de autossuficiência e
projetos que exijam espaço e discrição.

\subsubsection{Analise de Sensibilidade}

\begin{itemize}
\tightlist
\item
  Cenario otimista: GSS 6.4 (Pavimentação asfáltica total da PY13).
\item
  Cenario conservador: GSS 5.4 (Secas prolongadas afetando pastagens).
\end{itemize}

\subsubsection{}

\begin{itemize}
\tightlist
\item
  Dados territoriais MOPC.
\item
  Estatísticas demográficas INE\footnote{Instituto Nacional de Estadística (Instituto Nacional de Estatística), órgão oficial de dados demográficos e censitários.}.
\item
  Relatórios de recursos hídricos departamentais.
\end{itemize}}
\begin{table}[ht!]
\small \begin{tabular}{p{3.0cm}p{0.8cm}p{6.0cm}} \toprule
\textbf{Critério} & \textbf{Nota} & \textbf{Justificativa} \\ \midrule
A - Ameaças Estratégicas & 5.0 & - \\
B - Risco Social & 5.0 & - \\
C - Riscos Naturais & 7.0 & - \\
D - Autossuficiência & 7.0 & - \\
E - Institucional & 6.0 & - \\
\midrule \textbf{GSS FINAL} & \textbf{5.9} & \textbf{Seguro (Potencial moderado de resiliência).} \\ \bottomrule \end{tabular}
\end{table}
\subsubsection{Vulnerabilidades e
Mitigação}\label{vuln-Vaqueria-vulnerabilidades-e-mitigauxe7uxe3o}

\begin{itemize}
\item
  \textbf{Conectividade:} Rotas de acesso podem ser precárias em épocas
  de chuva intensa. Mitigação: Veículos off-road e manutenção de pontes
  locais.
\item
  \textbf{Comunicação:} Sinal de telefonia/\allowbreak internet instável em áreas
  profundas. Mitigação: Sistemas via satélite (Starlink).
\item
  INE\footnote{Instituto Nacional de Estadística (Instituto Nacional de Estatística), órgão oficial de dados demográficos e censitários.} Censo 2022: 10.498 habitantes.
\item
  Densidade: 9,1 hab/\allowbreak km2.
\item
  Homicídios: 3,5/\allowbreak 100k.
\item
  População distrital: 10.498 habitantes.
\item
  Densidade demografica: 9,1 hab/\allowbreak km2.
\item
  Conectividade viaria: 65/\allowbreak 100.
\item
  Exposicao hidrometeorologica: 25/\allowbreak 100.
\item
  Curto prazo: Dificuldade logística em períodos de safra.
\item
  Médio prazo: Mudanças no uso da terra impactando a disponibilidade de
  água.
\end{itemize}

\subsubsection{Diagnostico Integrado}\label{vuln-Vaqueria-diagnostico-integrado}

Vaquería destaca-se pela sua baixíssima densidade demográfica e grande
extensão territorial, oferecendo um ambiente de alta privacidade e
segurança contra ameaças externas. É um refúgio ideal para quem busca
isolamento com potencial produtivo.

\subsubsection{Por que sim}\label{vuln-Vaqueria-por-que-sim}

\begin{itemize}
\tightlist
\item
  Maior disponibilidade de terra per capita na região.
\item
  Ambiente rural extremamente tranquilo.
\end{itemize}

\subsubsection{Por que não}\label{vuln-Vaqueria-por-que-nuxe3o}

\begin{itemize}
\tightlist
\item
  Infraestrutura rodoviária e de serviços menos desenvolvida que
  vizinhos.
\end{itemize}

\subsubsection{Recomendação
Técnica}\label{vuln-Vaqueria-recomendauxe7uxe3o-tuxe9cnica}

Altamente recomendado para grandes propriedades de autossuficiência e
projetos que exijam espaço e discrição.

\begin{itemize}
\item
  Cenario otimista: GSS 6.4 (Pavimentação asfáltica total da PY13).
\item
  Cenario conservador: GSS 5.4 (Secas prolongadas afetando pastagens).
\item
  Dados territoriais MOPC.
\item
  Estatísticas demográficas INE\footnote{Instituto Nacional de Estadística (Instituto Nacional de Estatística), órgão oficial de dados demográficos e censitários.}.
\item
  Relatórios de recursos hídricos departamentais.
\end{itemize}

\subsubsection{Combustível}\label{vuln-Vaqueria-combustuxedvel}

\textbf{Referência:} PETROPAR /\allowbreak  postos locais (2024) \textbf{Tipo de
localidade:} Interior

\begin{longtable}[]{@{}lll@{}}
\toprule\noalign{}
Combustível & USD/\allowbreak litro & Gs/\allowbreak litro (aprox.) \\
\midrule\noalign{}
\endhead
\bottomrule\noalign{}
\endlastfoot
Gasolina 93 oct & 0.98 & 7,252 \\
Gasolina 97 oct (premium) & 1.08 & 7,992 \\
Diesel & 0.91 & 6,734 \\
\end{longtable}

\begin{quote}
Preços podem variar ±5\% conforme posto e sazonalidade. Chaco e interior
remoto apresentam maior variação.
\end{quote}

\subsubsection{Cobertura Celular}\label{vuln-Vaqueria-cobertura-celular}

\textbf{Fonte:} CONATEL PY /\allowbreak  operadoras (2024)

\begin{longtable}[]{@{}ll@{}}
\toprule\noalign{}
Parâmetro & Valor \\
\midrule\noalign{}
\endhead
\bottomrule\noalign{}
\endlastfoot
Cobertura 4G população (dept.) & 88\% \\
Cobertura 4G área rural & 72\% \\
Melhor operadora & Tigo \\
Qualidade rural & boa \\
\end{longtable}

\begin{quote}
Para áreas rurais fora do núcleo urbano, recomenda-se chip Tigo como
principal e Personal como backup.
\end{quote}

\subsubsection{Internet}\label{vuln-Vaqueria-internet}

\textbf{Fonte:} CONATEL /\allowbreak  Speedtest Ookla (2024)

\begin{longtable}[]{@{}ll@{}}
\toprule\noalign{}
Parâmetro & Valor \\
\midrule\noalign{}
\endhead
\bottomrule\noalign{}
\endlastfoot
Velocidade média download & 55 Mbps \\
Domicílios com internet (dept.) & 62\% \\
Tecnologia predominante & rádio/\allowbreak fibra \\
Opção rural & Starlink disponível (\textasciitilde USD 44/\allowbreak mês) \\
\end{longtable}

\subsubsection{Mercado Imobiliário e Terra
Rural}\label{vuln-Vaqueria-mercado-imobiliuxe1rio-e-terra-rural}

\textbf{Fonte:} INDERT /\allowbreak  Clasificados.com.py (2024)

\begin{longtable}[]{@{}ll@{}}
\toprule\noalign{}
Tipo & Referência \\
\midrule\noalign{}
\endhead
\bottomrule\noalign{}
\endlastfoot
Terra agrícola alta prod. (USD/\allowbreak ha) & 8,800 \\
Imóvel urbano (USD/\allowbreak m²) & 850 \\
Aluguel 2 quartos (USD/\allowbreak mês) & 350 \\
\end{longtable}

\begin{quote}
Valores de referência departamental. Localidades menores podem ter
preços 20--40\% abaixo da capital departamental. \#\#\# Saúde
\end{quote}

\textbf{Fonte:} MSPBS /\allowbreak  IPS Paraguay (2024-2026), consolidação
departamental e proxy local

\begin{longtable}[]{@{}ll@{}}
\toprule\noalign{}
Serviço & Disponibilidade \\
\midrule\noalign{}
\endhead
\bottomrule\noalign{}
\endlastfoot
USF /\allowbreak  Posto de Saúde & sim \\
Hospital Regional & não \\
IPS (seguro social) & sim \\
Farmácia & sim \\
Distância ao hospital de referência & \textasciitilde101 km (Coronel
Oviedo) \\
\end{longtable}

\textbf{Principais estabelecimentos:} USF local; referencia hospitalar
em Coronel Oviedo

\textbf{Observação para imigrantes:} Atendimento primario local ou em
raio curto; casos de maior complexidade seguem para Coronel Oviedo.
Cobertura privada continua recomendada para especialidades e urgências
de maior complexidade.
\subsubsection{Dossiê de Campo}\label{dossie-Vaqueria-dossiuxea-de-campo}
\begin{description}[style=nextline,leftmargin=0.5cm,font=\bfseries]
\item[Coordenadas]  24°55'S 55°45'W.
\item[Topografia]  Relevo de planície com ondulações suaves. Área de \textasciitilde 1.150 km².
\item[Alvos Estrategicos]  Grande extensão territorial com baixa densidade, favorável à dispersão. Sem alvos militares.
\item[Fallout]  Ventos predominantes do Norte.
\item[Populacao]  10.498 (Censo 2022 Final).
\item[Densidade]  \textasciitilde 9,1 hab/\allowbreak km² (Extremamente baixa, ideal para segurança e isolamento).
\item[Vias de Acesso]  Conexão com a Ruta PY13. Infraestrutura rodoviária em desenvolvimento.
\item[Servicos]  Saúde e educação básicas. Dependência de Coronel Oviedo para infraestrutura complexa.
\item[Hidrologia]  Baixo risco de inundações em grande escala. Ocorrências pontuais em pontes de madeira rurais.
\item[Sismicidade]  Nula.
\item[Solo]  Alfisois de boa profundidade, aptos para pecuária e agricultura extensiva.
\item[Agua]  Abundância de recursos hídricos superficiais e subterrâneos.
\item[Energia]  Rede ANDE\footnote{Administración Nacional de Electricidade (Administração Nacional de Eletricidade), companhia estatal de energia.} com cobertura parcial em áreas remotas.
\item[Producao]  Pecuária de corte, milho e mandioca.
\item[Seguranca]  Ambiente rural seguro, baixa densidade dificulta a atuação de grupos criminosos urbanos.
\item[Clima Social]  Estabilidade social baseada na produção agropecuária.
\end{description}
\textbf{Fonte climática:} NASA POWER Climatology API (período 2001-2020)
\textbf{Fonte luminosa:} estimativa world atlas

\paragraph{Irradiação Solar
(kWh/\allowbreak m²/\allowbreak dia)}\label{irradiauxe7uxe3o-solar-kwhmuxb2dia}

\begin{longtable}[]{@{}
  >{\raggedright\arraybackslash}p{(\linewidth - 24\tabcolsep) * \real{0.0746}}
  >{\raggedright\arraybackslash}p{(\linewidth - 24\tabcolsep) * \real{0.0746}}
  >{\raggedright\arraybackslash}p{(\linewidth - 24\tabcolsep) * \real{0.0746}}
  >{\raggedright\arraybackslash}p{(\linewidth - 24\tabcolsep) * \real{0.0746}}
  >{\raggedright\arraybackslash}p{(\linewidth - 24\tabcolsep) * \real{0.0746}}
  >{\raggedright\arraybackslash}p{(\linewidth - 24\tabcolsep) * \real{0.0746}}
  >{\raggedright\arraybackslash}p{(\linewidth - 24\tabcolsep) * \real{0.0746}}
  >{\raggedright\arraybackslash}p{(\linewidth - 24\tabcolsep) * \real{0.0746}}
  >{\raggedright\arraybackslash}p{(\linewidth - 24\tabcolsep) * \real{0.0746}}
  >{\raggedright\arraybackslash}p{(\linewidth - 24\tabcolsep) * \real{0.0746}}
  >{\raggedright\arraybackslash}p{(\linewidth - 24\tabcolsep) * \real{0.0746}}
  >{\raggedright\arraybackslash}p{(\linewidth - 24\tabcolsep) * \real{0.0746}}
  >{\raggedright\arraybackslash}p{(\linewidth - 24\tabcolsep) * \real{0.1045}}@{}}
\toprule\noalign{}
\begin{minipage}[b]{\linewidth}\raggedright
Jan
\end{minipage} & \begin{minipage}[b]{\linewidth}\raggedright
Fev
\end{minipage} & \begin{minipage}[b]{\linewidth}\raggedright
Mar
\end{minipage} & \begin{minipage}[b]{\linewidth}\raggedright
Abr
\end{minipage} & \begin{minipage}[b]{\linewidth}\raggedright
Mai
\end{minipage} & \begin{minipage}[b]{\linewidth}\raggedright
Jun
\end{minipage} & \begin{minipage}[b]{\linewidth}\raggedright
Jul
\end{minipage} & \begin{minipage}[b]{\linewidth}\raggedright
Ago
\end{minipage} & \begin{minipage}[b]{\linewidth}\raggedright
Set
\end{minipage} & \begin{minipage}[b]{\linewidth}\raggedright
Out
\end{minipage} & \begin{minipage}[b]{\linewidth}\raggedright
Nov
\end{minipage} & \begin{minipage}[b]{\linewidth}\raggedright
Dez
\end{minipage} & \begin{minipage}[b]{\linewidth}\raggedright
Média
\end{minipage} \\
\midrule\noalign{}
\endhead
\bottomrule\noalign{}
\endlastfoot
6.56 & 6.04 & 5.59 & 4.57 & 3.37 & 2.97 & 3.27 & 4.04 & 4.58 & 5.29 &
6.28 & 6.54 & \textbf{4.92} \\
\end{longtable}

\textbf{Inclinação solar recomendada:} 25° N (anual) · 35° N (inverno
jun-ago) · 15° N (verão nov-jan)

\paragraph{Precipitação (mm/\allowbreak mês)}\label{precipitauxe7uxe3o-mmmuxeas}

\begin{longtable}[]{@{}
  >{\raggedright\arraybackslash}p{(\linewidth - 24\tabcolsep) * \real{0.0704}}
  >{\raggedright\arraybackslash}p{(\linewidth - 24\tabcolsep) * \real{0.0704}}
  >{\raggedright\arraybackslash}p{(\linewidth - 24\tabcolsep) * \real{0.0704}}
  >{\raggedright\arraybackslash}p{(\linewidth - 24\tabcolsep) * \real{0.0704}}
  >{\raggedright\arraybackslash}p{(\linewidth - 24\tabcolsep) * \real{0.0704}}
  >{\raggedright\arraybackslash}p{(\linewidth - 24\tabcolsep) * \real{0.0704}}
  >{\raggedright\arraybackslash}p{(\linewidth - 24\tabcolsep) * \real{0.0704}}
  >{\raggedright\arraybackslash}p{(\linewidth - 24\tabcolsep) * \real{0.0704}}
  >{\raggedright\arraybackslash}p{(\linewidth - 24\tabcolsep) * \real{0.0704}}
  >{\raggedright\arraybackslash}p{(\linewidth - 24\tabcolsep) * \real{0.0704}}
  >{\raggedright\arraybackslash}p{(\linewidth - 24\tabcolsep) * \real{0.0704}}
  >{\raggedright\arraybackslash}p{(\linewidth - 24\tabcolsep) * \real{0.0704}}
  >{\raggedright\arraybackslash}p{(\linewidth - 24\tabcolsep) * \real{0.1549}}@{}}
\toprule\noalign{}
\begin{minipage}[b]{\linewidth}\raggedright
Jan
\end{minipage} & \begin{minipage}[b]{\linewidth}\raggedright
Fev
\end{minipage} & \begin{minipage}[b]{\linewidth}\raggedright
Mar
\end{minipage} & \begin{minipage}[b]{\linewidth}\raggedright
Abr
\end{minipage} & \begin{minipage}[b]{\linewidth}\raggedright
Mai
\end{minipage} & \begin{minipage}[b]{\linewidth}\raggedright
Jun
\end{minipage} & \begin{minipage}[b]{\linewidth}\raggedright
Jul
\end{minipage} & \begin{minipage}[b]{\linewidth}\raggedright
Ago
\end{minipage} & \begin{minipage}[b]{\linewidth}\raggedright
Set
\end{minipage} & \begin{minipage}[b]{\linewidth}\raggedright
Out
\end{minipage} & \begin{minipage}[b]{\linewidth}\raggedright
Nov
\end{minipage} & \begin{minipage}[b]{\linewidth}\raggedright
Dez
\end{minipage} & \begin{minipage}[b]{\linewidth}\raggedright
Total/\allowbreak ano
\end{minipage} \\
\midrule\noalign{}
\endhead
\bottomrule\noalign{}
\endlastfoot
134 & 138 & 121 & 148 & 174 & 93 & 73 & 51 & 104 & 186 & 185 & 172 &
\textbf{1581 mm} \\
\end{longtable}

\paragraph{Poluição Luminosa}\label{poluiuxe7uxe3o-luminosa}

\begin{longtable}[]{@{}ll@{}}
\toprule\noalign{}
Parâmetro & Valor \\
\midrule\noalign{}
\endhead
\bottomrule\noalign{}
\endlastfoot
Escala Bortle & 4 --- Céu rural-suburbano \\
Radiância artificial & 3.0 nW/\allowbreak cm²/\allowbreak sr \\
\end{longtable}
\paragraph{Solo (SoilGrids 2.0, média ponderada 0--30
cm)}

\begin{longtable}[]{@{}ll@{}}
\toprule\noalign{}
Parâmetro & Valor \\
\midrule\noalign{}
\endhead
\bottomrule\noalign{}
\endlastfoot
pH (H$_2$O) & 5.32 \\
Carbono orgânico (SOC) & 18.97 g/\allowbreak kg \\
Argila & 23.57 \% \\
Areia & 58.52 \% \\
Silte (calc.) & 17.9 \% \\
Densidade aparente & 1.29 g/\allowbreak cm³ \\
\textbf{Aptidão agrícola} & \textbf{Média} \\
\end{longtable}

Fonte: ISRIC SoilGrids 2.0 via WCS. Coords: -24.9167°, -55.75°. Média
ponderada camadas 0-5, 5-15, 15-30 cm.

\begin{itemize}
\tightlist
\item
  \textbf{Agua:} Abundância de recursos hídricos superficiais e
  subterrâneos.
\item
  \textbf{Energia:} Rede ANDE\footnote{Administración Nacional de Electricidade (Administração Nacional de Eletricidade), companhia estatal de energia.} com cobertura parcial em áreas remotas.
\item
  \textbf{Producao:} Pecuária de corte, milho e mandioca.
\end{itemize}
\subsubsection{Indicadores Sociais}

\textbf{Fontes:} PNUD 2020, INE\footnote{Instituto Nacional de Estadística (Instituto Nacional de Estatística), órgão oficial de dados demográficos e censitários.} Censo 2022, INE\footnote{Instituto Nacional de Estadística (Instituto Nacional de Estatística), órgão oficial de dados demográficos e censitários.} EPHC 2023, Ministerio
Público 2024\\
\textbf{Nota:} Valores marcados como \emph{(dept.)} referem-se ao
departamento de Caaguazú; valores \emph{(dist.)} são específicos deste
distrito. Dados marcados \emph{(est.)} são estimativas.

\begin{longtable}[]{@{}
  >{\raggedright\arraybackslash}p{(\linewidth - 4\tabcolsep) * \real{0.4231}}
  >{\raggedright\arraybackslash}p{(\linewidth - 4\tabcolsep) * \real{0.2692}}
  >{\raggedright\arraybackslash}p{(\linewidth - 4\tabcolsep) * \real{0.3077}}@{}}
\toprule\noalign{}
\begin{minipage}[b]{\linewidth}\raggedright
Indicador
\end{minipage} & \begin{minipage}[b]{\linewidth}\raggedright
Valor
\end{minipage} & \begin{minipage}[b]{\linewidth}\raggedright
Âmbito
\end{minipage} \\
\midrule\noalign{}
\endhead
\bottomrule\noalign{}
\endlastfoot
IDH (2020) & 0,715 & dept. \\
Ranking IDH nacional & 7/\allowbreak 18 & dept. \\
Esperança de vida & 73,1 anos & dept. \\
Escolaridade média & 7.9 anos & dist. \\
RNB per capita (USD PPA) & 9.000 & dept. \\
Pobreza monetária (\%) & 28,5\% & dept. \\
Pobreza extrema (\%) & 6,8\% & dept. \\
Índice de Gini & 0,447 & dept. \\
Acesso a água potável (\%) & 79,8\% & dept. \\
Acesso a saneamento (\%) & 78,3\% & dept. \\
Taxa de homicídios (est., /\allowbreak 100k hab) & \textasciitilde6--8 (estimativa,
próxima ou acima da média nacional) & dept. \\
Índice de segurança & médio & dept. \\
População (Censo 2022) & 10.498 hab. & dist. \\
Idade mediana & 26.0 anos & dist. \\
Taxa de fecundidade & N/\allowbreak D & dist. \\
\end{longtable}
\subsubsection{Combustível}

\textbf{Referência:} PETROPAR /\allowbreak  postos locais (2024) \textbf{Tipo de
localidade:} Interior

\begin{longtable}[]{@{}lll@{}}
\toprule\noalign{}
Combustível & USD/\allowbreak litro & Gs/\allowbreak litro (aprox.) \\
\midrule\noalign{}
\endhead
\bottomrule\noalign{}
\endlastfoot
Gasolina 93 oct & 0.98 & 7,252 \\
Gasolina 97 oct (premium) & 1.08 & 7,992 \\
Diesel & 0.91 & 6,734 \\
\end{longtable}

\begin{quote}
Preços podem variar ±5\% conforme posto e sazonalidade. Chaco e interior
remoto apresentam maior variação.
\end{quote}

\subsubsection{Cobertura Celular}

\textbf{Fonte:} CONATEL PY /\allowbreak  operadoras (2024)

\begin{longtable}[]{@{}ll@{}}
\toprule\noalign{}
Parâmetro & Valor \\
\midrule\noalign{}
\endhead
\bottomrule\noalign{}
\endlastfoot
Cobertura 4G população (dept.) & 88\% \\
Cobertura 4G área rural & 72\% \\
Melhor operadora & Tigo \\
Qualidade rural & boa \\
\end{longtable}

\begin{quote}
Para áreas rurais fora do núcleo urbano, recomenda-se chip Tigo como
principal e Personal como backup.
\end{quote}

\subsubsection{Internet}

\textbf{Fonte:} CONATEL /\allowbreak  Speedtest Ookla (2024)

\begin{longtable}[]{@{}ll@{}}
\toprule\noalign{}
Parâmetro & Valor \\
\midrule\noalign{}
\endhead
\bottomrule\noalign{}
\endlastfoot
Velocidade média download & 55 Mbps \\
Domicílios com internet (dept.) & 62\% \\
Tecnologia predominante & rádio/\allowbreak fibra \\
Opção rural & Starlink disponível (\textasciitilde USD 44/\allowbreak mês) \\
\end{longtable}

\subsubsection{Mercado Imobiliário e Terra
Rural}

\textbf{Fonte:} INDERT /\allowbreak  Clasificados.com.py (2024)

\begin{longtable}[]{@{}ll@{}}
\toprule\noalign{}
Tipo & Referência \\
\midrule\noalign{}
\endhead
\bottomrule\noalign{}
\endlastfoot
Terra agrícola alta prod. (USD/\allowbreak ha) & 8,800 \\
Imóvel urbano (USD/\allowbreak m²) & 850 \\
Aluguel 2 quartos (USD/\allowbreak mês) & 350 \\
\end{longtable}

\begin{quote}
Valores de referência departamental. Localidades menores podem ter
preços 20--40\% abaixo da capital departamental. \#\#\# Saúde
\end{quote}

\textbf{Fonte:} MSPBS /\allowbreak  IPS Paraguay (2024-2026), consolidação
departamental e proxy local

\begin{longtable}[]{@{}ll@{}}
\toprule\noalign{}
Serviço & Disponibilidade \\
\midrule\noalign{}
\endhead
\bottomrule\noalign{}
\endlastfoot
USF /\allowbreak  Posto de Saúde & sim \\
Hospital Regional & não \\
IPS (seguro social) & sim \\
Farmácia & sim \\
Distância ao hospital de referência & \textasciitilde101 km (Coronel
Oviedo) \\
\end{longtable}

\textbf{Principais estabelecimentos:} USF local; referencia hospitalar
em Coronel Oviedo

\textbf{Observação para imigrantes:} Atendimento primario local ou em
raio curto; casos de maior complexidade seguem para Coronel Oviedo.
Cobertura privada continua recomendada para especialidades e urgências
de maior complexidade.
\newpage
\secaoDiagnostico{Yhu}{\mapaConteudo{-24.9500}{-55.9333}}{Yhú é um polo em crescimento que combina boa infraestrutura com solo
altamente produtivo. Sua localização estratégica e densidade moderada o
tornam um ponto de equilíbrio para quem busca segurança e potencial
econômico.

\subsubsection{Por que sim}

\begin{itemize}
\tightlist
\item
  Conectividade regional sólida.
\item
  Solo de alta produtividade (Oxisois).
\end{itemize}

\subsubsection{Por que não}

\begin{itemize}
\tightlist
\item
  Pressão populacional superior à de vizinhos como Vaquería.
\end{itemize}

\subsubsection{Recomendação
Técnica}

Ideal para investimentos em silvicultura e agroindústria em ambiente
seguro.

\subsubsection{Analise de Sensibilidade}

\begin{itemize}
\tightlist
\item
  Cenario otimista: GSS 6.5 (Consolidação de Yhú como hub comercial).
\item
  Cenario conservador: GSS 5.7 (Crise no setor de exportação de grãos).
\end{itemize}

\subsubsection{}

\begin{itemize}
\tightlist
\item
  Dados populacionais validados (INE\footnote{Instituto Nacional de Estadística (Instituto Nacional de Estatística), órgão oficial de dados demográficos e censitários.} 2022).
\item
  Mapeamento de solos (MAG/\allowbreak INE\footnote{Instituto Nacional de Estadística (Instituto Nacional de Estatística), órgão oficial de dados demográficos e censitários.}).
\item
  Dados de conectividade (MOPC).
\end{itemize}}
\begin{table}[ht!]
\small \begin{tabular}{p{3.0cm}p{0.8cm}p{6.0cm}} \toprule
\textbf{Critério} & \textbf{Nota} & \textbf{Justificativa} \\ \midrule
A - Ameaças Estratégicas & 5.4 & - \\
B - Risco Social & 6.0 & - \\
C - Riscos Naturais & 7.0 & - \\
D - Autossuficiência & 6.0 & - \\
E - Institucional & 6.5 & - \\
\midrule \textbf{GSS FINAL} & \textbf{6.1} & \textbf{Seguro (Potencial moderado de resiliência).} \\ \bottomrule \end{tabular}
\end{table}
\subsubsection{Vulnerabilidades e
Mitigação}\label{vuln-Yhu-vulnerabilidades-e-mitigauxe7uxe3o}

\begin{itemize}
\item
  \textbf{Logística Florestal:} Dependência de grandes caminhões pode
  congestionar rotas em crises. Mitigação: Mapear rotas rurais para
  veículos leves.
\item
  \textbf{Recursos Hídricos:} Pressão sobre nascentes pelo agronegócio.
  Mitigação: Proteção de APP (Área de Preservação Permanente) na
  propriedade.
\item
  INE\footnote{Instituto Nacional de Estadística (Instituto Nacional de Estatística), órgão oficial de dados demográficos e censitários.} Censo 2022: 22.524 habitantes.
\item
  Preço da Terra: Gs 45M a 60M por hectare.
\item
  Homicídios: 3,5/\allowbreak 100k.
\item
  População distrital: 22.524 habitantes.
\item
  Densidade demografica: 33,2 hab/\allowbreak km2.
\item
  Conectividade viaria: 80/\allowbreak 100.
\item
  Exposicao hidrometeorologica: 30/\allowbreak 100.
\item
  Curto prazo: Sazonalidade climática afetando o setor florestal.
\item
  Médio prazo: Integração logística nacional conectando o norte ao sul
  via Yhú.
\end{itemize}

\subsubsection{Diagnostico Integrado}\label{vuln-Yhu-diagnostico-integrado}

Yhú é um polo em crescimento que combina boa infraestrutura com solo
altamente produtivo. Sua localização estratégica e densidade moderada o
tornam um ponto de equilíbrio para quem busca segurança e potencial
econômico.

\subsubsection{Por que sim}\label{vuln-Yhu-por-que-sim}

\begin{itemize}
\tightlist
\item
  Conectividade regional sólida.
\item
  Solo de alta produtividade (Oxisois).
\end{itemize}

\subsubsection{Por que não}\label{vuln-Yhu-por-que-nuxe3o}

\begin{itemize}
\tightlist
\item
  Pressão populacional superior à de vizinhos como Vaquería.
\end{itemize}

\subsubsection{Recomendação
Técnica}\label{vuln-Yhu-recomendauxe7uxe3o-tuxe9cnica}

Ideal para investimentos em silvicultura e agroindústria em ambiente
seguro.

\begin{itemize}
\item
  Cenario otimista: GSS 6.5 (Consolidação de Yhú como hub comercial).
\item
  Cenario conservador: GSS 5.7 (Crise no setor de exportação de grãos).
\item
  Dados populacionais validados (INE\footnote{Instituto Nacional de Estadística (Instituto Nacional de Estatística), órgão oficial de dados demográficos e censitários.} 2022).
\item
  Mapeamento de solos (MAG/\allowbreak INE\footnote{Instituto Nacional de Estadística (Instituto Nacional de Estatística), órgão oficial de dados demográficos e censitários.}).
\item
  Dados de conectividade (MOPC).
\end{itemize}

\subsubsection{Combustível}\label{vuln-Yhu-combustuxedvel}

\textbf{Referência:} PETROPAR /\allowbreak  postos locais (2024) \textbf{Tipo de
localidade:} Interior

\begin{longtable}[]{@{}lll@{}}
\toprule\noalign{}
Combustível & USD/\allowbreak litro & Gs/\allowbreak litro (aprox.) \\
\midrule\noalign{}
\endhead
\bottomrule\noalign{}
\endlastfoot
Gasolina 93 oct & 0.98 & 7,252 \\
Gasolina 97 oct (premium) & 1.08 & 7,992 \\
Diesel & 0.91 & 6,734 \\
\end{longtable}

\begin{quote}
Preços podem variar ±5\% conforme posto e sazonalidade. Chaco e interior
remoto apresentam maior variação.
\end{quote}

\subsubsection{Cobertura Celular}\label{vuln-Yhu-cobertura-celular}

\textbf{Fonte:} CONATEL PY /\allowbreak  operadoras (2024)

\begin{longtable}[]{@{}ll@{}}
\toprule\noalign{}
Parâmetro & Valor \\
\midrule\noalign{}
\endhead
\bottomrule\noalign{}
\endlastfoot
Cobertura 4G população (dept.) & 88\% \\
Cobertura 4G área rural & 72\% \\
Melhor operadora & Tigo \\
Qualidade rural & boa \\
\end{longtable}

\begin{quote}
Para áreas rurais fora do núcleo urbano, recomenda-se chip Tigo como
principal e Personal como backup.
\end{quote}

\subsubsection{Internet}\label{vuln-Yhu-internet}

\textbf{Fonte:} CONATEL /\allowbreak  Speedtest Ookla (2024)

\begin{longtable}[]{@{}ll@{}}
\toprule\noalign{}
Parâmetro & Valor \\
\midrule\noalign{}
\endhead
\bottomrule\noalign{}
\endlastfoot
Velocidade média download & 55 Mbps \\
Domicílios com internet (dept.) & 62\% \\
Tecnologia predominante & rádio/\allowbreak fibra \\
Opção rural & Starlink disponível (\textasciitilde USD 44/\allowbreak mês) \\
\end{longtable}

\subsubsection{Mercado Imobiliário e Terra
Rural}\label{vuln-Yhu-mercado-imobiliuxe1rio-e-terra-rural}

\textbf{Fonte:} INDERT /\allowbreak  Clasificados.com.py (2024)

\begin{longtable}[]{@{}ll@{}}
\toprule\noalign{}
Tipo & Referência \\
\midrule\noalign{}
\endhead
\bottomrule\noalign{}
\endlastfoot
Terra agrícola alta prod. (USD/\allowbreak ha) & 8,800 \\
Imóvel urbano (USD/\allowbreak m²) & 850 \\
Aluguel 2 quartos (USD/\allowbreak mês) & 350 \\
\end{longtable}

\begin{quote}
Valores de referência departamental. Localidades menores podem ter
preços 20--40\% abaixo da capital departamental. \#\#\# Saúde
\end{quote}

\textbf{Fonte:} MSPBS /\allowbreak  IPS Paraguay (2024-2026), consolidação
departamental e proxy local

\begin{longtable}[]{@{}ll@{}}
\toprule\noalign{}
Serviço & Disponibilidade \\
\midrule\noalign{}
\endhead
\bottomrule\noalign{}
\endlastfoot
USF /\allowbreak  Posto de Saúde & sim \\
Hospital Regional & não \\
IPS (seguro social) & sim \\
Farmácia & sim \\
Distância ao hospital de referência & \textasciitilde85 km (Coronel
Oviedo) \\
\end{longtable}

\textbf{Principais estabelecimentos:} USF local; referencia hospitalar
em Coronel Oviedo

\textbf{Observação para imigrantes:} Atendimento primario local ou em
raio curto; casos de maior complexidade seguem para Coronel Oviedo.
Cobertura privada continua recomendada para especialidades e urgências
de maior complexidade.
\subsubsection{Dossiê de Campo}\label{dossie-Yhu-dossiuxea-de-campo}
\begin{description}[style=nextline,leftmargin=0.5cm,font=\bfseries]
\item[Coordenadas]  24°57'S 55°56'W.
\item[Topografia]  Relevo predominantemente plano a suavemente ondulado. Área de \textasciitilde 678 km².
\item[Alvos Estrategicos]  Centro de serviços regional em expansão na Ruta PY13. Áreas de produção florestal e agrícola.
\item[Fallout]  Ventos predominantes do Norte/\allowbreak Nordeste.
\item[Populacao]  22.524 (Censo 2022 Final).
\item[Densidade]  \textasciitilde 33,2 hab/\allowbreak km².
\item[Vias de Acesso]  Cruzamento de rotas importantes conectando o departamento de Caaguazú ao norte do país.
\item[Servicos]  Rede de saúde em estruturação e comércio regional consolidado.
\item[Hidrologia]  Rio Yhú e seus afluentes. Risco moderado de inundações em áreas ribeirinhas baixas.
\item[Sismicidade]  Nula.
\item[Solo]  Oxisois de alta qualidade, aptos para reflorestamento e agricultura de grãos.
\item[Agua]  Recursos hídricos superficiais e subterrâneos abundantes.
\item[Energia]  Rede ANDE\footnote{Administración Nacional de Electricidade (Administração Nacional de Eletricidade), companhia estatal de energia.}.
\item[Producao]  Madeira, soja, milho e mandioca.
\item[Seguranca]  Zoneamento urbano e rural seguro. Taxa de homicídios acompanhando a média departamental baixa (3,5/\allowbreak 100k).
\item[Clima Social]  Sociedade tradicional em transição para o agronegócio mecanizado.
\end{description}
\textbf{Fonte climática:} NASA POWER Climatology API (período 2001-2020)
\textbf{Fonte luminosa:} estimativa world atlas

\paragraph{Irradiação Solar
(kWh/\allowbreak m²/\allowbreak dia)}\label{irradiauxe7uxe3o-solar-kwhmuxb2dia}

\begin{longtable}[]{@{}
  >{\raggedright\arraybackslash}p{(\linewidth - 24\tabcolsep) * \real{0.0746}}
  >{\raggedright\arraybackslash}p{(\linewidth - 24\tabcolsep) * \real{0.0746}}
  >{\raggedright\arraybackslash}p{(\linewidth - 24\tabcolsep) * \real{0.0746}}
  >{\raggedright\arraybackslash}p{(\linewidth - 24\tabcolsep) * \real{0.0746}}
  >{\raggedright\arraybackslash}p{(\linewidth - 24\tabcolsep) * \real{0.0746}}
  >{\raggedright\arraybackslash}p{(\linewidth - 24\tabcolsep) * \real{0.0746}}
  >{\raggedright\arraybackslash}p{(\linewidth - 24\tabcolsep) * \real{0.0746}}
  >{\raggedright\arraybackslash}p{(\linewidth - 24\tabcolsep) * \real{0.0746}}
  >{\raggedright\arraybackslash}p{(\linewidth - 24\tabcolsep) * \real{0.0746}}
  >{\raggedright\arraybackslash}p{(\linewidth - 24\tabcolsep) * \real{0.0746}}
  >{\raggedright\arraybackslash}p{(\linewidth - 24\tabcolsep) * \real{0.0746}}
  >{\raggedright\arraybackslash}p{(\linewidth - 24\tabcolsep) * \real{0.0746}}
  >{\raggedright\arraybackslash}p{(\linewidth - 24\tabcolsep) * \real{0.1045}}@{}}
\toprule\noalign{}
\begin{minipage}[b]{\linewidth}\raggedright
Jan
\end{minipage} & \begin{minipage}[b]{\linewidth}\raggedright
Fev
\end{minipage} & \begin{minipage}[b]{\linewidth}\raggedright
Mar
\end{minipage} & \begin{minipage}[b]{\linewidth}\raggedright
Abr
\end{minipage} & \begin{minipage}[b]{\linewidth}\raggedright
Mai
\end{minipage} & \begin{minipage}[b]{\linewidth}\raggedright
Jun
\end{minipage} & \begin{minipage}[b]{\linewidth}\raggedright
Jul
\end{minipage} & \begin{minipage}[b]{\linewidth}\raggedright
Ago
\end{minipage} & \begin{minipage}[b]{\linewidth}\raggedright
Set
\end{minipage} & \begin{minipage}[b]{\linewidth}\raggedright
Out
\end{minipage} & \begin{minipage}[b]{\linewidth}\raggedright
Nov
\end{minipage} & \begin{minipage}[b]{\linewidth}\raggedright
Dez
\end{minipage} & \begin{minipage}[b]{\linewidth}\raggedright
Média
\end{minipage} \\
\midrule\noalign{}
\endhead
\bottomrule\noalign{}
\endlastfoot
6.58 & 6.08 & 5.45 & 4.47 & 3.30 & 2.90 & 3.20 & 3.96 & 4.53 & 5.31 &
6.35 & 6.66 & \textbf{4.9} \\
\end{longtable}

\textbf{Inclinação solar recomendada:} 25° N (anual) · 35° N (inverno
jun-ago) · 15° N (verão nov-jan)

\paragraph{Precipitação (mm/\allowbreak mês)}\label{precipitauxe7uxe3o-mmmuxeas}

\begin{longtable}[]{@{}
  >{\raggedright\arraybackslash}p{(\linewidth - 24\tabcolsep) * \real{0.0704}}
  >{\raggedright\arraybackslash}p{(\linewidth - 24\tabcolsep) * \real{0.0704}}
  >{\raggedright\arraybackslash}p{(\linewidth - 24\tabcolsep) * \real{0.0704}}
  >{\raggedright\arraybackslash}p{(\linewidth - 24\tabcolsep) * \real{0.0704}}
  >{\raggedright\arraybackslash}p{(\linewidth - 24\tabcolsep) * \real{0.0704}}
  >{\raggedright\arraybackslash}p{(\linewidth - 24\tabcolsep) * \real{0.0704}}
  >{\raggedright\arraybackslash}p{(\linewidth - 24\tabcolsep) * \real{0.0704}}
  >{\raggedright\arraybackslash}p{(\linewidth - 24\tabcolsep) * \real{0.0704}}
  >{\raggedright\arraybackslash}p{(\linewidth - 24\tabcolsep) * \real{0.0704}}
  >{\raggedright\arraybackslash}p{(\linewidth - 24\tabcolsep) * \real{0.0704}}
  >{\raggedright\arraybackslash}p{(\linewidth - 24\tabcolsep) * \real{0.0704}}
  >{\raggedright\arraybackslash}p{(\linewidth - 24\tabcolsep) * \real{0.0704}}
  >{\raggedright\arraybackslash}p{(\linewidth - 24\tabcolsep) * \real{0.1549}}@{}}
\toprule\noalign{}
\begin{minipage}[b]{\linewidth}\raggedright
Jan
\end{minipage} & \begin{minipage}[b]{\linewidth}\raggedright
Fev
\end{minipage} & \begin{minipage}[b]{\linewidth}\raggedright
Mar
\end{minipage} & \begin{minipage}[b]{\linewidth}\raggedright
Abr
\end{minipage} & \begin{minipage}[b]{\linewidth}\raggedright
Mai
\end{minipage} & \begin{minipage}[b]{\linewidth}\raggedright
Jun
\end{minipage} & \begin{minipage}[b]{\linewidth}\raggedright
Jul
\end{minipage} & \begin{minipage}[b]{\linewidth}\raggedright
Ago
\end{minipage} & \begin{minipage}[b]{\linewidth}\raggedright
Set
\end{minipage} & \begin{minipage}[b]{\linewidth}\raggedright
Out
\end{minipage} & \begin{minipage}[b]{\linewidth}\raggedright
Nov
\end{minipage} & \begin{minipage}[b]{\linewidth}\raggedright
Dez
\end{minipage} & \begin{minipage}[b]{\linewidth}\raggedright
Total/\allowbreak ano
\end{minipage} \\
\midrule\noalign{}
\endhead
\bottomrule\noalign{}
\endlastfoot
134 & 138 & 121 & 148 & 174 & 93 & 73 & 51 & 104 & 186 & 185 & 172 &
\textbf{1581 mm} \\
\end{longtable}

\paragraph{Poluição Luminosa}\label{poluiuxe7uxe3o-luminosa}

\begin{longtable}[]{@{}ll@{}}
\toprule\noalign{}
Parâmetro & Valor \\
\midrule\noalign{}
\endhead
\bottomrule\noalign{}
\endlastfoot
Escala Bortle & 4 --- Céu rural-suburbano \\
Radiância artificial & 3.0 nW/\allowbreak cm²/\allowbreak sr \\
\end{longtable}
\paragraph{Solo (SoilGrids 2.0, média ponderada 0--30
cm)}

\begin{longtable}[]{@{}ll@{}}
\toprule\noalign{}
Parâmetro & Valor \\
\midrule\noalign{}
\endhead
\bottomrule\noalign{}
\endlastfoot
pH (H$_2$O) & 5.4 \\
Carbono orgânico (SOC) & 19.58 g/\allowbreak kg \\
Argila & 22.01 \% \\
Areia & 61.24 \% \\
Silte (calc.) & 16.7 \% \\
Densidade aparente & 1.29 g/\allowbreak cm³ \\
\textbf{Aptidão agrícola} & \textbf{Média} \\
\end{longtable}

Fonte: ISRIC SoilGrids 2.0 via WCS. Coords: -24.95°, -55.9333°. Média
ponderada camadas 0-5, 5-15, 15-30 cm.

\begin{itemize}
\tightlist
\item
  \textbf{Agua:} Recursos hídricos superficiais e subterrâneos
  abundantes.
\item
  \textbf{Energia:} Rede ANDE\footnote{Administración Nacional de Electricidade (Administração Nacional de Eletricidade), companhia estatal de energia.}.
\item
  \textbf{Producao:} Madeira, soja, milho e mandioca.
\end{itemize}
\subsubsection{Indicadores Sociais}

\textbf{Fontes:} PNUD 2020, INE\footnote{Instituto Nacional de Estadística (Instituto Nacional de Estatística), órgão oficial de dados demográficos e censitários.} Censo 2022, INE\footnote{Instituto Nacional de Estadística (Instituto Nacional de Estatística), órgão oficial de dados demográficos e censitários.} EPHC 2023, Ministerio
Público 2024\\
\textbf{Nota:} Valores marcados como \emph{(dept.)} referem-se ao
departamento de Caaguazú; valores \emph{(dist.)} são específicos deste
distrito. Dados marcados \emph{(est.)} são estimativas.

\begin{longtable}[]{@{}
  >{\raggedright\arraybackslash}p{(\linewidth - 4\tabcolsep) * \real{0.4231}}
  >{\raggedright\arraybackslash}p{(\linewidth - 4\tabcolsep) * \real{0.2692}}
  >{\raggedright\arraybackslash}p{(\linewidth - 4\tabcolsep) * \real{0.3077}}@{}}
\toprule\noalign{}
\begin{minipage}[b]{\linewidth}\raggedright
Indicador
\end{minipage} & \begin{minipage}[b]{\linewidth}\raggedright
Valor
\end{minipage} & \begin{minipage}[b]{\linewidth}\raggedright
Âmbito
\end{minipage} \\
\midrule\noalign{}
\endhead
\bottomrule\noalign{}
\endlastfoot
IDH (2020) & 0,715 & dept. \\
Ranking IDH nacional & 7/\allowbreak 18 & dept. \\
Esperança de vida & 73,1 anos & dept. \\
Escolaridade média & 7.3 anos & dist. \\
RNB per capita (USD PPA) & 9.000 & dept. \\
Pobreza monetária (\%) & 28,5\% & dept. \\
Pobreza extrema (\%) & 6,8\% & dept. \\
Índice de Gini & 0,447 & dept. \\
Acesso a água potável (\%) & 79,8\% & dept. \\
Acesso a saneamento (\%) & 78,3\% & dept. \\
Taxa de homicídios (est., /\allowbreak 100k hab) & \textasciitilde6--8 (estimativa,
próxima ou acima da média nacional) & dept. \\
Índice de segurança & médio & dept. \\
População (Censo 2022) & 22.524 hab. & dist. \\
Idade mediana & 26.0 anos & dist. \\
Taxa de fecundidade & N/\allowbreak D & dist. \\
\end{longtable}
\subsubsection{Combustível}

\textbf{Referência:} PETROPAR /\allowbreak  postos locais (2024) \textbf{Tipo de
localidade:} Interior

\begin{longtable}[]{@{}lll@{}}
\toprule\noalign{}
Combustível & USD/\allowbreak litro & Gs/\allowbreak litro (aprox.) \\
\midrule\noalign{}
\endhead
\bottomrule\noalign{}
\endlastfoot
Gasolina 93 oct & 0.98 & 7,252 \\
Gasolina 97 oct (premium) & 1.08 & 7,992 \\
Diesel & 0.91 & 6,734 \\
\end{longtable}

\begin{quote}
Preços podem variar ±5\% conforme posto e sazonalidade. Chaco e interior
remoto apresentam maior variação.
\end{quote}

\subsubsection{Cobertura Celular}

\textbf{Fonte:} CONATEL PY /\allowbreak  operadoras (2024)

\begin{longtable}[]{@{}ll@{}}
\toprule\noalign{}
Parâmetro & Valor \\
\midrule\noalign{}
\endhead
\bottomrule\noalign{}
\endlastfoot
Cobertura 4G população (dept.) & 88\% \\
Cobertura 4G área rural & 72\% \\
Melhor operadora & Tigo \\
Qualidade rural & boa \\
\end{longtable}

\begin{quote}
Para áreas rurais fora do núcleo urbano, recomenda-se chip Tigo como
principal e Personal como backup.
\end{quote}

\subsubsection{Internet}

\textbf{Fonte:} CONATEL /\allowbreak  Speedtest Ookla (2024)

\begin{longtable}[]{@{}ll@{}}
\toprule\noalign{}
Parâmetro & Valor \\
\midrule\noalign{}
\endhead
\bottomrule\noalign{}
\endlastfoot
Velocidade média download & 55 Mbps \\
Domicílios com internet (dept.) & 62\% \\
Tecnologia predominante & rádio/\allowbreak fibra \\
Opção rural & Starlink disponível (\textasciitilde USD 44/\allowbreak mês) \\
\end{longtable}

\subsubsection{Mercado Imobiliário e Terra
Rural}

\textbf{Fonte:} INDERT /\allowbreak  Clasificados.com.py (2024)

\begin{longtable}[]{@{}ll@{}}
\toprule\noalign{}
Tipo & Referência \\
\midrule\noalign{}
\endhead
\bottomrule\noalign{}
\endlastfoot
Terra agrícola alta prod. (USD/\allowbreak ha) & 8,800 \\
Imóvel urbano (USD/\allowbreak m²) & 850 \\
Aluguel 2 quartos (USD/\allowbreak mês) & 350 \\
\end{longtable}

\begin{quote}
Valores de referência departamental. Localidades menores podem ter
preços 20--40\% abaixo da capital departamental. \#\#\# Saúde
\end{quote}

\textbf{Fonte:} MSPBS /\allowbreak  IPS Paraguay (2024-2026), consolidação
departamental e proxy local

\begin{longtable}[]{@{}ll@{}}
\toprule\noalign{}
Serviço & Disponibilidade \\
\midrule\noalign{}
\endhead
\bottomrule\noalign{}
\endlastfoot
USF /\allowbreak  Posto de Saúde & sim \\
Hospital Regional & não \\
IPS (seguro social) & sim \\
Farmácia & sim \\
Distância ao hospital de referência & \textasciitilde85 km (Coronel
Oviedo) \\
\end{longtable}

\textbf{Principais estabelecimentos:} USF local; referencia hospitalar
em Coronel Oviedo

\textbf{Observação para imigrantes:} Atendimento primario local ou em
raio curto; casos de maior complexidade seguem para Coronel Oviedo.
Cobertura privada continua recomendada para especialidades e urgências
de maior complexidade.
